
%%%%%%%%%%%%%%%%%%%%%%%%%%%%%%%%%%
%\newpage
%\part{Preliminaries}
%%\section{Preliminaries}
%\label{partpreliminaries}
%This part will be divided as follows: In Section \ref{sectionROSSONCTs} we define the class of rank one symmetric spaces of noncompact type (\ROSSs). In Sections \ref{sectiongeometry1}-\ref{sectiongeometry2}, we define the class of hyperbolic metric spaces and study their geometry. In Section \ref{sectiondiscreteness}, we explore different notions of discreteness for groups of isometries of a metric space. In Section \ref{sectionclassification} we prove two classification theorems, one for isometries (Theorem \ref{classificationofisometries}) and one for semigroups of isometries (Theorem \ref{classificationofsemigroups}). Finally, in Section \ref{sectionlimitsets} we define and study the limit set of a semigroup of isometries.

\chapter{Algebraic hyperbolic spaces}
\label{sectionROSSONCTs}
%%%%%%%%%%%%%%%%%%%%%%%%%%%%%%%%%%

In this chapter we introduce our main objects of interest, algebraic hyperbolic spaces in finite and infinite dimensions. References for the theory of finite-dimensional algebraic hyperbolic spaces, which are often called \emph{rank one symmetric spaces of noncompact type}, include \cite{BridsonHaefliger,CFKP,MackayTyson}. Infinite-dimensional algebraic hyperbolic spaces, as well as some non-hyperbolic infinite-dimensional symmetric spaces, have been discussed in \cite{Duchesne}.


%\ignore{
%
%%%%%%%%%%%%%%%%%%%%%%%%%%%%%%%%%%%
%\bigskip
%\subsection{Infinite-dimensional Riemannian manifolds}
%%%%%%%%%%%%%%%%%%%%%%%%%%%%%%%%%%%
%
%The infinite-dimensional hyperbolic space is an infinite-dimensional Riemannian manifold, so we begin by reviewing the theory of infinite-dimensional Riemannian manifolds. A good reference for this subsection is \cite{Lang_differential_geometry}.
%
%\begin{definition} \label{definitionhilbertmanifold}
%A ($\CC^\infty$) \emph{Hilbert manifold} is a Banach manifold modelled on a real Hilbert space $(\HH,\|\cdot\|)$, i.e. a separable metrizable topological space $(X,\scrT)$ together with a collection of triples $(U_\alpha,V_\alpha,\Phi_\alpha)_{\alpha\in I}$ (called an \emph{atlas}), where
%\begin{itemize}
%\item[(I)] $(U_\alpha)_{\alpha\in I}$ is an open cover of $X$,
%\item[(II)] $(V_\alpha)_{\alpha\in I}$ are open subsets of $\HH$,
%\item[(III)] for each $\alpha\in I$, $\Phi_\alpha:U_\alpha\to V_\alpha$ is a homeomorphism, and
%\item[(IV)] for each $\alpha,\beta\in I$, the \emph{transition map}
%\begin{equation}
%\label{transition}
%\Phi_\beta\circ\Phi_\alpha^{-1}:\Phi_\alpha(U_\alpha\cap U_\beta)\to \Phi_\beta(U_\alpha\cap U_\beta)
%\end{equation}
%is a $\CC^\infty$ diffeomorphism. ()
%\end{itemize}
%For the most part we will not describe the atlases of our manifolds explicitly; it should not be too difficult to construct them. If $X$ is a Hilbert manifold, then for each $x\in X$ we may define the \emph{tangent space} $T_x X$ (see \cite[\6 II.1]{Lang_differential_geometry}), a topological vector space isomorphic to $\HH$. For each $\alpha\in I$ and for each $\xx\in V_\alpha$, there is a natural map $(\Phi_\alpha^{-1})'(\xx):\HH\to T_{\Phi_\alpha^{-1}(\xx)} X$ which is an isomorphism of topological vector spaces.
%\end{definition}
%\begin{definition} \label{definitionIDRM}
%An \emph{infinite-dimensional Riemannian manifold} is a Hilbert manifold equipped with a Riemannian structure, i.e. a field of inner products $( \lb \cdot , \cdot \rb_x )_{x\in X}$ such that:
%\begin{itemize}
%\item[(I)] for each $x\in X$, $\lb \cdot , \cdot \rb_x$ is an inner product (i.e. a positive definite symmetric bilinear form) on the tangent space $T_x X$,
%\item[(II)] for each $\alpha\in I$, the map
%\[
%(\xx,\uu,\vv) \mapsto \lb \uu,\vv\rb_{\xx,\alpha} := \big\lb (\Phi_\alpha^{-1})'(\xx)[\uu],(\Phi_\alpha^{-1})'(\xx)[\vv]\big\rb_{\Phi_\alpha^{-1}(\xx)}
%\]
%is $\CC^\infty$,
%\item[(III)] for each $\alpha\in I$ and for each $\xx\in V_\alpha$, there exists $C_{\xx,\alpha} > 0$ such that for all $u\in\HH$ 
%\[
%\frac{1}{C_{\xx,\alpha}} \| \uu \| \leq \| \uu \|_{\xx,\alpha} \leq C_{\xx,\alpha} \| \uu \|,
%\]
%where
%\[
%\| \uu \|_{\xx,\alpha} := \sqrt{\lb \uu,\uu \rb_{\xx,\alpha}}.
%\]
%\end{itemize}
%\end{definition}
%\ignore{
%\begin{observation}
%For each $x\in X$, it suffices to check condition (III) of Definition \ref{definitionIDRM} for a single $\alpha\in I$ for which $x\in U_\alpha$.
%\end{observation}
%\begin{proof}
%Let $\beta\in I$ be another index for which $x\in U_\beta$. Since the transition map \eqref{transition} is assumed to be a diffeomorphism, its derivative at $\Phi_\alpha(x)$ is an isomorphism of $\HH$ (in the Banach space sense). Letting
%\[
%C_{x,\beta} = K C_{x,\alpha},
%\]
%where $K$ is the bi-Lipschitz constant of this isomorphism, finishes the proof.
%\end{proof}
%}% end ignore
%
%Let $X$ be an infinite-dimensional Riemannian manifold.
%\begin{definition}
%Let $\gamma:[a,b]\to X$ be a piecewise $\CC^1$ map. The \emph{length} of $\gamma$ is
%\[
%\ell_X(\gamma) := \int_a^b \|\gamma'(t)\|_{\gamma(t)}dt.
%\]
%For two points $x,y\in X$, the \emph{distance} between $x$ and $y$ is
%\[
%\dist_X(x,y):= \inf_{\gamma: x \leadsto y} \ell_X(\gamma) \cdot
%\]
%Here the notation $\gamma:x\leadsto y$ means that $\gamma$ is a path from $x$ to $y$, i.e. $\gamma(a) = x$ and $\gamma(b) = y$.
%\end{definition}
%
%\begin{proposition}[{\cite[Proposition 6.1]{Lang_differential_geometry}}]
%$\dist_X:X\times X\to\Rplus$ is a metric on $X$ compatible with the topology $\scrT$.
%\end{proposition}
%
%\begin{example}
%Hilbert space $\HH$ can be considered by itself as an infinite-dimensional Riemannian manifold, with at each point the standard inner product
%\[
%\lb \uu,\vv\rb_\xx := \sum_{i = 1}^\infty u_i v_i.
%\]
%The atlas is just $I = \{1\}$, $U_1 = V_1 = \HH$, and $\Phi_1 = \id$, where $\id$ is the identity map.
%\end{example}
%
%\begin{example}
%\label{examplewhatH}
%The space $\what\HH := \HH\cup\{\infty\}$ can be given the structure of a Hilbert manifold, by letting $I = \{1,2\}$ and
%\begin{align*}
%U_1 &:= \HH & V_1 &:= \HH & \Phi_1 &:= \id\\
%U_2 &:= \what\HH\butnot\{\0\} & V_2 &:= \HH & \Phi_2(\xx) &:= \frac{\xx}{\|\xx\|^2}.
%\end{align*}
%The topology on $\what\HH$ is defined as follows: a subset $U \subset \what\HH$ is open if and only if $U \cap \HH$ is open and if
%\[
%\infty\in U \;\Rightarrow\; \text{$\HH \setminus U$ is bounded}.
%\]
%\end{example}
%\begin{warning}
%The topology on $\what\HH$ is not a one-point compactification. Indeed, $\what\HH$ with the topology defined above is not a compact space.
%\end{warning}
%%\begin{remark}
%%We shall study the map $\Phi_2$ and similar maps in Subsection \ref{subsectionconformal}.
%%\end{remark}
%
%\subsubsection{Geodesics}
%Let $X$ be an infinite-dimensional Riemannian manifold.
%
%\begin{definition}
%A piecewise $\CC^1$ map $\gamma:[a,b]\to X$ is a \emph{length-minimizing geodesic} if
%\[
%\ell_X(\gamma) = \dist_X(\gamma(a),\gamma(b)).
%\]
%It is \emph{travelling at unit speed} if $\|\gamma'(t)\|_{\gamma(t)} = 1$ for all $t\in(a,b)$.
%\end{definition}
%
%\begin{remark}
%In finite-dimensional Riemannian geometry, the Hopf--Rinow theorem asserts the existence of a length-minimizing geodesic connecting any two points on a Riemannian manifold. In infinite dimensions, it is possible that such a geodesic does not exist; see \cite{Atkin}. However, we will see (Subsection \ref{subsectiongeodesics}) that in any infinite-dimensional nonpositively curved manifold, any two points can be connected by a geodesic (in fact, a unique geodesic).
%\end{remark}
%
%\begin{proposition}
%\label{propositiongeodesicequivalent}
%Let $\gamma:[a,b]\to X$ be a continuous map. The following are equivalent:
%\begin{itemize}
%\item[(A)] $\gamma$ is a length-minimizing geodesic travelling at unit speed.
%\item[(B)] $\gamma$ is an isometric embedding, i.e. $\dist_X(\gamma(t_1),\gamma(t_2)) = |t_2 - t_1|$ for all $t_1,t_2\in[a,b]$.
%\end{itemize}
%\end{proposition}
%\begin{proof}[Proof of \text{(A)} \implies \text{(B)}] First, since $\gamma$ is travelling at unit speed, we have $\dist_X(\gamma(t_1),\gamma(t_2)) \leq |t_2 - t_1|$ for all $t_1,t_2\in[a,b]$. Thus
%\begin{align*}
%\dist_X(\gamma(t_1),\gamma(t_2))
%&\geq \dist_X(\gamma(a),\gamma(b)) - \dist_X(\gamma(a),\gamma(t_1)) - \dist_X(\gamma(t_2),\gamma(t_1)) \noreason \\
%&\geq \ell_X(\gamma) - (t_1 - a) - (b - t_2)\\
%&= (b - a) - (t_1 - a) - (b - t_2) \since{$\gamma$ is travelling at unit speed}\\
%&= t_2 - t_1.
%\end{align*}
%\end{proof}
%\begin{proof}[Proof of \text{(B)} \implies \text{(A)}] See \cite[Lemma 4]{GJR}.
%\end{proof}
%
%\begin{remark}
%Below, we will define a \emph{geodesic segment} in a metric space to be the image of an isometric embedding of an interval. By Proposition \ref{propositiongeodesicequivalent}, for a Riemannian manifold this is the same as defining a geodesic segment to be the image of a length-minimizing geodesic. (Here we use the fact that any piecewise $\CC^1$ map can be reparametrized to travel at unit speed.)
%\end{remark}
%
%
%\subsubsection{Isometries of infinite-dimensional Riemannian manifolds}
%
%\begin{definition}
%\label{definitionRiemannianisomorphism}
%Let $X$ and $Y$ be infinite-dimensional Riemannian manifolds, and let $\Psi:X\to Y$ be a $\CC^\infty$ diffeomorphism. $\Psi$ is an \emph{isomorphism} if 
%\[
%\lb \Psi'(x)[u],\Psi'(x)[v]\rb_{\Psi(x),Y} = \lb u,v\rb_{x,X}
%\]
%for all $x\in X$ and for all $u,v\in T_x X$.
%\end{definition}
%
%\begin{observation}
%\label{observationRiemannianisomorphism}
%Ever isomorphism of Riemannian manifolds $\Psi:X\to Y$ is also an \emph{isometry}, i.e.
%\begin{equation}
%\label{Riemannianisometry}
%\dist_Y(\Psi(x_1),\Psi(x_2)) = \dist_X(x_1,x_2) \all x_1,x_2\in X.
%\end{equation}
%\end{observation}
%
%The converse is also true, but much less trivial:
%
%\begin{theorem}[{\cite[Theorem 7]{GJR}}]
%\label{theoremmyerssteenrod}
%Let $X$ and $Y$ be infinite-dimensional Riemannian manifolds, and let $\Psi:X\to Y$ be a bijection. If $\Psi$ satisfies \eqref{Riemannianisometry}, then $\Psi$ is an isomorphism (and in particular is $\CC^\infty$).
%\end{theorem}
%
%We remark that in finite-dimensions, Theorem \ref{theoremmyerssteenrod} is due to S. B. Myers and N. E. Steenrod \cite{MyersSteenrod}.
%
%}% end ignore

%%%%%%%%%%%%%%%%%%%%%%%%%%%%%%%%%%
\bigskip
\section{The definition}
\label{subsectionROSSONCT}
%%%%%%%%%%%%%%%%%%%%%%%%%%%%%%%%%%

Finite-dimensional rank one symmetric spaces of noncompact type come in four flavors, corresponding to the classical division algebras $\R$, $\C$, $\Q$ (quaternions), and $\mathbb O$ (octonions).\Footnote{We denote the quaternions by $\Q$ in order to avoid confusion with the rank one symmetric space of noncompact type (defined over $\Q$) itself, which we will denote by $\H$. Be aware that $\Q$ should not be confused with the set of rational numbers.} The first three division algebras have corresponding rank one symmetric spaces of noncompact type of arbitrary dimension, but there is only one rank one symmetric space of noncompact type corresponding to the octonions; it occurs in dimension two (which corresponds to real dimension 16). Consequently, the octonion rank one symmetric space of noncompact type (known as the \emph{Cayley hyperbolic plane}\Footnote{Not to be confused with the \emph{Cayley plane}, a different mathematical object.}) does not have an infinite-dimensional analogue, while the other three classes do admit infinite-dimensional analogues.

%\begin{remark}
The rank one symmetric spaces of noncompact type corresponding to $\R$ have constant negative curvature. However, those corresponding to the other division algebras have variable negative curvature \cite[Lemmas 2.3, 2.7, 2.11]{Quint} (see also \cite[Corollary of Proposition 4]{Heintze}).
%\end{remark}

\begin{remark}
\label{remarkH2O}
In this monograph we will use the term ``algebraic hyperbolic spaces'' to refer to all rank one symmetric spaces of noncompact type \emph{except} the Cayley hyperbolic plane $\H_{\mathbb O}^2$, in order to avoid dealing with the complicated algebra of the octonions.\Footnote{The complications come from the fact that the octonions are not associative, thus making it somewhat unclear what it means to say that $\amsbb O^3$ is a vector space ``over'' the octonions, since in general $(\xx a) b \neq \xx (a b)$.} However, we feel confident that all the theorems regarding algebraic hyperbolic spaces in this monograph can be generalized to the Cayley hyperbolic plane (possibly after modifying the statements slightly). We leave this task to an algebraist.

For the reader interested in learning more about the Cayley hyperbolic plane, see \cite[pp.136-139]{Mostow}, \cite{SpringerVeldkamp}, or \cite{Allcock}; see also \cite{Baez} for an excellent introduction to octonions in general.
\end{remark}

Fix $\F\in\{\R,\C,\Q\}$ and an index set $J$, and let us construct an algebraic hyperbolic space (i.e. a rank one symmetric space of noncompact type) over the field $\F$ in dimension $\#(J)$. We remark that usually we will let $J = \Namer = \{1,2,\ldots\}$, but occasionally $J$ may be an uncountable set. Let
\[
\HH = \HH_\F^J := \left\{ \xx = (x_i)_{i\in J}\in\F^J \left\vert\;\; \sum_{i\in J} |x_i|^2 < \infty
\right.\right\},
\]
and for $\xx\in\HH$ let
\[
\|\xx\| := \left( \sum_{i\in J} |x_i|^2 \right)^{1/2}.
\]
We will think of $\HH$ as a right $\F$-module, so scalars will always act on the right.\Footnote{The advantage of this convention is that it allows operators to act on the left.} Note that
\[
\|\xx a\| = |a|\cdot \|\xx\| \all \xx\in \HH \all a\in\F.
\]
A \emph{sesquilinear form on $\HH$} is an $\R$-bilinear map $B(\cdot,\cdot):\HH\times\HH\to \F$ satisfying
\[
B(\xx a,\yy) = \wbar a B( \xx,\yy) \text{ and } B( \xx,\yy a) = B( \xx,\yy) a.\Footnote{In the case $\amsbb F = \amsbb C$, this disagrees with the usual convention; we follow here the convention of \cite[\63.3.1]{MackayTyson}.}
\]
Here and from now on $\wbar a$ denotes the conjugate of a complex or quaternionic number $a\in\F$; if $\F = \R$, then $\wbar a = a$.

A sesquilinear form is said to be \emph{skew-symmetric} if $B(\yy,\xx) = \wbar{B(\xx,\yy)}$. For example, the map
\[
B_\EE(\xx,\yy) := \sum_{i\in J} \wbar{x_i} y_i
\]
is a skew-symmetric sesquilinear form. Note that
\[
\EE(\xx) := B_\EE(\xx,\xx) = \|\xx\|^2.
\]

\section{The hyperboloid model}
\label{subsectionhyperboloidmodel}
Assume that $0\notin J$, and let
\[
\LL = \LL_\F^{J\cup\{0\}} = \HH_\F^{J\cup\{0\}} = \left\{\xx = (x_i)_{i\in J\cup\{0\}}\in\F^{J\cup\{0\}} \left\vert\sum_{i\in J\cup\{0\}} |x_i|^2 < \infty\right.\right\}.
\]
Consider the skew-symmetric sesquilinear form $B_\QQ:\LL\times\LL\to\F$ defined by
\[
B_\QQ(\xx,\yy) := -\wbar x_0 y_0 + \sum_{i\in J} \wbar x_i y_i
\]
and its associated quadratic form
\begin{equation}
\label{Qdef}
\QQ(\xx) := B_\QQ(\xx,\xx) = -|x_0|^2 + \sum_{i\in J} |x_i|^2.
\end{equation}
We observe that the form $\QQ$ is not positive definite, since $\QQ(\ee_0) = -1$.

\begin{remark}
If $\F = \R$, then the form $\QQ$ is called a \emph{Lorentzian} quadratic form, and the pair $(\LL,\QQ)$ is called a \emph{Minkowski space}.
\end{remark}

Let $\proj(\LL)$ denote the \emph{projectivization} of $\LL$, i.e. the quotient of $\LL\butnot\{\0\}$ under the equivalence relation $\xx\sim \xx a$ ($\xx\in\LL\butnot\{\0\}$, $a\in \F\butnot\{0\}$). Let
\[
\H = \H_\F^J := \{ [\xx]\in\proj(\LL_\F^{J\cup\{0\}}) : \QQ(\xx) < 0 \},
\]
and consider the map $\dist_\H: \H\times\H\to \Rplus$ defined by the equation
\begin{equation}
\label{distanceinL}
\cosh\dist_\H([\xx],[\yy]) = \frac{|B_\QQ(\xx,\yy)|}{\sqrt{|\QQ(\xx)|\cdot|\QQ(\yy)}|}, \;\;\; [\xx],[\yy]\in \H.
\end{equation}
Note that the map $\dist_\H$ is well-defined because the right hand side is invariant under multiplying $\xx$ and $\yy$ by scalars.

\begin{proposition}
\label{propositionROSSONCT}
The map $\dist_\H$ is a metric on $\H$ that is compatible with the natural topology, when viewed as a subspace of the quotient space $\proj(\LL)$. Moreover, for any two distinct points $[\xx],[\yy]\in\HH$ there exists a unique isometric embedding $\gamma:\R\to \H$ such that $\gamma(0) = [\xx]$ and $\gamma\circ\dist_\H([\xx],[\yy]) = [\yy]$.
\end{proposition}
\begin{remark}
The second sentence is the \emph{unique geodesic extension} property of $\H$. It holds more generally for Riemannian manifolds (cf. Remark \ref{remarkriemannian} below), but is an important distinguishing feature in the larger class of uniquely geodesic metric spaces.
\end{remark}
\begin{proof}[Proof of Proposition \ref{propositionROSSONCT}]
The key to the proof is the following lemma, which may also be deduced from the infinite-dimensional analogue of Sylvester's law of inertia \cite[Lemma 3]{Maddocks}.
\begin{lemma}
\label{lemmasylvester}
Fix $\zz\in\LL$ with $\QQ(\zz) < 0$, and let $\zz^\perp = \{\ww : B_\QQ(\zz,\ww) = 0\}$. Then $\QQ\given\zz^\perp$ is positive definite.
\end{lemma}
%\begin{remark}
%Lemma \ref{lemmasylvester} may be deduced from the infinite-dimensional analogue of Sylvester's law of inertia \cite[Lemma 3]{Maddocks}.
%\end{remark}
\begin{subproof}[Proof of Lemma \ref{lemmasylvester}]
By contradiction, suppose $\QQ(\yy) \leq 0$ for some $\yy\in\zz^\perp$. There exist $a,b\in\F$, not both zero, such that $y_0 a + z_0 b = 0$. But then
\[
0 < \QQ(\yy a + \zz b) = |a|^2\QQ(\yy) + |b|^2\QQ(\zz) \leq 0,
\]
which provides a contradiction.
\end{subproof}

Now fix $[\xx],[\yy],[\zz]\in\H$, and let $\xx,\yy,\zz\in\LL\butnot\{\0\}$ be representatives which satisfy
\[
B_\QQ(\xx,\zz) = B_\QQ(\yy,\zz) = \QQ(\zz) = -1.
\]
Then
\begin{align*}
\cosh\dist_\H([\xx],[\zz]) &= \frac{1}{\sqrt{1 - \QQ(\xx - \zz)}} &
\cosh\dist_\H([\yy],[\zz]) &= \frac{1}{\sqrt{1 - \QQ(\yy - \zz)}}\\
\sinh\dist_\H([\xx],[\zz]) &= \frac{\sqrt{\QQ(\xx - \zz)}}{\sqrt{1 - \QQ(\xx - \zz)}} &
\sinh\dist_\H([\yy],[\zz]) &= \frac{\sqrt{\QQ(\yy - \zz)}}{\sqrt{1 - \QQ(\yy - \zz)}}\cdot
\end{align*}
By the addition law for hyperbolic cosine we have
\[
\cosh(\dist_\H([\xx],[\zz]) + \dist_\H([\yy],[\zz])) = \frac{1 + \sqrt{\QQ(\xx - \zz)}\sqrt{\QQ(\yy - \zz)}}{\sqrt{1 - \QQ(\xx - \zz)}\sqrt{1 - \QQ(\yy - \zz)}} \cdot
\]
On the other hand, we have
\[
\cosh\dist_\H([\xx],[\yy]) = \frac{1}{\sqrt{1 - \QQ(\xx - \zz)}}\frac{1}{\sqrt{1 - \QQ(\yy - \zz)}}\left|-1 + B_\QQ(\xx - \zz,\yy - \zz)\right|.
\]
Since $\xx - \zz,\yy - \zz\in\zz^\perp$, the Cauchy--Schwartz inequality together with Lemma \ref{lemmasylvester} gives
\[
\left|-1 + B_\QQ(\xx - \zz,\yy - \zz)\right| \leq 1 + \sqrt{\QQ(\xx - \zz)}\sqrt{\QQ(\yy - \zz)},
\]
with equality if and only if $\xx - \zz$ and $\yy - \zz$ are proportional with a negative real constant of proportionality. This demonstrates the triangle inequality.

To show that $\dist_\H$ is compatible with the natural topology, it suffices to show that if $U$ is a neighborhood in the natural topology of a point $[\xx]\in\H$, then there exists $\epsilon > 0$ such that $B([\xx],\epsilon) \subset U$. Indeed, fix a representative $\xx\in [\xx]$; then there exists $\delta > 0$ such that $\|\yy - \xx\|\leq \delta$ implies $[\yy]\in U$. Now, given $[\yy]\in B([\xx],\epsilon)$, choose a representative $\yy\in [\yy]$ such that $\zz := \yy - \xx$ satisfies $B_\QQ(\xx,\zz) = 0$; this is possible since any representative $\yy\in[\yy]$ satisfies $B_\QQ(\xx,\yy)\neq 0$ by Lemma \ref{lemmasylvester}. Then
\[
\cosh\dist_\H([\xx],[\yy]) = \frac{|\QQ(\xx)|}{\sqrt{|\QQ(\xx)| \cdot |\QQ(\xx) + \QQ(\zz)|}} = \sqrt{\left|\frac{\QQ(\xx)}{\QQ(\xx) + \QQ(\zz)}\right|}\cdot
\]
So if $\dist_\H([\xx],[\yy]) \leq \epsilon$, then $\QQ(\zz)\leq \QQ(\xx)[1 - 1/\cosh(\epsilon)]$. By Lemma \ref{lemmasylvester}, there exists $C > 0$ such that $\|\zz\|^2 \leq C\QQ(\xx)[1 - 1/\cosh(\epsilon)]$. In particular, we may choose $\epsilon$ so that $C\QQ(\xx)[1 - 1/\cosh(\epsilon)] \leq \delta$, which completes the proof that $\dist_\H$ is compatible with the natural topology.

Now suppose that $\gamma:\R\to \H$ is an isometric embedding, and let $[\zz] = \gamma(0)$. Choose a representative $\zz\in\LL\butnot\{\0\}$ such that $\QQ(\zz) = -1$, and for each $t\in\R\butnot\{0\}$ choose a representative $\xx_t\in \LL\butnot\{\0\}$ such that $B_\QQ(\xx_t,\zz) = -1$. The preceding argument shows that for $t_1 < 0 < t_2$, $\xx_{t_1} - \zz$ and $\xx_{t_2} - \zz$ are proportional with a negative constant of proportionality. Together with \eqref{distanceinL}, this implies that
\begin{equation}
\label{geodesic}
\xx_t = \zz + \tanh(t)\ww
\end{equation}
for some $\ww\in \zz^\perp$ with $\QQ(\ww) = 1$. Conversely, direct calculation shows that the equation \eqref{geodesic} defines an isometric embedding $\gamma_{\zz,\ww}:\R\to\H$ via the formula $\gamma_{\zz,\ww}(t) = [\xx_t]$.
\end{proof}

\begin{definition}
\label{definitionROSSONCT}
%A \emph{rank one symmetric space of noncompact type (\ROSS)} is a pair $(\H_\F^J,\dist_\H)$, where $\F\in\{\R,\C,\Q\}$ and $J$ is a nonempty set such that $0\notin J$.
An \emph{algebraic hyperbolic space} is a pair $(\H_\F^J,\dist_\H)$, where $\F\in\{\R,\C,\Q\}$ and $J$ is a nonempty set such that $0\notin J$. 
\end{definition}


\begin{remark}
\label{remarkROSSONCTterminology}
In finite dimensions, the class of algebraic hyperbolic spaces is identical (modulo the Cayley hyperbolic plane, cf. Remark \ref{remarkH2O}) to the class of rank one symmetric spaces of noncompact type. This follows from the classification theorem for finite-dimensional symmetric spaces, see e.g. \cite[p.518]{Helgason}.\Footnote{In the notation of \cite{Helgason}, the spaces $\amsbb H_{\amsbb R}^p$, $\amsbb H_{\amsbb C}^p$, $\amsbb H_{\amsbb Q}^p$, and $\amsbb H_{\amsbb O}^2$ are written as $\SO(p,1)/\SO(p)$, $\SU(p,1)/\SU(p)$, $\Sp(p,1)/\Sp(p)$, and $(\mathfrak f_{4(-20)},\so(9))$, respectively.} It is not clear whether an analogous theorem holds in infinite dimensions (but see \cite{Duchesne3} for some results in this direction).
\end{remark}

%\begin{remark}
%\label{remarkCROSSONCT}
%Although, to agree with the literature, we have included the Cayley hyperbolic plane $\H_{\mathbb O}^2$ as an example of a \ROSS, we will not prove any theorems about this space. The reason for this is that an understanding of the geometry of the Cayley hyperbolic plane requires much more algebra than for the other \ROSSs, and we prefer not to spend time here on such a digression, especially since the example is finite-dimensional.
%
%We define a \emph{classical} \ROSS, or \emph{\CROSS}, to be a \ROSS\ which is not the Cayley hyperbolic plane.\Footnote{The distinction between \emph{classical} and \emph{exceptional} symmetric spaces comes from the corresponding distinction on the level of Lie groups. Roughly, exceptional Lie groups are those which do not belong to an infinite family of Lie groups.}
%\end{remark}

\begin{remark}
\label{remarkriemannian}
In finite dimensions, the metric $\dist_\H$ may be defined as the length metric associated to a certain Riemannian metric on $\H$; cf. \cite[\62.2]{Quint}. The same procedure works in infinite dimensions; cf. \cite{Lang_differential_geometry} for an exposition of the theory of infinite dimensional manifolds. Although a detailed account of the theory of infinite-dimensional Riemannian manifolds would be too much of a digression, let us make the following points:
\begin{itemize}
\item An infinite-dimensional analogue of the Hopf--Rinow theorem is false \cite{Atkin}, i.e. there exists an infinite-dimensional Riemannian manifold such that some two points on that manifold cannot be connected by a geodesic. However, if an infinite-dimensional Riemannian manifold $X$ is nonpositively curved, then any two points of $X$ can be connected by a unique geodesic as a result of the infinite-dimensional Cartan--Hadamard theorem \cite[IX, Theorem 3.8]{Lang_differential_geometry}; moreover, this geodesic is length-minimizing. In particular, if one takes a Riemannian manifolds approach to defining infinite-dimensional algebraic hyperbolic spaces, then the second assertion of Proposition \ref{propositionROSSONCT} follows from the Cartan--Hadamard theorem.
\item A bijection between two infinite-dimensional Riemannian manifolds is an isometry with respect to the length metric if and only if it is a diffeomorphism which induces an isometry on the Riemannian metric \cite[Theorem 7]{GJR}. This theorem is commonly known as the Myers--Steenrod theorem, as S. B. Myers and N. E. Steenrod proved its finite-dimensional version \cite{MyersSteenrod}. The difficult part of this theorem is proving that any bijection which is an isometry with respect to the length metric is differentiable. In the case of algebraic hyperbolic spaces, however, this follows directly from Theorem \ref{theoremisometries} below.
\end{itemize}
\end{remark}


%%%%%%%%%%%%%%%%%%%%%%%%%%%%%%%%%%
\bigskip
\section{Isometries of algebraic hyperbolic spaces}
\label{subsectionisometries}
%%%%%%%%%%%%%%%%%%%%%%%%%%%%%%%%%%

We define the \emph{group of isometries} of a metric space $(X,\dist)$ to be the group
\[
\Isom(X) := \{g:X\to X: \text{$g$ is a bijection and } \dist(g(x),g(y)) = \dist(x,y)\all x,y\in X\}.
\]
In this section we will compute the group of isometries of an arbitrary algebraic hyperbolic space. Fix $\F\in\{\R,\C,\Q\}$ and an index set $J$, and let $\HH = \HH_\F^J$, $\LL = \LL_\F^{J\cup\{0\}}$, and $\H = \H_\F^J$. We begin with the following observation:

\begin{observation}
Let $\O_\F(\LL;\QQ)$ denote the group of $\QQ$-preserving $\F$-linear automorphisms of $\LL$. Then for all $T\in\O_\F(\LL;\QQ)$, the map $[T]:\H\to\H$ defined by the equation
\begin{equation}
\label{[T]def}
[T]([\xx]) = [T\xx]
\end{equation}
is an isometry of $(\H,\dist_\H)$.
\end{observation}
\begin{proof}
The map $[T]$ is well-defined by the associativity property $T(\xx a) = (T\xx) a$. Since $T$ is $\QQ$-preserving and $\F$-linear, the polarization identity (the three versions cover the three cases when the base field $\F = \R$, $\C$, and $\Q$ respectively)
\[
B_\QQ(\xx,\yy) = \begin{cases}
\frac 14 [\QQ(\xx + \yy) - \QQ(\xx - \yy)] \\
\frac 14 [\QQ(\xx + \yy) - \QQ(\xx - \yy) - i \QQ(\xx + \yy i) + i\QQ(\xx - \yy i)] \\
\frac 14 \left[\QQ(\xx + \yy) - \QQ(\xx - \yy) + \sum\limits_{\ell = i,j,k}\big( -\ell \QQ(\xx + \yy \ell) + \ell\QQ(\xx - \yy \ell) \big)\right] 
\end{cases}
\] 
%\[
%B_\QQ(\xx,\yy) = \begin{cases}
%\frac 14 [\QQ(\xx + \yy) - \QQ(\xx - \yy)] & \F = \R\\
%\frac 14 [\QQ(\xx + \yy) - \QQ(\xx - \yy) - i \QQ(\xx + \yy i) + i\QQ(\xx - \yy i)] & \F = \C\\
%\frac 14 \left[\QQ(\xx + \yy) - \QQ(\xx - \yy) + \sum\limits_{\ell = i,j,k}\big( -\ell \QQ(\xx + \yy \ell) + \ell\QQ(\xx - \yy \ell) \big)\right] & \F = \Q
%\end{cases}
%\] 
guarantees that
\begin{equation}
\label{BQpreserving}
B_\QQ(T\xx,T\yy) = B_\QQ(\xx,\yy) \all \xx,\yy\in\HH.
\end{equation}
Comparing with \eqref{distanceinL} shows that $[T]$ is an isometry.
\end{proof}

The group $\O_\F(\LL;\QQ)$ is quite large. In addition to containing all maps of the form $T\oplus I$, where $T\in\O_\F(\HH;\EE)$ and $I:\F\to\F$ is the identity map, it also contains the so-called \emph{Lorentz boosts}
\begin{equation}
\label{lorentzboost}
T_{j,t}(\xx) = \left(\begin{cases}
x_i & i\neq 0,j\\
\cosh(t) x_j + \sinh(t) x_0 & i = j\\
\sinh(t) x_j + \cosh(t) x_0 & i = 0
\end{cases}\right)_{i\in J\cup\{0\}}, \;\; j\in J, t\in\R.
\end{equation}
We leave it as an exercise that $\O_\F(\HH;\EE)\oplus\{I\}$ and the Lorentz boosts in fact generate the group $\O_\F(\LL;\QQ)$.

\begin{observation}
\label{observationtransitivity}
The group
\[
\PO_\F(\LL;\QQ) = \{[T] : T\in\O_\F(\LL;\QQ)\} \leq \Isom(\H)
\]
acts transitively on $\H$.
\end{observation}
\begin{proof}
Let $\zero = [(1,\0)]$. The orbit of $\zero$ under $\PO_\F(\LL;\QQ)$ contains its image under the Lorentz boosts. Specifically, for every $t\in\R$ the orbit of $\zero$ contains the point $[(\cosh(t),\sinh(t),\0)]$. Applying maps of the form $[T\oplus I]$, $T\in \O_\F(\HH,\EE)$, shows that the orbit of $\zero$ is $\H$.
\end{proof}

We may ask the question of whether the group $\PO_\F(\LL;\QQ)$ is equal to $\Isom(\H)$ or is merely a subgroup. The answer turns out to depend on the division algebra $\F$:

\begin{theorem}
\label{theoremisometries}
If $\F\in\{\R,\Q\}$ then $\Isom(\H) = \PO_\F(\LL;\QQ)$. If $\F = \C$, then $\PO_\F(\LL;\QQ)$ is of index 2 in $\Isom(\H)$.
\end{theorem}
\begin{remark}
In finite dimensions, Theorem \ref{theoremisometries} is given as an exercise in \cite[Exercise II.10.21]{BridsonHaefliger}. Because of the importance of Theorem \ref{theoremisometries} to \thispaper, we provide a full proof.
\end{remark}
Before proving Theorem \ref{theoremisometries}, it will be convenient for us to introduce a group somewhat larger than $\O_\F(\LL;\QQ)$. Let $\Aut(\F)$ denote the group of automorphisms of $\F$ as an $\R$-algebra, i.e.
\[
\Aut(\F) = \left\{\sigma:\F\to\F \left|
\begin{split}
\text{$\sigma$ is an $\R$-linear bijection and}\\
\text{$\sigma(ab) = \sigma(a)\sigma(b)$ for all $a,b\in\F$}
\end{split}
\right.
\right\}.
\]
We will say that an $\R$-linear map $T:\LL\to\LL$ is \emph{$\F$-skew linear} if there exists $\sigma\in\Aut(\F)$ such that
\begin{equation}
\label{sigmaTdef}
T(\xx a) = T(\xx) \sigma(a) \text{ for all $\xx\in\HH$ and $a\in\F$}.
\end{equation}
The group of skew-linear bijections $T:\LL\to\LL$ which preserve $\QQ$ will be denoted $\O_\F^*(\LL;\QQ)$. For each $T$, the unique $\sigma\in\Aut(\F)$ satisfying \eqref{sigmaTdef} will be denoted $\sigma_T$. Note that the map $T\mapsto\sigma_T$ is a homomorphism.

\begin{warning*}
The associative law $(T\xx)a = T(\xx a)$ is \emph{not} valid for $T\in\O_\F^*(\LL;\QQ)$; rather, $T(\xx a) = (T\xx)\sigma_T(a)$ by \eqref{sigmaTdef}. Thus when discussing elements of $\O_\F^*(\LL;\QQ)$, we must be careful of parentheses.
\end{warning*}

\begin{example}
For each $\sigma\in\Aut(\F)$, the map
\[
\sigma^J(\xx) = \left(\sigma(x_i)\right)_{i\in J}
\]
is $\F$-skew-linear and $\QQ$-preserving, and $\sigma_{\sigma^J} = \sigma$.
\end{example}

\begin{observation}
\label{observationBQpreserving}
For $T\in\O_\F^*(\LL;\QQ)$,
\[
B_\QQ(T\xx,T\yy) = \sigma_T(B_\QQ(\xx,\yy)) \all \xx,\yy\in\LL.
\]
\end{observation}
\begin{proof}
By \eqref{BQpreserving}, the formula holds when $T\in\O_\F(\LL;\QQ)$, and direct calculation shows that it holds when $T = \sigma^J$ for some $\sigma\in\Aut(\F)$. Since $\O_\F^*(\LL;\QQ)$ is a semidirect product of the groups $\O_\F(\LL;\QQ)$ and $\{\sigma^J : \sigma\in\Aut(\F)\}$, this completes the proof.
\end{proof}

We observe that if $T\in\O_\F^*(\LL;\QQ)$, then \eqref{sigmaTdef} shows that $T$ preserves $\F$-lines, i.e. $T(\xx\F) = T(\xx)\F$ for all $\xx\in\LL\butnot\{\0\}$. Thus the equation \eqref{[T]def} defines a map $[T]:\H\to\H$, which is an isometry by Observation \ref{observationBQpreserving}. Thus if
\[
\PO_\F^*(\LL;\QQ) = \{[T] : T\in\O_\F^*(\LL;\QQ)\},
\]
then
\[
\PO_\F(\LL;\QQ) \leq \PO_\F^*(\LL;\QQ) \leq \Isom(\H).
\]
We are now ready to begin the
\begin{proof}[Proof of Theorem \ref{theoremisometries}]
The proof will consist of two parts. In the first, we show that $\PO_\F^*(\LL;\QQ) = \Isom(\H)$, and in the second we show that $\PO_\F(\LL;\QQ)$ is equal to $\PO_\F^*(\LL;\QQ)$ if $\F = \R,\Q$ and is of index $2$ in $\PO_\F^*(\LL;\QQ)$ if $\F = \C$.

Fix $g\in \Isom(\H)$; we claim that $g\in\PO_\F^*(\LL;\QQ)$. Let $\zz = (1,\0)$, and let $\zero = [\zz]$. By Observation \ref{observationtransitivity}, there exists $[T]\in \PO_\F(\LL;\QQ)$ such that $[T](\zero) = g(\zero)$. Thus, we may without loss of generality assume that $g(\zero) = \zero$.

We observe that $\zz^\perp = \HH$. Let $S(\HH)$ denote the unit sphere of $\HH$, i.e. $S(\HH) = \{\ww\in\HH : \QQ(\ww) = 1\}$. For each $\ww\in S(\HH)$, the embedding $\gamma_{\zz,\ww}:\R\to\H$ defined in the proof of Proposition \ref{propositionROSSONCT} is an isometry. By Proposition \ref{propositionROSSONCT}, its image under $g$ must also be an isometry. Specifically, there exists $f(\ww)\in S(\HH)$ such that
\begin{equation}
\label{tanht}
g([\zz + \tanh(t)\ww]) = [\zz + \tanh(t)f(\ww)] \all t\in\R.
\end{equation}
The fact that $g$ is a bijection implies that $f:S(\HH)\to S(\HH)$ is a bijection. Moreover, the fact that $g$ is an isometry means that for all $\ww_1,\ww_2\in S(\HH)$ and $t_1,t_2\in\R$, we have
\[
\dist([\zz + \tanh(t_1)\ww_1], [\zz + \tanh(t_2)\ww_2]) = \dist([\zz + \tanh(t_1)f(\ww_1)], [\zz + \tanh(t_2)f(\ww_2)]).
\]
Recalling that
\begin{align*}
\cosh\dist([\zz + \tanh(t_1)\ww_1] &, [\zz + \tanh(t_2)\ww_2])\\ &= \frac{|B_\QQ(\zz + \tanh(t_1)\ww_1, \zz + \tanh(t_2)\ww_2)|}{\sqrt{|\QQ(\zz + \tanh(t_1)\ww_1)|\cdot|\QQ(\zz + \tanh(t_2)\ww_2)|}}\\
&= \frac{|-1 + \tanh(t_1)\tanh(t_2)B_\QQ(\ww_1,\ww_2)|}{\sqrt{\left(1 - \tanh^2(t_1)\right)\left(1 - \tanh^2(t_2)\right)}},
\end{align*}
%\begin{align*}
%\cosh\dist([\zz + \tanh(t_1)\ww_1], [\zz + \tanh(t_2)\ww_2]) &= \frac{|B_\QQ(\zz + \tanh(t_1)\ww_1, \zz + \tanh(t_2)\ww_2)|}{\sqrt{|\QQ(\zz + \tanh(t_1)\ww_1)|\cdot|\QQ(\zz + \tanh(t_2)\ww_2)|}}\\
%&= \frac{|-1 + \tanh(t_1)\tanh(t_2)B_\QQ(\ww_1,\ww_2)|}{\sqrt{\left(1 - \tanh^2(t_1)\right)\left(1 - \tanh^2(t_2)\right)}},
%\end{align*}
we see that
\[
|-1 + \tanh(t_1)\tanh(t_2)B_\QQ(\ww_1,\ww_2)| = |-1 + \tanh(t_1)\tanh(t_2)B_\QQ(f(\ww_1),f(\ww_2))|.
\]
Write $\theta = \tanh(t_1)\tanh(t_2)$. Squaring both sides gives
%\begin{equation}
%\label{thetapolynomial}
%\theta^2 |B_\QQ(\ww_1,\ww_2)|^2 - 2\theta \Re[B_\QQ(\ww_1,\ww_2)] + 1 = \theta^2 |B_\QQ(f(\ww_1),f(\ww_2))|^2 - 2\theta \Re[B_\QQ(f(\ww_1),f(\ww_2))] + 1.
%\end{equation}
\begin{eqnarray}
\label{thetapolynomial}
\mathrm{LHS}^2 & = & \theta^2 |B_\QQ(\ww_1,\ww_2)|^2 - 2\theta \Re[B_\QQ(\ww_1,\ww_2)] + 1 \nonumber \\
|| &  & \\ 
\mathrm{RHS}^2 & = & \theta^2 |B_\QQ(f(\ww_1),f(\ww_2))|^2 - 2\theta \Re[B_\QQ(f(\ww_1),f(\ww_2))] + 1. \nonumber
\end{eqnarray}
%\begin{equation}
%\label{thetapolynomial}
%\theta^2 |B_\QQ(\ww_1,\ww_2)|^2 - 2\theta \Re[B_\QQ(\ww_1,\ww_2)] + 1 = \theta^2 |B_\QQ(f(\ww_1),f(\ww_2))|^2 - 2\theta \Re[B_\QQ(f(\ww_1),f(\ww_2))] + 1.
%\end{equation}
We observe that for $\ww_1,\ww_2\in S(\HH)$ fixed, \eqref{thetapolynomial} holds for all $-1 < \theta < 1$. In particular, taking the first and second derivatives and plugging in $\theta = 0$ gives
\begin{align} \label{represerved}
\Re[B_\QQ(\ww_1,\ww_2)] &= \Re[B_\QQ(f(\ww_1),f(\ww_2))]\\ \label{abspreserved}
|B_\QQ(\ww_1,\ww_2)| &= |B_\QQ(f(\ww_1),f(\ww_2))|.
\end{align}
Extend $f$ to a bijection $f:\HH\to\HH$ by letting $f(\0) = \0$ and $f(t\ww) = tf(\ww)$ for $t > 0$, $\ww\in S(\HH)$. We observe that \eqref{represerved} and \eqref{abspreserved} hold also for the extended version of $f$.

\begin{claim}
\label{claimflinear}
$f$ is $\R$-linear.
\end{claim}
\begin{subproof}
Fix $\ww_1,\ww_2\in\HH$ and $c_1,c_2\in\R$. By \eqref{represerved}, the maps
\[
\ww\mapsto \Re[B_\QQ(f(c_1\ww_1 + c_2\ww_2),f(\ww))] \text{ and } \ww\mapsto \Re[B_\QQ(c_1 f(\ww_1) + c_2 f(\ww_2),f(\ww))] 
\]
are identical. By the surjectivity of $f$ together with the Riesz representation theorem, this implies that $f(c_1\ww_1 + c_2\ww_2) = c_1 f(\ww_1) + c_2 f(\ww_2)$.
\end{subproof}
\begin{claim}
\label{claimFlines}
$f$ preserves $\F$-lines.
\end{claim}
\begin{subproof}
For each $\xx\in\HH\butnot\{\0\}$, the $\F$-line $\xx\F$ may be defined using the quantity $|B_\QQ|$ via the formula
\[
\xx\F = \left\{\yy\in\HH : \forall \ww\in \HH, \;\; |B_\QQ(\xx,\ww)| = 0 \;\;\Leftrightarrow\;\; |B_\QQ(\yy,\ww)| = 0\right\}.
\]
The claim therefore follows from \eqref{abspreserved}.
\end{subproof}
From Claim \ref{claimFlines}, we see that for all $\xx\in\HH\butnot\{\0\}$ and $a\in\F$, there exists $\sigma_\xx(a)\in\F$ such that
\[
f(\xx a) = f(\xx) \sigma_\xx(a).
\]
\begin{claim}
For $\xx,\yy\in\HH\butnot\{\0\}$,
\[
\sigma_\xx(a) = \sigma_\yy(a).
\]
\end{claim}
\begin{subproof}
By Claim \ref{claimflinear},
\[
[f(\xx) + f(\yy)] \sigma_{\xx + \yy}(a) = f(\xx a + \yy a) = f(\xx)\sigma_\xx(a) + f(\yy)\sigma_\yy(a).
\]
Rearranging, we see that
\[
f(\xx)[\sigma_{\xx + \yy}(a) - \sigma_\xx(a)] + f(\yy)[\sigma_{\xx + \yy}(a) - \sigma_\yy(a)] = 0.
\]
If $\xx$ and $\yy$ are linearly independent, then $\sigma_{\xx + \yy}(a) - \sigma_\xx(a) = 0$ and $\sigma_{\xx + \yy}(a) - \sigma_\yy(a) = 0$, so $\sigma_\xx(a) = \sigma_\yy(a)$. But the general case clearly follows from the linearly independent case.
\end{subproof}
For $a\in\F$, denote the common value of $\sigma_\xx(a)$ ($\xx\in\HH\butnot\{\0\}$) by $\sigma(a)$. Then
\begin{equation}
\label{fsigma}
f(\xx a) = f(\xx)\sigma(a) \all \xx\in\HH \all a\in\F.
\end{equation}
\begin{claim}
$\sigma\in\Aut(\F)$.
\end{claim}
\begin{subproof}
The $\R$-linearity of $\sigma$ follows from Claim \ref{claimflinear}, and the bijectivity of $\sigma$ follows from the bijectivity of $f$. Fix $\xx\in\HH\butnot\{\0\}$ arbitrary. For $a,b\in\F$,
\[
f(\xx)\sigma(ab) = f(\xx ab) = f(\xx a)\sigma(b) = f(\xx) \sigma(a)\sigma(b),
\]
which proves that $\sigma$ is a multiplicative homomorphism.
\end{subproof}
Thus $f\in\O_\F^*(\HH;\EE)$, and so $T = f\oplus I\in\O_\F^*(\LL;\QQ)$. But $[T] = g$ by \eqref{tanht}, so $g\in\PO_\F^*(\LL;\QQ)$. This completes the first part of the proof, namely that $\PO_\F^*(\LL;\QQ) = \Isom(\H)$.

To complete the proof, we need to show that $\PO_\F(\LL;\QQ)$ is equal to $\PO_\F^*(\LL;\QQ)$ if $\F = \R,\Q$ and is of index $2$ in $\PO_\F^*(\LL;\QQ)$ if $\F = \C$. If $\F = \R$, this is obvious. If $\F = \C$, it follows from the semidirect product structure $\O_\F^*(\LL;\QQ) = \O_\F(\LL;\QQ) \ltimes \{\sigma^J : \sigma\in\Aut(\F)\}$ together with the fact that $\Aut(\F) = \{I,z\mapsto \bar z\} \equiv\Z_2$.

If $\F = \Q$, then $\Aut(\F) = \{\Phi_a : a\in S(\Q)\}$, where $\Phi_a(b) = aba^{-1}$. Here $S(\F) = \{a\in\F : |a| = 1\}$. So $\O_\Q^*(\LL;\QQ) \neq \O_\Q(\LL;\QQ)$; nevertheless, we will show that $\PO_\Q^*(\LL;\QQ) = \PO_\Q(\LL;\QQ)$. Fix $[T]\in \PO_\Q^*(\LL;\QQ)$, and fix $a\in S(\Q)$ for which $\sigma_T = \Phi_a$. Consider the map
\begin{equation}
\label{Tadef}
T_a(\xx) = \xx a.
\end{equation}
We have $T_a\in\O_\Q^*(\LL;\QQ)$ and $\sigma_{T_a} = \Phi_a^{-1}$. Thus $\sigma_{T_a T} = \Phi_a \Phi_a^{-1} = I$, so $T_a T$ is $\F$-linear. But
\[
[T_a T] = [T],
\]
so $[T]\in\PO_\Q(\LL;\QQ)$. The completes the proof of Theorem \ref{theoremisometries}.
\end{proof}

\begin{remark}
Using algebraic language, the automorphisms $\Phi_a$ of $\Q$ are \emph{inner} automorphisms, while the automorphism $z\mapsto \bar z$ of $\C$ is an \emph{outer} automorphism. Although both inner and outer automorphisms contribute to the quotient $\O_\F^*(\LL;\QQ)/\O_\F(\LL;\QQ)$, only the outer automorphisms contribute to the quotient $\PO_\F^*(\LL;\QQ)/\PO_\F(\LL;\QQ)$. This explains why the index $\#(\PO_\F^*(\LL;\QQ)/\PO_\F(\LL;\QQ))$ is smaller when $\F = \Q$ than when $\F = \C$: although the group $\Aut(\Q)$ is much larger than $\Aut(\C)$, it consists entirely of inner automorphisms, while $\Aut(\C)$ has an outer automorphism.
\end{remark}

\begin{definition}
The \emph{bordification} of $\H$ is its closure relative to the topological space $\proj(\LL)$, i.e.
\[
\bord\H = \{[\xx] : \QQ(\xx) \leq 0\}.
\]
The \emph{boundary} of $\H$ is its topological boundary relative to $\proj(\LL)$, i.e.
\[
\del\H = \bord\H\butnot\H = \{[\xx] : \QQ(\xx) = 0\}.
\]
\end{definition}

The following is a corollary of Theorem \ref{theoremisometries}:

\begin{corollary}
\label{corollaryextension}
Every isometry of $\H$ extends uniquely to a homeomorphism of $\bord\H$.
\end{corollary}
\begin{proof}
If $T\in\O_\F^*(\LL;\QQ)$, then the formula \eqref{[T]def} defines a homeomorphism of $\bord\H$ which extends the action of $[T]$ on $\H$. The uniqueness is automatic.
\end{proof}
\begin{remark}
Corollary \ref{corollaryextension} can also be proven independently of Theorem \ref{theoremisometries} via the theory of hyperbolic metric spaces; cf. Lemma \ref{lemmaisometryextension} and Proposition \ref{propositionboundariesequivalent}.
\end{remark}

The following observation will be useful in the sequel:
\begin{observation}
\label{observationBQnonzero}
Fix $[\xx],[\yy]\in\bord\H$. Then
\[
B_\QQ(\xx,\yy) = 0 \;\;\Leftrightarrow\;\; [\xx] = [\yy] \in\del\H.
\]
\end{observation}
\begin{proof}
If either $[\xx]$ or $[\yy]$ is in $\H$, this follows from Lemma \ref{lemmasylvester}. Suppose that $[\xx],[\yy]\in\del\H$, and that $B_\QQ(\xx,\yy) = 0$. Then $\QQ$ is identically zero on $\xx\F + \yy\F$. Thus $(\xx\F + \yy\F)\cap\HH = \{\0\}$, and so $\xx\F + \yy\F$ is one-dimensional. This implies $[\xx] = [\yy]$.
\end{proof}

\section{Totally geodesic subsets of algebraic hyperbolic spaces}
\label{subsectiontotallygeodesic}

Given two pairs $(X,\bord X)$ and $(Y,\bord Y)$, where $X$ and $Y$ are metric spaces contained in the topological spaces $\bord X$ and $\bord Y$ (and dense in these spaces), an \emph{isomorphism} between $(X,\bord X)$ and $(Y,\bord Y)$ is a homeomorphism between $\bord X$ and $\bord Y$ which restricts to an isometry between $X$ and $Y$.

\begin{proposition}
\label{propositiontotallygeodesic}
Let $\K\leq\F$ be an $\R$-subalgebra, and let $V\leq\LL$ be a closed (right) $\K$-module such that
\begin{equation}
\label{totallygeodesic}
B_\QQ(\xx,\yy)\in\K \all \xx,\yy\in V.
\end{equation}
Then either $[V]\cap\H = \emptyset$ and $\#([V]\cap\bord\H)\leq 1$, or $([V]\cap\H,[V]\cap\bord\H)$ is isomorphic to an algebraic hyperbolic space together with its closure.
\end{proposition}
\begin{proof}~
\begin{itemize}
\item[Case 1:] $[V]\cap\H\neq\emptyset$. In this case, fix $[\zz]\in [V]\cap\H$, and let $\zz$ be a representative of $[\zz]$ with $\QQ(\zz) = -1$. By Lemma \ref{lemmasylvester}, $\QQ$ is positive-definite on $\zz^\perp$. We leave it as an exercise that the quadratic forms $\QQ\given\zz^\perp$ and $\EE\given\zz^\perp$ agree up to a bounded multiplicative error factor, which implies that $\zz^\perp$ is complete with respect to the norm $\sqrt\QQ$.

From \eqref{totallygeodesic}, we see that $(V\cap\zz^\perp,B_\QQ)$ is a $\K$-Hilbert space. By the usual Gram--Schmidt process, we may construct an orthonormal basis $(\ee_i)_{i\in J'}$ for $V\cap\zz^\perp$, thus proving that $V\cap\zz^\perp$ is isomorphic to $\HH_\K^{J'}$ for some set $J'$. Thus $V$ is isomorphic to $\LL_\K^{J'\cup\{0\}}$, and so $([V]\cap\H,[V]\cap\bord\H)$ is isomorphic to $(\H_\K^{J'},\bord\H_\K^{J'})$.

\item[Case 2:] $[V]\cap\H = \emptyset$. We need to show that $\#([V]\cap\bord\H)\leq 1$. By contradiction fix $[\xx],[\yy]\in[V]$ distinct, and let $\xx,\yy\in V$ be representatives. By Observation \ref{observationBQnonzero}, $B_\QQ(\xx,\yy)\neq 0$. On the other hand, $\QQ(\xx) = \QQ(\yy) = 0$ since $[\xx],[\yy]\in\del\H$. Thus $\QQ(\xx - \yy B(\xx,\yy)^{-1}) = -2 < 0$. On the other hand, $\xx - \yy B(\xx,\yy)^{-1}\in V$ by \eqref{totallygeodesic}. Thus $[\xx - \yy B(\xx,\yy)^{-1}]\in [V]\cap\H$, a contradiction.
\end{itemize}
\end{proof}

\begin{definition}
\label{definitiontotallygeodesic}
A \emph{totally geodesic subset} of an algebraic hyperbolic space $\H$ is a set of the form $[V]\cap\bord\H$, where $V$ is as in Proposition \ref{propositiontotallygeodesic}. A totally geodesic subset is \emph{nontrivial} if it contains an element of $\H$.
\end{definition}

\begin{remark}
As with Definition \ref{definitionROSSONCT}, the terminology ``totally geodesic'' is motivated here by the finite-dimensional situation, where totally geodesic subsets correspond precisely with the closures of those submanifolds which are totally geodesic in the sense of Riemannian geometry; see \cite[Proposition A.4 and A.7]{Quint}. However, note that we consider both the empty set and singletons in $\del\H$ to be totally geodesic.
\end{remark}

\begin{remark}
\label{remarkVrotation}
If $V\leq\LL$ is a closed $\K$-module satisfying \eqref{totallygeodesic}, then for each $a\in\F\butnot\{0\}$, $Va$ is a closed $a^{-1}\K a$-module satisfying \eqref{totallygeodesic} (with $\K = a^{-1}\K a$).
\end{remark}


\begin{lemma}
\label{lemmatotallygeodesic}
The intersection of any collection of totally geodesic sets is totally geodesic.
\end{lemma}
\begin{proof}
Suppose that $(S_\alpha)_{\alpha\in A}$ is a collection of totally geodesic sets, and suppose that $S = \bigcap_\alpha S_\alpha\neq\emptyset$. Fix $[\zz]\in S$, and let $\zz$ be a representative of $[\zz]$. Then for each $\alpha\in A$, there exist (cf. Remark \ref{remarkVrotation}) an $\R$-algebra $\K_\alpha$ and a closed $\K_\alpha$-subspace $V_\alpha \leq \LL$ satisfying \eqref{totallygeodesic} (with $\K = \K_\alpha$) such that $\zz\in V_\alpha$ and $S_\alpha = [V_\alpha]\cap\bord\H$. Let $\K = \bigcap_\alpha\K_\alpha$ and $V = \bigcap_\alpha V_\alpha$. Clearly, $V$ is a $\K$-module and satisfies \eqref{totallygeodesic}.

We have $[V]\cap\bord\H\subset S$. To complete the proof, we must show the converse direction. Fix $[\xx]\in S\butnot\{[\zz]\}$. By Observation \ref{observationBQnonzero}, there exists a representative $\xx$ of $[\xx]$ such that $B_\QQ(\zz,\xx) = 1$. Then for each $\alpha$, we may find $a_\alpha\in\F\butnot\{0\}$ such that $\xx a_\alpha\in V_\alpha$. We have
\[
a_\alpha = B_\QQ(\zz,\xx)a_\alpha = B_\QQ(\zz,\xx a_\alpha)\in\K_\alpha.
\]
Since $V_\alpha$ is a $\K_\alpha$-module, this implies $\xx\in V_\alpha$. Since $\alpha$ was arbitrary, $\xx\in V$, and so $[\xx]\in[V]\cap\bord\H$.
\end{proof}
\begin{remark}
Given $K\subset\bord\H$, Lemma \ref{lemmatotallygeodesic} implies that there exists a smallest totally geodesic set containing $K$. If we are only interested in the geometry of $K$, then by Proposition \ref{propositiontotallygeodesic} we can assume that this totally geodesic set is really our ambient space. In such a situation, we may without loss of generality suppose that there is no proper totally geodesic subset of $\bord\H$ which contains $K$. In this case we say that $K$ is \emph{irreducible}.
\end{remark}

\begin{warning*}
Although the intersection of any collection of totally geodesic sets is totally geodesic, it is not necessarily the case that the decreasing intersection of nontrivial totally geodesic sets is nontrivial; cf. Remark \ref{remarkparabolictorsion}.
\end{warning*}

%[Is Theorem \ref{theoremtotallygeodesic} being used? (Yes, but only the special case of Remark \ref{remarktotallygeodesic}.) Is it interesting by itself? If not, delete it but keep Lemma \ref{lemmaoperatornorm} (which you should move to another subsection).\internal]

The main reason that totally geodesic sets are relevant to our development is their relationship with the group of isometries. Specifically, we have the following:
\begin{theorem}
\label{theoremtotallygeodesic}
Let $(g_n)_1^\infty$ be a sequence in $\Isom(\H)$, and let
\begin{equation}
\label{fixedROSSONCT}
S = \left\{[\xx]\in\bord\H : g_n([\xx]) \tendsto n [\xx]\right\}.
\end{equation}
Then either $S\subset\del\H$ and $\#(S) = 2$, or $S$ is a totally geodesic set.
\end{theorem}
\begin{remark}
\label{remarktotallygeodesic}
An important example is the case where the sequence $(g_n)_1^\infty$ is constant, say $g_n = g$ for all $n$. Then $S$ is precisely the \emph{fixed point set} of $g$:
\[
S = \Fix(g) := \left\{[\xx]\in\bord\H : g([\xx]) = [\xx]\right\}.
\]
If $\H$ is finite-dimensional, then it is possible to reduce Theorem \ref{theoremtotallygeodesic} to this special case by a compactness argument.
\end{remark}
\begin{proof}[Proof of Theorem \ref{theoremtotallygeodesic}]
If $S = \emptyset$, then the statement is trivial. Suppose that $S\neq\emptyset$, and fix $[\zz]\in S$.

{\it Step 1: Choosing representatives $T_n$.} From the proof of Theorem \ref{theoremisometries}, we see that each $g_n$ may be written in the form $[T_n]$ for some $T_n\in\O_\F^*(\LL;\QQ)$. We have some freedom in choosing the representatives $T_n$; specifically, given $a_n\in S(\F)$ we may replace $T_n$ by $T_n T_{a_n}$, where $T_{a_n}$ is defined by \eqref{Tadef}.

Since $g_n([\zz])\to [\zz]$, there exist representatives $\zz_n$ of $g_n([\zz])$ such that $\zz_n\to \zz$. For each $n$, there is a unique representative $T_n$ of $g_n$ such that
\[
(T_n \zz) c_n = \zz_n \text{ for some $c_n\in\R\butnot\{0\}$.}
\]
Then
\[
(T_n \zz) c_n \to \zz.
\]
\begin{remark}
If $\F = \Q$, it may be necessary to choose $T_n\in\O_\F^*(\LL;\QQ)\butnot\O_\F(\LL;\QQ)$, despite the fact that each $g_n$ can be represented by an element of $\O_\F(\LL;\QQ)$.
\end{remark}

{\it Step 2: A totally geodesic set.}
Write $\sigma_n = \sigma_{T_n}$, and let
\begin{align*}
\K &= \{a\in \F : \sigma_n(a)\to a\}\\
V &= \left\{\xx\in\LL : T_n \xx \tendsto n \xx\right\}.
\end{align*}
Then $\K$ is an $\R$-subalgebra of $\F$, and $V$ is a $\K$-module. Given $\xx,\yy\in V$, by Observation \ref{observationBQpreserving} we have
\[
\sigma_n(B_\QQ(\xx,\yy)) = B_\QQ(T_n\xx,T_n\yy) \tendsto n B_\QQ(\xx,\yy),
\]
so $B(\xx,\yy)\in \K$. Thus $V$ satisfies \eqref{totallygeodesic}. If $V$ is closed, then the above observations show that $[V]\cap\bord\H$ is totally geodesic. However, this issue is a bit delicate:

\begin{claim}
If $\#([V]\cap\bord\H)\geq 2$, then $V$ is closed.
\end{claim}
\begin{subproof}
Suppose that $\#([V]\cap\bord\H)\geq 2$. The proof of Proposition \ref{propositiontotallygeodesic} shows that $[V]\cap\H\neq\emptyset$. Thus, there exists $\xx\in V$ for which $[\xx]\in\H$. In particular, $g_n([\xx])\to [\xx]$. Letting $\zero = [(1,\0)]$, we have
\[
\dist_\H(\zero,g_n(\zero)) \leq 2\dist_\H(\zero,[\xx]) + \dist_\H([\xx],g_n([\xx])) \tendsto n 2\dist_\H(\zero,[\xx]).
\]
In particular $\dist_\H(\zero,g_n(\zero))$ is bounded, say $\dist_\H(\zero,g_n(\zero)) \leq C$.

\begin{lemma}
\label{lemmaoperatornorm}
Fix $T\in\O_\F^*(\LL;\QQ)$, and let $\|T\|$ denote the operator norm of $T$. Then
\[
\|T\| = e^{\dist_\H(\zero,[T](\zero))}.
\]
\end{lemma}
\begin{subproof}
Write $T = T_{j,t}(A\oplus I)$, where $T_{j,t}$ is a Lorentz boost (cf. \eqref{lorentzboost}) and $A\in\O_\F^*(\HH;\EE)$. Then
\[
[T](\zero) = [T_{j,t}](\zero) = [(\cosh(t),\sinh(t),\0)].
\]
Here the second entry represents the $j$th coordinate. In particular,
\[
\cosh\dist_\H(\zero,[T](\zero)) = \frac{|B_\QQ((1,\0),(\cosh(t),\sinh(t),\0))|}{\sqrt{|\QQ(1,\0)|\cdot|\QQ(\cosh(t),\sinh(t),\0)|}} = \frac{\cosh(t)}{1} = \cosh(t).
\]
On the other hand,
\[
\|T\| = \|T_{j,t}\| = \left\|
\left[\begin{array}{ccc}
\cosh(t) & \sinh(t) & \\
\sinh(t) & \cosh(t) &\\
&& I
\end{array}\right] \right\| = e^t.
\]
This completes the proof.
%Since the matrix \eqref{coshsinh} is symmetric, it is diagonalizable by an orthogonal matrix, and so its operator norm is equal to the maximum of its eigenvalues, which is readily computed to be $e^t$.
\end{subproof}
\noindent Thus $\|T_n\| \leq e^C$ for all $n$, and so the sequence $(T_n)_1^\infty$ is equicontinuous. It follows that $V$ is closed.
\end{subproof}
\noindent Since $\#([V]\cap\bord\H)\leq 1$ implies that $[V]\cap\bord\H$ is totally geodesic, we conclude that $[V]\cap\bord\H$ is totally geodesic, regardless of whether or not $V$ is closed.

\begin{remark}
When $\#([V]\cap\bord\H)\leq 1$, there seems to be no reason to think that $V$ should be closed.
\end{remark}

{\it Step 3: Relating $S$ to $[V]\cap\bord\H$.}
The object of this step is to show that $S = [V]\cap\bord\H$ unless $S\subset\del\H$ and $\#(S)\leq 2$. For each $[\xx]\in S\butnot\{[\zz]\}$, let $\xx$ be a representative of $[\xx]$ such that $B_\QQ(\zz,\xx) = 1$; this is possible by Observation \ref{observationBQnonzero}. It is possible to choose a sequence of scalars $(a_n^{([\xx])})_{n = 1}^\infty$ in $\F\butnot\{\0\}$ such that $(T_n\xx) a_n^{([\xx])}\to \xx$. Let $a_n^{([\zz])} = c_n$. For $[\xx],[\yy]\in S$, we have
\begin{equation}
\label{anxany}
\begin{split}
\wbar{a_n^{([\xx])}}\sigma_{T_n}(B_\QQ(\xx,\yy)) a_n^{([\yy])}
&= \wbar{a_n^{([\xx])}}B_\QQ(T_n\xx,T_n\yy) a_n^{([\yy])} \hspace{.5 in} \text{(by Observation \ref{observationBQpreserving})}\\
&= B_\QQ((T_n\xx) a_n^{([\xx])},(T_n\yy) a_n^{([\yy])})\\
&\tendsto n B_\QQ(\xx,\yy).
\end{split}
\end{equation}
In particular,
\begin{equation}
\label{anxanyabs}
|a_n^{([\xx])}|\cdot|a_n^{([\yy])}| \tendsto n 1 \text{ whenever }B_\QQ(\xx,\yy)\neq 0.
\end{equation}
\begin{claim}
Unless $S\subset\del\H$ and $\#(S)\leq 2$, then for all $[\xx]\in S$ we have
\begin{equation}
\label{anx}
|a_n^{([\xx])}|\tendsto n 1.
\end{equation}
\end{claim}
\begin{subproof}
We first observe that it suffices to demonstrate \eqref{anx} for one value of $\xx$; if \eqref{anx} holds for $\xx$ and $[\yy]\neq[\xx]$, then $B_\QQ(\xx,\yy)\neq 0$ by Observation \ref{observationBQnonzero} and so \eqref{anxanyabs} implies $|a_n^{([\yy])}|\to 1$.

Now suppose that $S\nsubset\del\H$, and choose $[\xx]\in S\cap\H$. Then $B_\QQ(\xx,\xx)\neq 0$, and so \eqref{anxanyabs} implies \eqref{anx}.

Finally, suppose that $\#(S)\geq 3$, and choose $[\xx],[\yy],[\zz]\in S$ distinct. By \eqref{anxanyabs} together with Observation \ref{observationBQnonzero}, we have $|a_n^{([\xx])}|\cdot|a_n^{([\yy])}| \to 1$, $|a_n^{([\xx])}|\cdot|a_n^{([\zz])}| \to 1$, and $|a_n^{([\yy])}|\cdot|a_n^{([\zz])}| \to 1$. Multiplying the first two formulas and dividing by the third, we see that $|a_n^{([\xx])}|\to 1$.
\end{subproof}

For the remainder of the proof we assume that either $S\nsubset\del\H$ or $\#(S)\geq 3$.

Plugging $\zz = \xx$ into \eqref{anx}, we see that $c_n\to 1$. In particular, $[\zz]\in[V]\cap\bord\H$. Now fix $[\xx]\in S\butnot\{[\zz]\}$. Since $c_n\to 1$ and $B_\QQ(\zz,\xx) = 1$, \eqref{anxany} becomes
\[
a_n^{([\xx])} \to 1.
\]
Thus $\xx\in V$, and so $[\xx]\in [V]\cap\bord\H$.

%\ignore{
%
%
%
%For each $n$, let $T_n\in\O_\F^*(\LL;\QQ)$ \comdavid{define} be the unique representative of $g_n$ such that $T_n \xx c_n = \xx_n$ for some $c_n\in\F\butnot\{0\}$. Then $T_n \xx c_n \to \xx$.
%
%The set
%\[
%V = \left\{\yy\in\LL : T_n\yy \tendsto n \yy\right\}
%\]
%is an $\F$-linear subspace of $\LL$. Clearly, $[V]\cap \bord\H\subset S$. Conversely, fix $[\yy]\in S$, and let $\yy\in\LL\butnot\{\0\}$ be a representative of $[\yy]$. There exists a sequence of nonzero scalars $(a_n)_1^\infty$ such that
%\[
%T_n\yy a_n \tendsto n \yy.
%\]
%It follows that
%\begin{align*}
%B_\QQ(\xx,\yy) c_n a_n
%&= \wbar{c_n}B_\QQ(\xx,\yy) a_n \since{$c_n$ is real}\\
%&= B_\QQ(\xx c_n,\yy a_n)\\
%&= B_\QQ(T_n\xx c_n,T_n\yy a_n) \by{\eqref{BQpreserving}}\\
%&\tendsto n B_\QQ(\xx,\yy).
%\end{align*}
%In particular, if $[\xx]\neq[\yy]$, then $B_\QQ(\xx,\yy)\neq 0$ and so $c_n a_n\to 1$.
%
%\begin{itemize}
%\item[Case 1:] For some $[\xx]\in S$, $c_n\to 1$. In this case, $\xx\in V$, and for all $[\yy]\neq[\xx]$, the above argument shows that $\yy\in V$. Thus $S = [V]\cap\bord\H$.
%\item[Case 2:] For every $[\xx]\in S$, $c_n\not\to 1$. In this case, we must have $S\subset\del\H$, since $c_n = 1$ whenever $[\xx]\in\H$. By hypothesis, this implies $\#(S)\geq 3$. Choose $[\xx],[\yy],[\zz]\in S$ distinct, and let $\xx,\yy,\zz$ be representatives chosen as above. Then
%\begin{equation}
%\label{cnanbn}
%c_n a_n \to 1 \text{ and }c_n b_n \to 1,
%\end{equation}
%while
%\[
%\wbar{a_n} B_\QQ(\yy,\zz) b_n \to B_\QQ(\yy,\zz).
%\]
%(Note that $a_n$, $b_n$ may not be real.) Thus
%\[
%|a_n|\cdot |b_n| \to 1,
%\]
%and combining with \eqref{cnanbn} gives $c_n\to 1$. Thus we are reduced to the first case.
%\end{itemize}
%To complete the proof, we must show that the set $V$ has the desired properties (i)-(ii). 
%}% end ignore
\end{proof}




%%%%%%%%%%%%%%%%%%%%%%%%%%%%%%%%%%
\bigskip
\section{Other models of hyperbolic geometry}
\label{subsectionmodels}
%%%%%%%%%%%%%%%%%%%%%%%%%%%%%%%%%%

Fix $\F\in\{\R,\C,\Q\}$ and a set $J$, and let $\H = \H_\F^J$. The pair $(\H,\bord\H)$ is known as the \emph{hyperboloid model} of hyperbolic geometry (over the division algebra $\F$ and in dimension $\#(J)$). In this section we discuss two other important models of hyperbolic geometry. Note that the Poincar\'e ball model, which many of the figures of later chapters are drawn in, is not discussed here. References for this section include \cite{CFKP, Goldman}.

\subsection{The (Klein) ball model}
\label{subsubsectionballmodel}
Let
\[
\B = \B_\F^J = \{\xx\in \HH := \HH_\F^J : \|\xx\| < 1\},
\]
and let $\bord\B$ denote the closure of $\B$ relative to $\HH$.
\begin{observation}
\label{observationHequivB}
The map $e_{\B,\H}:\bord\B\to\bord\H$ defined by the equation
\[
e_{\B,\H}(\xx) = [(1,\xx)]
\]
is a homeomorphism, and $e_{\B,\H}(\B) = \H$. Thus if we let
\begin{equation}
\label{distB}
\cosh\dist_\B(\xx,\yy) = \cosh\dist_\H(e_{\B,\H}(\xx),e_{\B,\H}(\yy)) = \frac{|1 - B_\EE(\xx,\yy)|}{\sqrt{1 - \|\xx\|^2}\sqrt{1 - \|\yy\|^2}},
\end{equation}
then $e_{\B,\H}$ is an isomorphism between $(\B,\bord\B)$ and $(\H,\bord\H)$.
\end{observation}
The pair $(\B,\bord\B)$ is called the \emph{ball model} of hyperbolic geometry. It is often convenient for computations, especially those for which a single point plays an important role: by Observation \ref{observationtransitivity}, such a point can be moved to the origin $\0\in\B$ via an isomorphism of $(\B,\bord\B)$.

\begin{remark}
We should warn that the ball model $\B_\R^J$ of real hyperbolic geometry is not the same as the well-known Poincar\'e model, rather, it is the same as the Klein model.
\end{remark}

\begin{observation}
For all $T\in \O_\F^*(\HH;\EE)$, $T\given\B$ is an isometry which stabilizes $\0$.
\end{observation}

\begin{proposition}
\label{propositionIsomB}
In fact,
\[
\Stab(\Isom(\B);\0) = \{T\given\B : T\in\O_\F^*(\HH;\EE)\}.
\]
\end{proposition}
\begin{proof}
This is an immediate consequence of Theorem \ref{theoremisometries}.
\end{proof}


\subsection{The half-space model}
\label{subsubsectionE}

Now suppose $\F = \R$.\Footnote{The appropriate analogue of the half-space model when $\amsbb F\in\{\amsbb C,\amsbb Q\}$ is the \emph{paraboloid model}; see e.g. \cite[Chapter 4]{Goldman}.} Assume that $1\in J$, and let
\[
\E = \E^J = \left\{\left.\xx\in\HH := \HH_\F^J \right\vert x_1 > 0 \right\}.
\]
We will view $\E$ as resting inside the larger space
\[
\what\HH := \HH \cup\{\infty\}.
\]
The topology on $\what\HH$ is defined as follows: a subset $U \subset \what\HH$ is open if and only if
\[
\text{$U \cap \HH$ is open and ($\infty\in U$ \implies $\HH \setminus U$ is bounded)}.
\]
The boundary and closure of $\E$ will be subsets of $\what\HH$ according to the topology defined above, i.e.
\begin{align*}
\del\E &= \{\xx\in\HH : x_1 = 0\}\cup\{\infty\}\\
\bord\E &= \{\xx\in\HH : x_1 \geq 0\}\cup\{\infty\}.
\end{align*}
\begin{proposition}
\label{propositionHequivE}
The map $e_{\E,\H}:\bord\E\to\bord\H$ defined by the formula
\begin{equation}
\label{paraboloidembedding}
e_{\E,\H}(\xx) = \begin{cases}
\left[\left(\begin{cases}
2x_i & i\neq 0,1\\
1 + \|\xx\|^2 & i = 0\\
1 - \|\xx\|^2 & i = 1
\end{cases}\right)_{i\in J\cup\{0\}}\right] & \xx\neq\infty\\
[(1,-1,\0)] & \xx = \infty
\end{cases}
\end{equation}
is a homeomorphism, and $e_{\E,\H}(\E) = \H$. Thus if we let
\begin{equation}
\label{distE}
\cosh\dist_\E(\xx,\yy) = \cosh\dist_\H(e_{\E,\H}(\xx),e_{\E,\H}(\yy)) = 1 + \frac{\|\yy - \xx\|^2}{2x_1 y_1},
\end{equation}
then $e_{\E,\H}$ is an isomorphism between $(\E,\bord\E)$ and $(\H,\bord\H)$.
\end{proposition}
\begin{proof}
For $\xx\in\bord\E\butnot\{\infty\}$,
\[
\QQ(e_{\E,\H}(\xx)) = -(1 + \|\xx\|^2)^2 + (1 - \|\xx\|^2)^2 + \sum_{i\in J\butnot\{1\}}(2x_i)^2 = -4x_1^2.
\]
It follows that $e_{\E,\H}(\E) \subset\H$ and $e_{\E,\H}(\del\E)\subset\del\H$. Calculation verifies that the map
\begin{equation}
\label{iotainverse}
e_{\H,\E}([\xx]) = \begin{cases}\left(\begin{cases}
x_i/2 & i\neq 1\\
\sqrt{-\QQ(\xx)}/2 & i = 1
\end{cases}\right)_{i\in J} &\text{ if }x_0 + x_1 = 2 \\
\infty &\text{ if } \xx = (1,-1,\0)
\end{cases}
\end{equation}
is both a left and a right inverse of $e_{\E,\H}$. Notice that it is defined in a way such that for each $[\xx]\in\bord\H$, there is a unique representative $\xx$ of $[\xx]$ for which the formula \eqref{iotainverse} makes sense. We leave it to the reader to verify that $e_{\E,\H}$ and $e_{\H,\E}$ are both continuous, and that \eqref{distE} holds.
\end{proof}

The point $\infty\in\del\E$, corresponding to the point $[(1,-1,\0)]\in\del\H$, plays a special role in the half-space model. In fact, the half-space model can be thought of as an attempt to understand the geometry of hyperbolic space when a single point on the boundary is fixed. Consequently, we are less interested in the set of all isometries of $\E$ than simply the set of all isometries which fix $\infty$.

\begin{observation}[Poincar\'e extension]
\label{observationpoincareextension}
Let $\BB = \del\E\butnot\{\infty\} = \HH_\R^{J\butnot\{1\}}$, and let $g:\BB\to\BB$ be a \emph{similarity}, i.e. a map of the form
\[
g(\xx) = \lambda T\xx + \bb,
\]
where $\lambda > 0$, $T\in\O_\R(\BB;\EE)$, and $\bb\in\BB$. Then the map $\what g:\bord\E\to\bord\E$ defined by the formula
\begin{equation}
\label{poincareextension}
\what g(\xx) = \begin{cases}
(\lambda x_1,g(\pi(\xx))) & \xx\neq \infty \\
\infty & \xx = \infty
\end{cases}
\end{equation}
is an isomorphism of $(\E,\bord\E)$; in particular, $\what g\given\E$ is an isometry of $\E$. Here $\pi:\HH\to\BB$ is the natural projection.
\end{observation}
\begin{proof}
This is immediate from \eqref{distE}.
\end{proof}
The isometry $\what g$ defined by \eqref{poincareextension} is called the \emph{Poincar\'e extension} of $g$ to $\E$.

\begin{remark}
Intuitively we shall think of the number $x_1$ as representing the \emph{height} of a point $\xx\in\bord\E$. Then \eqref{poincareextension} says that if $g:\BB\to\BB$ is an isometry, then the Poincar\'e extension of $g$ is an isometry of $\E$ which preserves the heights of points.
\end{remark}

\begin{proposition}
\label{propositionIsomE}
For all $g\in\Isom(\E)$ such that $g(\infty) = \infty$, there exists a similarity $h:\BB\to\BB$ such that $g = \what h$.
\end{proposition}
\begin{proof}
By Theorem \ref{theoremisometries}, there exists $T\in\O(\LL;\QQ)$ such that $[T] = e_{\E,\H}\circ g\circ e_{\E,\H}^{-1}$. This gives an explicit formula for $g$, and one must check that if $[T]$ preserves $[(1,-1,\0)]$, then $g$ is a Poincar\'e extension.
%\[
%[T] = e_{\E,\H}\circ g\circ e_{\E,\H}^{-1}.
%\]
%Since $[T](e_{\E,\H}(\infty)) = e_{\E,\H}(\infty)$, we have $T(1,-1,\0) = \lambda^{-1}(1,-1,\0)$ for some $\lambda\in\R\butnot\{0\}$. It follows that for all $\xx\in\LL$,
%\[
%x_0 + x_1 = -B_\QQ(\xx,(1,-1,\0)) = -\BB_\QQ(T\xx,T(1,-1,\0)) = -\lambda^{-1}\BB_\QQ(T\xx,(1,-1,\0)) = \lambda^{-1}((T\xx)_0 + (T\xx)_1).
%\]
%Together with \eqref{iotainverse}, this shows that $T$ multiplies heights by $\lambda$. The conclusion now follows from \eqref{distE}.
\end{proof}

\subsection{Transitivity of the action of $\Isom(\H)$ on $\del\H$}
Using the ball and half-plane models of hyperbolic geometry, it becomes easy to prove the following assertion:

\begin{proposition}
\label{propositiontransitivity}
If $\F = \R$, the group $\Isom(\H)$ acts triply transitively on $\del\H$.
\end{proposition}
This complements the fact that $\Isom(\H)$ acts transitively on $\H$ (Observation \ref{observationtransitivity}).
\begin{proof}
By Observation \ref{observationHequivB} and Proposition \ref{propositionHequivE}, we may switch between models as convenient. It is clear that $\Isom(\B)$ acts transitively on $\del\B$, and that $\Stab(\Isom(\E);\infty)$ acts doubly transitively on $\del\E\butnot\{\infty\}$. Therefore given any triple $(\xi_1,\xi_2,\xi_3)$, we may conjugate to $\B$, conjugate $\xi_1$ to a standard point, conjugate to $\E$ while conjugating $\xi_1$ to $\infty$, and then conjugate $\xi_2,\xi_3$ to standard points.
\end{proof}


We end this chapter with a convention:

\begin{convention}
\label{conventionH}
When $\alpha$ is a cardinal number, $\H_\F^\alpha$ will denote $\H_\F^J$ for any set $J$ of cardinality $\alpha$, but particularly $J = \{1,\ldots,n\}$ if $\alpha = n\in\Namer$ and $J = \Namer$ if $\alpha = \#(\Namer)$. Moreover, $\H_\F^\infty$ will always be used to denote $\H_\F^{\#(\Namer)} = \H_\F^{\Namer}$, the unique (up to isomorphism) infinite-dimensional separable algebraic hyperbolic space defined over $\F$. Finally, real hyperbolic spaces will be denoted without using $\R$ as a subscript, e.g. $\H^\infty = \H_\R^\infty$, $\B^J = \B_\R^J$, $\HH^\alpha = \HH_\R^\alpha$.
\end{convention}



%\ignore{
%There are several models of hyperbolic geometry, which are equivalent as (infinite-dimensional) Riemannian manifolds but which reflect different aspects of hyperbolic geometry. The models we will be interested in are the ball model ($\B$), the half-space model ($\H$), the hyperboloid model ($\L$ for Lorentzian), and the Klein model ($\K$). When we do not wish to specify a model we will write $\H$ (for ``hyperbolic'').
%
%\subsubsection{The ball model, the half-space model, and the Klein model}
%The ball model is the set
%\[
%\B := \ \{ \xx\in\HH : \|\xx\| < 1 \} 
%\]
%together with the Riemannian metric
%\[
%\lb \uu,\vv \rb_{\xx,\B} \ := \ \frac{4 \lb \uu,\vv \rb}{(1 - \|\xx\|^2)^2} ~\cdot
%\]
%The half-space model is the set
%\[
%\H \ := \ \{ \xx\in\HH : x_1 > 0 \}
%\]
%together with the Riemannian metric
%\[
%\lb \uu,\vv \rb_{\xx,\H} \ := \ \frac{\lb \uu,\vv \rb}{x_1^2} ~\cdot
%\]
%The Klein model is the set $\K := \B$ together with the Riemannian metric
%\[
%\lb \uu,\vv \rb_{\xx,\K} \ := \ \frac{\lb \uu,\vv \rb}{1 - \|\xx\|^2} - \frac12\cdot\frac{\sum_{i,j = 1}^\infty x_i x_j [u_i v_j + u_j v_i]}{(1 - \|\xx\|^2)^2}
%\]
%\begin{remark}
%In most references, the $(d + 1)$-dimensional half-space model is defined to be the set $\{\xx\in\R^{d + 1} :x_{d + 1} > 0\}$. When $d = \infty$, this does not make sense since there is no $\infty$th coordinate. Thus we have decided to use the first coordinate instead.
%\end{remark}
%
%Note that the topological boundaries of $\H$ and $\B$ are also Hilbert manifolds (although it requires slightly more work to come up with coordinate charts):
%\begin{align*}
%\del\B = \del \K &= \{ \xx\in\HH : \|\xx\| = 1 \} , \\
%\del\H &= \{ \xx\in\HH : x_1 = 0 \} \cup \{\infty\}.
%\end{align*}
%Note that we have taken the boundary of $\H$ with respect to the Hilbert manifold $\what\HH$ defined in Example \ref{examplewhatH}. Henceforth we shall always respect this convention.
%
%We note that the closures $\bord\B$, $\bord\H$, and $\bord\K$ are not Hilbert manifolds, but they are Hilbert manifolds with boundary (see \cite[\6 II.4]{Lang_differential_geometry}). For the purposes of \thispaper, we will be content with considering them as topological subspaces $\what\HH$.
%
%Finally, let us say a word about the geometric significance of $\B$, $\H$, and $\K$. The ball model is best if you want to figure out what the world looks like if you are ``at a point in $X$''; the half-space model is best if you are ``at a point on the boundary $\del X$''. The Klein model is the best for arguments which involve convexity, since geodesic segments are just (Euclidean) straight lines (see Subsection \ref{subsectiongeodesics} below).
%
%\subsubsection{The hyperboloid model}
%Now we consider a more complicated model of hyperbolic geometry, the hyperboloid model. The main advantage of this model is that isometries of hyperbolic space extend to \emph{linear} maps of the ambient space (Theorem \ref{theoremlinearextension} below).
%
%\subsubsection{Equivalence of models}
%We end this subsection with the following theorem which demonstrates that our four different models of hyperbolic geometry are all isomorphic as Riemannian manifolds.
%\begin{theorem}
%\label{equivalenceofmodels}
%Let
%\begin{align} \label{eBH}
%e_{\B,\H}(\xx) &= -\ee_1 + 2\frac{\xx + \ee_1 \;}{\| \xx + \ee_1 \|^2} \\ \label{eLB}
%e_{\L,\B}(\xx) &= \frac{(x_1,x_2,\ldots)}{1 + x_0} \\ \label{eLK}
%e_{\L,\K}(\xx) &= \frac{(x_1,x_2,\ldots)}{x_0}.
%\end{align}
%Then $e_{\B,\H}:\B\to\H$, $e_{\L,\B}:\L\to\B$, and $e_{\L,\K}:\L\to\K$ are isomorphisms of Riemannian manifolds as defined in Definition \ref{definitionRiemannianisomorphism}. In particular, by Observation \ref{observationRiemannianisomorphism} they are also isometries. Furthermore, each map $e_{X,Y}$ extends uniquely to a homeomorphism between $\bord X$ and $\bord Y$.
%\end{theorem}
%The theorem can be proven by direct calculation; cf. \cite{CFKP}.
%
%\begin{remark}
%By composition and inversion we may define $e_{X,Y}$ for any distinct $X,Y\in\{\B,\H,\K,\L\}$; for example, $e_{\K,\H} = e_{\B,\H}\circ e_{\L,\B}\circ e_{\L,\K}^{-1}$. Such a map would also satisfy the conclusion of Theorem \ref{equivalenceofmodels}.
%\end{remark}
%
%\begin{remark}[Formula for the continuous extension to the boundary]
%The same formula \eqref{eBH} may be used to extend $e_{\B,\H}$ to the closure $\bord\B$. On the other hand, the formula \eqref{eLK} defines a map $e_{\L,\K}:\LL\to \what\H$ which is homogeneous, meaning that $e_{\L,\K}(t\xx) = e_{\L,\K}(\xx)$ for all $t\neq 0$ and $\xx\in\LL$. Projectivizing gives a map $e_{\L,\K}:\proj(\LL)\to\what\HH$, and restricting to $\del\L$ gives a map $e_{\L,\K}:\del\L\to\del\K$. It can be verified that combining this map with the map $e_{\L,\K}:\L\to\K$ gives a continuous map $e_{\L,\K}:\bord\L\to\bord\K$. More surprisingly, combining the map $e_{\L,\K}:\del\L\to\del\K$ with the map $e_{\L,\B}:\L\to\B$ gives a continuous map $e_{\L,\B}:\bord\L\to\bord\B$.
%%Since the formula \eqref{eLK} is homogeneous, it is well-defined on $\proj(\LL)$ and therefore accurately describes the continuous extension of $e_{\L,\K}$ to $\cl{\L}$. On the other hand, \eqref{eLB} is not homogeneous. In fact, the continuous extensions of $e_{\L,\B}$ and $e_{\L,\K}$ agree on $\del\L$, so \eqref{eLK} can be used to describe the action of $e_{\L,\B}$ on $\del\L$.
%\end{remark}
%}% end ignore

%\ignore{
%%%%%%%%%%%%%%%%%%%%%%%%%%%%%%%%%%%
%\bigskip
%\subsection{Geodesics and distance in hyperbolic space}
%\label{subsectiongeodesics}
%%%%%%%%%%%%%%%%%%%%%%%%%%%%%%%%%%%
%
%
%
%\begin{proposition}[Hyperbolic law of cosines] \label{lawofcosines}
%Let $x,y,z\in\H$, and let
%\begin{align*}
%A &= \dist(y,z)\\
%B &= \dist(x,z)\\
%C &= \dist(x,y)\\
%\gamma &= \ang_z(x,y).
%\end{align*}
%Then
%\[
%\cosh(C) = \cosh(A)\cosh(B) - \sinh(A)\sinh(B)\cos(\gamma).
%\]
%\end{proposition}
%[Proof?\internal]
%
%}% end ignore



%\ignore{
%%%%%%%%%%%%%%%%%%%%%%%%%%%%%%%%%%%
%\bigskip
%\subsection{Conformal maps and Liouville's Theorem}
%\label{subsectionliouville}
%%%%%%%%%%%%%%%%%%%%%%%%%%%%%%%%%%%
%
%\selfnote{This section should be deleted, but maybe some parts should be kept.}
%
%
%To define conformal maps we first need the notion of a similarity.
%\begin{definition}
%Let $T:\HH\to\HH$ be a bounded linear operator. $T$ is a \emph{similarity} if it can be written as the product of a positive real number and a linear isometry. The positive number is called the \emph{scaling constant} of $T$, and will be denoted $\|T\|$.
%
%An affine map $A:\HH\to\HH$ is a \emph{similarity} if its linear part $\xx\mapsto A(\xx) - A(\0)$ is a similarity.
%
%The group of similarities of $\HH$ will be denoted $\Sim(\HH)$.
%\end{definition}
%
%\begin{definition}
%Let $X$ and $Y$ be infinite-dimensional Riemannian manifolds, and let $f:X \to Y$ be a diffeomorphism. We say that $f$ is \emph{conformal} if for each $x\in X$, $f'(x)$ is a similarity.
%\end{definition}
%
%As in the finite-dimensional case, our first example of a conformal map is the inversion with respect to a sphere.
%
%\begin{definition}
%Fix $\pp\in\mathcal{H}$ and $\alpha > 0$, and let $S(\pp,\alpha)$ denote the sphere around $\pp$ of radius $\alpha$. The \emph{inversion with respect to the sphere $S(\pp,\alpha)$} is the map
%\[
%\inv_{\pp,\alpha} : \xx \mapsto \alpha^2 \frac{\xx - \pp \;}{\|\xx - \pp\|^2} + \pp .
%\]
%We make the conventions that $\inv_\pp = \inv_{\pp,1}$ and $\inv = \inv_\0$.
%\end{definition}
%\begin{observation}
%\[
%\Fix(\inv_{\pp,\alpha}) = S(\pp,\alpha),
%\]
%which justifies the name ``inversion with respect to $S(\pp,\alpha)$''.
%\end{observation}
%Note that as defined $\inv_{\pp,\alpha}: \mathcal{H}\setminus \{\pp\} \to \mathcal{H}\setminus \{\pp\}$; however, we are more interested in its unique continuous extension to $\what\HH$ which is defined by sending $\pp \mapsto \infty$ and $\infty \mapsto \pp$.
%
%It may be verified that $\inv_{\pp,\alpha}: \mathcal{H}\setminus \{\pp\} \to \mathcal{H}\setminus \{\pp\}$ is a conformal map and that the following formulae for the scaling constant and ``direction'' of $\inv_{\pp,\alpha}'(z)$ analogous to the finite-dimensional case are still valid:
%\[
%\| \inv_{\pp,\alpha}'(\xx) \| = \frac{\alpha^2}{\| \xx - \pp \|^2} ~ \text{and}
%\]
%\[
%\frac{\inv_{\pp,\alpha}'(\xx)}{\| \inv_{\pp,\alpha}'(\xx) \|} = \id - 2P_{\xx - \pp} ~,
%\]
%where, $P_{\xx - \pp}:\HH \to \HH$ is orthogonal projection onto the linear subspace spanned by $\xx - \pp$.
%
%%%%%%%%%%%%%%%%%%%%%%%%%%%%%%%%%%%
%\subsubsection{Liouville's Theorem}
%%%%%%%%%%%%%%%%%%%%%%%%%%%%%%%%%%%
%
%The following theorem generalizes the classical result known as Liouville's theorem, which tells us that for $d\geq 3$, any conformal diffeomorphism betwen two subsets of $\R^d$ is the restriction of a M\"obius transformation.
%
%%It was mentioned in \cite{MSU} that the finite-dimensional proof
%%could be extended to the infinite-dimensional Hilbert space
%%setting by following the proof in Section 5.2 of \cite{B}. It
%%turns out that the proof there is somewhat incomplete and we shall
%%give an exposition of the proof in what follows. We note our proof
%%will essentially follow \cite{B} which claims to follow
%%Nevanlinna's original argument in \cite{Nevanlinna}. This will be done elsewhere.
%
%\begin{theorem}[Liouville's theorem in Hilbert space]\label{Liouville}
%Let $U,V \subset \mathcal{H}$ be nonempty open connected sets and let $\phi : U \to V$ be a conformal diffeomorphism. Then one of the following cases holds: 
%\begin{itemize}
%\item[(Nonlinear)] $\phi$ is the composition of an inversion and an affine similarity, or
%%There exist $\lambda >0$, $\pp,\qq\in\HH$, and $T\in\O(\HH)$ such that
%%\[
%%\phi(\xx) = \lambda T[\inv_\pp(\xx)]+\qq
%%\]
%\item[(Linear)] $\phi$ is an affine similarity.
%%There exist $\lambda >0$, $\pp\in\HH$, and $T\in\O(\HH)$ such that
%%\[
%%\phi(\xx) = \lambda T[\xx]+\pp.
%%\]
%\end{itemize}
%\end{theorem}
%Note that in either case the map $\phi$ extends uniquely to a conformal map $\what\phi : \what\HH \to \what\HH$. As in the finite-dimensional case, the map $\what\phi$ is called a \emph{M\"obius transformation} and we denote the class of such maps by $\Mob(\what\HH)$. To avoid clutter we avoid decorating the extension whenever it is clear.
%
%\begin{remark}
%Theorem \ref{Liouville} follows from the observation (see \cite{Huff} or \cite{Zorich}) that R.~Nevanlinna's proof of the finite-dimensional Liouville's theorem \cite{Nevanlinna} extends perfectly well to infinite dimensions.\Footnote{We would like to note, however, that there is a small omission towards the end of Nevanlinna's proof, at the place where it breaks into a few cases. There is a relatively easy argument to fill in the gap; however, the change is significant because the new argument is the only place in which the positive definiteness of the inner product is used. In fact, if we consider transformations which are conformal with respect to an indefinite quadratic form, then we end up needing a third case in Theorem \ref{Liouville}: ``double inversion''. Such considerations will be discussed in detail in \cite{DS_Lorentzian_Kleinian}.}
%\end{remark}
%
%%
%%there either there 
%%exists a unique quadruple $(\lambda,x,y,M)$ with $\lambda>0$, $x,y\in\HH$ and $M$ a 
%%bijective linear isometry such that
%%\[
%%\phi(\xx) = \lambda M(\inv_\pp(z))+y ,
%%\]
%%or $\phi$ is of the form $\phi(z)=\lambda M(z) + y$. As in the finite-dimensional case,
%%the map $\phi$ is called a M\"obius transformation and we denote the class of such maps by $\Mob(\what\HH)$.
%
%
%
%\begin{remark}
%It follows from the formula for the derivative of an inversion given above that for
%any given conformal map $\phi$ of the form $\phi(\xx) = \lambda T(\inv_\pp(\xx))+\qq$ we have
%\[
%\| \phi'(\xx) \| = \frac{\lambda}{ \|\xx - \pp\|^2}
%\]
%for every $\xx\in\mathcal{H}$. Note that if $\phi$ is of the 
%form $\phi(\xx)=\lambda T(\xx) + \pp$ we have $\| \phi'(\xx)\| = \lambda$.
%\end{remark}
%
%\begin{observation}
%\label{observationLiouvillecases}
%The nonlinear case corresponds to the case $g(\infty)\neq \infty$ and the linear case corresponds to the case $g(\infty) = \infty$.
%\end{observation}
%
%}% end ignore

%\ignore{
%
%%%%%%%%%%%%%%%%%%%%%%%%%%%%%%%%%%%
%\bigskip
%\subsection{The equivalence of models}
%%%%%%%%%%%%%%%%%%%%%%%%%%%%%%%%%%%
%
%Consider $\B,\H \subset \what\HH$. Recall that the topology on $\HH$ was defined as follows: a subset $U \subset \what\HH$ is open if and only if $U \cap \HH$ is open and $\HH \setminus U$ is bounded if $\infty\in U$. In what follows, $\del \B$, $\bord\B$, $\del \H$ and $\bord\H$ refer to the boundary and closure with respect to the topology on $\what\HH$. We then have the following descriptions:
%
%\begin{remark}
%We mention here the beautiful exposition of models and their equivalence in Cannon et al \cite{cannonfloyd} so that the interested reader may extend their computations to the Hilbertian setting. 
%For example, just as in the finite-dimensional situation, the conformal map $\inv_{-\ee_1,\sqrt2} : \what\HH \to \what\HH$ is a homeomorphism between $\bord\H$ and $\bord\B$ and an isometry between $\H$ and $\B$ with their respective metrics.
%For the purposes of this work, we do not talk about the boundary in the Lorentzian setting.
%\end{remark}
%
%To give some flavor for these computations let us make the following observations. Note that $\xx,\uu,\ee_1$ will refer to vectors in Hilbert space; though when we write $(x)_1$ we refer to the first coordinate of the vector $\xx$, e.g. $(\ee_2)_2 = 1$. 
%
%\begin{observation}
%The map $e_{\L,\B} : \L \to \B$ defined by
%\[
%e_{\L,\B}(\xx) = \frac{(x_1,x_2,\ldots)}{1 + x_0} 
%\]
%is an isometry between $\L$ and $\B$.
%\end{observation}
%
%%\begin{proof}
%%Note that the derivative of $e_{\L,\B} : \L \to \B$ is given by
%%\[
%%e_{\L,\B}'(a,x)[(b,u)] = \frac{u}{1+a} - \frac{bx}{(1+a)^2} \cdot
%%\]
%%Now we compute
%%\[
%%\aligned
%%\| e_{\L,\B}'(a,x)[(b,u)] \|^2_{e_{\L,\B}(a,x), \B} &= \Bigl\| \frac{u}{1+a} - \frac{bx}{(1+a)^2} \Bigr\|^2_{\frac{x}{1+a},\B} \\
%%&= 4 \frac{\Bigl\| \frac{u}{1+a} - \frac{bx}{(1+a)^2} \Bigr\|^2 }{\Bigl[ 1 - \left\| \frac{x}{(1+a)^2} \right\|^2 \Bigr]^2} \\
%%&= \frac{\Bigl[ \frac{\| u \|^2}{(1+a)^2} -\frac{2ab^2}{(1+a)^3} + \frac{b^2(b^2 -1)}{(1+a)^4} \Bigr]}{\Bigl[ 1 - \frac{b^2 +1}{(1+a)^2} \Bigr]^2}\\
%%&= 4 \frac{\Bigl[ \frac{\| u \|^2}{(1+a)^2} + \frac{b^2(-2a + a -1)}{(1+a)^3} \Bigr]}{\Bigl[ 1 - \frac{a-1}{a+1} \Bigr]^2} \\{b^2(1+a)
%%&= 4 \frac{\Bigl[ \frac{\| u \|^2}{(1+a)^2} - \frac{(1+a)^3} \Bigr]}{\Bigl[ \frac{2}{1+a} \Bigr]^2} \\
%%&= \|u\|^2 - b^2 \\
%%&= \| (b,u) \|^2_{(a,x),\L} \ \cdot
%%\endaligned
%%\]
%%
%%\end{proof}
%
%\begin{observation}
%The map $e_{\B,\H} : \B \to \H$ defined by
%\[
%e_{\B,\H}(\xx) = -\ee_1 + 2\frac{\xx + \ee_1 \;}{\| \xx + \ee_1 \|^2} 
%\]
%is an isometry between $\B$ and $\H$.
%\end{observation}
%
%%\begin{proof}
%%Note that the derivative of $e_{\B,\H} : \B \to \H$ is given by
%%\[
%%e_{\B,\H}'(x)[\uu] = 2\frac{u}{\| x + e_1 \|^2} - 4 \frac{\lb x + e_1, u \rb}{\| x + e_1 \|^4}(x + e_1) \cdot
%%\]
%%As before we now compute
%%\[
%%\aligned
%%\| e_{\B,\H}'(x)[\uu] \|^2_{e_{\B,\H}(x),\H} &= \frac{\Bigl\| 2\frac{u}{\| x + e_1 \|^2} - 4 \frac{\lb x + e_1, u \rb (x + e_1)}{\| x + e_1 \|^4} \Bigr\|^2}{\Bigl(-1 + 2 \frac{x_1 + 1}{\| x + e_1 \|^2} \Bigr)^2}\\
%%&= \frac{\Bigl\| 2u - 4 \frac{\lb x + e_1, u \rb (x + e_1)}{\| x + e_1 \|^2} \Bigr\|^2}{\(2(x_1 + 1) - \| x + e_1 \|^2\)^2}\\
%%&= \frac{4 \|u\|^2 - \frac{16 \lb x + e_1, u \rb^2}{\| x + e_1 \|^2} + \frac{16 \lb x + e_1, u \rb^2}{\| x + e_1 \|^2}}{\( 2x_1 + 2 - [\|x\|^2 + 2x_1 +1] \)^2} \\
%%&= \frac{4 \| u \|^2}{\( 1 - \| x \|^2 \)^2} \\
%%&= \| u \|^2_{x,\B}
%%\endaligned
%%\]
%%
%%\end{proof}
%
%}% end ignore


%\ignore{
%\begin{theorem}
%\label{theoremMobiusisometry}
%Let $X$ be either $\B$, $\H$, $\K$, or $\L$, and let $g\in\Isom(X)$. Then $g$ admits a unique continuous extension to $\bord X$. Moreover,
%\begin{itemize}
%\item[(i)] If $X = \B$ or $X = \H$, then $g$ admits a unique conformal extension to $\what\HH$. Conversely, any M\"obius transformation of $\what\HH$ preserving $X$ acts isometrically on $X$.
%\item[(ii)] If $X = \L$, then $g$ admits a unique linear extension to $\LL$. This extension preserves the quadratic form \text{\eqref{Qdef}}.
%\end{itemize}
%\end{theorem}
%In each case the extensions of $g$ will also be denoted by $g$.
%\begin{proof}~
%\begin{itemize}
%\item[(i)] Since the metrics $\lb \cdot,\cdot \rb_{\cdot,\B}$ and $\lb \cdot,\cdot \rb_{\cdot,\H}$ are both conformal to the Euclidean metric, it follows from Proposition \ref{propositionRiemannianisomorphism} that any $g\in\Isom(X)$ is a conformal isomorphism of $X$ viewed as a subset of $\HH$. Then by Liouville's theorem $g$ has a unique conformal extension to $\what\HH$.
%
%The ``conversely'' part follows from direct calculation.
%\item[(ii)] Fix $g\in\Isom(\L)$ and let $h = e_{\L,\B}\circ g\circ e_{\B,\L}\in \Isom(\B)$. Then by (i), $h$ extends uniquely to a M\"obius transformation of $\what\HH$. Now direct calculation shows that $g$ is a linear map preserving $Q$.
%
%The ``conversely'' part is clear from \eqref{hyperboloidmetric}.
%\end{itemize}
%Finally, if $X = \K$, then any $g\in\Isom(X)$ extends continuously to $\bord\K$ by conjugating to any of the other models.
%\end{proof}
%}% end ignore

%\ignore{
%
%\begin{proposition}
%\label{propositionOLQ}
%\begin{equation}
%\label{OLQ}
%\Isom(\H) \equiv \PO^*(\LL;\QQ).
%\end{equation}
%\end{proposition}
%\begin{proof}~
%[Add in proof.\internal]
%Fix $g\in\Isom(\L)$. By \eqref{}, we have
%\begin{equation}
%\label{isometryBQ}
%B_\QQ(g(\vv_1),g(\vv_2)) = B_\QQ(\vv_1,\vv_2)
%\end{equation}
%for all $\vv_1,\vv_2\in\L$.
%\begin{claim}
%\label{claimlinearisometry}
%If a linear combination $\sum_n c_n \vv_n$ is zero with $\vv_n\in\L$, then $\sum_n c_n g(\vv_n) = 0$.
%\end{claim}
%\begin{subproof}
%For all $\vv\in\L$, by \eqref{isometryBQ} we have
%\[
%B_\QQ(\sum_n c_n g(\vv_n),\vv) = B_\QQ(\sum_n c_n \vv_n,g^{-1}(\vv)) = 0.
%\]
%By linearity, $B_\QQ(\sum_n c_n g(\vv_n),\vv) = 0$ for all $\vv\in\LL$. Letting $\vv$ range over the basis vectors, \eqref{Qdef} implies that $\sum_n c_n g(\vv_n) = 0$.
%\end{subproof}
%Now define $\w g:\LL\to\LL$ by letting $\w g(\sum_n c_n \vv_n) = \sum_n c_n g(\vv_n)$ whenever $\vv_n\in\L$. Claim \ref{claimlinearisometry} implies that $\w g$ is well-defined. Clearly, $\w g$ is linear and extends $g$. Direct calculation shows that \eqref{isometryBQ} holds when $g$ is replaced by $\w g$; i.e. $\w g$ is a linear $Q$-isometry. This completes the proof.
%\end{proof}
%
%}% end ignore



%%% Old definition of \Mob(\what\HH)
%\begin{definition}
%\[
%\Mob(\what\HH) := \{ g : \K \to \K \vert \ g \ \emph{preserves the Riemannian metric $\lb \cdot ~,~ \cdot \rb_{\K}$} \} \cdot
%\]
%\end{definition}


%\ignore{
%\subsubsection{Orbits of $\0$ and $\infty$}
%
%\begin{proposition}\label{propositiontransitivity}
%The group $\Isom(\H)$ acts transitively on $\H$ and triply transitively on $\del \H$.
%\end{proposition}
%
%The proof of Proposition \ref{propositiontransitivity} will use the following lemma:
%\begin{lemma}
%\label{lemmaOHtransitive}
%The group $\O(\HH)$ acts transitively on $S(\0,1)$.
%\end{lemma}
%\begin{proof}
%Take $\pp\in S(\0,1)$ and extend to an orthonormal basis $\mathcal{B} := \{ \bb_1 = \pp, \bb_2, \ldots \}$. This is possible since every maximal orthonormal set is an orthonormal basis. Now define $T:\HH\to\HH$ by $T[\xx] = \sum_{i = 1}^\infty x_i \bb_i$. Then $T\in\O(\HH)$ and $T(\ee_1) = \pp$. Thus the orbit of $\ee_1$ is all of $S(\0,1)$.
%\end{proof}
%\begin{proof}[Proof of Proposition \ref{propositiontransitivity}]
%The proof splits into three parts. First we show that $\Isom(\H)$ acts transitively on $\H$. Next we show that $\Isom(\B)$ acts transitively on $\del \B$ and by equivalence of models we get the same result for $\Isom(\H)$ acting on $\del \H$. Finally, we show that $\Isom(\H)$ acts triply transitively on $\del \H$.
%\begin{itemize}
%\item[1.][$\Isom(\H)$ acts transitively on $\H$]
%For any $\qq\in\H$, let $\lambda = q_1$, let $\pp = (0,q_2,q_3, \ldots)$, and then let $g(\xx) := \lambda \xx + \pp$. By Proposition \ref{propositionstabilizer} we have $g\in\Stab(\Isom(\H);\infty)$; thus $g\in\Isom(\H)$. On the other hand, direct calculation shows that $g(\ee_1) = \qq$. Therefore the action is transitive on $\H$. 
%\item[2.][$\Isom(\B)$ acts transitively on $\del \B$]
%By Proposition \ref{propositionstabilizer} we have $\Stab(\B;\0) = \O(\HH)$, which by Lemma \ref{lemmaOHtransitive} acts transitively on $S(\0,1) = \del\B$.
%\item[3.][$\Isom(\H)$ acts triply transitively on $\del \H$]
%Fix $\xi_1,\xi_2,\xi_3\in\del \H$ distinct. Then by transitivity of $\Isom(\H)$ acting on $\del\H$ there exists $g\in\Isom(\H)$ such that $g(\infty) = \xi_1$. Let $\eta_i := g^{-1}(\xi_i)$ for $i = 1,2,3$, so that $\eta_1 = \infty$. Then $\eta_2,\eta_3\in\EE$, so by Lemma \ref{lemmaOHtransitive}, we may choose $T\in\O(\EE)$ such that $T[\ee_2] = \frac{\eta_3 - \eta_2}{\| \eta_3 - \eta_2 \|}$. Let $\lambda = \| \eta_3 - \eta_2 \|$, let $\pp = \eta_2$, and let $h(\xx) = \lambda T[\xx] + \pp$ for $\xx\in\EE$. Then $h\in\Sim(\EE)$. Let $\what h\in\Sim(\HH)$ be the Poincar\'e extension of $h$. Then by Proposition \ref{propositiontransitivity} we have $\what h\in\Stab(\Isom(\H);\infty)\leq\Isom(\H)$. On the other hand, direct calculation shows that
%\[
%\what h(\0) = \eta_2, \hspace{.5 in} \what h(\ee_2) = \eta_3.
%\]
%Letting $j = g\what h$, we have
%\[
%j(\infty) = \xi_1,\hspace{.5 in}j(\0) = \xi_2,\hspace{.5 in}j(\ee_2) = \xi_3.
%\]
%Thus $\Isom(\H)$ acts triply transitively on $\del \H$.
%\end{itemize}
%\end{proof}
%
%}% end ignore




%\ignore{
%\begin{remark} If we restrict to the subclass of $\Mob(\what\HH)$ defined by
%\[
%\Mob^*(\what\HH) := \{ \ g\in\Mob(\what\HH) \ \vert \ \emph{$g$ is the composition of finitely many inversions}\} ,
%\]
%then we have
%\[
%\Mob^*(\what\HH) \subsetneq \Mob(\what\HH) .
%\] 
%\end{remark}
%This is in contrast to finite dimensions etc.
%\begin{proof}
%Let $g(\xx) = \lambda T\circ \inv_\pp (\xx) + \qq$ for some $(\lambda,\pp,\qq,T)$ as in Liouville's Theorem. Then $g$ an be written as a finite composition of inversions if and only if $\Fix(T)$ has finite codimension. For example, the shift map on $\ell^2(\Z)$ cannot be written as a finite composition of inversions. We give an indication of why this is true. Say $\Fix(T)$ has finite codimension, then one can find a finite-dimensional subspace $V$ such that the entire map is the Poincar\'e extension of its restriction $V$. This reduces us to the finite-dimensional statement \cite{Ratcliffe} that every M\"obius map is a composition of finitely many inversions which may then be re-extended. On the other hand if one computes the composition of two inversions it can be shown that $\Fix(T)$ has codimension $1$ and so composing finitely many inversions only adds one finitely many times to the codimension. 
%\end{proof}
%}% end ignore

%\ignore{
%\begin{remark}
%One cannot make sense of \emph{orientation-preserving} transformations in infinite dimensions as one cannot define a meaningful notion of orientation. If one wanted to define orientation-preserving via the kernel of a continuous homomorphism $\mathcal{O}: \O(\HH) \to \Z_2$ one would easily fall into a trap. (Here $\Z_2$ is the group with two elements.) For example, any reflection in a hyperplane on $\ell^2(\Z)$ would be orientation-preserving. For a construction of such a map, take the commutator of the shift map squared and the map that switches consecutive pairs of nonnegative coordinates, i.e. $0$ and $1$, $2$ and $3$, etc.
%
%Indeed, for $\uu,\vv\in\HH$, let $H_{\uu,\vv}:= \{ \xx + \vv : \xx\in\uu^{\perp} \}$ be the hyperplane determined by $\uu$ and $\vv$ and let $r_{\uu,\vv}$ be reflection in this hyperplane given by $\xx \mapsto \left(\id - 2P_\uu \right)(\yy - \xx) + \yy$, where $P_\uu$ is the projection onto the hyperplane $\uu^\perp$. Then $\O(r_{\uu,\vv})$ can be shown to equal $I$, i.e. be orientation-preserving.
%\end{remark}
%
%}% end ignore


%\draftnewpage % Organization for the next section
%\begin{itemize}
%\item[A.1] Define CAT(-1) space
%\item[subsub] Define $\R$-trees
%\item[B.2] Define Gromov hyperbolic metric space
%\item[subsub] Examples of hyperbolic metric spaces
%\item[C] The boundary of a hyperbolic space
%\item[subsub] Extending the Gromov product and Busemann function to the boundary
%\item[subsub] A topology on $\bord X$ (Requires (I))
%\item[E] Gromov product in \CROSSs
%\item[F] Gromov boundary of a \CROSS\ (requires (C, including the topology part))
%\item[I] Different (meta)metrics on the boundary
%\item[subsub] Visual (spherical / Bourdon) (requires (E) for motivation) (requires (C, not including topology part) to prove Obs 3.64)
%\item[subsub] Standing assumptions
%\item[subsub] Euclidean / Hamenstadt (requires (C, topology part))
%\item[*] Sullivan's formula
%\item SECTION BREAK
%\item[G] Rips condition
%\item[H] CAT(-1) implies regularly geodesic (requires (G), (I))
%\item[J] Derivatives
%\item[subsub] Of metametrics
%\item[subsub] Of maps
%\item[subsub] Dynamical
%\item[subsub] In the half-space model
%\item[K] Generalized polar coordinates
%\item[L] Geometry of shadows
%\item[*] SSL (metric version)
%\item[*] Bounded Distortion Lemma
%\item[*] Big Shadows Lemma
%\item[*] Intersecting Shadows Lemma
%\item[*] Diameter of Shadows Lemma
%\end{itemize}
%\draftnewpage

%%%%%%%%%%%%%%%%%%%%%%%%%%%%%%%%%%
%%\draftnewpage
%%\section{$\R$-trees, CAT(-1) spaces, and Gromov hyperbolic metric spaces} \label{sectiongeometry1}
%%%%%%%%%%%%%%%%%%%%%%%%%%%%%%%%%%
%%
%%In this section we review the theory of ``negative curvature'' in general metric spaces. A good reference for this subject is \cite{BridsonHaefliger}. We begin by defining the class of \emph{$\R$-trees}, the main class of examples we will talk about in this monograph other than the class of \ROSSs, which we will discuss in more detail in Section \ref{sectionRtrees}. Next we will define CAT(-1) spaces, which are geodesic metric spaces whose triangles are ``thinner'' than the corresponding triangles in two-dimensional real hyperbolic space $\H^2$. Both \ROSSs\ and $\R$-trees are examples of CAT(-1) spaces. The next level of generality considers Gromov hyperbolic metric spaces. After defining these spaces, we proceed to define the boundary $\del X$ of a hyperbolic metric space $X$, introducing the family of so-called \emph{visual metrics} on the bordification $\bord X := X\cup\del X$. We show that the bordification of a \CROSS\ $X$ is isomorphic to its closure $\bord X$ defined in Section \ref{sectionROSSONCTs}; under this isomorphism, the visual metric on $\del \B^\alpha$ is proportional to the Euclidean metric.
%%
%%%%%%%%%%%%%%%%%%%%%%%%%%%%%%%%%%
%%\bigskip
%%\subsection{Graphs and $\R$-trees}
%%\label{subsectiongraphs}
%%%%%%%%%%%%%%%%%%%%%%%%%%%%%%%%%%
%%
%%To motivate the definition of $\R$-trees we begin by defining simplicial trees, which requires first defining graphs.
%%
%%\begin{definition}
%%\label{definitiongraphmetrization}
%%A \emph{weighted undirected graph} is a triple $(V,E,\ell)$, where $V$ is a nonempty set, $E\subset V\times V\butnot\{(x,x) : x\in V\}$ is invariant under the map $(x,y)\mapsto (y,x)$, and $\ell:E\to(0,\infty)$ is also invariant under $(x,y)\to (y,x)$. (If $\ell \equiv 1$, the graph is called \emph{unweighted}, and can be denoted simply $(V,E)$.) The graph is called \emph{connected} if for all $x,y\in V$, there exist $x = z_0,z_1,\ldots,z_n = y$ such that $(z_i,z_{i + 1}) \in E$ for all $i = 0,\ldots,n - 1$. If $(V,E,\ell)$ is connected, then the \emph{path metric} on $V$ is the metric
%%\begin{equation}
%%\label{pathmetric}
%%\dist_{E,\ell}(x,y) := \inf\left\{ \sum_{i = 0}^{n - 1} \ell(z_i,z_{i + 1}) : z_0 = x,\; z_n = y,\; (z_i,z_{i + 1})\in E \all i = 0,\ldots,n - 1\right\}.
%%\end{equation}
%%The \emph{geometric realization} of the graph $(V,E,\ell)$ is the metric space
%%\[
%%X = X(V,E,\ell) = \left(V\cup\bigcup_{(v,w)\in E} [0,\ell(v,w)]\right)/\sim,
%%\]
%%where $\sim$ represents the following identifications:
%%\begin{align*}
%%v &\sim ((v,w),0) \all (v,w)\in E\\
%%((v,w),t) &\sim ((w,v),\ell(v,w) - t) \all (v,w)\in E \all t\in[0,\ell(v,w)]
%%\end{align*}
%%and the metric $\dist$ on $X$ is given by
%%\[
%%\dist\big(((v_0,v_1),t),((w_0,w_1),s)\big) = \min\{ |t - i\ell(v_0,v_1)| + \dist(v_i,w_j) + |s - j\ell(w_0,w_1)| : i = 0,1, j = 0,1\}.
%%\]
%%(The geometric realization of a graph is sometimes also called a graph. In the sequel, we shall call it a \emph{geometric graph}.)
%%\end{definition}
%%
%%\begin{example}[The Cayley graph of a group]
%%\label{examplecayleygraph}
%%Let $\Gamma$ be a group, and let $E_0\subset \Gamma$ be a generating set. (In most circumstances $E_0$ will be finite; there is an exception in Example \ref{examplenotstronglydiscrete} below.) Assume that $E_0 = E_0^{-1}$. The \emph{Cayley graph} of $\Gamma$ with respect to the generating set $E_0$ is the unweighted graph $(\Gamma,E)$, where
%%\begin{equation}
%%\label{Edef}
%%(\gamma,\beta)\in E \;\;\Leftrightarrow\;\; \gamma^{-1}\beta\in E_0.
%%\end{equation}
%%More generally, if $\ell_0:E_0\to (0,\infty)$ satisfies $\ell_0(g^{-1}) = \ell_0(g)$, the \emph{weighted Cayley graph} of $\Gamma$ with respect to the pair $(E_0,\ell_0)$ is the graph $(\Gamma,E,\ell)$, where $E$ is defined by \eqref{Edef}, and
%%\begin{equation}
%%\label{elldef}
%%\ell(\gamma,\beta) = \ell_0(\gamma^{-1}\beta).
%%\end{equation}
%%The path metric of a Cayley graph is called a \emph{Cayley metric}.
%%\end{example}
%%
%%\begin{remark}
%%\label{remarknaturalaction}
%%The equations \eqref{Edef}, \eqref{elldef} guarantee that for each $\gamma\in\Gamma$, the map $\Gamma\ni\beta\to \gamma\beta\in\Gamma$ is an isometry of $\Gamma$ with respect to any Cayley metric. This isometry extends in a unique way to an isometry of the geometric Cayley graph $X = X(\Gamma,E,\ell)$. The map sending $\gamma$ to this isometry is a homomorphism from $\Gamma$ to $\Isom(X)$, and is called the \emph{natural action} of $\Gamma$ on $X$.
%%\end{remark}
%%
%%\begin{remark}
%%\label{remarkpathmetricuniversal}
%%The path metric \eqref{pathmetric} satisfies the following \emph{universal property}: If $Y$ is a metric space and if $\phi:V\to Y$ satisfies $\dist(\phi(v),\phi(w)) \leq \ell(v,w) \all (v,w)\in E$, then $\dist(\phi(v),\phi(w)) \leq \dist(v,w) \all v,w\in V$.% In particular, any two Cayley metrics defined by finite generating sets are bi-Lipschitz equivalent. Thus, whenever $\Gamma$ is finitely generated we call any 
%%\end{remark}
%%
%%\begin{remark}
%%\label{remarkgeodesicmetric}
%%The main difference between the metric space $(V,\dist_{E,\ell})$ and the geometric graph $X = X(V,E,\ell)$ is that the latter is a \emph{geodesic} metric space. A metric space $X$ is said to be \emph{geodesic} if for every $p,q\in X$, there exists an isometric embedding $\pi:[t,s]\to X$ such that $\pi(t) = p$ and $\pi(s) = q$, for some $t,s\in\R$. The set $\pi([t,s])$ is denoted $\geo pq$ and is called a \emph{geodesic segment connecting $p$ and $q$}. The map $\pi$ is called a \emph{parameterization} of the geodesic segment $\geo pq$. (Note that although $\geo qp = \geo pq$, $\pi$ is not a parameterization of $\geo qp$.)
%%
%%Warning: Although we may denote any geodesic segment connecting $p$ and $q$ by $\geo pq$, such a geodesic segment is not necessarily unique. A geodesic metric space $X$ is called \emph{uniquely geodesic} if for every $p,q\in X$, the geodesic segment connecting $p$ and $q$ is unique.
%%\end{remark}
%%
%%\begin{notation}
%%\label{notationgeopqt}
%%If $\pi:[0,t_0]\to X$ is a parameterization of the geodesic segment $\geo pq$, then for each $t\in[0,t_0]$, $\geo pq_t$ denotes the point $\pi(t)$, i.e. the unique point on the geodesic segment $\geo pq$ such that $\dist(p,\geo pq_t) = t$.
%%\end{notation}
%%
%%We are now ready to define the class of simplicial trees. Let $(V,E,\ell)$ be a weighted undirected graph. A \emph{cycle} in $(V,E,\ell)$ is a finite sequence of distinct vertices $v_1,\ldots,v_n\in V$, with $n\geq 3$, such that
%%\begin{equation}
%%\label{cycle}
%%(v_1,v_2), (v_2,v_3), \ldots, (v_{n - 1},v_n), (v_n,v_1)\in E.
%%\end{equation}
%%\begin{definition}
%%\label{definitionweightedtree}
%%A \emph{simplicial tree} is the geometric realization of a weighted undirected graph with no cycles. A \emph{$\Z$-tree} (or \emph{unweighted simplicial tree}, or just \emph{tree}) is the geometric realization of an unweighted undirected graph with no cycles.
%%\end{definition}
%%
%%\begin{example}
%%Let $\F_2(\Z)$ denote the free group on two elements $\gamma_1,\gamma_2$. Let $E_0 = \{\gamma_1,\gamma_1^{-1},\gamma_2,\gamma_2^{-1}\}$. The geometric Cayley graph of $\F_2(\Z)$ with respect to the generating set $E_0$ is an unweighted simplicial tree.
%%\end{example}
%%
%%\begin{example}
%%\label{exampletreetriangle}
%%Let $V = \{\wbar C,\wbar p,\wbar q,\wbar r\}$, and fix $\ell_{\wbar p},\ell_{\wbar q},\ell_{\wbar r} > 0$. Let $E = \{(\wbar C,x), (x,\wbar C) : x = \wbar p,\wbar q,\wbar r\}$, and let $\ell(\wbar C,x) = \ell(x,\wbar C) = \ell_x$. The geometric realization of the graph $(V,E,\ell)$ is a simplicial tree; see Figure \ref{figureRtree}. It will be denoted $\wbar \Delta = \Delta(\wbar p,\wbar q,\wbar r)$, and will be called a \emph{tree triangle}. For $x,y\in\{\wbar p,\wbar q,\wbar r\}$ distinct, the distance between $x$ and $y$ is given by
%%\[
%%\dist(x,y) = \ell_x + \ell_y.
%%\]
%%Solving for $\ell_{\wbar p}$ in terms of $\dist(\wbar p,\wbar q),\dist(\wbar p,\wbar r),\dist(\wbar q,\wbar r)$ gives
%%\begin{equation}
%%\label{gromovprecursor}
%%\ell_{\wbar p} = \dist(\wbar p,\wbar C) = \frac12[\dist(\wbar p,\wbar q) + \dist(\wbar p,\wbar r) - \dist(\wbar q,\wbar r)].
%%\end{equation}
%%\end{example}
%%
%%\begin{definition}
%%\label{definitionRtree}
%%A metric space $X$ is an \emph{$\R$-tree} if for all $p,q,r\in X$, there exists a tree triangle $\wbar\Delta = \Delta(\wbar p,\wbar q,\wbar r)$ and an isometric embedding $\iota:\wbar\Delta\to X$ sending $\wbar p,\wbar q,\wbar r$ to $p,q,r$, respectively.
%%\end{definition}
%%\begin{definition}
%%\label{definitioncenterRtree}
%%Let $X$ be an $\R$-tree, fix $p,q,r\in X$, and let $\iota:\wbar\Delta\to X$ be as above. The point $C = C(p,q,r) := \iota(\wbar C)$ is called the \emph{center} of the geodesic triangle $\Delta = \Delta(p,q,r)$.
%%\end{definition}
%%
%%As the name suggests, every simplicial tree is an $\R$-tree; the converse does not hold; see e.g. \cite[Example on p.50]{Chiswell}. Before we can prove that every simplicial tree is an $\R$-tree, we will need a lemma:
%%
%%\begin{lemma}[Cf. {\cite[p.29]{Chiswell}}]
%%\label{lemmaRtreeequivalent}
%%Let $X$ be a metric space. The following are equivalent:
%%\begin{itemize}
%%\item[(A)] $X$ is an $\R$-tree.
%%\item[(B)] There exists a collection of geodesics $\GG$, with the following properties:
%%\begin{itemize}
%%\item[(BI)] For each $x,y\in X$, there is a geodesic $\geo xy\in\GG$ connecting $x$ and $y$.
%%\item[(BII)] Given $\geo xy\in\GG$ and $z,w\in \geo xy$, we have $\geo zw\in\GG$, where $\geo zw$ is interpreted as the set of points in $\geo xy$ which lie between $z$ and $w$.
%%\item[(BIII)] Given $x_1,x_2,x_3\in X$ distinct and geodesics $\geo{x_1}{x_2},\geo{x_1}{x_3},\geo{x_2}{x_3}\in\GG$, at least one pair of the geodesics $\geo{x_i}{x_j}$, $i\neq j$, has a nontrivial intersection. More precisely, there exist distinct $i,j,k\in\{1,2,3\}$ such that
%%\[
%%\geo{x_i}{x_j} \cap \geo{x_i}{x_k} \propersupset \{x_i\}.
%%\]
%%\end{itemize}
%%\end{itemize}
%%\end{lemma}
%%\begin{proof}[Proof of \text{(A)} \implies \text{(B)}]
%%Note that (BI) and (BII) are true for any uniquely geodesic metric space. Given $x_1,x_2,x_3$ distinct, let $C$ be the center. Then $x_i\neq C$ for some $i$; without loss of generality $x_1\neq C$. Then
%%\[
%%\geo{x_1}{x_2} \cap \geo{x_1}{x_3} = \geo{x_1}{C} \propersupset \{x_1\}.
%%\]
%%\end{proof}
%%\begin{proof}[Proof of \text{(B)} \implies \text{(A)}]
%%We first show that given $x_1,x_2,x_3\in X$ and $\geo{x_1}{x_2},\geo{x_1}{x_3},\geo{x_2}{x_3}\in\GG$, the intersection $\bigcap_{i\neq j} \geo{x_i}{x_j}$ is nonempty. Indeed, suppose not. For $i = 2,3$ let $\gamma_i:[0,\dist(x_1,x_i)]\to X$ be a parameterization of $\geo{x_1}{x_i}$, and let
%%\[
%%t_1 = \max\{t\geq 0 : \gamma_2(t) = \gamma_3(t)\}.
%%\]
%%By replacing $x$ with $\gamma_2(t_1) = \gamma_3(t_1)$ and using (BII), we may without loss of generality assume that $t_1 = 0$, or equivalently that $\geo{x_1}{x_2}\cap \geo{x_1}{x_3} = \{x_1\}$. Similarly, we may without loss of generality assume that $\geo{x_2}{x_1}\cap \geo{x_2}{x_3} = \{x_2\}$ and $\geo{x_3}{x_1}\cap\geo{x_3}{x_2} = \{x_3\}$. But then (BIII) implies that $x_1,x_2,x_3$ cannot be all distinct. This immediately implies that $\bigcap_{i\neq j} \geo{x_i}{x_j} \neq \emptyset$.
%%
%%To complete the proof, we must show that $X$ is uniquely geodesic. Indeed suppose that for some $x_1,x_2\in X$, there is more than one geodesic connecting $x_1$ and $x_2$. Let $\geo{x_1}{x_2}\in\GG$ be a geodesic connecting $x_1$ and $x_2$, and let $\geo{x_1}{x_2}'$ be another geodesic connecting $x_1$ and $x_2$. Then there exists $x_3\in\geo{x_1}{x_2}'\butnot\geo{x_1}{x_2}$. By the above paragraph, there exists $w\in\bigcap_{i\neq j} \geo{x_i}{x_j}$. Since $w\in \geo{x_i}{x_3}$, we have
%%\begin{equation}
%%\label{dxiw}
%%\dist(x_i,w) \leq \dist(x_i,x_3).
%%\end{equation}
%%On the other hand, since $w\in \geo{x_1}{x_2}$ and $x_3\in\geo{x_1}{x_2}'$, we have
%%\[
%%\dist(x_1,x_2) = \dist(x_1,w) + \dist(x_2,w) \leq \dist(x_1,x_3) + \dist(x_2,x_3) = \dist(x_1,x_2).
%%\]
%%It follows that equality holds in \eqref{dxiw}, i.e. $\dist(x_i,w) = \dist(x_i,x_3)$. Since $w\in \geo{x_i}{x_3}$, this implies $w = x_3$. But then $x_3 = w\in \geo{x_1}{x_2}$, a contradiction.
%%\end{proof}
%%
%%%\begin{lemma}
%%%\label{lemmaRtreeequivalent}
%%%Let $X$ be a metric space. The following are equivalent:
%%%\begin{itemize}
%%%\item[(A)] $X$ is an $\R$-tree.
%%%\item[(B)] There exists a collection of geodesics $(\geo xy)_{x,y\in X}$, with the following properties:
%%%\begin{itemize}
%%%\item[(BI)] For each $x,y\in X$, $\geo xy$ is a geodesic between $x$ and $y$.
%%%\item[(BII)] Given $z,w\in \geo xy$, $\geo zw\subset \geo xy$.
%%%\item[(BIII)] Given $x_1,x_2,x_3\in X$ distinct, at least one pair of the geodesics $\geo{x_i}{x_j}$, $i\neq j$, has a nontrivial intersection. More precisely, there exist distinct $i,j,k\in\{1,2,3\}$ such that
%%%\[
%%%\geo{x_i}{x_j} \cap \geo{x_i}{x_k} \propersupset \{x_i\}.
%%%\]
%%%\end{itemize}
%%%\end{itemize}
%%%\end{lemma}
%%%\begin{proof}[Proof of \text{(A)} \implies \text{(B)}]
%%%Note that (BI) and (BII) are true for any uniquely geodesic metric space. Given $x_1,x_2,x_3$ distinct, let $C$ be the center. Then $x_i\neq C$ for some $i$; without loss of generality $x_1\neq C$. Then
%%%\[
%%%\geo{x_1}{x_2} \cap \geo{x_1}{x_3} = \geo{x_1}{C} \propersupset \{x_1\}.
%%%\]
%%%\end{proof}
%%%\begin{proof}[Proof of \text{(B)} \implies \text{(A)}]
%%%We first show that given $x_1,x_2,x_3\in X$, the intersection $\bigcap_{i\neq j} \geo{x_i}{x_j}$ is nonempty. Indeed, suppose not. For $i = 2,3$ let $\gamma_i:[0,\dist(x_1,x_i)]\to X$ be a parameterization of $\geo{x_1}{x_i}$, and let
%%%\[
%%%t_1 = \max\{t\geq 0 : \gamma_2(t) = \gamma_3(t)\}.
%%%\]
%%%By replacing $x$ with $\gamma_2(t_1) = \gamma_3(t_1)$ and using (BII), we may without loss of generality assume that $t_1 = 0$, or equivalently that $\geo{x_1}{x_2}\cap \geo{x_1}{x_3} = \{x_1\}$. Similarly, we may without loss of generality assume that $\geo{x_2}{x_1}\cap \geo{x_2}{x_3} = \{x_2\}$ and $\geo{x_3}{x_1}\cap\geo{x_3}{x_2} = \{x_3\}$. But then (BIII) implies that $x_1,x_2,x_3$ cannot be all distinct. This immediately implies that $\bigcap_{i\neq j} \geo{x_i}{x_j} \neq \emptyset$.
%%%
%%%To complete the proof, we must show that $X$ is uniquely geodesic. Indeed suppose that for some $x_1,x_2\in X$, there is a geodesic $\geo{x_1}{x_2}'$ which is not in the collection. Then there exists $x_3\in\geo{x_1}{x_2}'\butnot\geo{x_1}{x_2}$. By the above paragraph, there exists $w\in\bigcap_{i\neq j} \geo{x_i}{x_j}$. Since $w\in \geo{x_i}{x_3}$, we have
%%%\begin{equation}
%%%\label{dxiw}
%%%\dist(x_i,w) \leq \dist(x_i,x_3).
%%%\end{equation}
%%%On the other hand, since $w\in \geo{x_1}{x_2}$ and $x_3\in\geo{x_1}{x_2}'$, we have
%%%\[
%%%\dist(x_1,x_2) = \dist(x_1,w) + \dist(x_2,w) \leq \dist(x_1,x_3) + \dist(x_2,x_3) = \dist(x_1,x_2).
%%%\]
%%%It follows that equality holds in \eqref{dxiw}, i.e. $\dist(x_i,w) = \dist(x_i,x_3)$. Since $w\in \geo{x_i}{x_3}$, this implies $w = x_3$. But then $x_3 = w\in \geo{x_1}{x_2}$, a contradiction.
%%%\end{proof}
%%
%%\begin{corollary}
%%Every simplicial tree is an $\R$-tree.
%%\end{corollary}
%%\begin{proof}
%%Let $X = X(V,E,\ell)$ be a simplicial tree, and let $\GG$ be the collection of all geodesics; then (BI) and (BII) both hold. By contradiction, suppose that there exist points $x_1,x_2,x_3\in X$ such that $\geo{x_i}{x_j}\cap\geo{x_i}{x_k} = \{x_i\}$ for all distinct $i,j,k\in\{1,2,3\}$. Then the path $\geo{x_1}{x_2}\cup\geo{x_2}{x_3}\cup\geo{x_3}{x_1}$ is equal to the union of the edges of a cycle of the graph $(V,E,\ell)$. This is a contradiction.
%%\end{proof}
%%
%%We shall investigate $\R$-trees in more detail in Section \ref{sectionRtrees}, where we will give various examples of $\R$-trees together with groups acting isometrically on them.
%%
%%
%%
%%%\begin{corollary}
%%%Let $X$ be a tree, and fix $p,q,r\in X$.
%%%\begin{itemize}
%%%\item[(i)] If $\geo pq\cap \geo qr = \{q\}$, then $\geo pr = \geo pq \cup \geo qr$
%%%\item[(ii)] There exists $C\in X$ such that $\geo pq\cap \geo pr = \geo pC$.
%%%\item[(iii)] $\geo pq\cap \geo qr\cap \geo rp \neq \emptyset$.
%%%\end{itemize}
%%%\end{corollary}
%%
%%
%%
%%
%%
%%
%%%\begin{definition}[Cf. {\cite[p.29]{Chiswell}}]
%%%A metric space $X$ is an \emph{$\R$-tree} (or \emph{metric tree}) if the following hold:
%%%\begin{itemize}
%%%\item[(I)] $X$ is a geodesic metric space.
%%%\item[(II)] Given three points $p,q,r\in X$, if $\geo pq\cap \geo qr = \{q\}$, then $\geo pr = \geo pq \cup \geo qr$.
%%%\item[(III)] Given three points $p,q,r\in X$, there exists $C$ such that $\geo pq\cap \geo pr = \geo pC$.
%%%\end{itemize}
%%%\end{definition}
%%
%%%\begin{definition}
%%%A metric space $X$ is an \emph{$\R$-tree} (or \emph{metric tree}) if for any geodesic triangle $\Delta = \Delta(p,q,r)$ in $X$, there is a (unique) point
%%%\[
%%%C(p,q,r)\in\geo pq\cap \geo qr\cap \geo rp.
%%%\]
%%%The point $C(p,q,r)$ is called the \emph{center} of the triangle $\Delta$; see Figure \ref{figureRtree}.
%%%\end{definition}
%%
%%
%%%\begin{proposition}
%%%Let $(V,E)$ be an undirected graph. The metrization $X(V,E)$ is an $\R$-tree if and only if $(V,E)$ is a tree.
%%%\end{proposition}
%%
%%\begin{figure}
%%\begin{center}
%%\begin{tikzpicture}[line cap=round,line join=round,>=triangle 45,x=1.0cm,y=1.0cm]
%%\clip(0.22,0.49) rectangle (5.46,4.0);
%%\draw (0.8,0.8)-- (2.76,1.84);
%%\draw (2.76,1.84)-- (2.84,3.74);
%%\draw (2.76,1.84)-- (4.74,2);
%%\begin{scriptsize}
%%\fill [color=black] (0.8,0.8) circle (1.2pt);
%%\draw[color=black] (0.68,0.6) node {$p$};
%%\fill [color=black] (2.76,1.84) circle (1.2pt);
%%\draw[color=black] (2.06,1.98) node {$C(p,q,r)$};
%%\fill [color=black] (2.84,3.74) circle (1.2pt);
%%\draw[color=black] (3.02,3.86) node {$q$};
%%\fill [color=black] (4.74,2) circle (1.2pt);
%%\draw[color=black] (5.04,2.04) node {$r$};
%%\end{scriptsize}
%%\end{tikzpicture}
%%\caption{A geodesic triangle in an $\R$-tree.}
%%\label{figureRtree}
%%\end{center}
%%\end{figure}
%%
%%%\begin{figure}
%%%\centerline{\mbox{\includegraphics[scale = 0.7]{Rtree.png}}}
%%%\caption{A geodesic triangle in an $\R$-tree.}
%%%\label{figureRtree}
%%%\end{figure}
%%
%%%There are several equivalent definitions of $\R$-trees; cf. \cite[p.29]{Chiswell}. One which we will find particularly useful is the following:
%%
%%
%%
%%%
%%%\begin{remark}
%%%This definition disagrees slightly with other sources (e.g. \cite{Chiswell}) that would consider the set $V$ to be the $\Z$-tree rather than $X$. Hopefully this will not cause too much confusion.
%%%\end{remark}
%%
%%
%%%\begin{remark}
%%%\label{remarksimplicialtree}
%%%An $\R$-tree $X$ is a simplicial tree if and only if there exists a discrete set $V\subset X$ with the following properties:
%%%\begin{itemize}
%%%\item[(I)] If $x,y,z\in V$, then $C(x,y,z)\in V$.
%%%\item[(II)] $X = \bigcup_{x,y\in V}\geo xy$.
%%%\end{itemize}
%%%Moreover, the simplicial tree is unweighted if and only if $\dist(x,y)\in\N$ for all $x,y\in V$.
%%%\end{remark}
%%%\begin{proof}
%%%First, suppose that $X$ is a simplicial tree, the geometric realization of the weigted undirected graph $(V,E,\ell)$. Fix $x,y,z\in V$, and let $C = C(x,y,z)$. If $C\in \geo pq\butnot\{p,q\}$ for some $(p,q)\in E$, we would have $\#(\del B(C,r)) = 2$ for all sufficiently small $r$, but the fact that $C$ is the center of the triangle $\Delta(x,y,z)$ implies that $\#(\del B(C,r)) \geq 3$ for all sufficiently small $r$, unless $C\in \{x,y,z\}$. Either is a contradiction, so we conclude that $C\in V$, demonstrating (I). (II) is obvious from the definition of a simplicial tree.
%%%
%%%Conversely, suppose that we have a discrete set $V$ satisfying (I) and (II). Let $E = \{(x,y)\in V: \geo xy\cap V = \{x,y\}\}$, let $\ell((x,y)) = \dist(x,y)$, and let $Y$ be the geometric realization of the weighted undirected graph $(V,E,\ell)$. We can define an embedding from $Y$ to $X$ by letting $\phi(((x,y),t)) = \geo xy_t$. This embedding is clearly $1$-Lipschitz; to show that it is isometric, take $x,y\in Y$ and [Continue\internal]
%%%\end{proof}
%%
%%
%%
%%%%%%%%%%%%%%%%%%%%%%%%%%%%%%%%%%
%%\bigskip
%%\subsection{CAT(-1) spaces}
%%\label{subsectionCAT}
%%%%%%%%%%%%%%%%%%%%%%%%%%%%%%%%%%
%%
%%The following definitions have been modified from \cite[p.158]{BridsonHaefliger}, to which the reader is referred for more details.
%%
%%A \emph{geodesic triangle} in $X$ consists of three points $p,q,r\in X$ (the \emph{vertices} of the triangle) together with a choice of three geodesic segments $\geo pq$, $\geo qr$, and $\geo rp$ joining them (the \emph{sides}). Such a geodesic triangle will be denoted $\Delta(p,q,r)$, although we note that this could cause ambiguity if $X$ is not uniquely geodesic. Although formally $\Delta(p,q,r)$ is an element of $X^3\times \P(X)^3$, we will sometimes identify $\Delta(p,q,r)$ with the set $\geo pq \cup \geo qr\cap \geo rp \subset X$, writing $x\in\Delta(p,q,r)$ if $x\in \geo pq \cup \geo qr\cup \geo rp$.
%%
%%A triangle $\wbar{\Delta} = \Delta(\wbar p,\wbar q,\wbar r)$ in $\H^2$ is called a \emph{comparison triangle} for $\Delta = \Delta(p,q,r)$ if $\dist(\wbar p, \wbar q) = \dist(p,q)$, $\dist(\wbar q,\wbar r) = \dist(q, r)$, and $\dist(\wbar p,\wbar r) = \dist(p,r)$. Any triangle admits a comparison triangle, unique up to isometry. For any point $x\in\geo pq$, we define its \emph{comparison point} $\wbar x\in\geo{\wbar p}{\wbar q}$ to be the unique point such that $\dist(\wbar x,\wbar p) = \dist(x,p)$ and $\dist(\wbar x,\wbar q) = \dist(x,q)$. In the notation above, the comparison point of $\geo pq_t$ is equal to $\geo{\wbar p}{\wbar q}_t$ for all $t\in[0,\dist(p,q)] = [0,\dist(\wbar p,\wbar q)]$. For $x\in \geo qr$ and $x\in \geo rp$, the comparison point is defined similarly.
%%
%%Let $X$ be a metric space and let $\Delta$ be a geodesic triangle in $X$. We say that $\Delta$ \emph{satisfies the CAT(-1) inequality} if for all $x,y\in\Delta$,
%%\begin{equation}
%%\label{CAT}
%%\dist(x, y) \leq \dist(\wbar x, \wbar y),
%%\end{equation}
%%where $\wbar x$ and $\wbar y$ are any\Footnote{The comparison points $\wbar x$ and $\wbar y$ may not be uniquely determined if either $x$ or $y$ lies on two sides of the triangle simultaneously.} comparison points for $x$ and $y$, respectively. Intuitively, $\Delta$ satisfies the CAT(-1) inequality if it is ``thinner'' than its comparison triangle $\wbar{\Delta}$.
%%
%%\begin{definition}
%%\label{definitionCAT}
%%$X$ is a \emph{CAT(-1) space} if it is a geodesic metric space and if all of its geodesic triangles satisfy the CAT(-1) inequality.
%%\end{definition}
%%%Note that it follows from this definition that CAT(-1) spaces are uniquely geodesic.
%%\begin{observation}[{\cite[Proposition II.1.4(1)]{BridsonHaefliger}}]
%%\label{observationCATuniquelygeo}
%%CAT(-1) spaces are uniquely geodesic.
%%\end{observation}
%%\begin{proof}
%%Let $X$ be a CAT(-1) space, and suppose that two points $p,q\in X$ are connected by two geodesic segments $\geo pq$ and $\geo pq'$. Fix $t\in [0,\dist(p,q)]$ and let $x = \geo pq_t$, $x' = \geo pq_t'$. Consider the triangle $\Delta(p,q,x)$ determined by the geodesic segments $\geo pq'$, $\geo px$, and $\geo xq$, and a comparison triangle $\Delta(\wbar p,\wbar q,\wbar x)$. Then $x$ and $x'$ have the same comparison point $\wbar x$, so by the CAT(-1) inequality
%%\[
%%\dist(x,x')\leq \dist(\wbar x,\wbar x) = 0,
%%\]
%%and thus $x = x'$. Since $t$ was arbitrary, it follows that $\geo pq = \geo pq'$. Since $\geo pq'$ was arbitrary, $\geo pq$ is the unique geodesic segment connecting $p$ and $q$.
%%\end{proof}
%%
%%\subsubsection{Examples of CAT(-1) spaces}
%%\label{subsubsectionexamplesCAT}
%%In this text we concentrate on two main examples of CAT(-1) spaces: \ROSSs\ and $\R$-trees. We therefore begin by proving the following result which implies that \ROSSs\ are CAT(-1):
%%
%%\begin{proposition}
%%\label{propositioncurvatureCAT}
%%Any Riemannian manifold (finite- or infinite-dimensional) with sectional curvature bounded above by $-1$ is a CAT(-1) space.
%%\end{proposition}
%%\begin{proof}
%%The finite-dimensional case is proven in \cite[Theorem II.1A.6]{BridsonHaefliger}. The infinite-dimensional follows upon augmenting the finite-dimensional proof with the infinite-dimensional Cartan--Hadamard theorem \cite[IX, Theorem 3.8]{Lang_differential_geometry} to guarantee surjectivity of the exponential map.%\comdavid{ Check the BH more carefully to see what precisely do they need...}
%%\end{proof}
%%
%%Since \ROSSs\ have sectional curvature bounded between $-4$ and $-1$ (e.g. \cite[Corollary of Proposition 4]{Heintze}; see also \cite[Lemmas 2.3, 2.7, and 2.11]{Quint}% for \CROSSs, \cite[(5)]{AravindaFarrell} for the Cayley hyperbolic plane
%%), the following corollary is immediate:
%%
%%\begin{corollary}
%%\label{corollaryROSSONCTCAT}
%%Every \ROSS\ is a CAT(-1) space.
%%\end{corollary}
%%
%%\begin{remark}
%%One can prove Corollary \ref{corollaryROSSONCTCAT} without using the full strength of Proposition \ref{propositioncurvatureCAT}. Indeed, any geodesic triangle in a \ROSS\ is isometric to a geodesic triangle in $\H_\F^2$ for some $\F\in\{\R,\C,\Q\}$. Since $\H_\F^2$ is finite-dimensional, thinness of its geodesic triangles follows from the finite-dimensional version of Proposition \ref{propositioncurvatureCAT}.
%%\end{remark}
%%
%%%\ignore{
%%%
%%%In fact, the above remark is true as well for infinite-dimensional Riemannian manifolds:
%%%
%%%\begin{theorem}[Infinite-dimensional Cartan--Hadamard theorem]
%%%Let $(X,g)$ be a simply-connected infinite-dimensional Riemannian manifold with nonpositive sectional curvature. Then for each $x\in X$, the exponential map $\exp:T_x X\to X$ is a diffeomorphism.
%%%\end{theorem}
%%%
%%%\begin{proof}
%%%Let $h$ be the pullback of $g$ under the exponential map. It follows from \cite{}\internal that $h\geq e$, where $e$ is the Euclidean metric on $T_x X\equiv \H$.
%%%
%%%We claim first that $\exp$ is a local diffeomorphism. By the inverse function theorem, it suffices to show that $D\exp(\vv)$ is invertible for all $\vv\in T_x X$. Indeed, fix suc a $\vv$, and for each $t\in [0,1]$ let $T_t = D\exp(t\vv)$. Let's assume for simplicity that $X = \H$, although in fact this is only locally true. Then $t\mapsto T_t$ is a continuous map from $[0,1]$ to the set of semiexpanding operators on $\H$. (The semiexpandingness comes from the theorem from \cite{}\intermal.) The continuity is with respect to the norm topology. Moreover, $T_0 = I$.
%%%
%%%Let $S = \{t\in[0,1]: T_t \text{ is invertible}\}$. To complete the proof, we need to show that $S$ is a clopen set. To see that $S$ is closed, suppose $S\ni t_n\to t$. Fix $\ee\in \H$. For each $n\in\N$, let $\vv_n = T_{t_n}^{-1}[\ee]$; we have $\|\vv_n\| \leq \|\ee\|$ since $T_{t_n}$ is semiexpanding. Then
%%%\[
%%%\|T_t[\vv_n] - \ee\| = \|(T_t - T_{t_n})[\vv_n]\| \leq \|T_t - T_{t_n}\| \cdot \|\ee\| \to 0,
%%%\]
%%%so $\ee\in \cl{T_t[\H]}$. Since $T_t$ is semiexpanding, $T_t[\H]$ is closed and so $\ee\in T_t[\H]$. Since $\ee$ was arbitrary, $T_t$ is invertible.
%%%
%%%To see that $S$ is open, suppose $t_n\to t\in S$. Then for sufficiently large $n$, $\|T_{t_n} - T_t\| < 1/\|T_t^{-1}\|$, which implies that the series
%%%\[
%%%T_t^{-1}\sum_{n = 0}^\infty (1 - T_{t_n}T_t^{-1})^n
%%%\]
%%%converges. Direct calculation shows that the sum of this series is an inverse of $T_{t_n}$, so $t_n\in S$ for all sufficiently large $n$.
%%%
%%%The conclusion now follows upon applying Lemma \ref{lemmacartanhadamard}.\end{proof}
%%%\begin{lemma}
%%%\label{lemmacartanhadamard}
%%%Let $(X,g)$ and $(Y,h)$ be complete infinite-dimensional Riemannian manifolds, and suppose that $\Phi:Y\to X$ is a local diffeomorphism which is also a local isometry. If $X$ is simply-connected, then $\Phi$ is a diffeomorphism (and thus a global isometry).
%%%\end{lemma}
%%%\begin{proof}
%%%We claim first that $\Phi$ has the homotopy lifting property for smooth maps: If $f:[0,1]\to X$ is smooth, and $y\in \Phi^{-1}(f(0))$, then there is a unique smooth map $g:[0,1]\to Y := T_x X$ such that $\Phi\circ g = f$. Indeed, consider the set $S = \{t\in[0,1]: \exists g:[0,t]\to Y \text{ a lift of $f$}\}$. Clearly $0\in S$; we need to show $1\in S$, so it suffices to show that $S$ is clopen. Suppose first that $S\ni t_n\to t$. Then there is a lift $g:\CO 0t\to Y$. The length of the path $g$ is the same as the length of the path $\Phi\circ g = f\given \CO 0t$, which is finite since $f$ is smooth. It follows that the sequence $(g(t_n))_1^\infty$ is Cauchy with respect to the metric $\dist_h$ for any sequence $t_n\to t$. Since $\dist_h\geq\dist_e$, the sequence is also Cauchy with respect to $\dist_e$, and since $\dist_e$ is complete it converges to a point $z\in Y$. Letting $g(t) = z$ therefore gives a lift $g:[0,t]\to Y$, demonstrating that $t\in S$.
%%%
%%%Now suppose that $t\in S$. Since $\Phi$ is a local diffeomorphism, there is a neighborhood $U$ of $g(t)$ such that $\Phi$ is a diffeomorphism on $U$. Let $t_0 > t$ be close enough to $t$ so that $f(t')\in \Phi(U)$ for all $t < t' < t_0$. Then there exists a lift $g:\CO 0{t_0}\to Y$, and it follows that $\CO 0{t_0}\subset S$.
%%%
%%%Thus $S$ is a clopen set, and so $1\in S$, completing the proof.
%%%\end{proof}
%%%
%%%}% end ignore
%%
%%
%%
%%%Now we move on to our second example of a CAT(-1) space, namely $\R$-trees.
%%%
%%%\begin{definition}
%%%A metric space $X$ is an \emph{$\R$-tree} (or \emph{metric tree}) if for any geodesic triangle $\Delta = \Delta(p,q,r)$ in $X$, there is a (unique) point
%%%\[
%%%C(p,q,r)\in\geo pq\cap \geo qr\cap \geo rp.
%%%\]
%%%The point $C(p,q,r)$ is called the \emph{center} of the triangle $\Delta$; see Figure \ref{figureRtree}.
%%%\end{definition}
%%%
%%%
%%%\begin{figure}
%%%\begin{center}
%%%\begin{tikzpicture}[line cap=round,line join=round,>=triangle 45,x=1.0cm,y=1.0cm]
%%%\clip(0.22,0.28) rectangle (5.46,4.2);
%%%\draw (0.8,0.8)-- (2.76,1.84);
%%%\draw (2.76,1.84)-- (2.84,3.74);
%%%\draw (2.76,1.84)-- (4.74,2);
%%%\begin{scriptsize}
%%%\fill [color=black] (0.8,0.8) circle (1.2pt);
%%%\draw[color=black] (0.68,0.6) node {$p$};
%%%\fill [color=black] (2.76,1.84) circle (1.2pt);
%%%\draw[color=black] (2.06,1.98) node {$C(p,q,r)$};
%%%\fill [color=black] (2.84,3.74) circle (1.2pt);
%%%\draw[color=black] (3.02,3.86) node {$q$};
%%%\fill [color=black] (4.74,2) circle (1.2pt);
%%%\draw[color=black] (5.04,2.04) node {$r$};
%%%\end{scriptsize}
%%%\end{tikzpicture}
%%%\caption{A geodesic triangle in an $\R$-tree.}
%%%\label{figureRtree}
%%%\end{center}
%%%\end{figure}
%%%
%%%%\begin{figure}
%%%%\centerline{\mbox{\includegraphics[scale = 0.7]{Rtree.png}}}
%%%%\caption{A geodesic triangle in an $\R$-tree.}
%%%%\label{figureRtree}
%%%%\end{figure}
%%%
%%\begin{observation}
%%\label{observationRtreesCAT}
%%$\R$-trees are CAT(-1).
%%\end{observation}
%%\begin{proof}
%%First of all, an argument similar to the proof of Observation \ref{observationCATuniquelygeo} shows that $\R$-trees are uniquely geodesic, justifying Figure \ref{figureRtree}. In particular, if $\Delta(p,q,r)$ is a geodesic triangle in an $\R$-tree and if $C = C(p,q,r)$ then $\geo pq = \geo pC\cup\geo qC$, $\geo qr = \geo qC\cup\geo rC$, and $\geo rp = \geo rC\cup\geo pC$. It follows that any two points $x,y\in \Delta$ share a side in common, without loss of generality say $x,y\in\geo pq$. Then
%%\[
%%\dist(x,y) = \dist(p,q) - \dist(x,p) - \dist(y,q) = \dist(\wbar p,\wbar q) - \dist(\wbar x,\wbar p) - \dist(\wbar y,\wbar q) \leq \dist(\wbar x,\wbar y).
%%\]
%%\end{proof}
%%
%%In a sense $\R$-trees are the ``most negatively curved spaces''; although we did not define the notion of a CAT($\kappa$) space, $\R$-trees are CAT($\kappa$) for every $\kappa\in\R$.
%%%
%%%There are several equivalent definitions of $\R$-trees; cf. \cite[p.29\internal]{Chiswell}. One which we will find particularly useful is the following:
%%%
%%%\begin{lemma}
%%%\label{lemmaRtreeequivalent}
%%%Let $X$ be a metric space. The following are equivalent:
%%%\begin{itemize}
%%%\item[(A)] $X$ is an $\R$-tree.
%%%\item[(B)] There exists a collection of geodesics $(\geo xy)_{x,y\in X}$, with the following properties:
%%%\begin{itemize}
%%%\item[(BI)] For each $x,y\in X$, $\geo xy$ is a geodesic between $x$ and $y$.
%%%\item[(BII)] Given $z,w\in \geo xy$, $\geo zw\subset \geo xy$.
%%%\item[(BIII)] Given $x_1,x_2,x_3\in X$ distinct, at least one pair of the geodesics $\geo{x_i}{x_j}$, $i\neq j$, have a nontrivial intersection. More precisely, there exist distinct $i,j,k\in\{1,2,3\}$ such that
%%%\[
%%%\geo{x_i}{x_j} \cap \geo{x_i}{x_k} \propersupset \{x_i\}.
%%%\]
%%%\end{itemize}
%%%\end{itemize}
%%%\end{lemma}
%%%\begin{proof}[Proof of \text{(A)} \implies \text{(B)}]
%%%Note that (BI) and (BII) are true for any uniquely geodesic metric space. Given $x_1,x_2,x_3$ distinct, let $C$ be the center. Then $x_i\neq C$ for some $i$; without loss of generality $x_1\neq C$. Then
%%%\[
%%%\geo{x_1}{x_2} \cap \geo{x_1}{x_3} = \geo{x_1}{C} \propersupset \{x_1\}.
%%%\]
%%%\end{proof}
%%%\begin{proof}[Proof of \text{(B)} \implies \text{(A)}]
%%%We first show that given $x_1,x_2,x_3\in X$, the intersection $\bigcap_{i\neq j} \geo{x_i}{x_j}$ is nonempty. Indeed, suppose not. For $i = 2,3$ let $\gamma_i:[0,\dist(x_1,x_i)]\to X$ be a parameterization of $\geo{x_1}{x_i}$, and let
%%%\[
%%%t_1 = \max\{t\geq 0 : \gamma_2(t) = \gamma_3(t)\}.
%%%\]
%%%By replacing $x$ with $\gamma_2(t_1) = \gamma_3(t_1)$ and using (BII), we may without loss of generality assume that $t_1 = 0$, or equivalently that $\geo{x_1}{x_2}\cap \geo{x_1}{x_3} = \{x_1\}$. Similarly, we may without loss of generality assume that $\geo{x_2}{x_1}\cap \geo{x_2}{x_3} = \{x_2\}$ and $\geo{x_3}{x_1}\cap\geo{x_3}{x_2} = \{x_3\}$. But then (BIII) implies that $x_1,x_2,x_3$ cannot be all distinct. This immediately implies that $\bigcap_{i\neq j} \geo{x_i}{x_j} \neq \emptyset$.
%%%
%%%To complete the proof, we must show that $X$ is uniquely geodesic. Indeed suppose that for some $x_1,x_2\in X$, there is a geodesic $\geo{x_1}{x_2}'$ which is not in the collection. Then there exists $x_3\in\geo{x_1}{x_2}'\butnot\geo{x_1}{x_2}$. By the above paragraph, there exists $w\in\bigcap_{i\neq j} \geo{x_i}{x_j}$. Since $w\in \geo{x_i}{x_3}$, we have
%%%\begin{equation}
%%%\label{dxiw}
%%%\dist(x_i,w) \leq \dist(x_i,x_3).
%%%\end{equation}
%%%On the other hand, since $w\in \geo{x_1}{x_2}$ and $x_3\in\geo{x_1}{x_2}'$, we have
%%%\[
%%%\dist(x_1,x_2) = \dist(x_1,w) + \dist(x_2,w) \leq \dist(x_1,x_3) + \dist(x_2,x_3) = \dist(x_1,x_2).
%%%\]
%%%It follows that equality holds in \eqref{dxiw}, i.e. $\dist(x_i,w) = \dist(x_i,x_3)$. Since $w\in \geo{x_i}{x_3}$, this implies $w = x_3$. But then $x_3 = w\in \geo{x_1}{x_2}$, a contradiction.
%%%\end{proof}
%%%
%%%\begin{definition}
%%%\label{definitionZtree}
%%%A \emph{$\Z$-tree} is an $\R$-tree $X$ together with a set $V\subset X$ (the \emph{vertex set}) with the following properties:
%%%\begin{itemize}
%%%\item[(I)] $\dist(x,y)\in\N$ for all $x,y\in V$
%%%\item[(II)] If $x,y\in V$, $z\in \geo xy$, and $\dist(x,z)\in\N$, then $z\in V$
%%%\item[(III)] $X = \bigcup_{x,y\in V}\geo xy$.
%%%\end{itemize}
%%%\end{definition}
%%%
%%%\begin{remark}
%%%This definition disagrees slightly with other sources (e.g. \cite{Chiswell}) that would consider the set $V$ to be the $\Z$-tree rather than $X$. Hopefully this will not cause too much confusion.
%%%\end{remark}
%%%
%%
%%
%%
%%%%%%%%%%%%%%%%%%%%%%%%%%%%%%%%%%
%%\bigskip
%%\subsection{Gromov hyperbolic metric spaces}
%%%%%%%%%%%%%%%%%%%%%%%%%%%%%%%%%%
%%
%%We now come to the theory of Gromov hyperbolic metric spaces. In a sense, Gromov hyperbolic metric spaces are those which are ``approximately $\R$-trees''. A good reference for this subsection is \cite{Vaisala}.
%%
%%% In this subsection we review the theory of hyperbolic metric spaces and their boundaries. There is a vast literature on this subject \cite{BH, KapovichBenakli, Oshika, Va}. We would like to single out \cite{BridsonHaefliger}, \cite{BonkSchramm}, and \cite{Vaisala} as exceptions that work in not necessarily proper settings and make reference to infinite-dimensional hyperbolic space.
%%
%%For any three numbers $d_{pq},d_{qr},d_{rp}\geq 0$ satisfying the triangle inequality, there exists an $\R$-tree $X$ and three points $p,q,r\in X$ such that $\dist(p,q) = d_{pq}$, etc. Thus in some sense looking at triples ``does not tell you'' that you are looking at an $\R$-tree. Now let us look at quadruples. A quadruple $(p,q,r,s)$ in an $\R$-tree $X$ looks something like Figure \ref{figurequadruple}. Of course, the points $p,q,r,s\in X$ could be arranged in any order. However, let us consider them the way that they are arranged in Figure \ref{figurequadruple} and note that
%%
%%
%%\begin{figure}
%%\begin{center}
%%\begin{tikzpicture}[line cap=round,line join=round,>=triangle 45,x=1.0cm,y=1.0cm]
%%%\clip(0.56,0.9) rectangle (7.44,4.9);
%%\clip(0.56,1.2) rectangle (7.44,4.8);
%%\draw (2.02,3.3)-- (2.8,2.68);
%%\draw (2.8,2.68)-- (1.56,1.46);
%%\draw (2.8,2.68)-- (5.5,2.64);
%%\draw (5.5,2.64)-- (6.58,4.54);
%%\draw (5.5,2.64)-- (6.3,1.46);
%%\begin{scriptsize}
%%\fill [color=black] (2.02,3.3) circle (1.2pt);
%%\draw[color=black] (1.83,3.3) node {$q$};
%%\fill [color=black] (1.56,1.46) circle (1.2pt);
%%\draw[color=black] (1.44,1.3) node {$p$};
%%\fill [color=black] (6.58,4.54) circle (1.2pt);
%%\draw[color=black] (6.61,4.72) node {$r$};
%%\fill [color=black] (6.3,1.46) circle (1.2pt);
%%\draw[color=black] (6.16,1.3) node {$s$};
%%\end{scriptsize}
%%\end{tikzpicture}
%%\caption{A quadruple of points in an $\R$-tree.}
%%\label{figurequadruple}
%%\end{center}
%%\end{figure}
%%
%%%\begin{figure}
%%%\centerline{\mbox{\includegraphics[scale = 0.7]{quadruple.png}}}
%%%\caption{A quadruple of points in an $\R$-tree.}
%%%\label{figurequadruple}
%%%\end{figure}
%%
%%\begin{equation}
%%\label{centersequal}
%%C(p,q,r) = C(p,q,s).
%%\end{equation}
%%In order to write this equality in terms of distances, we need some way of measuring the distance from the vertex of a geodesic triangle to its center.
%%\begin{observation}
%%If $\Delta(p,q,r)$ is a geodesic triangle in an $\R$-tree then $\dist(p,C(p,q,r))$ is equal to
%%\begin{equation}
%%\label{gromovproduct}
%%\lb q|r\rb_p := \frac12[\dist(p,q) + \dist(p,r) - \dist(q,r)].
%%\end{equation}
%%\end{observation}
%%The expression $\lb q|r\rb_p$ is called the \emph{Gromov product} of $q$ and $r$ with respect to $p$, and it makes sense in any metric space. It can be thought of as measuring the ``defect in the triangle inequality''; indeed, the triangle inequality is exactly what assures that $\lb q|r\rb_p\geq 0$ for all $p,q,r\in X$.
%%
%%Now \eqref{centersequal} implies that
%%\[
%%\lb q|r\rb_p = \lb q|s\rb_p \leq \lb r|s\rb_p.
%%\]
%%(The last inequality does not follow from \eqref{centersequal} but it may be seen from Figure \ref{figurequadruple}.) However, since the arrangement of points was arbitrary we do not know which two Gromov products will be equal and which one will be larger. An inequality which captures all possibilities is
%%\begin{equation}
%%\label{gromovinRtrees}
%%\lb q|r\rb_p \geq \min(\lb q|s\rb_p,\lb r|s\rb_p).
%%\end{equation}
%%Now, as mentioned before, we will define hyperbolic metric spaces as those which are ``approximately $\R$-trees''. Thus they will satisfy \eqref{gromovinRtrees} with an asymptotic.
%%\begin{definition}
%%\label{definitiongromovhyperbolic}
%%A metric space $X$ is called \emph{hyperbolic} (or \emph{Gromov hyperbolic}) if for every four points $x,y,z,w\in X$ we have 
%%\begin{equation}
%%\label{gromov}
%%\lb x|z\rb_w \gtrsim_\plus \min(\lb x|y\rb_w,\lb y|z\rb_w),
%%\end{equation}
%%We will refer to \eqref{gromov} as \emph{Gromov's inequality}.
%%\end{definition}
%%From the above discussion, every $\R$-tree is Gromov hyperbolic with an implied constant of $0$ in \eqref{gromov}. (This can also be deduced from Proposition \ref{propositionCATimpliesGromov} below.)
%%
%%Note that many authors require $X$ to be a geodesic metric space in order to be hyperbolic; we do not. If $X$ is a geodesic metric space, then the condition of hyperbolicity can be reformulated in several different ways, including the \emph{thin triangles condition}; for details, see \cite[\6 III.H.1]{BridsonHaefliger} or Subsection \ref{subsectionrips} below.
%%
%%It will be convenient for us to make a list of several identities satisfied by the Gromov product. For each $z\in X$, let $\busemann_z$ denote the \emph{Busemann function}
%%\begin{equation}
%%\label{busemanndef}
%%\busemann_z(x,y) := \dist(z,x) - \dist(z,y).
%%\end{equation}
%%
%%
%%\begin{proposition}
%%\label{propositionbasicidentities}
%%The Gromov product and Busemann function satisfy the following identities and inequalities:
%%\begin{align*}
%%\tag{a} \lb x|y\rb_z &= \lb y|x\rb_z \\
%%\tag{b} \dist(y,z) &= \lb y|x\rb_z + \lb z|x\rb_y\\
%%\tag{c} 0\leq \lb x|y\rb_z &\leq \min(\dist(x,z),\dist(y,z))\\
%%\tag{d} \lb x|y\rb_z &\leq \lb x|y\rb_w + \dist(z,w)\\
%%\tag{e} \lb x|y\rb_w &\leq \lb x|z\rb_w + \dist(y,z)\\
%%\tag{f} |\busemann_x(z,w)| &\leq \dist(z,w)\\
%%\tag{g} \lb x|y\rb_z &= \lb x|y\rb_w + \frac12[\busemann_x(z,w) + \busemann_y(z,w)]\\
%%\tag{h} \lb x|y\rb_z &= \frac12[\dist(x,z) + \busemann_y(z,x)]\\
%%\tag{j} \busemann_x(y,z) &= \lb z|x\rb_y - \lb y|x\rb_z \\
%%\tag{k} \lb x|y\rb_z &= \lb x|y\rb_w + \dist(z,w) - \lb x|z\rb_w - \lb y|z\rb_w\\% (i) omitted because it causes confusion with Roman numerals - (v) and (x) will also be omitted
%%\tag{l} \lb x|y\rb_w &= \lb x|z\rb_w + \frac12[\busemann_w(y,z) - \busemann_x(y,z)]
%%\end{align*}
%%\end{proposition}
%%The proof is a straightforward computation. We remark that (a)-(e) may be found in \cite[Lemma 2.8]{Vaisala}.% We will refer to (f) as the \emph{change of viewpoint formula}.
%%
%%\begin{figure}
%%\begin{center}
%%\begin{tikzpicture}[line cap=round,line join=round,>=triangle 45,x=1.0cm,y=1.0cm]
%%\clip(-1.54,-0.63) rectangle (5.02,3.6);
%%\draw (-0.68,-0.48)-- (1.02,0.66);
%%\draw (1.02,0.66)-- (1.76,1.86);
%%\draw (1.76,1.86)-- (2.06,3.26);
%%\draw (1.76,1.86)-- (4.3,2.42);
%%\draw (1.9621294832739777,3.244449991852601)-- (1.6662459503620535,1.8663069374640353);
%%\draw (1.9621294832739777,3.244449991852601)-- (1.920651867546605,3.2029723761252282);
%%\draw (1.9621294832739777,3.244449991852601)-- (1.9850941143735463,3.194322722245245);
%%\draw (1.6662459503620535,1.8663069374640353)-- (1.645781680327839,1.9358422158191166);
%%\draw (1.6662459503620535,1.8663069374640353)-- (1.712620501296158,1.919132510577037);
%%\draw (1.811554396247592,1.7618958525911042)-- (4.322472096464471,2.330897992557095);
%%\draw (1.811554396247592,1.7618958525911042)-- (1.8414665959029357,1.8362946416296355);
%%\draw (1.811554396247592,1.7618958525911042)-- (1.8757672524481015,1.7219591198124153);
%%\draw (4.322472096464471,2.330897992557095)-- (4.258520422692855,2.3649896261428003);
%%\draw (4.322472096464471,2.330897992557095)-- (4.284871880487056,2.2595837949659954);
%%\begin{scriptsize}
%%\draw [fill=black] (-0.68,-0.48) circle (.75pt);
%%\draw[color=black] (-0.9047548310928106,-0.48743844483680343) node {$x$};
%%\draw [fill=black] (2.06,3.26) circle (.75pt);
%%\draw[color=black] (2.1685467646948235,3.4849673412078808) node {$y$};
%%\draw[color=black] (1.3983054592304987,2.7983786868297873) node {$\lb z|x\rb_y$};
%%\draw [fill=black] (4.3,2.42) circle (.75pt);
%%\draw[color=black] (4.408133565880492,2.6512525466059103) node {$z$};
%%\draw[color=black] (3.083998303865608,1.701190403090175) node {$\lb y|x\rb_z$};
%%\end{scriptsize}
%%\end{tikzpicture}
%%\caption{An illustration of (b) of Proposition \ref{propositionbasicidentities} in an $\R$-tree.}
%%\end{center}
%%\end{figure}
%%
%%\subsubsection{Examples of Gromov hyperbolic metric spaces}
%%
%%\begin{proposition}
%%\label{propositionCATimpliesGromov}
%%Every CAT(-1) space (and in particular every \ROSS) is Gromov hyperbolic. In fact, if $X$ is a CAT(-1) space then for every four points $x,y,z,w\in X$ we have 
%%\begin{equation}
%%\label{gromovforCAT}
%%e^{-\lb x|z\rb_w} \leq e^{-\lb x|y\rb_w} + e^{-\lb y|z\rb_w},
%%\end{equation}
%%and so $X$ satisfies \eqref{gromov} with an implied constant of $\log(2)$.
%%\end{proposition}
%%Proposition \ref{propositionCATimpliesGromov} will be proven below in Subsection \ref{subsectionGromovinH}.
%%\begin{remark}
%%The first assertion of Proposition \ref{propositionCATimpliesGromov}, namely, that CAT(-1) spaces are Gromov hyperbolic, is \cite[Proposition III.H.1.2]{BridsonHaefliger}. The inequality \eqref{gromovforCAT} in the case where $x,y,z\in\del X$ and $w\in X$ can be found in \cite[Th\'eor\`eme 2.5.1]{Bourdon}.
%%\end{remark}
%%
%%\begin{definition}
%%\label{definitionstronglyhyperbolic}
%%A space $X$ satisfying the conclusion of Proposition \ref{propositionCATimpliesGromov} is said to be \emph{strongly hyperbolic}.
%%\end{definition}
%%Note that
%%\[
%%\text{$\R$-tree \implies CAT(-1) \implies Strongly hyperbolic \implies Hyperbolic.}
%%\]
%%
%%A large class of examples of hyperbolic metric spaces which are not CAT(-1) is furnished by the Cayley graphs of finitely presented groups. Indeed, we have the following:
%%
%%\begin{theorem}[{\cite[p.78]{Gromov4}}, \cite{Olshanskii}; see also \cite{Champetier}]
%%\label{theoremolshanskii}
%%Fix $k\geq 2$ and an alphabet $A = \{a_1^{\pm 1},a_2^{\pm 1},\cdots,a_k^{\pm 1}\}$. Fix $i\in\Namer$ and a sequence of positive integers $(n_1,\cdots,n_i)$. Let $N = N(k,i,n_1,\cdots,n_i)$ be the number of group presentations $G = \lb a_1,\cdots, a_k|r_1,\cdots,r_i\rb$ such that $r_1,\cdots,r_i$ are reduced words in the alphabet $A$ such that the length of $r_j$ is $n_j$ for $j = 1,2,\cdots,i$. If $N_h$ is the number of groups in this collection whose Cayley graphs are hyperbolic and if $n = \min(n_1,\cdots,n_i)$ then $\lim_{n\to\infty}N_h/N = 1$.
%%\end{theorem}
%%
%%This theorem says that in some sense, ``almost every'' finitely presented group is hyperbolic.
%%
%%%\begin{remark}
%%%A similar theorem was asserted without proof by M. Gromov in \cite[pp.78-79]{Gromov4}.
%%%\end{remark}
%%
%%If one has a hyperbolic metric space $X$, there are two ways to get another hyperbolic metric space from $X$, one trivial and one nontrivial.
%%
%%\begin{observation}
%%\label{observationsubspace}
%%Any subspace of a hyperbolic metric space is hyperbolic. Any subspace of a strongly hyperbolic metric space is strongly hyperbolic.
%%\end{observation}
%%
%%To describe the other method we need to define the notion of a \emph{quasi-isometric embedding}.
%%\begin{definition}
%%\label{definitionquasiisometry}
%%Let $(X_1,\dist_1)$ and $(X_2,\dist_2)$ be metric spaces. A map $\Phi:X_1\to X_2$ is a \emph{quasi-isometric embedding} if for every $x,y\in X_1$
%%\[
%%\dist_2(\Phi(x),\Phi(y)) \asymp_{\plus,\times} \dist_1(x,y).
%%\]
%%A quasi-isometric embedding $\Phi$ is called a \emph{quasi-isometry} if its image $\Phi(X_1)$ is cobounded in $X_2$, that is, if there exists $R > 0$ such that $\Phi(X_1)$ is $R$-dense in $X_2$, meaning that $\min_{x\in X_2} \dist(x,\Phi(X_1)) \leq R$. In this case, the spaces $X_1$ and $X_2$ are said to be \emph{quasi-isometric}.
%%\end{definition}
%%
%%\begin{theorem}[{\cite[Theorem III.H.1.9]{BridsonHaefliger}}]
%%\label{theoremquasiisometric}
%%Any geodesic metric space which can be quasi-isometrically embedded into a geodesic hyperbolic metric space is also a hyperbolic metric space.
%%\end{theorem}
%%
%%\begin{remark}
%%Theorem \ref{theoremquasiisometric} is not true if the hypothesis of geodesicity is dropped. For example, $\R$ is quasi-isometric to $\Rplus\times\{0\}\cup\{0\}\times\Rplus\subset\R^2$, but the former is hyperbolic and the latter is not.
%%\end{remark}
%%
%%There are many more examples of hyperbolic metric spaces which we will not discuss; cf. the list in \6\ref{subsubsectionlist}.
%%
%%
%%\bigskip
%%\subsection{The boundary of a hyperbolic metric space}
%%\label{subsectionboundary}
%%
%%In this subsection we define the Gromov boundary of a hyperbolic metric space $X$. The construction will depend on a distinguished point $\zero\in X$, but the resulting space will be independent of which point is chosen. If $X$ is an $\R$-tree, then the boundary of $X$ will turn out to be the set of infinite branches through $X$, i.e. the set of all isometric embeddings $\pi:\Rplus\to X$ sending $0$ to $\zero$, where $\zero\in X$ is a distinguished fixed point. If $X$ is a \CROSS, then the boundary of $X$ will turn out to be isomorphic to the space $\del X$ defined in Section \ref{sectionROSSONCTs}.
%%
%%To motivate the definition of the boundary, suppose that $X$ is an $\R$-tree. An infinite branch through $X$ can be approximated by finite branches which agree on longer and longer segments. Suppose that $(\geo\zero{x_n})_1^\infty$ is a sequence of geodesic segments. For each $n,m\in\Namer$, the length of the intersection of $\geo\zero{x_n}$ and $\geo\zero{x_m}$ is equal to $\dist(\zero,C(\zero,x_n,x_m))$, which in turn is equal to $\lb x_n|x_m\rb_\zero$. Thus, the sequence $(\geo\zero{x_n})_1^\infty$ converges to an infinite geodesic if and only if
%%\begin{equation}
%%\label{gromovsequence}
%%\lb x_n|x_m\rb_\zero \tendsto{n,m} \infty.
%%\end{equation}
%%(Cf. Figure \ref{figuregromovseq}.) The formula \eqref{gromovsequence} is reminiscent of the definition of a Cauchy sequence. This intuition will be made explicit in Subsection \ref{subsectionmetametrics}, where we will introduce a \emph{metametric} on $X$ with the property that a sequence in $X$ satisfies \eqref{gromovsequence} if and only if it is Cauchy with respect to this metametric.
%%%\ignore{
%%%To motivate the definition of the boundary, consider the inequality \eqref{gromovforCAT}. It indicates that the expression
%%%\[
%%%\Dist_w(x,y) := b^{-\lb x|y\rb_w}
%%%\]
%%%satisfies a triangle inequality. Note, however, that it is not a metric since $\Dist_w(x,x) > 0$. Now recall that the completion of a metric space can be defined as the set of Cauchy sequences modulo their equivalence. We can do the same thing with our near-metric $\Dist_w$ (actually called a \emph{metametric}; see Subection \internalmissingref\ below).
%%%}% end ignore
%%
%%\begin{figure}
%%\begin{center}
%%\begin{tikzpicture}[line cap=round,line join=round,>=triangle 45,x=1.0cm,y=1.0cm]
%%\clip(-2.74,-1.0) rectangle (7.98,3.16);
%%\draw (-1.8,-0.78)-- (0.52,0.72);
%%\draw (0.52,0.72)-- (0.7,2.3);
%%\draw (0.52,0.72)-- (2.28,0.52);
%%\draw (2.28,0.52)-- (2.96,-0.64);
%%\draw (2.28,0.52)-- (3.46,1.62);
%%\draw (3.46,1.62)-- (3.7,2.48);
%%\draw (3.46,1.62)-- (4.78,2.16);
%%\draw (4.78,2.16)-- (5.38,2.92);
%%\draw (4.78,2.16)-- (6.8,2.36);
%%\begin{scriptsize}
%%\fill [color=black] (-1.8,-0.78) circle (1pt);
%%\draw[color=black] (-2,-0.86) node {$\zero$};
%%\fill [color=black] (0.7,2.3) circle (1pt);
%%\draw[color=black] (0.9,2.48) node {$x_1$};
%%\fill [color=black] (2.96,-0.64) circle (1pt);
%%\draw[color=black] (3.32,-0.72) node {$x_2$};
%%\fill [color=black] (3.7,2.48) circle (1pt);
%%\draw[color=black] (3.96,2.74) node {$x_3$};
%%\fill [color=black] (5.38,2.92) circle (1pt);
%%\draw[color=black] (5.74,3.08) node {$x_5$};
%%\fill [color=black] (6.8,2.36) circle (1pt);
%%\draw[color=black] (7.36,2.4) node {$\ldots$};
%%\end{scriptsize}
%%\end{tikzpicture}
%%\caption{A Gromov sequence in an $\R$-tree.}
%%\label{figuregromovseq}
%%\end{center}
%%\end{figure}
%%
%%
%%\begin{definition}
%%\label{definitiongromovsequence}
%%A sequence $(x_n)_1^\infty$ in $X$ for which \eqref{gromovsequence} holds is called a \emph{Gromov sequence}. Two Gromov sequences $(x_n)_1^\infty$ and $(y_n)_1^\infty$ are called \emph{equivalent} if
%%\[
%%\lb x_n|y_n\rb_\zero \tendsto n \infty,
%%\]
%%or equivalently if
%%\[
%%\lb x_n|y_m\rb_\zero \tendsto{n,m} \infty.
%%\]
%%In this case, we write $(x_n)_1^\infty \sim (y_n)_1^\infty$. It is readily verified using Gromov's inequality that $\sim$ is an equivalence relation on the set of Gromov sequences in $X$. We will denote the class of sequences equivalent to a given sequence $(x_n)_1^\infty$ by $[(x_n)_1^\infty]$.
%%\end{definition}
%%
%%\begin{definition}
%%\label{definitiongromovboundary}
%%The \emph{Gromov boundary} of $X$ is the set of Gromov sequences modulo equivalence. It is denoted $\del X$. The \emph{Gromov closure} or \emph{bordification} of $X$ is the disjoint union $\bord X := X\cup\del X$. 
%%\end{definition}
%%
%%\begin{remark}
%%\label{remarkambiguityboundary}
%%If $X$ is a \CROSS, then this notation causes some ambiguity, since it is not clear whether $\del X$ represents the Gromov boundary of $X$, or rather the topological boundary of $X$ as in Section \ref{sectionROSSONCTs}. This ambiguity will be resolved in \6\ref{subsubsectionboundariesequivalent} below when it is shown that the two bordifications are isomorphic.
%%\end{remark}
%%
%%\begin{remark}
%%\label{remarkambiguityboundary2}
%%In the literature, the ideal boundary of a hyperbolic metric space is often taken to be the set of equivalence classes of geodesic rays under asymptotic equivalence, rather than the set of equivalence classes of Gromov sequences (e.g. \cite[p.427]{BridsonHaefliger}). If $X$ is proper and geodesic, then these two notions are equivalent \cite[Lemma III.H.3.13]{BridsonHaefliger}, but in general they may be different.
%%\end{remark}
%%
%%\begin{remark}
%%\label{remarkbasepoint1}
%%By (d) of Proposition \ref{propositionbasicidentities}, the concepts of Gromov sequence and equivalence do not depend on the basepoint $\zero$. In particular, the Gromov boundary $\del X$ is independent of $\zero$.
%%\end{remark}
%%
%%\subsubsection{Extending the Gromov product and the Busemann function to the boundary}
%%We now wish to extend the Gromov product and Busemann function to the boundary ``by continuity''. Fix $\xi,\eta\in\del X$ and $z\in X$. Ideally, we would like to define $\lb \xi|\eta\rb_z$ to be
%%\begin{equation}
%%\label{gromovproductdef}
%%\lim_{n,m\to\infty}\lb x_n|y_m\rb_z ,
%%\end{equation}
%%where $(x_n)_1^\infty\in\xi$ and $(y_m)_1^\infty\in\eta$. (The definition would then have to be shown independent of which sequences were chosen.) The naive definition \eqref{gromovproductdef} does not work, because the limit \eqref{gromovproductdef} does not necessarily exist:
%%\begin{example}
%%\label{examplegromovnotcont}
%%Let
%%\[
%%X = \{\xx\in\R^2: x_2\in[0,1]\}
%%\]
%%be interpreted as a subspace of $\R^2$ with the $L^1$ metric. Then $X$ is a hyperbolic metric space, since it contains the cobounded hyperbolic metric space $\R\times\{0\}$. Its Gromov boundary consists of two points $-\infty$ and $+\infty$, which are the limits of $\xx$ as $x_1$ approaches $-\infty$ or $+\infty$, respectively. Let $\yy = (0,1)$ and $\zz = (1,0)$. Then for all $\xx\in X$, $\lb \xx|\yy\rb_\zz = x_2$. In particular, we can find a sequence $\xx_n\to +\infty$ such that $\lim_{n\to\infty} \lb \xx_n|\yy\rb_\zz$ does not exist.
%%\end{example}
%%
%%Fortunately, the limit \eqref{gromovproductdef} ``exists up to a constant'':
%%\begin{lemma}
%%\label{lemmawelldefined}
%%Let $(x_n)_1^\infty$ and $(y_m)_1^\infty$ be Gromov sequences, and fix $y,z\in X$. Then
%%\begin{align} \label{welldefined1}
%%\liminf_{n,m\to\infty}\lb x_n|y_m\rb_z &\asymp_\plus \limsup_{n,m\to\infty}\lb x_n|y_m\rb_z\\ \label{welldefined2}
%%\liminf_{n\to\infty}\lb x_n|y\rb_z &\asymp_\plus \limsup_{n\to\infty}\lb x_n|y\rb_z ,
%%\end{align}
%%with equality if $X$ is strongly hyperbolic.
%%\end{lemma}
%%Note that except for the statement about strongly hyperbolic spaces, this lemma is simply \cite[Lemma 5.6]{Vaisala}.
%%\begin{proof}[Proof of Lemma \ref{lemmawelldefined}]
%%Fix $n_1,n_2,m_1,m_2\in\Namer$. By Gromov's inequality
%%\[
%%\lb x_{n_1}|y_{m_1}\rb_z \gtrsim_\plus \min(\lb x_{n_2}|y_{m_2}\rb_z,\lb x_{n_1}|x_{n_2}\rb_z,\lb y_{m_1}|y_{m_2}\rb_z).
%%\]
%%Taking the liminf over $n_1,m_1$ and the limsup over $n_2,m_2$ gives
%%\begin{align*}
%%\liminf_{n,m\to\infty}\lb x_n|y_m \rb_z &\gtrsim_\plus \min\left(\limsup_{n,m\to\infty}\lb x_n|y_m \rb_z,\liminf_{n_1,n_2\to\infty}\lb x_{n_1}|x_{n_2}\rb_z,\liminf_{m_1,m_2\to\infty}\lb y_{m_1}|y_{m_2}\rb_z\right) \noreason\\
%%&=_\pt \limsup_{n,m\to\infty}\lb x_n|y_m\rb_z, \since{$(x_n)_1^\infty$ and $(y_m)_1^\infty$ are Gromov}
%%\end{align*}
%%demonstrating \eqref{welldefined1}. On the other hand, suppose that $X$ is strongly hyperbolic. Then by \eqref{gromovforCAT} we have
%%\[
%%\exp\big(-\lb x_{n_1}|y_{m_1}\rb_z\big) \leq \exp\big(-\lb x_{n_2}|y_{m_2}\rb_z\big) + \exp\big(-\lb x_{n_1}|x_{n_2}\rb_z\big) + \exp\big(-\lb y_{m_1}|y_{m_2}\rb_z\big);
%%\]
%%taking the limsup over $n_1,m_1$ and the liminf over $n_2,m_2$ gives
%%\begin{align*}
%%\exp&\big(-\liminf_{n,m\to\infty}\lb x_n|y_m\rb_z\big)\\
%%&\leq \exp\big(-\limsup_{n,m\to\infty}\lb x_n|y_m \rb_z\big) + \exp\big(-\liminf_{n_1,n_2\to\infty}\lb x_{n_1}|x_{n_2}\rb_z\big) + \exp\big(-\liminf_{m_1,m_2\to\infty}\lb y_{m_1}|y_{m_2}\rb_z\big) \noreason\\
%%&= \exp\big(-\limsup_{n,m\to\infty}\lb x_n|y_m\rb_z\big), \since{$(x_n)_1^\infty$ and $(y_m)_1^\infty$ are Gromov}
%%\end{align*}
%%demonstrating equality in \eqref{welldefined1}. The proof of \eqref{welldefined2} is similar and will be omitted.
%%\end{proof}
%%
%%\begin{remark}
%%Many of the statements in this monograph concerning strongly hyperbolic metric spaces are in fact valid for all hyperbolic metric spaces satisfying the conclusion of Lemma \ref{lemmawelldefined}.
%%\end{remark}
%%
%%Now that we know that it does not matter too much whether we replace the limit in \eqref{gromovproductdef} by a liminf or a limsup, we make the following definition without fear:
%%
%%\begin{definition}
%%For $\xi,\eta\in\del X$ and $y,z\in X$, let
%%\begin{align} \label{gromovboundarydef1}
%%\lb \xi|\eta\rb_z &:= \inf\left\{\liminf_{n,m\to\infty}\lb x_n|y_m\rb_z:(x_n)_1^\infty\in\xi,(y_m)_1^\infty\in\eta\right\}\\ \label{gromovboundarydef2}
%%\lb \xi|y\rb_z &:= \lb y|\xi\rb_z := \inf\left\{\liminf_{n\to\infty}\lb x_n|y\rb_z:(x_n)_1^\infty\in\xi\right\}\\ \label{gromovboundarydef3}
%%\busemann_\xi(y,z) &= \lb z|\xi\rb_y - \lb y|\xi\rb_z.
%%\end{align}
%%\end{definition}
%%
%%As a corollary of Lemma \ref{lemmawelldefined}, we have the following:
%%
%%\begin{lemma}
%%\label{lemmaboundaryasymptotic}
%%Fix $\xi,\eta\in\del X$ and $y,z\in X$. For all $(x_n)_1^\infty\in\xi$ and $(y_m)_1^\infty\in\eta$ we have 
%%\begin{align} \label{boundaryasymptotic1}
%%\lb x_n|y_m\rb_z &\tendsto{n,m,\plus} \lb \xi|\eta\rb_z\\ \label{boundaryasymptotic2}
%%\lb x_n|y\rb_z &\tendsto{\;\; n,\plus \;\;} \lb \xi|y\rb_z\\ \label{boundaryasymptotic3}
%%\busemann_{x_n}(y,z) &\tendsto{\;\; n,\plus \;\;} \busemann_\xi(y,z),
%%\end{align}
%%(cf. Convention \ref{conventiontendston}), with exact limits if $X$ is strongly hyperbolic.
%%\end{lemma}
%%Note that except for the statement about strongly hyperbolic spaces, this lemma is simply \cite[Lemma 5.11]{Vaisala}.
%%\begin{proof}[Proof of Lemma \ref{lemmaboundaryasymptotic}]
%%Suppose that for each $i = 1,2$, we are given $(x_n^{(i)})_1^\infty\in \xi$ and $(y_m^{(i)})_1^\infty\in\eta$. Let
%%\[
%%x_n = \begin{cases}
%%x_{n/2}^{(1)} & n \text{ even}\\
%%x_{(n + 1)/2}^{(2)} & n \text{ odd}
%%\end{cases},
%%\]
%%and define $y_m$ similarly. It may be verified using Gromov's inequality that $(x_n)_1^\infty\in \xi$ and $(y_m)_1^\infty\in\eta$. Applying Lemma \ref{lemmawelldefined}, we have
%%\[
%%\min_{i = 1}^2 \min_{j = 1}^2 \liminf_{n,m\to\infty} \lb x_n^{(i)}|y_m^{(j)}\rb_z \asymp_\plus \max_{i = 1}^2 \max_{j = 1}^2 \limsup_{n,m\to\infty} \lb x_n^{(i)}|y_m^{(j)}\rb_z .
%%\]
%%In particular,
%%\begin{align*}
%%\liminf_{n,m\to\infty} \lb x_n^{(1)} | y_n^{(1)} \rb_z
%%&\lesssim_\plus \limsup_{n,m\to\infty} \lb x_n^{(1)} | y_n^{(1)} \rb_z\\
%%&\lesssim_\plus \liminf_{n,m\to\infty} \lb x_n^{(2)} | y_n^{(2)} \rb_z\\
%%&\lesssim_\plus \limsup_{n,m\to\infty} \lb x_n^{(2)} | y_n^{(2)} \rb_z
%%\lesssim_\plus \liminf_{n,m\to\infty} \lb x_n^{(1)} | y_n^{(1)} \rb_z.
%%\end{align*}
%%Taking the infimum over all $(x_n^{(2)})_1^\infty\in\xi$ and $(y_m^{(2)})_1^\infty\in\eta$ gives \eqref{boundaryasymptotic1}. A similar argument gives \eqref{boundaryasymptotic2}. Finally, \eqref{boundaryasymptotic3} follows from \eqref{boundaryasymptotic2}, \eqref{gromovboundarydef3}, and (j) of Proposition \ref{propositionbasicidentities}.
%%
%%If $X$ is strongly hyperbolic, then all error terms are equal to zero, demonstrating that the limits converge exactly.
%%\end{proof}
%%\begin{remark}
%%In the sequel, the statement that ``if $X$ is strongly hyperbolic, then all error terms are zero'' will typically be omitted from our proofs.
%%\end{remark}
%%
%%A simple but useful consequence of Lemma \ref{lemmaboundaryasymptotic} is the following:
%%
%%\begin{corollary}
%%\label{corollaryboundaryasymptotic}
%%The formulas of Proposition \ref{propositionbasicidentities} together with Gromov's inequality hold for points on the boundary as well, if the equations and inequalities there are replaced by additive asymptotics.
%% If $X$ is strongly hyperbolic, then we may keep the original formulas without adding an error term.
%%\end{corollary}
%%\begin{proof}
%%For each identity, choose a Gromov sequence representing each element of the boundary which appears in the formula. Replace each occurrence of this element in the formula by the general term of the chosen sequence. This yields a sequence of formulas, each of which is known to be true. Take a subsequence on which each term in these formulas converges. Taking the limit along this subsequence again yields a true formula, and by Lemma \ref{lemmaboundaryasymptotic} we may replace each limit term by the term which it stood for, with only bounded error in doing so, and no error if $X$ is strongly hyperbolic. Thus the formula holds as an additive asymptotic, and holds exactly if $X$ is strongly hyperbolic.
%%\end{proof}
%%
%%\begin{remark}
%%\label{remarkacde}
%%In fact, (a), (c), (d), and (e) of Proposition \ref{propositionbasicidentities} hold in $\bord X$ in the usual sense, i.e. as exact formulas without additive constants. 
%%\end{remark}
%%\begin{proof}
%%These are the identities where there is at most one Gromov product on each side of the formula. For each element of the boundary, we may simply replace each occurence of that element with the general term of an arbitrary Gromov sequence, take the liminf, and then take the infimum over all Gromov sequences.
%%\end{proof}
%%
%%\begin{observation}
%%\label{observationyzinfinity}
%%$\lb x|y\rb_z = \infty$ if and only if $x = y\in\del X$.
%%\end{observation}
%%\begin{proof}
%%This follows directly from \eqref{gromovboundarydef1} and \eqref{gromovboundarydef2}.
%%\end{proof}
%%
%%\subsubsection{A topology on $\bord X$}
%%\label{subsubsectiontopologyclX}
%%One can endow the bordification $\bord X = X\cup\del X$ with a topological structure $\scrT$ as follows: Given $S\subset\bord X$, write $S\in\scrT$ (i.e. call $S$ open) if
%%\begin{itemize}
%%\item[(I)] $S\cap X$ is open, and
%%\item[(II)] For each $\xi\in S\cap\del X$ there exists $t\geq 0$ such that $N_t(\xi)\subset S$, where
%%\[
%%N_t(\xi) := N_{t,\zero}(\xi) := \{y\in \bord X:\lb y|\xi\rb_\zero > t\}
%%\]
%%\end{itemize}
%%
%%\begin{remark}
%%The topology $\scrT$ may equivalently be defined to be the unique topology on $\bord X$ satisfying:
%%\begin{itemize}
%%\item[(I)] $\scrT\given X$ is compatible with the metric $\dist$, and
%%\item[(II)] For each $\xi\in\del X$, the collection
%%\begin{equation}
%%\label{Ntxibase}
%%\{N_t(\xi) : t \geq 0\}
%%\end{equation}
%%is a neighborhood base for $\scrT$ at $\xi$.
%%\end{itemize}
%%\end{remark}
%%
%%\begin{remark}
%%\label{remarkNtxiopen}
%%It follows from Lemma \ref{lemmalowersemicontinuous} below that the sets $N_t(\xi)$ are open in the topology $\scrT$.
%%\end{remark}
%%
%%\begin{remark}
%%\label{remarkbasepoint2}
%%By (d) of Proposition \ref{propositionbasicidentities} (cf. Remark \ref{remarkacde}), we have $N_{t,x}(\xi) \supset N_{t + \dist(x,y),y}(\xi)$ for all $x,y\in X$, $\xi\in\del X$, and $t\geq 0$. Thus the topology $\scrT$ is independent of the basepoint $\zero$.
%%\end{remark}
%%
%%The topology $\scrT$ is quite nice. In fact, we have the following:
%%
%%\begin{proposition}\label{propositionclX}
%%The topological space $(\bord X,\scrT)$ is completely metrizable. If $X$ is proper and geodesic, then $\bord X$ (and thus also $\del X$) is compact. If $X$ is separable, then $\bord X$ (and thus also $\del X$) is separable.
%%\end{proposition}
%%\begin{remark}
%%If $X$ is proper and geodesic, then Proposition \ref{propositionclX} is \cite[Exercise III.H.3.18(4)]{BridsonHaefliger}.
%%\end{remark}
%%\begin{proof}[Proof of Proposition \ref{propositionclX}]
%%We delay the proof of the complete metrizability of $\bord X$ until Subsection \ref{subsectionmetametrics}, where we will introduce a class of compatible complete metrics on $\bord X$ which are important from a geometric point of view, the so-called \emph{visual} metrics.
%%
%%Since $X$ is dense in $\bord X$, the separability of $X$ implies the separability of $\bord X$. Moreover, since $\bord X$ is metrizable (as we will show in Subsection \ref{subsectionmetametrics}), the separability of $\bord X$ implies the separability of $\del X$.
%%
%%Finally, assume that $X$ is proper and geodesic; we claim that $\bord X$ is compact. Let $(x_n)_1^\infty$ be a sequence in $X$. If $(x_n)_1^\infty$ contains a bounded subsequence, then since $X$ is proper it contains a convergent subsequence. Thus we assume that $(x_n)_1^\infty$ contains no bounded subsequence, i.e. $\dox{x_n}\to \infty$.
%%
%%For each $n\in\Namer$ and $t\geq 0$ let
%%\[
%%x_{n,t} = \geo\zero{x_n}_{t\wedge\dox{x_n}},\Footnote{Here and from now on $A\wedge B = \min(A,B)$ and $A\vee B = \max(A,B)$.}
%%\]
%%where $\geo\zero{x_n}$ is any geodesic connecting $\zero$ and $x_n$. Since $X$ is proper, there exists a sequence $(n_k)_1^\infty$ such that for each $t\geq 0$, the sequence $(x_{n_k,t})_1^\infty$ is convergent, say
%%\[
%%x_{n_k,t}\tendsto k x_t.
%%\]
%%It is readily verified that the map $t\mapsto x_t$ is an isometric embedding from $\Rplus$ to $X$. Thus there exists a point $\xi\in\del X$ such that $x_t\to \xi$. We claim that $x_{n_k}\to \xi$. Indeed, for each $t\geq 0$,
%%\[
%%\limsup_{k\to\infty}\Dist(x_{n_k},x_{n_k,t}) \asymp_\times \limsup_{k\to\infty}\Dist(x_{n_k,t},x_t) \asymp_\times \limsup_{k\to\infty} \Dist(x_t,\xi) \asymp_\times b^{-t},
%%\]
%%and so the triangle inequality gives
%%\[
%%\limsup_{k\to\infty}\Dist(x_{n_k},\xi) \lesssim_\times b^{-t}.
%%\]
%%Letting $t\to\infty$ shows that $x_{n_k}\to\xi$.
%%\end{proof}
%%
%%\begin{observation}
%%\label{observationboundaryconvergence}
%%A sequence $(x_n)_1^\infty$ in $\bord X$ converges to a point $\xi\in \del X$ if and only if
%%\begin{equation}\label{boundaryconvergence}
%%\lb x_n|\xi\rb_\zero \tendsto n \infty.
%%\end{equation}
%%\end{observation}
%%
%%\begin{observation}
%%\label{observationboundaryconvergence2}
%%A sequence $(x_n)_1^\infty$ in $X$ converges to a point $\xi\in \del X$ if and only if $(x_n)_1^\infty$ is a Gromov sequence and $(x_n)_1^\infty\in\xi$.
%%\end{observation}
%%
%%% Not sure what this corollary is saying nor how it is used later.
%%%\ignore{
%%%\begin{corollary}
%%%\label{corollarygromovboundarydef34}
%%%We may rewrite \eqref{gromovboundarydef1} and \eqref{gromovboundarydef2} as
%%%\begin{align} \label{gromovboundarydef1B}
%%%\lb \xi|\eta\rb_z &= \liminf_{x\to\xi,y\to\eta}\lb x|y\rb_z\\ \label{gromovboundarydef2B}
%%%\lb \xi|y\rb_z &:= \lb y|\xi\rb_z = \liminf_{x\to\xi}\lb x|y\rb_z
%%%\end{align}
%%%\end{corollary}
%%%}% end ignore
%%
%%
%%We now investigate the continuity properties of the Gromov product and Busemann function.
%%
%%\begin{lemma}[Near-continuity of the Gromov product and Busemann function]
%%\label{lemmanearcontinuity}
%%The functions $(x,y,z)\mapsto \lb x|y\rb_z$ and $(x,z,w)\mapsto \busemann_x(z,w)$ are \underline{nearly continuous} in the following sense: Suppose that $(x_n)_1^\infty$ and $(y_n)_1^\infty$ are sequences in $\bord X$ which converge to points $x_n\to x\in\bord X$ and $y_n\to y\in\bord X$. Suppose that $(z_n)_1^\infty$ and $(w_n)_1^\infty$ are sequences in $X$ which converge to points $z_n\to z\in X$ and $w_n\to w\in X$. Then
%%\begin{align} \label{continuousextension1}
%%\lb x_n|y_n\rb_{z_n} &\tendsto{n,\plus} \lb x|y\rb_z\\ \label{continuousextension2}
%%\busemann_{x_n}(z_n,w_n) &\tendsto{n,\plus} \busemann_x(z,w),
%%\end{align}
%%with $\tendsto n$ if $X$ is strongly hyperbolic.
%%\end{lemma}
%%\begin{proof}
%%In the proof of \eqref{continuousextension1}, there are three cases:
%%\begin{itemize}
%%\item[Case 1:] $x,y\in X$. In this case, \eqref{continuousextension1} follows directly from (d) and (e) of Proposition \ref{propositionbasicidentities}.
%%\item[Case 2:] $x,y\in\del X$. In this case, for each $n\in\Namer$, choose $\what x_n\in X$ such that either
%%\begin{itemize}
%%\item[(1)] $\what x_n = x_n$ (if $x_n\in X$), or
%%\item[(2)] $\lb \what x_n|x_n\rb_z\geq n$ (if $x_n\in\del X$).
%%\end{itemize}
%%Choose $\what y_n$ similarly. Clearly, $\what x_n\to x$ and $\what y_n \to y$. By Observation \ref{observationboundaryconvergence2}, $(\what x_n)_1^\infty\in x$ and $(\what y_n)_1^\infty\in y$. Thus by Lemma \ref{lemmaboundaryasymptotic},
%%\begin{equation}
%%\label{whatxwhaty}
%%\lb \what x_n|\what y_n\rb_z \tendsto{n,\plus} \lb x|y\rb_z.
%%\end{equation}
%%Now by Gromov's inequality and (e) of Proposition \ref{propositionbasicidentities}, either
%%\begin{align*} \tag{1}
%%\lb \what x_n|\what y_n\rb_z &\asymp_\plus \lb x_n|y_n\rb_{z_n} \text{ or }\\ \tag{2}
%%\lb \what x_n|\what y_n\rb_z &\gtrsim_\plus n,
%%\end{align*}
%%with which asymptotic is true depending on $n$. But for $n$ sufficiently large, \eqref{whatxwhaty} ensures that the (2) fails, so (1) holds.
%%\item[Case 3:] $x\in X$, $y\in\del X$, or vice-versa. In this case, a straightforward combination of the above arguments demonstrates \eqref{continuousextension1}.
%%\end{itemize}
%%Finally, note that \eqref{continuousextension2} is an immediate consequence of \eqref{continuousextension1}, \eqref{gromovboundarydef3}, and (j) of Proposition \ref{propositionbasicidentities}.
%%\end{proof}
%%
%%Although Lemma \ref{lemmanearcontinuity} is generally sufficient for applications, we include the following lemma which reassures us that the Gromov product does behave somewhat regularly even on an ``exact'' level.
%%
%%\begin{lemma}
%%\label{lemmalowersemicontinuous}
%%The function $(x,y,z)\mapsto\lb x|y\rb_z$ is lower semicontinuous on $\bord X\times\bord X\times X$.
%%\end{lemma}
%%\begin{proof}
%%Since $\bord X$ is metrizable, it is enough to show that if $x_n\to x$, $y_n\to y$, and $z_n\to z$, then
%%\[
%%\liminf_{n\to\infty}\lb x_n|y_n\rb_{z_n} \geq \lb x|y\rb_z.
%%\]
%%Now fix $\epsilon > 0$.
%%\begin{claim}
%%For each $n\in\N$, there exist points $\what x_n,\what y_n\in X$ satisfying:
%%\begin{align} \label{LSC1}
%%\lb \what x_n|\what y_n\rb_{z_n} &\leq \lb x_n|y_n\rb_{z_n} + \epsilon,\\ \label{LSC2}
%%\lb \what x_n|x_n\rb_\zero &\geq n, \text{ or }\what x_n = x_n\in X,\\ \label{LSC3}
%%\lb \what y_n|y_n\rb_\zero &\geq n, \text{ or }\what y_n = y_n\in X.
%%\end{align}
%%\end{claim}
%%\begin{subproof}
%%Suppose first that $x_n,y_n\in\del X$. By the definition of $\lb x_n|y_n\rb_{z_n}$, there exist $(x_{n,k})_1^\infty\in x_n$ and $(y_{n,\ell})_1^\infty\in y_n$ such that
%%\[
%%\liminf_{k,\ell\to\infty} \lb x_{n,k}|y_{n,\ell}\rb_{z_n} \leq \lb x_n|y_n\rb_{z_n} + \epsilon/2.
%%\]
%%It follows that there exist arbitrarily large\Footnote{Here, of course, ``arbitrarily large'' means that $\min(k,\ell)$ can be made arbitrarily large.} pairs $(k,\ell)\in\N^2$ such that the points $\what x_n := x_{n,k}$ and $\what y_n := y_{n,\ell}$ satisfy \eqref{LSC1}. Since \eqref{LSC2} and \eqref{LSC3} are satisfied for all sufficiently large $(k,\ell)\in\N^2$, this completes the proof. Finally, if either $x_n\in X$, $y_n\in X$, or both, a straightforward adaptation of the above argument yields the claim.
%%\end{subproof}
%%The equations \eqref{LSC2} and \eqref{LSC3}, together with Gromov's inequality, imply that $\what x_n\to x$ and $\what y_n\to y$. Now suppose that $x,y\in \del X$. Then by Observation \ref{observationboundaryconvergence2}, $(\what x_n)_1^\infty\in x$ and $(\what y_n)_1^\infty\in y$. So by the definition of $\lb x|y\rb_z$, we have
%%\begin{align*}
%%\lb x|y\rb_z &\leq \liminf_{n\to\infty}\lb \what x_n|\what y_n\rb_z \by{the definition of $\lb x|y\rb_z$}\\
%%&= \liminf_{n\to\infty}\lb \what x_n|\what y_n\rb_{z_n} \by{(d) of Proposition \ref{propositionbasicidentities}}\\
%%&\leq \liminf_{n\to\infty}\lb x_n|y_n\rb_{z_n} + \epsilon. \by{\eqref{LSC1}}
%%\end{align*}
%%Letting $\epsilon$ tend to zero completes the proof. A similar argument applies to the case where $x\in X$, $y\in X$, or both.
%%\end{proof}
%%
%%
%%%\begin{proof}[Proof of Remark \ref{remarkNtxiopen}]
%%%The fact that $N_t(\xi)\cap X$ is open follows from (g) of Proposition \ref{propositionbasicidentities} - recall that we do not have to add an error term when moving this formula to the boundary (Remark \ref{remarkace}). Fix $\eta\in %N_t(\xi)\cap\del X$. Then $\lb\xi|\eta\rb_\zero > t$; fix $t < \w t < \lb\xi|\eta\rb_\zero$.
%%%\begin{claim}
%%%\label{claimNtxiopen}
%%%There exists $n\in\N$ such that $\lb y|\xi\rb_\zero\geq \w t$ for all $y\in N_n(\eta)\cap X$.
%%%\end{claim}
%%%\begin{proof}
%%%Suppose not. Then we have a sequence $y_n\to \eta$ so that $\lb y_n|\xi\rb_\zero < \w t$ for all $n$. For each $n\in\N$, by \eqref{gromovboundarydef2B}, there exists $x_n\in X$ such that $\lb x_n|y_n\rb_\zero < \w t$ and $\lb %x_n|\xi\rb_\zero > n$. Then $x_n\to \xi$. Thus by \eqref{gromovboundarydef1B} we have
%%%\[
%%%\w t < \lb\xi|\eta\rb_\zero \leq \liminf_{m,n\to\infty}\lb x_m|y_n\rb_\zero
%%%\leq \liminf_{n\to\infty}\lb x_n|y_n\rb_\zero \leq \w t,
%%%\]
%%%a contradiction.
%%%\QEDmod\end{proof}
%%%Let $n$ be as in Claim \ref{claimNtxiopen}, and let
%%%\[
%%%U = \bord X \butnot \cl{X\butnot N_n(\eta)}.
%%%\]
%%%Clearly, $U$ is open and $\eta\in U$. Fix $z\in U$. If $z\in X$, then Claim \ref{claimNtxiopen} implies that $\lb z|\xi\rb_\zero \geq \w t > t$, so $z\in N_t(\xi)$. If $z\in\del X$, then
%%%\[
%%%\lb z|\xi\rb_\zero = \liminf_{x\to\xi,y\to z}\lb x|y\rb_\zero \geq \inf_{y\in N_n(\eta)}\liminf_{x\to\xi}\lb x|y\rb_\zero = \inf_{y\in N_n(\eta)}\lb y|\xi\rb_\zero \geq \w t > t,
%%%\]
%%%so again $z\in N_t(\xi)$. Thus $U\subset N_t(\xi)$, and $N_t(\xi)$ contains a neighborhood of $\eta$.
%%%\end{proof}
%%
%%
%%\begin{lemma}
%%\label{lemmaisometryextension}
%%If $g$ is an isometry of $X$, then it extends in a unique way to a continuous map $\w g:\bord X\to\bord X$.
%%\end{lemma}
%%\begin{proof}
%%This follows more or less directly from Remarks \ref{remarkbasepoint1} and \ref{remarkbasepoint2}; details are left to the reader.
%%\end{proof}
%%%\ignore{
%%%\begin{proof}
%%%It is clear that $g$ sends Gromov sequences to Gromov sequences, and equivalent Gromov sequences to equivalent Gromov sequences. Thus the formula
%%%\[
%%%\del g([(x_n)_1^\infty]) = [(g(x_n))_1^\infty]
%%%\]
%%%defines a map $\del g:\del X\to \del X$. Let $\w g:\bord X\to\bord X$ be the combination of $g$ and $\del g$. To show that $\w g$ is continuous, suppose that $U\subset\bord X$ is open, and let $V = \w g^{-1}(U)$. Then $V\cap X = g^{-1}(U\cap X)$ is open. Fix $\xi\in V\cap \del X$. Then there exists $t > 0$ such that $N_t(\del g(\xi)) \subset U$.
%%%\begin{claim}
%%%$\w g(N_{t + \dogo g})(\xi) \subset N_t(\del g(\xi))$.
%%%\end{claim}
%%%\begin{subproof}
%%%Suppose $y\in N_{t + \dogo g}(\xi)$, i.e. $\lb y|\xi\rb_\zero > t + \dogo g$. Suppose first that $y\in\del X$. Then
%%%\begin{align*}
%%%\lb \w g(y) | \w g(\xi)\rb_{\w g(\zero)}
%%%&= \inf\left\{\liminf_{n,m\to\infty} \lb y_n | x_m\rb_{g(\zero)} : (y_n)_1^\infty\in\w g(y), (x_m)_1^\infty\in\w g(\xi)\right\}\\
%%%&= \inf\left\{\liminf_{n,m\to\infty} \lb g(y_n) | g(x_m)\rb_{g(\zero)} : (y_n)_1^\infty\in y, (x_m)_1^\infty\in\xi\right\}\\
%%%&= \inf\left\{\liminf_{n,m\to\infty} \lb y_n | x_m\rb_\zero : (y_n)_1^\infty\in y, (x_m)_1^\infty\in\xi\right\}\\
%%%&= \lb y|\xi\rb_\zero > t + \dogo g.
%%%\end{align*}
%%%On the other hand, by (d) of Proposition \ref{propositionbasicidentities} (cf. Remark \ref{remarkacde}), we have $\lb \w g(y) | \w g(\xi)\rb_{\w g(\zero)} \leq \lb \w g(y) | \w g(\xi)\rb_\zero + \dogo g$. Thus $\lb \w g(y) | \w g(\xi)\rb_\zero > t$, i.e. $\w g(y)\in N_t(\w g(\xi))$. A similar argument works if $y\in X$. \comdavid{This proof seems too detailed - can we say it in words, like in the proof of Corollary \ref{corollaryboundaryasymptotic}?}
%%%\end{subproof}
%%%
%%%Thus $N_{t + \dogo g}(\xi) \subset \w g^{-1}(U) = V$. Since $\xi$ was arbitrary, this demonstrates that $V$ is open.
%%%\end{proof}
%%%}% end ignore
%%In the sequel we will omit the tilde from the extended map $\w g$.
%%
%%
%%
%%%%%%%%%%%%%%%%%%%%%%%%%%%%%%%%%%%%
%%\bigskip
%%\subsection{The Gromov product in \CROSSs} \label{subsectionGromovinH}
%%%%%%%%%%%%%%%%%%%%%%%%%%%%%%%%%%%%
%%
%%In this subsection we analyze the Gromov product in a \CROSS\ $X$. We prove Proposition \ref{propositionCATimpliesGromov} which states that CAT(-1) spaces are strongly hyperbolic, and then we show that the Gromov boundary of $X$ is isomorphic to its topological boundary, justifying Remark \ref{remarkambiguityboundary}.
%%
%%In what follows, we will switch between the hyperboloid model $\H = \H_\F^\alpha$ and the ball model $\B = \B_\F^\alpha$ according to convenience. In the following lemma, $\del\B$ and $\bord\B$ denote the topological boundary and closure of $\B$ as defined in Section \ref{sectionROSSONCTs}, not the Gromov boundary and closure as defined above.
%%
%%\begin{figure}
%%\begin{center}
%%\definecolor{qqwuqq}{rgb}{0.0,0.39215686274509803,0.0}
%%\begin{tikzpicture}[line cap=round,line join=round,>=triangle 45,x=1.0cm,y=1.0cm]
%%\clip(-3.40,-3.1) rectangle (3.40,3.13);
%%\draw [shift={(0.0,-0.0)},color=qqwuqq,fill=qqwuqq,fill opacity=0.1] (0,0) -- (27.07967008819517:0.5963691863565421) arc (27.07967008819517:60.832386620422206:0.5963691863565421) -- cycle;
%%\draw(0.0,0.0) circle (3.0cm);
%%\draw (0.0,-0.0)-- (2.6711231603651164,1.3656870293596084);
%%\draw (0.0,-0.0)-- (1.4620984777924115,2.619593106044737);
%%\draw (1.4620984777924115,2.619593106044737)-- (2.6711231603651164,1.3656870293596084);
%%\begin{scriptsize}
%%\draw [fill=black] (1.4620984777924115,2.619593106044737) circle (.75pt);
%%\draw[color=black] (1.5997838888502904,2.7999980589688112) node {$\xx$};
%%\draw [fill=black] (2.6711231603651164,1.3656870293596084) circle (.75pt);
%%\draw[color=black] (2.9117960988346834,1.4482279032273154) node {$\yy$};
%%\draw [fill=black] (0.0,-0.0) circle (.75pt);
%%\draw[color=black] (-0.2290816159764389,-0.12160581857100267) node {$\0$};
%%\draw[color=qqwuqq] (0.6040198381009913,0.57476336778553957) node {$\theta$};
%%\end{scriptsize}
%%\end{tikzpicture}
%%\caption{If $\F = \R$ and $\xx,\yy\in\del\B$, then $e^{-\lb\xx|\yy\rb_\0} = \frac12\|\yy - \xx\| = \sin(\theta/2)$, where $\theta$ denotes the angle $\ang_\0(\xx,\yy)$ drawn in the figure.}
%%\end{center}
%%\end{figure}
%%
%%
%%\begin{lemma}
%%\label{lemmagromovextension}
%%The Gromov product $(\xx,\yy,\zz)\mapsto \lb \xx|\yy\rb_\zz:\B\times\B\times\B\to\Rplus$ extends uniquely to a continuous function $(\xx,\yy,\zz)\mapsto \lb \xx|\yy\rb_\zz:\bord\B\times\bord\B\times\B\to[0,\infty]$. Moreover, the extension satisfies the following:
%%\begin{itemize}
%%\item[(i)] $\lb\xx|\yy\rb_{\zz} = \infty$ if and only if $\xx = \yy\in\del\B$.
%%\item[(ii)] For all $\xx,\yy\in\bord\B$,
%%\begin{equation}
%%\label{euclideanball1}
%%e^{-\lb\xx|\yy\rb_\0} \geq \frac{1}{\sqrt 8}\|\yy - \xx\|.
%%\end{equation}
%%If $\F = \R$ and $\xx,\yy\in\del\B$, then
%%\begin{equation}
%%\label{euclideanball2}
%%e^{-\lb\xx|\yy\rb_\0} = \frac12\|\yy - \xx\|.
%%\end{equation}
%%%\item[(iii)]
%%%\begin{equation}
%%%\label{Gromovasymptotic}
%%%\lb \xx|\yy\rb_\0 \asymp_\plus \min\left(\dist_\B(\0,\xx),\dist_\B(\0,\yy),-\log\ang(\xx,\yy)\right)\text{ for $\xx,\yy\in\bord\B$}.
%%%\end{equation}
%%%Here by convention $\dist_\B(\0,\xx) = \infty$ if $\xx\in\del\B$.
%%\end{itemize}
%%\end{lemma}
%%%\ignore{
%%%\begin{proof}
%%%Define the function $F:\Rplus\times\Rplus\times[0,\pi]\to\Rplus$ by
%%%\[
%%%F(A,B,\gamma) = \frac{\exp\left[\cosh^{-1}\left(\cosh(A)\cosh(B) - \sinh(A)\sinh(B)\cos(\gamma)\right)\right]}{e^A e^B}.
%%%\]
%%%Then by Proposition \ref{lawofcosines},we have
%%%\begin{align*}
%%%e^{-2\lb \xx|\yy\rb_\zero} &= F(\dist_\B(\zero,\xx),\dist_\B(\zero,\yy),\ang(\xx,\yy))
%%%%\\&= F(2\tanh^{-1}(\|\xx\|),2\tanh^{-1}(\|\yy\|),\cos^{-1}(\lb\xx,\yy\rb/(\|\xx\|\cdot\|\yy\|)))
%%%\end{align*}
%%%for all $\xx,\yy\in\B$. Now since $\lim_{t\to\infty}\frac{e^t}{\cosh(t)} = 2$, we have
%%%\begin{align*}
%%%\lim_{B\to\infty}F(A,B,\gamma)
%%%&= \lim_{B\to\infty}\frac{2\left(\cosh(A)\cosh(B) - \sinh(A)\sinh(B)\cos(\gamma)\right)}{e^A [2\cosh(B)]}\\
%%%&= \lim_{B\to\infty}\frac{\cosh(A) - \sinh(A)\tanh(B)\cos(\gamma)}{e^A}\\
%%%&= \frac{\cosh(A) - \sinh(A)\cos(\gamma)}{e^A}
%%%\end{align*}
%%%and 
%%%\[
%%%\lim_{A\to\infty}\frac{\cosh(A) - \sinh(A)\cos(\gamma)}{e^A}
%%%= \frac{1 - \cos(\gamma)}{2}
%%%= \sin^2(\gamma/2).
%%%\]
%%%Thus if we let
%%%\[
%%%\what F(A,B,\gamma) :=
%%%\begin{cases}
%%%F(A,B,\gamma) & A,B < \infty\\
%%%\frac{\cosh(A) - \sinh(A)\cos(\gamma)}{e^A} & A < B = \infty\\
%%%\frac{\cosh(B) - \sinh(B)\cos(\gamma)}{e^B} & B < A = \infty\\
%%%\sin^2(\gamma/2) & A = B = \infty
%%%\end{cases}
%%%\]
%%%then $\what F:[0,\infty]\times[0,\infty]\times[0,\pi]\to[0,\infty]$ is a continuous function.\Footnote{Technically, the calculations above do not prove the continuity of $\what F$; however, this continuity is easily verified.} Thus, letting
%%%\[
%%%\lb \xx|\yy\rb_\zero := -\frac12\log\left(\what F\bigl(\dist_\B(\zero,\xx),\dist_\B(\zero,\yy),\ang(\xx,\yy)\bigr)\right)
%%%\]
%%%defines a continuous extension of the Gromov product to the boundary.
%%%
%%%We now prove (i)-(iii):
%%%\begin{itemize}
%%%\item[(i)] By inspection, we see that $\what F(A,B,\gamma) = 0$ if and only if $A = B = \infty$ and $\sin^2(\gamma/2) = 0$, i.e. $\gamma = 0$. Thus, in order to have $\lb \xx|\yy\rb_\zero = \infty$ we must have $\dist_\B(\zero,\xx) = \dist_\B(\zero,\yy) = \infty$ and $\ang(\xx,\yy) = 0$, or in other words $\xx,\yy\in\del\B$ and $\xx = \yy$.
%%%\item[(ii)] If $\xx,\yy\in\del\B$, then $\dist_\B(\zero,\xx) = \dist_\B(\zero,\yy) = \infty$. Thus
%%%\[
%%%e^{-2\lb\xx|\yy\rb_\zero} = \what F\bigl(\infty,\infty,\ang(\xx,\yy)\bigr) = \sin^2\bigl(\ang(\xx,\yy)/2\bigr)
%%%\]
%%%and thus
%%%\[
%%%e^{-\lb\xx|\yy\rb_\zero} = \sin\bigl(\ang(\xx,\yy)/2\bigr)
%%%= \frac12\|\yy - \xx\|
%%%\]
%%%(see Figure \internalmissingref).
%%%\item[(iii)] For all $A,B\in\Rplus$ and $\gamma\in[0,\pi]$, using the asymptotic $\cosh(x)\asymp_\times e^x$ three times we have
%%%\begin{align*}
%%%\what F(A,B,\gamma)
%%%&\asymp_\times \frac{\cosh(A)\cosh(B) - \sinh(A)\sinh(B)\cos(\gamma)}{\cosh(A)\cosh(B)}\\
%%%&=_\pt 1 - \tanh(A)\tanh(B)\cos(\gamma)\\
%%%&\asymp_\times \max\left(1 - \tanh(A),1 - \tanh(B),1 - \cos(\gamma)\right)\\
%%%&\asymp_\times \max\left(e^{-2A},e^{-2B},\gamma^2\right).
%%%\end{align*}
%%%Thus
%%%\[
%%%-\frac12\log(\what F(A,B,\gamma)) \asymp_\plus \min(A,B,-\log(\gamma)),
%%%\]
%%%which implies \eqref{Gromovasymptotic}.
%%%\end{itemize}
%%%\end{proof}
%%%}% end ignore
%%\begin{proof}
%%We begin by making some computations in the hyperboloid model $\H$. For $[\xx],[\yy]\in\bord\H$ and $[\zz]\in\H$, let
%%\[
%%\alpha_{[\zz]}([\xx],[\yy]) = \frac{|\QQ(\zz)|\cdot |B_\QQ(\xx,\yy)|}{|B_\QQ(\xx,\zz)|\cdot |B_\QQ(\yy,\zz)|}\in\Rplus.
%%\]
%%By \eqref{distanceinL}, for $[\xx],[\yy],[\zz]\in\H$ we have
%%\begin{equation}
%%\label{lorentzlawofcosh}
%%\alpha_{[\zz]}([\xx],[\yy]) = \frac{\cosh\dist_\H([\xx],[\yy])}{\cosh\dist_\H([\xx],[\zz])\cosh\dist_\H([\yy],[\zz])}\cdot
%%\end{equation}
%%Let $\DD = \{(A,B,C)\in\Rplus^3: \cosh(A)\cosh(B)C \geq 1\}$, and define $F: \DD\to\Rplus$ by
%%\[
%%F(A,B,C) = \frac{\exp\left[\cosh^{-1}\left(\cosh(A)\cosh(B)C\right)\right]}{e^A e^B}\cdot
%%\]
%%Then by \eqref{lorentzlawofcosh}, we have
%%\[
%%e^{-2\lb [\xx]|[\yy]\rb_{[\zz]}} = F\left(\dist_\H([\zz],[\xx]),\dist_\H([\zz],[\yy]),\alpha_{[\zz]}([\xx],[\yy])\right)
%%\]
%%for all $[\xx],[\yy],[\zz]\in\H$. Now since $\lim_{t\to\infty}\frac{e^t}{\cosh(t)} = 2$, we have for all $A\geq 0$ and $C > 0$
%%\begin{align*}
%%\lim_{B\to\infty}F(A,B,C)
%%&= \lim_{B\to\infty}\frac{2\left(\cosh(A)\cosh(B)C\right)}{e^A e^B}\\
%%&= \frac{\cosh(A)}{e^A} C
%%\end{align*}
%%and 
%%\[
%%\lim_{A\to\infty}\frac{\cosh(A)}{e^A} C
%%= C/2.
%%\]
%%Let $\what \DD$ be the closure of $\DD$ relative to $[0,\infty]^2\times \Rplus$, i.e. $\what \DD = \DD\cup \left([0,\infty]^2\times \Rplus\butnot \Rplus^3\right)$. If we let
%%\[
%%\what F(A,B,C) :=
%%\begin{cases}
%%F(A,B,C) & A,B < \infty\\
%%\frac{\cosh(A)}{e^A} C & A < B = \infty\\
%%\frac{\cosh(B)}{e^B} C & B < A = \infty\\
%%C/2 & A = B = \infty
%%\end{cases}
%%\]
%%then $\what F: \what \DD\to\Rplus$ is a continuous function.\Footnote{Technically, the calculations above do not prove the continuity of $\what F$; however, this continuity is easily verified using standard methods.} Thus, letting
%%\[
%%\lb [\xx]|[\yy]\rb_{[\zz]} := -\frac12\log\what F\left(\dist_\H([\zz],[\xx]),\dist_\H([\zz],[\yy]),\alpha_{[\zz]}([\xx],[\yy])\right)
%%\]
%%defines a continuous extension of the Gromov product to $\bord\H\times\bord\H\times\H$.
%%
%%We now prove (i)-(ii):
%%\begin{itemize}
%%\item[(i)] Using the inequality $e^t/2 \leq \cosh(t) \leq e^t$, it is easily verified that
%%\begin{equation}
%%\label{C4}
%%\what F(A,B,C) \geq C/4
%%\end{equation}
%%for all $(A,B,C)\in\what\DD$. In particular, if $\what F(A,B,C) = 0$ then $C = 0$. Thus if $\lb [\xx]|[\yy]\rb_{[\zz]} = \infty$ then $\alpha_{[\zz]}([\xx],[\yy]) = 0$; since $[\zz]\in\H$ we have $B_\QQ(\xx,\yy) = 0$, and by Observation \ref{observationBQnonzero} we have $[\xx] = [\yy]\in\del\H$. Conversely, if $[\xx] = [\yy]\in\del\H$, then $\dist_\H([\zz],[\xx]) = \dist_\H([\zz],[\yy]) = \infty$ and $\alpha_{[\zz]}([\xx],[\yy]) = 0$, so $\lb [\xx]|[\yy]\rb_{[\zz]} = -\frac12 \log\what F(\infty,\infty,0) = \infty$.
%%
%%%\item[(ii)] If $\xx,\yy\in\del\B$, then $\dist_\B(\zz,\xx) = \dist_\B(\zz,\yy) = \infty$. Thus
%%%\[
%%%e^{-2\lb\xx|\yy\rb_{\zz}} = \what F\bigl(\infty,\infty,\alpha_{\zz}(\xx,\yy)\bigr) = \alpha_{\zz}(\xx,\yy)/2.
%%%\]
%%%Letting $\zz = \ee_0$ and normalizing $\xx$ and $\yy$ so that $x_0 = y_0 = 1$, we have $B_\QQ(\xx,\zz) = B_\QQ(\yy,\zz) = \QQ(\zz) = 1$; thus
%%%\[
%%%\alpha_{\zz}(\xx,\yy) = B_\QQ(\xx,\yy) = \cdots
%%%\]
%%\item[(ii)] Recall that in $\B$, $\zero = [(1,\0)]$. For $\xx,\yy\in\B$,
%%\begin{align*}
%%\alpha_\zero(e_{\B,\H}(\xx),e_{\B,\H}(\yy)) &= \frac{|\QQ(1,\0)|\cdot|B_\QQ((1,\xx),(1,\yy))|}{|B_\QQ((1,\0),(1,\xx))|\cdot |B_\QQ((1,\0),(1,\yy))|} \noreason\\
%%&= |1 - B_\EE(\xx,\yy)|\\
%%&\geq 1 - \Re B_\EE(\xx,\yy) &\text{(with equality if $\F = \R$)}\\
%%&\geq \frac12 [\|\xx\|^2 + \|\yy\|^2] - \Re B_\EE(\xx,\yy) &\text{(with equality if $\xx,\yy\in\del\B$)}\\
%%&= \frac12\|\yy - \xx\|^2.
%%\end{align*}
%%Combining with \eqref{C4} gives
%%\[
%%e^{-2\lb \xx|\yy\rb_\0} \geq \frac14\alpha_\zero(e_{\B,\H}(\xx),e_{\B,\H}(\yy)) \geq \frac18\|\yy - \xx\|^2.
%%\]
%%If $\F = \R$ and $\xx,\yy\in\del\B$, then
%%\[
%%e^{-2\lb \xx|\yy\rb_\0} = \frac12\alpha_\zero(e_{\B,\H}(\xx),e_{\B,\H}(\yy)) = \frac14\|\yy - \xx\|^2.
%%\]
%%\end{itemize}
%%\end{proof}
%%
%%
%%We now prove Proposition \ref{propositionCATimpliesGromov}, beginning with the following lemma:
%%\begin{lemma}
%%\label{lemmaHstrongly}
%%If $\F = \R$ then $\B$ is strongly Gromov hyperbolic.
%%\end{lemma}
%%\begin{proof}
%%By the transitivity of the isometry group (Observation \ref{observationtransitivity}), it suffices to check \eqref{gromovforCAT} for the special case $w = \zero$. So let us fix $x,y,z\in\B$, and by contradiction suppose that
%%\[
%%e^{-\lb x|z\rb_\zero} > e^{-\lb x|y\rb_\zero} + e^{-\lb y|z\rb_\zero},
%%\]
%%or equivalently that
%%\[
%%1 > e^{\lb x|z\rb_\zero - \lb x|y\rb_\zero} + e^{\lb x|z\rb_\zero - \lb y|z\rb_\zero}.
%%\]
%%Clearly, the above inequality implies that $x\neq z$ and $y\neq\zero$. Now let $\gamma_1$ and $\gamma_2$ be the unique bi-infinite geodesics extending the geodesic segments $\geo xz$ and $\geo\zero y$, respectively. Let $x_\infty,z_\infty\in\del\B$ be the appropriate endpoints of $\gamma_1$, and let $y_\infty$ be the endpoint of $\gamma_2$ which is closer to $y$ than to $\zero$. (See Figure \ref{figureHstrongly}.) For each $t\in\Rplus$, let
%%\[
%%x_t = \geo x{x_\infty}_t\in\gamma_1,
%%\]
%%and let $y_t\in\gamma_2$, $z_t\in\gamma_1$ be defined similarly. 
%%
%%
%%\begin{figure}
%%\begin{center}
%%\begin{tikzpicture}[line cap=round,line join=round,>=triangle 45,x=0.4402812927177227cm,y=0.44028129271787425cm]
%%\clip(-9.71,-7.2) rectangle (11,7.5);
%%\draw(0,0) ellipse (3.11cm and 3.11cm);
%%\draw (5,5)-- (-5,-5);
%%\draw [shift={(12.49,9.91)}] plot[domain=3.35:4.27,variable=\t]({1*14.67*cos(\t r)+0*14.67*sin(\t r)},{0*14.67*cos(\t r)+1*14.67*sin(\t r)});
%%\begin{scriptsize}
%%\fill [color=black] (0,0) circle (1pt);
%%\draw[color=black] (0.45,-0.58) node {$\zero$};
%%\fill [color=black] (5,5) circle (1pt);
%%\draw[color=black] (5.9,4.96) node {$y_\infty$};
%%\fill [color=black] (-1.85,6.83) circle (1pt);
%%\draw[color=black] (-2.5,7.14) node {$x_\infty$};
%%\fill [color=black] (6.29,-3.33) circle (1pt);
%%\draw[color=black] (7.2,-3.6) node {$z_\infty$};
%%\fill [color=black] (-0.97,4.09) circle (1pt);
%%\draw[color=black] (-1.53,3.88) node {$x$};
%%\fill [color=black] (-1.36,5.08) circle (1pt);
%%\draw[color=black] (-1.84,4.81) node {$x_t$};
%%\fill [color=black] (4.34,-2.28) circle (1pt);
%%\draw[color=black] (3.99,-2.7) node {$z$};
%%\fill [color=black] (4.91,-2.64) circle (1pt);
%%\draw[color=black] (4.93,-3.3) node {$z_t$};
%%\fill [color=black] (3.12,3.12) circle (1pt);
%%\draw[color=black] (3.67,2.63) node {$y_t$};
%%\fill [color=black] (2,2) circle (1pt);
%%\draw[color=black] (2.42,1.7) node {$y$};
%%\end{scriptsize}
%%\end{tikzpicture}
%%\caption{If Gromov's inequality fails for the quadruple $x,y,z,\zero$, then it also fails for the quadruple $x_\infty,y_\infty,z_\infty,\zero$.}
%%\label{figureHstrongly}
%%\end{center}
%%\end{figure}
%%
%%
%%We observe that
%%\begin{align*}
%%\frac{\del}{\del t}\left[\lb x_t|z_t\rb_\zero - \lb x_t|y_t\rb_\zero\right]
%%&= \frac12 \frac{\del}{\del t}\bigl[\dist_\B(\zero,z_t) + \dist_\B(x_t,y_t) - \dist_\B(x_t,z_t) - \dist_\B(\zero,y_t)\bigr]\\
%%&= \frac12 \frac{\del}{\del t}\bigl[\dist_\B(\zero,z_t) + \dist_\B(x_t,y_t) - 2t - t\bigr]\\
%%&\leq \frac12 \frac{\del}{\del t}\bigl[t + 2t - 2t - t\bigr] = 0,
%%\end{align*}
%%i.e. the expression $\lb x_t|z_t\rb_\zero - \lb x_t|y_t\rb_\zero$ is nonincreasing with respect to $t$. Taking the limit as $t$ approaches infinity, we have
%%\[
%%\lb x_\infty|z_\infty\rb_\zero - \lb x_\infty|y_\infty\rb_\zero \leq \lb x|z\rb_\zero - \lb x|y\rb_\zero
%%\]
%%and a similar argument shows that
%%\[
%%\lb x_\infty|z_\infty\rb_\zero - \lb y_\infty|z_\infty\rb_\zero \leq \lb x|z\rb_\zero - \lb y|z\rb_\zero.
%%\]
%%Thus
%%\[
%%1 > e^{\lb x_\infty|z_\infty\rb_\zero - \lb x_\infty|y_\infty\rb_\zero} + e^{\lb x_\infty|z_\infty\rb_\zero - \lb y_\infty|z_\infty\rb_\zero}
%%\]
%%or equivalently,
%%\[
%%e^{-\lb x_\infty|z_\infty\rb_\zero} > e^{-\lb x_\infty|y_\infty\rb_\zero} + e^{-\lb y_\infty|z_\infty\rb_\zero}.
%%\]
%%But by \eqref{euclideanball2}, if we write $x_\infty = \xx$, $y_\infty = \yy$, and $z_\infty = \zz$, then
%%\[
%%\frac12\|\zz - \xx\| > \frac12\|\yy - \xx\| + \frac12\|\zz - \yy\|.
%%\]
%%This is a contradiction.
%%\end{proof}
%%
%%We are now ready to prove
%%\begin{repproposition}{propositionCATimpliesGromov}
%%Every CAT(-1) space is strongly hyperbolic.
%%\end{repproposition}
%%\begin{proof}
%%Let $X$ be a CAT(-1) space, and fix $x,y,z,w\in X$. By \cite[Proposition II.1.11]{BridsonHaefliger}, there exist $\wbar x,\wbar y,\wbar z,\wbar w\in \H^2$ such that
%%\begin{align*}
%%\dist(x,y) &= \dist(\wbar x,\wbar y) &
%%\dist(y,z) &= \dist(\wbar y,\wbar z)\\
%%\dist(z,w) &= \dist(\wbar z,\wbar w) &
%%\dist(w,x) &= \dist(\wbar w,\wbar x)\\
%%\dist(x,z) &\leq \dist(\wbar x,\wbar z) &
%%\dist(y,w) &\leq \dist(\wbar y,\wbar w).
%%\end{align*}
%%It follows that
%%\[
%%e^{-\lb x|z\rb_w} \leq e^{-\lb \wbar x|\wbar z\rb_{\wbar w}} \leq e^{-\lb \wbar x|\wbar y\rb_{\wbar w}} + e^{-\lb \wbar y|\wbar z\rb_{\wbar w}} \leq e^{-\lb x|y\rb_w} + e^{-\lb y|z\rb_w}.
%%\]
%%%\ignore{
%%%
%%%
%%%
%%%Let $\Delta(\wbar x,\wbar y, \wbar w) \subset \B^2$ be a comparison triangle for $\Delta(x,y,w)$, Then the geodesic $\geo{\wbar y}{\wbar w}$ can be extended to a triangle $\Delta(\wbar y,\wbar z,\wbar w) \subset \B^2$ which is a comparison triangle for $\Delta(y,z,w)$ and which has the following property: the points $\wbar x$ and $\wbar z$ lie on opposite sides of the unique bi-infinite geodesic $\gamma$ which is an extension of $\geo{\wbar y}{\wbar w}$. Let $\wbar p$ be the unique intersection point of $\geo{\wbar x}{\wbar z}$ and $\gamma$. We now divide into three cases:
%%%\begin{itemize}
%%%\item[Case 1:] $\wbar p\in \geo{\wbar y}{\wbar w}$. In this case, we can draw a point $p\in\geo yw$ whose comparison point is $\wbar p$. Then by the CAT(-1) inequality
%%%\[
%%%\dist(x,z) \leq \dist(x,p) + \dist(p,z) \leq \dist(\wbar x, \wbar p) + \dist(\wbar p,\wbar z) = \dist(\wbar x, \wbar z).
%%%\]
%%%On the other hand, because of the way that the points $\wbar x$, $\wbar y$, $\wbar z$, and $\wbar w$ were constructed, all other distances are equal in the original and comparison figures. This and the fact that the comparison figure satisfies \eqref{gromovforCAT} implies that $(x,y,z,w)$ satisfies \eqref{gromovforCAT}.
%%%
%%%
%%%\item[Case 2:] $\wbar w\in\geo{\wbar p}{\wbar y}$. In this case, for each $t\in [0,\dist(\wbar w,\wbar p)]$ let $\wbar w_t = \geo{\wbar w}{\wbar p}_t$. Consider the inequality
%%%\begin{equation}
%%%\label{1leq}
%%%1 \leq e^{-\lb \wbar z|\wbar y\rb_{\wbar w_t}} + e^{-\lb \wbar x|\wbar y\rb_{\wbar w_t}}.
%%%\end{equation}
%%%By Lemma \ref{lemmaHstrongly}, \eqref{1leq} holds when $t = \dist(\wbar w,\wbar p)$, since then $\wbar w_t = \wbar p$ and thus $\lb \wbar x|\wbar z\rb_{\wbar w_t} = 0$. On the other hand, one checks that the right hand side of \eqref{1leq} is decreasing with respect to $t$. This implies that \eqref{1leq} holds when $t = 0$. Thus
%%%\[
%%%e^{-\lb x|z\rb_w} \leq 1 \leq e^{-\lb \wbar z|\wbar y\rb_{\wbar w}} + e^{-\lb \wbar x|\wbar y\rb_{\wbar w}} = e^{-\lb z|y\rb_w} + e^{-\lb x|y\rb_w}.
%%%\]
%%%\item[Case 3:] $\wbar y\in\geo{\wbar p}{\wbar w}$. This case follows from Case 2, since the equation \eqref{gromovforCAT} is invariant under switching $y$ and $w$.
%%%\end{itemize}
%%%}% end ignore
%%\end{proof}
%%
%%\bigskip
%%\subsubsection{The Gromov boundary of a \CROSS}
%%\label{subsubsectionboundariesequivalent}
%%
%%Again let $X = \H = \H_\F^\alpha$ be a \CROSS. By Proposition \ref{propositionCATimpliesGromov}, $X$ is a (strongly) hyperbolic metric space. (If $\F = \R$, we can use Lemma \ref{lemmaHstrongly}.) In particular, $X$ has a Gromov boundary, defined in Subsection \ref{subsectionboundary}. On the other hand, $X$ also has a topological boundary, defined in Section \ref{sectionROSSONCTs}. For this subsection only, we will write
%%\begin{align*}
%%\del_G X &= \text{ Gromov boundary of $X$}, \\
%%\del_T X &= \text{ topological boundary of $X$}.
%%\end{align*}
%%We will now show that this distinction is in fact unnecessary.
%%\begin{proposition}
%%\label{propositionboundariesequivalent}
%%The identity map $\id:X\to X$ extends uniquely to a homeomorphism $\what\id:X\cup\del_G X\to X\cup\del_T X$. Thus the pairs $(X,X\cup\del_G X)$ and $(X,X\cup\del_T X)$ are isomorphic in the sense of Subsection \ref{subsectiontotallygeodesic}.
%%\end{proposition}
%%
%%\begin{proof}[Proof of Proposition \ref{propositionboundariesequivalent}]
%%By Observation \ref{observationHequivB} and Proposition \ref{propositionHequivE}, it suffices to consider the case where $X = \B = \B_\F^\alpha$ is the ball model. Fix $\xi\in\del_G\B$. By definition, $\xi = [(\xx_n)_1^\infty]$ for some Gromov sequence $(\xx_n)_1^\infty$. By \eqref{euclideanball1}, the sequence $(\xx_n)_1^\infty$ is Cauchy in the metric $\|\cdot - \cdot\|$. Thus $\xx_n\to\xx$ for some $\xx\in\bord\B$; since $(\xx_n)_1^\infty$ is a Gromov sequence, we have
%%\[
%%\lb \xx|\xx\rb_\0 = \lim_{n,m\to\infty}\lb \xx_n|\xx_m\rb_\0 = \infty,
%%\]
%%and thus $\xx\in\del_T\B$ by (i) of Lemma \ref{lemmagromovextension}. Let
%%\[
%%\what\id(\xi) = \xx.
%%\]
%%To see that the map $\what\id$ is well-defined, note that if $(\yy_n)_1^\infty$ is another Gromov sequence equivalent to $(\xx_n)_1^\infty$, and if $\yy_n\to\yy\in\del_T \B$, then
%%\[
%%\lb \xx|\yy\rb_\0 = \lim_{n\to\infty}\lb \xx_n|\yy_n\rb_\0 = \infty,
%%\]
%%and so by (i) of Lemma \ref{lemmagromovextension} we have $\xx = \yy$.
%%
%%We next claim that $\what\id:\del_G\B\to\del_T\B$ is a bijection. To demonstrate injectivity, we note that if $\what\id(\xi) = \what\id(\eta) = \xx$, then by (i) of Lemma \ref{lemmagromovextension}
%%\[
%%\lim_{n\to\infty}\lb \xx_n|\yy_n\rb_\0 = \lb \xx|\xx\rb_\0 = \infty,
%%\]
%%where $(\xx_n)_1^\infty$ and $(\yy_n)_1^\infty$ are Gromov sequences representing $\xi$ and $\eta$, respectively. Thus $(\xx_n)_1^\infty$ and $(\yy_n)_1^\infty$ are equivalent, and so $\xi = \eta$.
%%
%%To demonstrate surjectivity, we observe that for $\xx\in\del_T\B$, we have
%%\[
%%\what\id\left(\left[\left(\frac{n - 1}{n}\xx\right)_1^\infty\right]\right) = \xx.
%%\]
%%Finally, we must demonstrate that $\what\id$ is a homeomorphism, or in other words that the topology defined in \6\ref{subsubsectiontopologyclX} is the usual topology on $\bord\B$ (i.e. the topology inherited from $\HH$). It suffices to demonstrate the following:
%%\begin{claim}
%%For any $\xx\in\del_T\B$, the collection \eqref{Ntxibase} (with $\xi = \xx$) is a neighborhood base of $\xx$ with respect to the usual topology.
%%\end{claim}
%%\begin{subproof}
%%By \eqref{euclideanball1}, we have
%%\[
%%N_t(\xx) \subset B(\xx,\sqrt 8 e^{-t}).
%%\]
%%On the other hand, the continuity of the Gromov product on $\bord\B$ guarantees that $N_t(\xx)$ contains a neighborhood of $\xx$ with respect to the usual topology.
%%\end{subproof}
%%
%%% Ignored since it's not clear whether we'll have metrizability at this point
%%%\ignore{
%%%Since both topologies are metrizable, it suffices to show that for any sequence $(\xx_n)_1^\infty$ in $\bord\B$ and for any $\xx\in\bord\B$, $\xx_n\to\xx$ in the former topology if and only if $\xx_n\to\xx$ in the latter topology. If $\xx\in\B$, this is clear. Suppose that $\xx\in\del \B$. Then $\xx_n\to\xx$ in the topology $\scrT$ if and only if
%%%\[
%%%\lb \xx_n|\xx\rb_\zero \tendsto n \infty,
%%%\]
%%%and by (iii) of Lemma \ref{lemmagromovextension}, this is equivalent to
%%%\[
%%%\dist_\B(\zero,\xx_n)\tendsto n \infty, \hspace{.5 in} \ang(\xx,\xx_n)\tendsto n 0,
%%%\]
%%%i.e. $\xx_n\to\xx$ in the usual topology.
%%%}% end ignore
%%\end{proof}
%%
%%In the sequel, the following will be useful:
%%
%%\begin{figure}
%%\begin{center}
%%\begin{tikzpicture}[line cap=round,line join=round,>=triangle 45,x=1.0cm,y=1.0cm]
%%\clip(-3.8,-0.0) rectangle (5.57,4.41);
%%\draw (-3.0,0.0)-- (3.0,0.0);
%%\draw [dash pattern=on 5pt off 5pt] (-1.52,2.22)-- (-1.52,0.0);
%%\draw [dash pattern=on 5pt off 5pt] (1.76,1.48)-- (1.7800000000000002,0.0);
%%\begin{scriptsize}
%%%\draw[color=black] (0.05368624806072046,-0.1412911554774228) node {$a$};
%%\draw [fill=black] (0.0,4.0) circle (.75pt);
%%\draw[color=black] (0.13316978586974992,4.270045192923707) node {$\infty$};
%%\draw [fill=black] (-1.52,2.22) circle (.75pt);
%%\draw[color=black] (-1.7346933526424422,2.5810200144818323) node {$\xx$};
%%\draw [fill=black] (1.76,1.48) circle (.75pt);
%%\draw[color=black] (1.961291155477427,1.7663137519392818) node {$\yy$};
%%\draw[color=black] (-1.834047774903729,1.1304454494670468) node {$x_1$};
%%\draw[color=black] (2.0209038088341993,0.8721239515877014) node {$y_1$};
%%\end{scriptsize}
%%\end{tikzpicture}
%%\caption{The value of the Busemann function $\busemann_\infty(\xx,\yy)$ depends on the heights of the points $\xx$ and $\yy$.}
%%\end{center}
%%\end{figure}
%%
%%\begin{proposition}
%%\label{propositionbusemannE}
%%Let $\E = \E^\alpha$ be the half-space model of a real \ROSS. For $\xx,\yy\in\E$,
%%\[
%%\busemann_\infty(\xx,\yy) = -\log(x_1/y_1).
%%\]
%%\end{proposition}
%%\begin{proof}
%%By \eqref{distE} we have
%%\begin{align*}
%%e^{\busemann_\infty(\xx,\yy)}
%%= \lim_{\zz\to\infty} \frac{\exp\dist_\E(\zz,\xx)}{\exp\dist_\E(\zz,\yy)}
%%= \lim_{\zz\to\infty} \frac{\cosh\dist_\E(\zz,\xx)}{\cosh\dist_\E(\zz,\yy)}
%%= \lim_{\zz\to\infty} \frac{1 + \frac{\|\zz - \xx\|^2}{2 x_1 z_1}}{1 + \frac{\|\zz - \yy\|^2}{2 y_1 z_1}}
%%= \lim_{\zz\to\infty} \frac{\left(\frac{\|\zz - \xx\|^2}{2 x_1 z_1}\right)}{\left(\frac{\|\zz - \yy\|^2}{2 y_1 z_1}\right)}
%%= \frac{y_1}{x_1}\cdot
%%\end{align*}
%%\end{proof}
%%
%%
%%
%%%%%%%%%%%%%%%%%%%%%%%%%%%%%%%%%%%%
%%\bigskip
%%\subsection{Metrics and metametrics on $\bord X$}
%%\label{subsectionmetametrics}
%%%%%%%%%%%%%%%%%%%%%%%%%%%%%%%%%%%%
%%
%%\subsubsection{General theory of metametrics}
%%
%%\begin{definition}
%%\label{definitionmetametric}
%%Recall that a \emph{metric} on a set $Z$ is a map $\Dist:Z\times Z\to\Rplus$ which satisfies:
%%\begin{itemize}
%%\item[(I)] Reflexivity: $\Dist(x,x) = 0$.
%%\item[(II)] Reverse reflexivity: $\Dist(x,y) = 0$ \implies $x = y$.
%%\item[(III)] Symmetry: $\Dist(x,y) = \Dist(y,x)$.
%%\item[(IV)] Triangle inequality: $\Dist(x,z) \leq \Dist(x,y) + \Dist(y,z)$.
%%\end{itemize}
%%Now we can define a \emph{metametric} on $Z$ to be a map $\Dist: Z\times Z\to\Rplus$ which satisfies (II), (III), and (IV), but not necessarily (I). This concept is not to be confused with the more common notion of a \emph{pseudometric}, which satisfies (I), (III), and (IV), but not necessarily (II). The term ``metametric'' was introduced by J. V\"ais\"al\"a in \cite{Vaisala}.
%%
%%If $\Dist$ is a metametric, we define its \emph{domain of reflexivity} to be the set $Z_\refl := \{x\in Z:\Dist(x,x) = 0\}$.\Footnote{In the terminology of \cite[p.19]{Vaisala}, the domain of reflexivity is the set of ``small points''.} Obviously, $\Dist$ restricted to its domain of reflexivity is a metric.
%%\end{definition}
%%
%%As in metric spaces, a sequence $(x_n)_1^\infty$ in a metametric space $(Z,\Dist)$ is called \emph{Cauchy} if $\Dist(x_n,x_m)\tendsto{n,m} 0$, and \emph{convergent} if there exists $x\in Z$ such that $\Dist(x_n,x)\to 0$. (However, see Remark \ref{remarktopologymetametrics} below.) The metametric space $(Z,\Dist)$ is called \emph{complete} if every Cauchy sequence is convergent. Using these definitions, the standard proof of the Banach contraction principle immediately yields the following:
%%
%%\begin{theorem}[Banach contraction principle for metametric spaces]\label{banachcontractionprinciple}
%%Let $(Z,\Dist)$ be a complete metametric space. Fix $0 < \lambda < 1$. If $g:Z\to Z$ satisfies
%%\[
%%\Dist(g(z),g(w)) \leq \lambda\Dist(z,w) \all z,w\in Z,
%%\]
%%then there exists a unique point $z\in Z$ so that $g(z) = z$. Moreover, for all $w\in Z$, we have $g^n(w)\to z$ with respect to the metametric $\Dist$.
%%\end{theorem}
%%%\begin{proof}
%%%Any standard proof of the Banach contraction principle (see e.g. \cite{}\internal) does not use condition 1 in the proof, so it goes through verbatim to metametric spaces.
%%%\end{proof}
%%
%%\begin{observation}
%%The fixed point coming $z$ coming from Theorem \ref{banachcontractionprinciple} must lie in the domain of reflexivity $Z_\refl$.
%%\end{observation}
%%\begin{proof}
%%\[
%%\Dist(z,z) = \Dist(g(z),g(z)) \leq \lambda \Dist(z,z),
%%\]
%%and thus $\Dist(z,z) = 0$.
%%\end{proof}
%%
%%Recall that a metric is said to be \emph{compatible} with a topology if that topology is equal to the topology induced by the metric. We now generalize this concept by introducing the notion of compatibility between a topology and a metametric.
%%
%%\begin{definition}
%%\label{definitioncompatiblemetametric}
%%Let $(Z,\Dist)$ be a metametric space. A topology $\scrT$ on $Z$ is \emph{compatible} with the metametric $\Dist$ if for every $\xi\in Z_\refl$, the collection
%%\begin{equation}
%%\label{balls}
%%\left\{B_\Dist(\xi,r) := \{y\in Z : \Dist(\xi,y) < r\} : r > 0\right\}
%%\end{equation}
%%forms a neighborhood base for $\scrT$ at $\xi$.
%%\end{definition}
%%
%%Note that unlike a metric, a metametric may have multiple topologies with which it is compatible.\Footnote{The topology considered in \cite[p.19]{Vaisala} is the finest topology compatible with a given metametric.} The metametric is viewed as determining a neighborhood base for points in the domain of reflexivity; neighborhood bases for other points must arise from some other structure. In the case we are interested in, namely the case where the underlying space for the metametric is the Gromov closure of a hyperbolic metric space $X$, the topology on the complement of the domain of reflexivity will come from the original metric $\dist$ on $X$.
%%
%%\begin{remark}
%%\label{remarktopologymetametrics}
%%If $(Z,\Dist)$ is a metametric space with a compatible topology $\scrT$, then there are two notions of what it means for a sequence $(x_n)_1^\infty$ in $Z$ to converge to a point $x\in Z$: the sequence may converge with respect to the topology $\scrT$, or it may converge with respect to the metametric (i.e. $\Dist(x_n,x)\to 0$). The relation between these notions is as follows: $x_n\to x$ with respect to the metametric $\Dist$ if and only if both of the following hold: $x_n\to x$ with respect to the topology $\scrT$, and $x\in Z_\refl$.
%%\end{remark}
%%
%%\begin{remark}
%%If a metametric $\Dist$ on a set $Z$ is compatible with a topology $\scrT$, then the metric $\Dist\given Z_\refl$ is compatible with the topology $\scrT\given Z_\refl$. However, the converse does not necessarily hold.
%%\end{remark}
%%
%%
%%For the remainder of this section, we fix a hyperbolic metric space $X$, and we let $\scrT$ be the topology on $\bord X$ introduced in \6\ref{subsubsectiontopologyclX}. We will consider various metrics and metametrics on $\bord X$ which are compatible with the topology $\scrT$.
%%
%%\subsubsection{The visual metametric based at a point $\notzero\in X$}
%%The first metametric that we will consider is designed to emulate the Euclidean or ``spherical'' metric on the boundary of the ball model $\B$. Recall from Lemma \ref{lemmagromovextension} that
%%\begin{repequation}{euclideanball2}
%%\frac12\|\yy - \xx\| = e^{-\lb \xx|\yy\rb_\0} \text{ for all $\xx,\yy\in\del\B$.}
%%\end{repequation}
%%The metric $(\xx,\yy)\mapsto \frac12\|\yy - \xx\|$ can be thought of as ``seen from $\0$''. The expression on the right hand side makes sense if $\xx,\yy\in\bord\B$, and defines a metametric on $\bord\B$. Moreover, the formula can be generalized to an arbitrary strongly hyperbolic metric space:
%%\begin{observation}
%%\label{observationDist}
%%If $X$ is a strongly hyperbolic metric space, then for each $\notzero\in X$ the map $\Dist_\notzero:\bord X\times\bord X\to \Rplus$ defined by
%%\begin{equation}
%%\label{distancedef}
%%\Dist_\notzero(x,y) := e^{-\lb x|y\rb_\notzero}
%%\end{equation}
%%is a complete metametric on $\bord X$. This metametric is compatible with the topology $\scrT$; moreover, its domain of reflexivity is $\del X$.
%%\end{observation}
%%\begin{proof}
%%Reverse reflexivity and the fact that $(\bord X)_\refl = \del X$ follow directly from Observation \ref{observationyzinfinity}; symmetry follows from (a) of Proposition \ref{propositionbasicidentities} together with Corollary \ref{corollaryboundaryasymptotic}; the triangle inequality follows from the definition of strong hyperbolicity together with Corollary \ref{corollaryboundaryasymptotic}.
%%
%%To show that $\Dist_\notzero$ is complete, suppose that $(x_n)_1^\infty$ is a Cauchy sequence in $X$. Applying \eqref{distancedef}, we see that $\lb x_n | x_m\rb_\notzero \tendsto{n,m} \infty$, i.e. $(x_n)_1^\infty$ is a Gromov sequence. Letting $\xi = [(x_n)_1^\infty]$, we have $x_n\to \xi$ in the $\Dist_{b,\notzero}$ metametric. Thus every Cauchy sequence in $X$ converges in $\bord X$. Since $X$ is dense in $\bord X$, a standard approximation argument shows that $\bord X$ is complete.
%%
%%Given $\xi\in (\bord X)_\refl = \del X$, the collection \eqref{balls} is equal to the collection \eqref{Ntxibase}, and is therefore a neighborhood base for $\scrT$ at $\xi$. Thus $\Dist_\notzero$ is compatible with $\scrT$.
%%\end{proof}
%%
%%Next, we drop the assumption that $X$ is strongly hyperbolic. Fix $b > 1$ and $\notzero\in X$, and consider the function
%%\begin{equation}
%%\label{Distdef}
%%\Dist_{b,\notzero}(x,y) = \inf_{(x_i)_0^n} \sum_{i = 0}^{n - 1} b^{-\lb x_i | x_{i + 1}\rb_\notzero},
%%\end{equation}
%%where the infimum is taken over finite sequences $(x_i)_0^n$ satisfying $x_0 = x$ and $x_n = y$.
%%
%%\begin{proposition}
%%\label{propositionDist}
%%If $b > 1$ is sufficiently close to $1$, then for each $\notzero\in X$, the function $\Dist_{b,\notzero}$ defined by \eqref{Distdef} is a complete metametric on $\bord X$ satisfying the following inequality:
%%\begin{equation}
%%\label{distanceasymptotic}
%%b^{-\lb x|y\rb_\notzero} / 4 \leq \Dist_{b,\notzero}(x,y) \leq b^{-\lb x|y\rb_\notzero}.
%%\end{equation}
%%This metametric is compatible with the topology $\scrT$; moreover, its domain of reflexivity is $\del X$.
%%\end{proposition}
%%We will refer to $\Dist_{b,\notzero}$ as the ``visual (meta)metric from the point $\notzero$ with respect to the parameter $b$''.
%%
%%\begin{remark}
%%The metric $\Dist_{b,\notzero}\given\del X$ has been referred to in the literature as the \emph{Bourdon metric}.
%%\end{remark}
%%
%%\begin{remark}
%%The first part of Proposition \ref{propositionDist} is \cite[Propositions 5.16 and 5.31]{Vaisala}.
%%\end{remark}
%%
%%% Since the explicit form of the metametric $\Dist_{b,\notzero}$ will be important in what follows, we include the proof of Proposition \ref{propositionDist}.
%%\begin{proof}[Proof of Proposition \ref{propositionDist}]
%%Let $\delta \geq 0$ be the implied constant in Gromov's inequality, and fix $1 < b \leq 2^{1/\delta}$. Then raising $b^{-1}$ to the power of both sides of Gromov's inequality gives
%%\[
%%b^{-\lb x|z\rb_\notzero} \leq 2\max\left(b^{-\lb x|y\rb_w}, b^{-\lb y|z\rb_w}\right),
%%\]
%%i.e. the function
%%\[
%%(x,y)\mapsto b^{-\lb x|y\rb_\notzero}
%%\]
%%satisfies the ``weak triangle inequality'' of \cite{Schroeder}. A straightforward adaptation of the proof of \cite[Theorem 1.2]{Schroeder} demonstrates \eqref{distanceasymptotic}. Condition (II) of being a metametric and the equality $(\bord X)_\refl = \del X$ now follow from Observation \ref{observationyzinfinity}. Conditions (III) and (IV) of being a metametric are immediate from \eqref{Distdef}.
%%
%%The argument for completeness is the same as in the proof of Observation \ref{observationDist}.
%%
%%Finally, given $\xi\in (\bord X)_\refl = \del X$, we observe that although the collections \eqref{balls} and \eqref{Ntxibase} are no longer equal, \eqref{distanceasymptotic} guarantees that the filters they generate are equal, which is enough to show that $\Dist_{b,\notzero}$ is compatible with $\scrT$.
%%\end{proof}
%%
%%\begin{remark}
%%\label{remarkDist}
%%If $X$ is strongly hyperbolic, then Proposition \ref{propositionDist} holds for all $1 < b\leq e$; moreover, the metametric $\Dist_{e,\notzero}$ is equal to the metametric $\Dist_\notzero$ defined in Observation \ref{observationDist}.
%%\end{remark}
%%
%%\begin{remark}
%%\label{remarkDistRtree}
%%If $(X,\dist)$ is an $\R$-tree, then for all $t > 0$, $(X,t\dist)$ is also an $\R$-tree and is therefore strongly hyperbolic (by Observation \ref{observationRtreesCAT} and Proposition \ref{propositionCATimpliesGromov}). It follows that Proposition \ref{propositionDist} holds for all $b > 1$.
%%\end{remark}
%%
%%
%%For the remainder of this section, we fix $b > 1$ close enough to $1$ so that Proposition \ref{propositionDist} holds.
%%
%%% Not clear whether it's used; also, needs a new proof.
%%%\ignore{
%%%\begin{proposition}
%%%\label{propositionDistBasymp}
%%%For all $\xx,\yy\in\bord\B$,
%%%\[
%%%\Dist_\0(\xx,\yy) \asymp_\times \max(1 - \|\xx\|,1 - \|\yy\|,\|\yy - \xx\|).
%%%\]
%%%\end{proposition}
%%%\begin{proof}
%%%Follows directly from [(iii) of Lemma \ref{lemmagromovextension}\internal] and the fact that $e^{-\dist_\B(\0,\xx)} \asymp_\times 1 - \|\xx\|$ fpr all $\xx\in\bord\B$.
%%%\end{proof}
%%%}% end ignore
%%
%%\subsubsection{The visual metric on $\bord X$}
%%
%%Although the metametric $\Dist_{b,\notzero}$ has the advantage of being directly linked to the Gromov product via \eqref{distanceasymptotic}, it is sometimes desirable to put a metric on $\bord X$, not just a metametric. We show now that such a metric can be constructed which agrees with $\Dist_{b,\notzero}$ on $\del X$.
%%
%%In the following proposition, we use the convention that $\dist(x,y) = \infty$ if $x,y\in\bord X$ and either $x\in\del X$ or $y\in\del X$.
%%
%%\begin{proposition}
%%\label{propositionwbarDist}
%%Fix $\notzero\in X$, and for all $x,y\in \bord X$ let
%%\[
%%\wbar\Dist_{b,\notzero}(x,y) = \min\big(\log(b)\dist(x,y), \Dist_{b,\notzero}(x,y)\big).
%%\]
%%Then $\wbar\Dist = \wbar\Dist_{b,\notzero}$ is a complete metric on $\bord X$ which agrees with $\Dist = \Dist_{b,\notzero}$ on $\del X$ and induces the topology $\scrT$.
%%\end{proposition}
%%As an immediate consequence we have the following result which was promised in \6\ref{subsubsectiontopologyclX}:
%%\begin{corollary}
%%The topological space $(\bord X,\scrT)$ is completely metrizable.
%%\end{corollary}
%%\begin{proof}[Proof of Proposition \ref{propositionwbarDist}]
%%Let us show that $\wbar\Dist$ is a metric. Conditions (I)-(III) are obvious. To demonstrate the triangle inequality, fix $x,y,z\in\bord X$.
%%\begin{itemize}
%%\item[(1)] If $\wbar\Dist(x,y) = \log(b)\dist(x,y)$ and $\wbar\Dist(y,z) = \log(b)\dist(y,z)$, then $\wbar\Dist(x,z) \leq \log(b)\dist(x,z) \leq \log(b)\dist(x,y) + \log(b)\dist(y,z) = \wbar\Dist(x,y) + \wbar\Dist(y,z)$. Similarly, if $\wbar\Dist(x,y) = \Dist(x,y)$ and $\wbar\Dist(y,z) = \Dist(y,z)$, then $\wbar\Dist(x,z) \leq \Dist(x,z) \leq \Dist(x,y) + \Dist(y,z) = \wbar\Dist(x,y) + \wbar\Dist(y,z)$.
%%\item[(2)] If $\wbar\Dist(x,y) = \log(b)\dist(x,y)$ and $\wbar\Dist(y,z) = \Dist(y,z)$, fix $\epsilon > 0$, and let $y = y_0,y_1,\ldots,y_n = z$ be a sequence such that
%%\[
%%\sum_{i = 0}^{n - 1}b^{-\lb y_i,y_{i + 1}\rb_\notzero} \leq \Dist(y,z) + \epsilon.
%%\]
%%Let $x_i = y_i$ for $i\geq 1$ but let $x_0 = x$. Then by (e) of Proposition \ref{propositionbasicidentities} and the inequality
%%\[
%%b^{-t} \leq s\log(b) + b^{-(t + s)} \hspace{.5 in} (s,t \geq 0),
%%\]
%%we have
%%\[
%%b^{-\lb x|y_1\rb_\notzero} \leq \log(b)\dist(x,y) + b^{-\lb y|y_1\rb_\notzero}.
%%\]
%%It follows that
%%\begin{align*}
%%\wbar\Dist(x,z) \leq \Dist(x,z) &\leq \sum_{i = 0}^{n - 1} b^{-\lb x_i|x_{i + 1}\rb_\notzero}\\
%%&= b^{-\lb x|y_1\rb_\notzero} + \sum_{i = 1}^{n - 1} b^{-\lb y_i|y_{i + 1}\rb_\notzero}\\
%%&\leq \log(b)\dist(x,y) + b^{-\lb y|y_1\rb_\notzero} + \sum_{i = 1}^{n - 1} b^{-\lb y_i|y_{i + 1}\rb_\notzero}\\
%%&\leq \log(b)\dist(x,y) + \Dist(y,z) + \epsilon = \wbar\Dist(x,y) + \wbar\Dist(y,z) + \epsilon.
%%\end{align*}
%%Taking the limit as $\epsilon$ goes to zero finishes the proof.
%%
%%\item[(3)] The third case is identical to the second.
%%\end{itemize}
%%
%%If a sequence $(x_n)_1^\infty$ is Cauchy with respect to $\wbar\Dist$, then Ramsey's theorem (for example) guarantees that some subsequence is Cauchy with respect to either $\dist$ or $\Dist$. This subsequence converges with respect to that metametric, and therefore also with respect to $\wbar\Dist$. It follows that the entire sequence converges. Thus $\wbar\Dist$ is complete.
%%
%%Finally, to show that $\wbar\Dist$ induces the topology $\scrT$, suppose that $U\subset\del X$ is open in $\scrT$, and fix $x\in U$. If $x\in X$, then $B_\dist(x,r)\subset U$ for some $r > 0$. On the other hand, by the triangle inequality $\Dist(x,y) \geq \frac12\Dist(x,x) > 0$ for all $y\in\bord X$. Letting $\w r = \min(r, \frac12\Dist(x,x))$, we have $B_{\wbar\Dist}(x,\w r) \subset B_\dist(x,r) \subset U$. If $x\in\del X$, then $N_t(x)\subset U$ for some $t\geq 0$; letting $C$ be the implied constant of \eqref{distanceasymptotic}, we have $B_{\wbar\Dist}(x,e^{-t}/C) = B_\Dist(x,e^{-t}/C) \subset N_t(x) \subset U$. Thus $U$ is open in the topology generated by the $\wbar\Dist$ metric. The converse direction is similar but simpler, and will be omitted.
%%\end{proof}
%%
%%\begin{remark}
%%\label{remarkDistunifcont}
%%The proof of Proposition \ref{propositionwbarDist} actually shows more, namely that
%%\[
%%\Dist(x,z) \leq \Dist(x,y) + \wbar\Dist(y,z) \all x,y,z\in\bord X.
%%\]
%%Since $\Dist(x,x) = b^{-\dox x} = \inf_{y\in\bord X} \Dist(x,y)$, plugging in $x = z$ gives
%%\[
%%b^{-\dox x} \leq b^{-\dox y} + \wbar\Dist(x,y) \all x,y\in\bord X.
%%\]
%%\end{remark}
%%
%%\begin{remark}
%%Although the metric $\wbar\Dist$ is convenient since it induces the correct topology on $\bord X$, it is not a generalization of the Euclidean metric on the closure of a \CROSS. Indeed, when $X = \B^2$, then $\wbar\Dist$ is not bi-Lipschitz equivalent to the Euclidean metric on $\bord X$.
%%\end{remark}
%%
%%
%%
%%
%%\subsubsection{The visual metametric based at a point $\xi\in\del X$}
%%Our final metametric is supposed to generalize the Euclidean metric on the boundary of the half-space model $\E$. This metric should be thought of as ``seen from the point $\infty$''.
%%
%%\begin{notation}
%%If $X$ is a hyperbolic metric space and $\xi\in\del X$, then let $\EE_\xi := \bord X\butnot\{\xi\}$.
%%\end{notation}
%%
%%Since we have not yet introduced a formula analogous to \eqref{euclideanball2} for the Euclidean metric on $\del\E\butnot\{\infty\}$, we will instead motivate the visual metametric based at a point $\xi\in\del X$ by considering a sequence $(\notzero_n)_1^\infty$ in $X$ converging to $\xi$, and taking the limits of their visual metametrics.
%%
%%
%%
%%In fact, $\Dist_{b,\notzero_n}(y_1,y_2)\to 0$ for every $y_1, y_2\in\EE_\xi$. Some normalization is needed.
%%
%%
%%\begin{lemma}
%%\label{lemmaeuclideanmetametric}
%%Fix $\zero\in X$, and suppose $\notzero_n\to\xi\in\del X$. Then for each $y_1, y_2\in\EE_\xi$,
%%\[
%%b^{\dox{\notzero_n}} \Dist_{b,\notzero_n} (y_1,y_2) \tendsto{n,\times} b^{-[\lb y_1 | y_2 \rb_\zero - \sum_{i = 1}^2 \lb y_i | \xi \rb_\zero]} .
%%\]
%%with $\tendsto n$ if $X$ is strongly hyperbolic.
%%\end{lemma}
%%\begin{proof}
%%\begin{align*}
%%b^{\dox{\notzero_n}} \Dist_{b,\notzero_n}(y_1,y_2)
%%&\asymp_\times b^{-[\lb y_1 | y_2 \rb_{\notzero_n} - \dox{\notzero_n}]} \by{\eqref{distanceasymptotic}}\\
%%&\asymp_\times b^{-[\lb y_1 | y_2 \rb_\zero - \sum_{i = 1}^2 \lb y_i | \notzero_n \rb_\zero]} \by{(k) of Proposition \ref{propositionbasicidentities}}\\
%%&\tendsto{n,\times} b^{-[\lb y_1 | y_2 \rb_\zero - \sum_{i = 1}^2 \lb y_i | \xi \rb_\zero]}. \by{Lemma \ref{lemmanearcontinuity}}
%%\end{align*}
%%In each step, equality holds if $X$ is strongly hyperbolic.
%%\end{proof}
%%
%%We can now construct the visual metametric based at a point $\xi\in\del X$.
%%
%%\begin{proposition}
%%\label{propositioneuclideanmetametric}
%%For each $\zero\in X$ and $\xi\in\del X$, there exists a complete metametric $\Dist_{b,\xi,\zero}$ on $\EE_\xi$ satisfying
%%\begin{equation}
%%\label{euclideanmetametric}
%%\Dist_{b,\xi,\zero}(y_1,y_2) \asymp_\times b^{-[\lb y_1 | y_2 \rb_\zero - \sum_{i = 1}^2 \lb y_i | \xi \rb_\zero]} \ ,
%%\end{equation}
%%with equality if $X$ is strongly hyperbolic. The metametric $\Dist_{b,\xi,\zero}$ is compatible with the topology $\scrT\given\EE_\xi$; moreover, a set $S\subset\EE_\xi$ is bounded in the metametric $\Dist_{b,\xi,\zero}$ if and only if $\xi\notin\cl S$.
%%\end{proposition}
%%\begin{remark}
%%The metric $\Dist_{b,\xi,\zero}\given\EE_\xi\cap\del X$ has been referred to in the literature as the \emph{Hamenst\"adt metric}.
%%\end{remark}
%%
%%\begin{proof}[Proof of Proposition \ref{propositioneuclideanmetametric}]
%%Let
%%\[
%%\Dist_{b,\xi,\zero}(y_1,y_2) = \limsup_{\notzero\to\xi}b^{\dox{\notzero}}\Dist_\notzero(y_1,y_2) := \sup\left\{\limsup_{n\to\infty} b^{\dox{\notzero_n}}\Dist_{\notzero_n}(y_1,y_2) : \notzero_n \tendsto n \xi\right\}
%%\]
%%Since the class of metametrics is closed under suprema and limits, it follows that $\Dist_{b,\xi,\zero}$ is a metametric. The asymptotic \eqref{euclideanmetametric} follows from Lemma \ref{lemmaeuclideanmetametric}.
%%
%%For the remainder of this proof, we write $\Dist = \Dist_{b,\zero}$ and $\Dist_\xi = \Dist_{b,\xi,\zero}$.
%%
%%For all $x\in\EE_\xi$,
%%\begin{equation}
%%\label{euclideancomparison}
%%\Dist_\xi(\zero,x) \asymp_\times b^{-[\lb \zero|x\rb_\zero - \lb \zero|\xi\rb_\zero - \lb x|\xi\rb_\zero]} = b^{\lb x|\xi\rb_\zero} \asymp_\times \frac{1}{\Dist(x,\xi)},
%%\end{equation}
%%with equality if $X$ is strongly hyperbolic. It follows that for any set $S\subset\EE_\xi$, the function $\Dist_\xi(\zero,\cdot)$ is bounded on $S$ if and only if the function $\Dist(\cdot,\xi)$ is bounded from below on $S$. This demonstrates that $S$ is bounded in the $\Dist_\xi$ metametric if and only if $\xi\notin\cl S$.
%%
%%
%%Let $(x_n)_1^\infty$ be a Cauchy sequence with respect to $\Dist_\xi$. Since $\Dist \lesssim_\times \Dist_\xi$, it follows that $(x_n)_1^\infty$ is also Cauchy with respect to the metametric $\Dist$, so it converges to a point $x\in\bord X$ with respect to $\Dist$. If $x\in\EE_\xi$, then we have
%%\[
%%\Dist_\xi(x_n,x) \asymp_\times b^{\lb x_n | \xi\rb_\zero + \lb x | \xi\rb_\zero} \Dist(x_n,x) \tendsto{n,\times} b^{2\lb x | \xi\rb_\zero} 0 = 0.
%%\]
%%On the other hand, if $x = \xi$, then the sequence $(x_n)_1^\infty$ is unbounded in the $\Dist_\xi$ metametric, which contradicts the fact that it is Cauchy. Thus $\Dist_\xi$ is complete.
%%
%%Finally, given $\eta\in (\EE_\xi)_\refl = \EE_\xi\cap\del X$, consider the filters $\FF_1$ and $\FF_2$ generated by the collections $\{B_{\Dist}(\eta,r) : r > 0\}$ and $\{B_{\Dist_\xi}(\eta,r) : r > 0\}$, respectively. Since $\Dist \lesssim_\times \Dist_\xi$, we have $\FF_2\subset\FF_1$. Conversely, since $B_{\Dist_\xi}(\eta,1)$ is bounded in the $\Dist_\xi$ metametric, its closure does not contain $\xi$, and so the function $\lb \cdot | \xi\rb_\zero$ is bounded on this set. Thus $\Dist \asymp_{\times,\eta} \Dist_\xi$ on $B_{\Dist_\xi}(\eta,1)$. Letting $C$ be the implied constant of the asymptotic, we have $B_{\Dist_\xi}(\eta,\min(r,1)) \subset B_{\Dist}(\eta,Cr)$, which demonstrates that $\FF_1\subset\FF_2$. Thus $\Dist_\xi$ is compatible with the topology $\scrT\given\EE_\xi$.
%%\end{proof}
%%
%%From Lemma \ref{lemmaeuclideanmetametric} and Proposition \ref{propositioneuclideanmetametric} it immediately follows that
%%\begin{equation}
%%\label{euclideanmetrictendsasymp}
%%b^{d(\zero,\notzero_n)} \Dist_{b,\notzero_n} (y_1,y_2) \tendsto{n,\times} \Dist_{b,\xi,\zero}(y_1,y_2) 
%%\end{equation}
%%whenever $(\notzero_n)_1^\infty\in\xi$.
%%
%%\begin{remark}
%%It is not clear whether a result analogous to Proposition \ref{propositionwbarDist} holds for the metametric $\Dist_{b,\xi,\zero}$. A straightforward adaptation of the proof of Proposition \ref{propositioneuclideanmetametric} does not work, since
%%\[
%%b^{\dox{\notzero_n}}\wbar\Dist_{b,\notzero_n}(x,y)
%%= \min(b^{\dox{\notzero_n}}\dist(x,y)) , b^{\dox{\notzero_n}}\Dist_{b,\notzero_n}(x,y)
%%\tendsto{n,\times} \min(\infty, \Dist_{b,\xi,\zero}(x,y))
%%= \Dist_{b,\xi,\zero}(x,y).
%%\]
%%\end{remark}
%%
%%We finish this section by describing the relation between the visual metametric based at $\infty$ and the Euclidean metric on the boundary of the half-space model $\E$:
%%
%%\begin{proposition}[Cf. Figure \ref{figurehamenstadt}]
%%\label{propositionhamenstadt}
%%Let $X = \E = \E^\alpha$, let $\zero = (1,\0)\in X$, and fix $\xx,\yy\in \EE_\infty = \E\cup\BB$. We have
%%\begin{equation}
%%\label{hamenstadt}
%%\Dist_{e,\infty,\zero}(\xx,\yy) \asymp_\times \max(x_1,y_1,\|\yy - \xx\|),
%%\end{equation}
%%with equality if $\xx,\yy\in\BB = \del\E\butnot\{\infty\}$.
%%\end{proposition}
%%\begin{proof}
%%First suppose that $\xx,\yy\in\E$. By (h) of Proposition \ref{propositionbasicidentities},
%%\begin{align*}
%%\Dist_{e,\infty,\zero}(\xx,\yy)
%%&=_\pt \exp\left(\frac12\big(\dist(\xx,\yy) + \busemann_\infty(\zero,x) + \busemann_\infty(\zero,\yy)\big)\right)\\
%%&=_\pt \sqrt{x_1 y_1} \exp\left(\frac12\left(\cosh^{-1}\left(1 + \frac{\|\yy - \xx\|^2}{2 x_1 y_1}\right)\right)\right) \by{\eqref{distE} and Proposition \ref{propositionbusemannE}}\\
%%&\asymp_\times \sqrt{x_1 y_1} \sqrt{1 + \frac{\|\yy - \xx\|^2}{2 x_1 y_1}} \since{$e^{t/2} \asymp_\times \sqrt{\cosh(t)}$}\\
%%&=_\pt \sqrt{x_1 y_1 + \|\yy - \xx\|^2}\\
%%&\asymp_\times \max(\sqrt{x_1 y_1}, \|\yy - \xx\|).
%%\end{align*}
%%Since $\sqrt{x_1 y_1}\leq \max(x_1,y_1)$, this demonstrates the $\lesssim$ direction of \eqref{hamenstadt}. Since $y_1 \leq x_1 + \|\yy - \xx\|$ and $x_1 \leq y_1 + \|\yy - \xx\|$, we have
%%\[
%%\max(x_1,y_1) \lesssim_\times \max(\min(x_1,y_1),\|\yy - \xx\|) \leq \max(\sqrt{x_1 y_1}, \|\yy - \xx\|)
%%\]
%%which demonstrates the reverse inequality. Thus \eqref{hamenstadt} holds for $\xx,\yy\in\E$; a continuity argument demonstrates \eqref{hamenstadt} for $\xx,\yy\in\EE_\infty$.
%%
%%If $\xx,\yy\in\BB$, then
%%\begin{align*}
%%\Dist_{e,\infty,\zero}(\xx,\yy)
%%&= \lim_{a,b\to 0} \sqrt{a b} \exp\left(\frac12\left(\cosh^{-1}\left(1 + \frac{\|\yy - \xx\|^2}{2 a b}\right)\right)\right) \noreason\\
%%&= \lim_{a,b\to 0} \sqrt{a b} \sqrt{2\left(1 + \frac{\|\yy - \xx\|^2}{2 a b}\right)} \since{$\lim_{t\to\infty} e^{t/2}/\sqrt{2\cosh(t)} = 1$}\\
%%&= \lim_{a,b\to 0} \sqrt{2\left(ab + \frac{\|\yy - \xx\|^2}{2}\right)} = \sqrt{\|\yy - \xx\|^2} = \|\yy - \xx\|.\noreason
%%\end{align*}
%%\end{proof}
%%
%%\begin{figure}
%%\begin{center}
%%\begin{tikzpicture}[line cap=round,line join=round,>=triangle 45,x=1.0cm,y=1.0cm]
%%\clip(-3.801,-0.18) rectangle (3.96,4.50);
%%\draw (-3.0,0.0)-- (3.0,0.0);
%%\draw (0.44800000000000006,-0.0)-- (1.76,1.48);
%%\draw [dashed] (-1.52,2.22)-- (-1.52,0);
%%\draw [dashed] (1.76,1.48)-- (1.78,0);
%%\draw (0.44800000000000006,-0.0)-- (-1.52,2.22);
%%\begin{scriptsize}
%%%\draw[color=black] (0.05368624806072046,-0.1412911554774228) node {$a$};
%%%\draw [fill=black] (0.0,4.0) circle (.75pt);
%%\draw[color=black] (0.13316978586974992,4.270045192923707) node {$\infty$};
%%\draw [fill=black] (-1.52,2.22) circle (.75pt);
%%\draw[color=black] (-1.7346933526424422,2.5810200144818323) node {$\xx$};
%%\draw [fill=black] (1.76,1.48) circle (.75pt);
%%\draw[color=black] (1.961291155477427,1.7663137519392818) node {$\yy$};
%%%\draw[color=black] (-1.7346933526424422,-0.1810329243819375) node {$H$};
%%%\draw[color=black] (-1.834047774903729,1.1304454494670468) node {$d$};
%%\draw [fill=black] (1.78,-6.938893903907228) circle (1.5pt);
%%%\draw[color=black] (1.961291155477427,-0.16116203992968017) node {$I$};
%%%\draw[color=black] (2.0209038088341993,0.8721239515877014) node {$e$};
%%\end{scriptsize}
%%\end{tikzpicture}
%%\caption{The Hamenst\"adt distance $\Dist_{e,\infty,\zero}(\xx,\yy)$ between two points $\xx,\yy\in\E$ is coarsely asymptotic to the maximum of the following three quantities: $x_1$, $y_1$, and $\|\yy - \xx\|$. Equivalently, $\Dist_{e,\infty,\zero}(\xx,\yy)$ is coarsely asymptotic to the length of the shortest path which both connects $\xx$ and $\yy$ and touches $\BB$.}
%%\label{figurehamenstadt}
%%\end{center}
%%\end{figure}
%%
%%\begin{corollary}[Cf. {\cite[Fig. 5]{Sullivan_entropy}}]
%%\label{corollarysullivansformula}
%%For $\xx,\yy\in\E$,
%%\[
%%e^{\dist(\xx,\yy)} \asymp_\times \frac{\max(x_1^2,y_1^2,\|\yy - \xx\|^2)}{x_1 y_1}\cdot
%%\]
%%\end{corollary}
%%\begin{proof}
%%We have
%%\[
%%\max(x_1^2,y_1^2,\|\yy - \xx\|^2) \asymp_\times \Dist_{e,\infty,\zero}(\xx,\yy)^2
%%= \exp\big(\dist(\xx,\yy) + \busemann_\infty(\zero,\xx) + \busemann_\infty(\zero,\yy)\big)
%%= x_1 y_1 e^{\dist(\xx,\yy)};
%%\]
%%rearranging completes the proof.
%%\end{proof}
%%
%%%[NOTE: Start talking about ``compatible with the topology'' metametrics in the sense introduced above. Note that we don't care about this compatibility for the Euclidean metametric (Observation \ref{observationxibounded} is more important), but it is important for the Bourdon metametric.]
%%
%%
%%
%%
%%% Note: Some deleted material here may be found in Frink's Lemma.tex.
%%
%%
%%
%%%%%%%%%%%%%%%%%%%%%%%%%%%%%%%%%%%%
%%\draftnewpage
%%\section{More about the geometry of hyperbolic metric spaces}\label{sectiongeometry2}
%%%%%%%%%%%%%%%%%%%%%%%%%%%%%%%%%%%%
%%
%%In this section we discuss various topics regarding the geometry of hyperbolic metric spaces, including metric derivatives, the Rips condition for hyperbolicity, construction of geodesic rays and lines in CAT(-1) spaces, ``shadows at infinity'', and some functions which we call ``generalized polar coordinates''. We start by introducing some conventions to apply in the remainder of the paper.
%%
%%%\setcounter{subsubsection}{0}
%%\subsection{Gromov triples}
%%%{\flushleft\textbf{Standing assumptions for Section \ref{sectiongeometry2}.}}
%%\label{standingassumptions2}
%%
%%The following definition is made for convenience of notation:
%%
%%\begin{definition}
%%\label{definitiongromovtriple}
%%A \emph{Gromov triple} is a triple $(X,\zero,b)$, where $X$ is a hyperbolic metric space, $\zero\in X$, and $b > 1$ is close enough to $1$ to guarantee for every $\notzero\in X$ the existence of a visual metametric $\Dist_{b,\notzero}$ via Proposition \ref{propositionDist} above.
%%\end{definition}
%%
%%\begin{notation}
%%Let $(X,\zero,b)$ be a Gromov triple. Given $\notzero\in X$ and $\xi\in\del X$, we will let $\Dist_\notzero = \Dist_{b,\notzero}$ be the metametric defined in Proposition \ref{propositionDist}, we will let $\wbar\Dist_\notzero = \wbar\Dist_{b,\notzero}$ be the metric defined in Proposition \ref{propositionwbarDist}, and we will let $\Dist_{\xi,\notzero} = \Dist_{b,\xi,\notzero}$ be the metametric defined in Proposition \ref{propositioneuclideanmetametric}. If $\notzero = \zero$, then we use the further shorthand $\Dist = \Dist_\zero$, $\wbar\Dist = \wbar\Dist_\zero$, and $\Dist_\xi = \Dist_{\xi,\zero}$.
%%
%%We will denote the diameter of a set $S$ with respect to the metametric $\Dist$ by $\Diam(S)$.
%%\end{notation}
%%
%%\begin{convention}
%%\label{conventiontriples}
%%For the remainder of the paper, with the exception of Section \ref{sectiondiscreteness}, all statements should be assumed to be universally quantified over Gromov triples $(X,\zero,b)$ unless context indicates otherwise.
%%\end{convention}
%%
%%\begin{convention}
%%For the remainder of the paper, whenever we make statements of the form ``Let $X = Y$'', where $Y$ is a hyperbolic metric space, we implicitly want to ``beef up'' $X$ into a Gromov triple $(X,\zero,b)$ whose underlying hyperbolic metric space is $Y$. For general $Y$, this may be done arbitrarily, but if $Y$ is strongly hyperbolic, we want to set $b = e$, and if $Y$ is a \CROSS, then we want to set $\zero = [(1,\0)]$, $\zero = \0$, or $\zero = (1,\0)$ depending on whether $Y$ is the hyperboloid model $\H$, the ball model $\B$, or the half-space model $\E$, respectively.
%%
%%For example, when saying ``Let $X = \H = \H^\infty$'', we really mean ``Let $X = \H = \H^\infty$, let $\zero = [(1,\0)]$, and let $b = e$.''
%%\end{convention}
%%
%%\begin{convention}
%%\label{conventionstandard}
%%The term ``Standard Case'' will always refer to the finite-dimensional situation where $X = \H^d$ for some $2\leq d < \infty$.
%%\end{convention}
%%
%%%%%%%%%%%%%%%%%%%%%%%%%%%%%%%%%%%%
%%\bigskip
%%\subsection{Derivatives}\label{subsectionderivatives}% Dynamical and non-dynamical
%%%%%%%%%%%%%%%%%%%%%%%%%%%%%%%%%%%%
%%
%%\subsubsection{Derivatives of metametrics}
%%
%%Let $(Z,\scrT)$ be a perfect topological space, and let $\Dist_1$ and $\Dist_2$ be two metametrics on $Z$. The \emph{metric derivative} of $\Dist_1$ with respect to $\Dist_2$ is the function $\Dist_1/\Dist_2:Z\to[0,\infty]$ defined by
%%\[
%%\frac{\Dist_1}{\Dist_2}(z) := \lim_{w\to z}\frac{\Dist_1(z,w)}{\Dist_2(z,w)},
%%\]
%%assuming the limit exists. If the limit does not exist, then we can speak of the \emph{upper} and \emph{lower} derivatives; these will be denoted $(\Dist_1/\Dist_2)^*(z)$ and $(\Dist_1/\Dist_2)_*(z)$, respectively. Note that the chain rule for metric derivatives takes the following form:
%%\[
%%\frac{\Dist_1}{\Dist_3} = \frac{\Dist_1}{\Dist_2}\frac{\Dist_2}{\Dist_3},
%%\]
%%assuming all limits exist.
%%
%%We proceed to calculate the derivatives of the metametrics introduced in Subsection \ref{subsectionmetametrics}:
%%
%%\begin{observation}
%%\label{observationGMVT}
%%Fix $y_1,y_2\in\bord X$.
%%\begin{itemize}
%%\item[(i)] For all $\notzero_1,\notzero_2\in X$, we have
%%\begin{equation}
%%\label{GMVT1}
%%\frac{\Dist_{\notzero_1}(y_1,y_2)}{\Dist_{\notzero_2}(y_1,y_2)} \asymp_\times b^{-\frac12[\busemann_{y_1}(\notzero_1,\notzero_2) + \busemann_{y_2}(\notzero_1,\notzero_2)]}.
%%\end{equation}
%%\item[(ii)] For all $\notzero\in X$ and $\xi\in\del X$, we have
%%\begin{equation}
%%\label{GMVT2}
%%\frac{\Dist_{\xi,\notzero}(y_1,y_2)}{\Dist_\notzero(y_1,y_2)} \asymp_\times b^{-[\lb y_1|\xi\rb_\notzero + \lb y_2|\xi\rb_\notzero]}.
%%\end{equation}
%%\item[(iii)] For all $\notzero_1,\notzero_2\in X$ and $\xi\in\del X$, we have
%%\begin{equation}
%%\label{GMVT3}
%%\frac{\Dist_{\xi,\notzero_1}(y_1,y_2)}{\Dist_{\xi,\notzero_2}(y_1,y_2)} \asymp_\times b^{\busemann_\xi(\notzero_1,\notzero_2)}.
%%\end{equation}
%%\end{itemize}
%%In each case, equality holds if $X$ is strongly hyperbolic.
%%\end{observation}
%%\begin{proof}
%%(i) follows from (g) of Proposition \ref{propositionbasicidentities}, while (ii) is immediate from \eqref{distanceasymptotic} and \eqref{euclideanmetametric}. (iii) follows from \eqref{euclideanmetrictendsasymp}.
%%\end{proof}
%%
%%Combining with Lemma \ref{lemmanearcontinuity} yields the following:
%%
%%\begin{corollary}
%%\label{corollarymetricderivative}
%%Suppose that $\bord X$ is perfect. Fix $y\in\bord X$.
%%\begin{itemize}
%%\item[(i)] For all $\notzero_1,\notzero_2\in X$, we have
%%\begin{equation}
%%\label{derivative1}
%%\left(\frac{\Dist_{\notzero_1}}{\Dist_{\notzero_2}}\right)^*(y) \asymp_\times \left(\frac{\Dist_{\notzero_1}}{\Dist_{\notzero_2}}\right)_*(y) \asymp_\times b^{-\busemann_y(\notzero_1,\notzero_2)}.
%%\end{equation}
%%\item[(ii)] For all $\notzero\in X$ and $\xi\in\del X$, we have
%%\begin{equation}
%%\label{derivative2}
%%\left(\frac{\Dist_{\xi,\notzero}}{\Dist_\notzero}\right)^*(y) \asymp_\times \left(\frac{\Dist_{\xi,\notzero}}{\Dist_\notzero}\right)_*(y) \asymp_\times b^{-2\lb y|\xi\rb_\notzero}.
%%\end{equation}
%%\item[(iii)] For all $\notzero_1,\notzero_2\in X$ and $\xi\in\del X$, we have
%%\begin{equation}
%%\label{derivative3}
%%\left(\frac{\Dist_{\xi,\notzero_1}}{\Dist_{\xi,\notzero_2}}\right)^*(y) \asymp_\times \left(\frac{\Dist_{\xi,\notzero_1}}{\Dist_{\xi,\notzero_2}}\right)_*(y) \asymp_\times b^{\busemann_\xi(\notzero_1,\notzero_2)}.
%%\end{equation}
%%\end{itemize}
%%In each case, equality holds if $X$ is strongly hyperbolic.
%%\end{corollary}
%%\begin{remark}
%%In case $\bord X$ is not perfect, \eqref{derivative1} - \eqref{derivative3} may be taken as definitions. We will ignore the issue henceforth.
%%\end{remark}
%%Combining Observation \ref{observationGMVT} with Corollary \ref{corollarymetricderivative} yields the following:
%%\begin{proposition}[Geometric mean value theorem]
%%\label{propositionGMVT}
%%Fix $y_1,y_2\in\bord X$.
%%\begin{itemize}
%%\item[(i)] For all $\notzero_1,\notzero_2\in X$, we have
%%\[
%%\frac{\Dist_{\notzero_1}(y_1,y_2)}{\Dist_{\notzero_2}(y_1,y_2)} \asymp_\times \left(\frac{\Dist_{\notzero_1}}{\Dist_{\notzero_2}}(y_1)\frac{\Dist_{\notzero_1}}{\Dist_{\notzero_2}}(y_2)\right)^{1/2}.
%%\]
%%\item[(ii)] For all $\notzero\in X$ and $\xi\in\del X$, we have
%%\[
%%\frac{\Dist_{\xi,\notzero}(y_1,y_2)}{\Dist_\notzero(y_1,y_2)} \asymp_\times \left(\frac{\Dist_{\xi,\notzero}}{\Dist_\notzero}(y_1)\frac{\Dist_{\xi,\notzero}}{\Dist_\notzero}(y_2)\right)^{1/2}.
%%\]
%%\item[(iii)] For all $\notzero_1,\notzero_2\in X$ and $\xi\in\del X$, we have
%%\[
%%\frac{\Dist_{\xi,\notzero_1}(y_1,y_2)}{\Dist_{\xi,\notzero_2}(y_1,y_2)} \asymp_\times \left(\frac{\Dist_{\xi,\notzero_1}}{\Dist_{\xi,\notzero_2}}(y_1)\frac{\Dist_{\xi,\notzero_1}}{\Dist_{\xi,\notzero_2}}(y_2)\right)^{1/2}.
%%\]
%%\end{itemize}
%%In each case, equality holds if $X$ is strongly hyperbolic.
%%\end{proposition}
%%
%%\subsubsection{Derivatives of maps}
%%\label{subsubsectionderivativemaps}
%%As before, let $(Z,\scrT)$ be a perfect topological space, and now fix just one metametric $\Dist$ on $Z$. For any map $g:Z\to Z$, the \emph{metric derivative} of $G$ is the function $g':Z\to(0,\infty)$ defined by
%%\[
%%g'(z) := \frac{\Dist\circ g}{\Dist}(z) = \lim_{w\to z}\frac{\Dist(g(z),g(w))}{\Dist(z,w)}\cdot
%%\]
%%If the limit does not exist, the upper and lower metric derivatives will be denoted $\overline g'$ and $\underline g'$, respectively.
%%
%%\begin{remark}
%%\label{remarkconfusion}
%%To avoid confusion, in what follows $g'$ will always denote the derivative of an isometry $g\in\Isom(X)$ with respect to the metametric $\Dist = \Dist_{b,\zero}$, rather than with respect to any other metametric.
%%\end{remark}
%%
%%
%%\begin{proposition}
%%\label{propositionderivativebusemann}
%%For all $g\in\Isom(X)$,
%%\begin{align*}
%%\overline g'(y) \asymp_\times \underline g'(y) &\asymp_\times b^{-\busemann_y(g^{-1}(\zero),\zero)} \all y\in\bord X\\
%%\frac{\Dist(g(y_1),g(y_2))}{\Dist(y_1,y_2)} &\asymp_\times \left(\overline g'(y_1) \overline g'(y_2)\right)^{1/2} \all y_1,y_2\in\bord X,
%%\end{align*}
%%with equality if $X$ is strongly hyperbolic.
%%\end{proposition}
%%\begin{proof}
%%This follows from (i) of Corollary \ref{corollarymetricderivative}, (i) of Proposition \ref{propositionGMVT}, and the fact that $\Dist\circ g = \Dist_{g^{-1}(\zero)}$.
%%\end{proof}
%%
%%\begin{corollary}
%%\label{corollaryproductoneorig}
%%For any distinct $y_1,y_2\in\Fix(g)\cap\del X$ we have
%%\[
%%\overline g'(y_1) \overline g'(y_2) \asymp_\times 1,
%%\]
%%with equality if $X$ is strongly hyperbolic.
%%\end{corollary}
%%
%%The next proposition shows the relation between the derivative of an isometry $g\in\Isom(X)$ at a point $\xi\in\Fix(g)$ and the action on the metametric space $(\EE_\xi,\Dist_\xi)$:
%%
%%
%%
%%\begin{proposition}
%%\label{propositioneuclideansimilarity}
%%Fix $g\in\Isom(X)$ and $\xi\in\Fix(g)$. Then for all $y_1,y_2\in\EE_\xi$,
%%\[
%%\frac{\Dist_\xi(g(y_1), g(y_2))}{\Dist_\xi(y_1,y_2)} \asymp_\times \frac{1}{g'(\xi)},
%%\]
%%with equality if $X$ is strongly hyperbolic.
%%\end{proposition}
%%\begin{proof}
%%\begin{align*}
%%\frac{\Dist_\xi(g(y_1), g(y_2))}{\Dist_\xi(y_1,y_2)} &=_\pt \frac{\Dist_{\xi,g^{-1}(\zero)}(y_1, y_2)}{\Dist_{\xi,\zero}(y_1,y_2)} \\
%%&\asymp_\times b^{-\busemann_\xi(\zero,g^{-1}(\zero))} \by{\eqref{GMVT3}} \\
%%&\asymp_\times 1/g'(\xi). \by{Proposition \ref{propositionderivativebusemann}}
%%\end{align*}
%%\end{proof}
%%
%%\begin{figure}
%%\begin{center}
%%\begin{tikzpicture}[line cap=round,line join=round,>=triangle 45,x=1.0cm,y=1.0cm]
%%\clip(-5.8,-0.36) rectangle(6.01,5.53);
%%\draw (-5.0,0.0)-- (5.0,0.0);
%%\draw [<->, dash pattern=on 4pt off 2pt] (1.35,1.12)-- (1.3452937374574478,0.0);
%%\draw (0.0,5.0)-- (0.0,0.0);
%%\draw (0.0,0.0)-- (4.246442867487024,3.5745642370946995);
%%\begin{scriptsize}
%%\draw[color=black] (0.057057757531416,-0.26212992254098266) node {$g_-$};
%%\draw [fill=black] (0.0,5.0) circle (.75pt);
%%\draw[color=black] (0.04808489249598222,5.322254959074134) node {$g_+ = \infty$};
%%\draw [fill=black] (1.35,1.12) circle (.75pt);
%%\draw[color=black] (1.17,1.53) node {$g^{-1}(x)$};
%%\draw [fill=black] (0,0.47635850682121) circle (.75pt);
%%\draw[color=black] (-0.5,0.47635850682121) node {$g^{-1}(\zero)$};
%%\draw[color=black] (2.5272758356701829,0.7026278596223968) node {$\mathrm{height}(g^{-1}(x))$};
%%\draw [fill=black] (0.0,0.0) circle (.75pt);
%%%\draw[color=black] (-0.3070507823268488,-0.1848867062821242) node {$D$};
%%\draw [fill=black] (4.246442867487024,3.5745642370946995) circle (.75pt);
%%\draw[color=black] (4.4036034520894525,3.6885775833822156) node {$x$};
%%\draw [fill=black] (0.0,2.0) circle (.75pt);
%%\draw[color=black] (0.2618688112016899,2.2273323702788814) node {$\zero$};
%%\end{scriptsize}
%%\end{tikzpicture}
%%\caption{The derivative of $g$ at $\infty$ is equal to the reciprocal of the dilatation ratio of $g$. In particular, $\infty$ is an attracting fixed point if and only if $g$ is expanding, and $\infty$ is a repelling fixed point if and only if $g$ is contracting.}
%%\end{center}
%%\end{figure}
%%
%%
%%\begin{remark}
%%Proposition \ref{propositioneuclideansimilarity} can be interpreted as a geometric mean value theorem for the action of $g$ on the metametric space $(\EE_\xi,\Dist_\xi)$. Specifically, it tells us that the derivative of $g$ on this metametric space is identically $1/g'(\xi)$.
%%\end{remark}
%%
%%\begin{remark}
%%If $g'(\xi) = 1$, then Proposition \ref{propositioneuclideansimilarity} tells us that the bi-Lipschitz constant of $g$ is independent of $g$, and that $g$ is an isometry if $X$ is strongly hyperbolic. This special case will be important in Section \ref{sectionparabolic}.
%%\end{remark}
%%
%%\begin{example}
%%\label{examplegprimeinfty}
%%Suppose that $X = \E = \E^\alpha$ is the half-space model of a real \ROSS, let $\BB = \del\E\butnot\{\infty\}$, let $g(\xx) = \lambda T(\xx) + \bb$ be a similarity of $\BB$, and consider the Poincar\'e extension $\what g\in\Isom(\E)$ defined in Observation \ref{observationpoincareextension}. Clearly $\what g$ acts as a similarity on the metametric space $(\EE_\infty,\Dist_\infty)$ in the following sense: For all $y_1,y_2\in \EE_\infty$,
%%\[
%%\Dist_\infty(g(y_1),g(y_2)) = \lambda \Dist_\infty(y_1,y_2).
%%\]
%%Comparing with Proposition \ref{propositioneuclideansimilarity} shows that $g'(\infty) = 1/\lambda$.
%%\end{example}
%%
%%
%%
%%
%%\subsubsection{The dynamical derivative} \label{subsubsectiondynamicalderivative}
%%We can interpret Corollary \ref{corollarymetricderivative} as saying that the metric derivative is well-defined only up to an asymptotic in a general hyperbolic metric space (although it is perfectly well defined in a strongly hyperbolic metric space). Nevertheless, if $\xi$ is a fixed point of the isometry $g$, then we can iterate in order to get arbitrary accuracy.
%%
%%\begin{proposition}
%%\label{propositiondynamicalderivative}
%%Fix $g\in\Isom(X)$ and $\xi\in\Fix(g)$. Then
%%\[
%%g'(\xi) := \lim_{n\to\infty} \big( (\overline{g^n})'(\xi) \bigr)^{1/n} = \lim_{n\to\infty} \bigl( (\underline{g^n})'(\xi) \bigr)^{1/n}.
%%\]
%%Furthermore
%%\[
%%\underline g'(\xi)\leq g'(\xi) \leq \overline g'(\xi).
%%\]
%%\end{proposition}
%%The number $g'(\xi)$ will be called the \emph{dynamical derivative} of $g$ at $\xi$.
%%\begin{proof}[Proof of Proposition \ref{propositiondynamicalderivative}]
%%The limits converge due to the submultiplicativity and supermultiplicativity of the expressions inside the radicals, respectively. To see that they converge to the same number, note that by Corollary \ref{corollarymetricderivative}
%%\[
%%\lim_{n\to\infty} \left(\frac{(\overline{g^n})'(\xi) }{(\underline{g^n})'(\xi) }\right)^{1/n} \leq \lim_{n\to\infty} C^{1/n} = 1
%%\]
%%for some constant $C$ independent of $n$.
%%\end{proof}
%%\begin{remark}
%%Let $\beta_\xi$ denote the \emph{Busemann quasicharacter} of \cite[p.14]{CCMT}. Then $\beta_\xi$ is related to the dynamical derivative via the following formula: $g'(\xi) = b^{-\beta_\xi(g)}$.
%%\end{remark}
%%Note that although the dynamical derivative is ``well-defined'', it is not necessarily the case that the chain rule holds for any two $g,h\in\Stab(\Isom(X);\xi)$ (although it must hold up to a multiplicative coarse asymptotic). For a counterexample see \cite[Example 3.12]{CCMT}. Note that this counterexample includes the possibility of two elementa $g,h\in\Stab(\Isom(X);\xi)$ such that $g'(\xi) = h'(\xi) = 1$ but $(gh)'(\xi)\neq 1$. A sufficient condition for the chain rule to hold exactly is given in \cite[Corollary 3.9]{CCMT}.
%%
%%Despite the failure of the chain rule, the following ``iteration'' version of the chain rule holds:
%%
%%\begin{proposition}
%%\label{propositionderivativeiteration}
%%Fix $g\in\Isom(X)$ and $\xi\in\Fix(g)$. Then
%%\[
%%(g^n)'(\xi) = [g'(\xi)]^n \all n\in\Z.
%%\]
%%In particular
%%\begin{equation}
%%\label{derivativeiteration2}
%%(g^{-1})'(\xi) = \frac{1}{g'(\xi)}\cdot
%%\end{equation}
%%\end{proposition}
%%\begin{proof}
%%The only difficulty lies in establishing \eqref{derivativeiteration2}:
%%\begin{align*}
%%(g^{-1})'(\xi) = \lim_{n\to\infty} \big((\overline{g^{-n}})'(\xi)\big)^{1/n}
%%&= \exp_{1/b} \lim_{n\to\infty} \frac{1}{n} \busemann_\xi(\zero,g^n(\zero))\\
%%&= \exp_{1/b} \lim_{n\to\infty} \frac{1}{n} \busemann_\xi(g^{-n}(\zero),\zero)\\
%%&= \exp_{1/b} \left( -\lim_{n\to\infty} \frac{1}{n} \busemann_\xi(\zero,g^{-n}(\zero)) \right) \\
%%&= \frac 1{g'(\xi)}\cdot
%%\end{align*}
%%\end{proof}
%%
%%Combining with Corollary \ref{corollaryproductoneorig} yields the following:
%%
%%\begin{corollary}
%%\label{corollaryproductone}
%%For any distinct $y_1,y_2\in\Fix(g)\cap\del X$ we have
%%\[
%%g'(y_1) g'(y_2) = 1.
%%\]
%%\end{corollary}
%%
%%
%%
%%
%%We end this subsection with the following result relating the dynamical derivative with the Busemann function:
%%
%%\begin{proposition}
%%\label{propositionbusemannpreserved}
%%Fix $g\in\Isom(X)$ and $\xi\in\Fix(g)$. Then for all $x\in X$ and $n\in\Z$,
%%\[
%%\busemann_\xi(x , g^{-n}(x)) \asymp_\plus n\log_b g'(\xi).
%%\]
%%with equality if $X$ is strongly hyperbolic.
%%\end{proposition}
%%\begin{proof}
%%If $x = \zero$, then
%%\begin{align*}
%%b^{-\busemann_\xi(\zero , g^{-n}(\zero))}
%%&\asymp_\times b^{\busemann_\xi(g^{-n}(\zero),\zero)}\\
%%&\asymp_\times (\overline{g^n})'(\xi) \by{Proposition \ref{propositionderivativebusemann}}\\
%%&\asymp_\times (g^n)'(\xi) = (g'(\xi))^n.
%%\end{align*}
%%For the general case, we note that
%%\begin{align*}
%%\busemann_\xi(x , g^{-n}(x))
%%&\asymp_\plus \busemann_\xi(x,\zero) + \busemann_\xi(\zero , g^{-n}(\zero)) + \busemann_\xi(g^{-n}(\zero), g^{-n}(x))\\
%%&\asymp_\plus \busemann_\xi(x,\zero) + n\log_b g'(\xi) + \busemann_\xi(\zero,x)\\
%%&\asymp_\plus n\log_b g'(\xi).
%%\end{align*}
%%\end{proof}
%%
%%
%%
%%% Not clear whether it's used; also, needs a new proof.
%%%\ignore{
%%%\begin{proposition}
%%%\label{propositionDistHasymp}
%%%For all $\xx,\yy\in\bord\E\butnot\{\infty\}$,
%%%\[
%%%\Dist_{\infty,\ee_1}(\xx,\yy) \asymp_\times \max(x_1,y_1,\|\yy - \xx\|).
%%%\]
%%%\end{proposition}
%%%\begin{proof}
%%%Let $e_{\E,\B}:\E\to\B$ be the Cayley transform. For each $n\in\N$ let $z_n = n\ee_1$. Then
%%%\begin{align*}
%%%\Dist_{\infty,\ee_1}(\xx,\yy)
%%%&\asymp_\times \limsup_{n\to\infty} e^{\dist_\E(\ee_1,n\ee_1)}\Dist_{n\ee_1}(\xx,\yy) \by{\eqref{euclideanmetrictendsasymp}} \\
%%%&=_\pt \limsup_{n\to\infty} n \Dist_{\ee_1}(\xx/n,\yy/n) \\
%%%&=_\pt \limsup_{n\to\infty} n \Dist_\0(e_{\E,\B}(\xx/n),e_{\E,\B}(\yy/n)) \\
%%%&\asymp_\times \limsup_{n\to\infty} n \max\big(1 - \|e_{\E,\B}(\xx/n)\|, 1 - \|e_{\E,\B}(\yy/n)\|, \|e_{\E,\B}(\yy/n) - e_{\E,\B}(\xx/n)\|\big)\noreason\\
%%%&\by{Proposition \ref{propositionDistBasymp}}\\
%%%&\asymp_\times \limsup_{n\to\infty} n \max\big(x_1/n,y_1/n, \|\yy/n - \xx/n\|\big) \since{$e_{\E,\B}$ is locally bi-Lipschitz}\\
%%%&=_\pt \max(x_1,y_1,\|\yy - \xx\|).
%%%\end{align*}
%%%\end{proof}
%%%}% end ignore
%%
%%%\ignore{
%%%\subsection{Sullivan's formula}
%%%\selfnote{How does this compare with the main theorem of the Generalized Polar Coordinates section? In fact, that section might be a better fit for this theorem than where it was previously...}
%%%
%%%Consider two points $\xx = (0,z)$ and $\yy = (t,1)$ in the upper half plane $\E^2$. In \cite[Fig. 5]{Sullivan_entropy}, Sullivan gave the following asymptotic formula for the hyperbolic distance between $\xx$ and $\yy$:
%%%\begin{equation}
%%%\label{sullivansformula}
%%%\dist_\E(\xx,\yy) \asymp_\plus \log\left(1 + \left(\frac{t}{z}\right)^2\right) + \log(z) \text{ if $z\geq 1$.}
%%%\end{equation}
%%%This formula can be generalized to the setting of hyperbolic metric spaces. Fix a distinguished point $\xi\in\del X$. For any two points $x,y\in X$, we define the \emph{height of $x$ relative to $y$} to be
%%%\begin{equation}
%%%\label{page37}
%%%\height_y(x) := b^{\busemann_\xi(x,y)},
%%%\end{equation}
%%%and the \emph{relative displacement of $x$ relative to $y$} to be
%%%\[
%%%\text{disp}_y(x) := \Dist_{\xi,y}(x,y).
%%%\]
%%%If $X = \E^{d + 1}$, then $\height_\yy(\xx) = x_{d + 1}/y_{d + 1}$ and $\text{disp}_\yy(\xx) \asymp_\times \max(1,\|\yy - \xx\|/y_{d + 1})$. [Continue\internal]
%%%
%%%
%%%
%%%\begin{lemma}[Sullivan's formula\comtushar{ DSC04299}]
%%%\label{lemmasullivansformula}
%%%\[
%%%b^{\dist(x,y)} \asymp_\times \frac{\max \Bigl[ 1,\text{disp}_y(x)^2, \height_y(x)^2 \Bigr]}{\height_y(x)} =\frac{\max \Bigl[ 1, \Dist_{\xi,y}(x, y)^2, b^{2 \busemann_\xi (x,y)} \Bigr]}{b^{\busemann_\xi (x,y)}}
%%%\]
%%%\end{lemma}
%%%\begin{proof}
%%%The logarithm of the right hand side base $b$ is coarsely asymptotic to
%%%\begin{align*}
%%%\max(-\busemann_\xi(x,y) , \busemann_\xi(x,y) , 2\lb x|\xi\rb_y - \busemann_\xi(x,y))
%%% = \max(|\busemann_\xi(x,y)| , \dist(x,y)) = \dist(x,y).
%%%\end{align*}
%%%
%%%
%%%\end{proof}
%%%
%%%
%%%
%%%\begin{remark}
%%%\eqref{sullivansformula} can be deduced from Lemma \ref{lemmasullivansformula} as follows: If $\xx = (0,z)$ and $\yy = (t,1)$, then $\text{disp}_\yy(\xx) = t$ and $\height_\yy(\xx) = z$, so
%%%\[
%%%b^{\dist_\E(\xx,\yy)} \asymp_\times \frac{\max \Bigl[ 1,t^2, z^2 \Bigr]}{z} = z \max\left(1, \left(\frac tz\right)^2, \left(\frac 1z\right)^2\right).
%%%\]
%%%If in addition $z\geq 1$, then
%%%\[
%%%\dist_\E(\xx,\yy) \asymp_\plus \log(z) + \log\max\left(1, \left(\frac tz\right)^2\right) \asymp_\plus \log(z) + \log\left(1 + \left(\frac tz\right)^2\right).
%%%\]
%%%\end{remark}
%%%
%%%}% end ignore
%%
%%
%%%%%%%%%%%%%%%%%%%%%%%%%%%%%%%%%%%%
%%\bigskip
%%\subsection{The Rips condition}\label{subsectionrips}
%%%%%%%%%%%%%%%%%%%%%%%%%%%%%%%%%%%%
%%
%%In this subsection, in addition to assuming that $X$ is a hyperbolic metric space (cf. \6\ref{standingassumptions2}), we assume that $X$ is geodesic. Recall (Subsection \ref{subsectionCAT}) that $\geo xy$ denotes the geodesic segment connecting two points $x,y\in X$.
%%
%%
%%\begin{proposition}~
%%\label{propositionrips}
%%\begin{itemize}
%%\item[(i)] For all $x,y,z\in X$,
%%\[
%%\dist(z,\geo xy) \asymp_\plus \lb x|y\rb_z.
%%\]
%%\item[(ii)] (Rips' thin triangles condition) For all $x,y_1,y_2\in X$ and $z\in\geo{y_1}{y_2}$, we have
%%\[
%%\min_{i = 1}^2 \dist(z,\geo x{y_i}) \asymp_\plus 0.
%%\]
%%\end{itemize}
%%\end{proposition}
%%
%%\begin{figure}
%%\begin{center}
%%\begin{tikzpicture}[line cap=round,line join=round,>=triangle 45,x=1.0cm,y=1.0cm]
%%\clip(-6.66,-3.12) rectangle (6.24,3.31);
%%\draw(0.0,0.0) circle (3.0cm);
%%\draw [shift={(3.001597531764653,1.959673853196985)}] plot[domain=2.7288669538659853:3.9124414341882043,variable=\t]({1.0*2.037582193265268*cos(\t r)+-0.0*2.037582193265268*sin(\t r)},{0.0*2.037582193265268*cos(\t r)+1.0*2.037582193265268*sin(\t r)});
%%\draw (0.0,-0.0)-- (1.309924961105806,0.823898739896739);
%%\begin{scriptsize}
%%\draw[color=black] (0.75,0.23) node {$\nwarrow$};
%%\draw (.85,0.3) node[anchor=north west] {$\asymp_\plus \langle x,y \rangle_z$};
%%\draw [fill=black] (0.0,-0.0) circle (.75pt);
%%\draw[color=black] (-0.19452753572169953,-0.12825951916697592) node {$z$};
%%\draw [fill=black] (1.135109224427748,2.7769636383321687) circle (.75pt);
%%\draw[color=black] (1.252879052282513,3.0161065168421803) node {$x$};
%%\draw [fill=black] (1.54,0.54) circle (.75pt);
%%\draw[color=black] (1.7519847722839659,0.5372147741682952) node {$y$};
%%\end{scriptsize}
%%\end{tikzpicture}
%%\caption{An illustration of Proposition \ref{propositionrips}(i).}
%%\end{center}
%%\end{figure}
%%
%%
%%\noindent In fact, the thin triangles condition is equivalent to hyperbolicity; see e.g. \cite[Proposition III.H.1.22]{BridsonHaefliger}.
%%\begin{proof}~
%%\begin{itemize}
%%\item[(i)] By the intermediate value theorem, there exists $w\in\geo xy$ such that $\lb x|z\rb_w = \lb y|z\rb_w$. Applying Gromov's inequality gives $\lb x|z\rb_w = \lb y|z\rb_w \lesssim_\plus \lb x|y\rb_w = 0$. Now (k) of Proposition \ref{propositionbasicidentities} shows that $\dist(z,\geo xy) \leq \dist(z,w) \lesssim_\plus \lb x|y\rb_z$. The other direction is immediate, since for each $w\in\geo xy$, we have $\lb x|y\rb_w = 0$, and so (d) of Proposition \ref{propositionbasicidentities} gives $\lb x|y\rb_w \leq \dist(z,w)$.
%%\item[(ii)] This is immediate from (i), Gromov's inequality, and the equation $\lb y_1|y_2\rb_z = 0$.
%%\end{itemize}
%%\end{proof}
%%
%%The next lemma demonstrates the correctness of the intuitive notion that if two points are close to each other, then the geodesic connecting them should not be very large.
%%
%%\begin{lemma}
%%\label{lemmageodesicdiameter}
%%Fix $x_1,x_2\in \bord X$. We have
%%\[
%%\Diam(\geo{x_1}{x_2}) \asymp_\times \Dist(x_1,x_2).
%%\]
%%\end{lemma}
%%\begin{proof}
%%It suffices to show that if $y\in\geo{x_1}{x_2}$, then
%%\[
%%\Dist(y,\{x_1,x_2\}) \lesssim_\times \Dist(x_1,x_2).
%%\]
%%Indeed, by the thin triangles condition, we may without loss of generality suppose that $\dist(y, \geo{\zero}{x_1}) \asymp_\plus 0$. Write $\dist(y,z) \asymp_\plus 0$ for some $z\in \geo\zero{x_1}$. Then
%%\[
%%\Dist(x_1,y) \asymp_\times \Dist(x_1,z) = e^{-\dox z} \asymp_\times e^{-\dox y} \leq e^{-\dist(\zero,\geo{x_1}{x_2})} \asymp_\times e^{-\lb x_1|x_2\rb_\zero} \asymp_\times \Dist(x_1,x_2).
%%\]
%%\end{proof}
%%
%%
%%%\ignore{
%%%
%%%\comdavid{I don't remember what the point of the following lemma is... maybe it's a first step towards proving the ``finite sets roughly isometrically embed into $\R$-trees'' theorem?}
%%%
%%%\begin{lemma}
%%%Suppose that four points $x_1,x_2,y_1,y_2$ satisfy
%%%\[
%%%\dist(x_1,y_i) \asymp_\plus \dist(x_2,y_i) \;\; i = 1,2
%%%\]
%%%and
%%%\[
%%%\dist(x_1,x_2) \asymp_\plus \dist(x_1,y_1) + \dist(x_2,y_1).
%%%\]
%%%Then $\dist(y_1,y_2) \asymp_\plus 0$.
%%%\end{lemma}
%%%}% end ignore
%%
%%
%%
%%%%%%%%%%%%%%%%%%%%%%%%%%%%%%%%%%%%
%%\bigskip
%%\subsection{Geodesics in CAT(-1) spaces} \label{subsectiongeodesicsCAT}
%%%%%%%%%%%%%%%%%%%%%%%%%%%%%%%%%%%%
%%
%%\begin{observation}
%%Any isometric embedding $\pi:\CO t\infty\to X$ extends uniquely to a continuous map $\pi:[t,\infty]\to \bord X$. Similarly, any isometric embedding $\pi:(-\infty,+\infty)\to X$ extends uniquely to a continuous map $\pi:[-\infty,+\infty]\to \bord X$. Abusing terminology, we will also call the extended maps ``isometric embeddings''.
%%\end{observation}
%%
%%
%%
%%\begin{definition}
%%\label{definitiongeodesic}
%%Fix $x\in X$ and $\xi,\eta\in\del X$.
%%\begin{itemize}
%%\item A \emph{geodesic ray} connecting $x$ and $\xi$ is the image of an isometric embedding $\pi:[0,\infty]\to X$ satisfying
%%\[
%%\pi(0) = x,\hspace{.5 in}\pi(\infty) = \xi.
%%\]
%%\item A \emph{geodesic line} or \emph{bi-infinite geodesic} connecting $\xi$ and $\eta$ is the image of an isometric embedding $\pi:[-\infty,+\infty]\to X$ satisfying
%%\[
%%\pi(-\infty) = \xi,\hspace{.5 in}\pi(+\infty) = \eta.
%%\]
%%\end{itemize}
%%When we do not wish to distinguish between geodesic segments (cf. Subsection \ref{subsectionCAT}), geodesic rays, and geodesic lines, we shall simply call them geodesics. For $x,y\in\bord X$, any geodesic connecting $x$ and $y$ will be denoted $\geo xy$.
%%\end{definition}
%%
%%\begin{notation}
%%\label{notationgeopqt2}
%%Extending Notation \ref{notationgeopqt}, if $\geo x\xi$ is the image of the isometric embedding $\pi:[0,\infty]\to X$, then for $t\in[0,\infty]$ we let $\geo x\xi_t = \pi(t)$, i.e. $\geo x\xi_t$ is the unique point on the geodesic ray $\geo x\xi$ such that $\dist(x,\geo x\xi_t) = t$.
%%\end{notation}
%%
%%%\begin{definition}
%%%A sequence of geodesics $(\geo{x_n}{y_n})_1^\infty$ \emph{converges} to a geodesic $\geo xy$ if
%%%\begin{itemize}
%%%\item[(I)] $x_n\to x$ and $y_n\to y$, and
%%%\item[(II)] Suppose that $(n_k)_1^\infty$ and $(z_k)_1^\infty$ are sequences such that $z_k\in \geo{x_{n_k}}{y_{n_k}}$. Then there is a sequence $(k_\ell)_1^\infty$ and a point $z\in\geo xy$ so that $z_{n_{k_\ell}}\to z%$.
%%%\end{itemize}
%%%\end{definition}
%%
%%The main goal of this subsection is to prove the following:
%%
%%\begin{proposition}
%%\label{propositiongeodesicconvergence}
%%Suppose that $X$ is a complete CAT(-1) space. Then:
%%\begin{itemize}
%%\item[(i)] For any two distinct points $x,y\in\bord X$, there is a unique geodesic $\geo xy$ connecting them.
%%\item[(ii)] Suppose that $(x_n)_1^\infty$ and $(y_n)_1^\infty$ are sequences in $\bord X$ which converge to points $x_n\to x\in\bord X$ and $y_n\to y\in\bord X$, with $x\neq y$. Then $\geo{x_n}{y_n}\to\geo xy$ in the Hausdorff metric on $(\bord X,\wbar\Dist)$. If $x = y$, then $\geo{x_n}{y_n}\to \{x\}$ in the Hausdorff metric.% following sense: Suppose that $(z_n)_1^\infty$ is a sequence such that $z_n\in \geo{x_n}{y_n}$. Then there is an increasing sequence $(n_k)_1^\infty$ and a point $z\in\geo xy$ so that $z_{n_k}\tendsto k z$.
%%\end{itemize}
%%\end{proposition}
%%
%%\begin{definition}
%%\label{definitionregularlygeodesic}
%%A hyperbolic metric space $X$ satisfying the conclusion of Proposition \ref{propositiongeodesicconvergence} will be called \emph{regularly geodesic}.
%%\end{definition}
%%
%%\begin{remark}
%%The existence of a geodesic connecting any two points in $\bord X$ was proven in \cite[Proposition 0.2]{Buyalo} under the weaker hypothesis that $X$ is a Gromov hyperbolic complete CAT(0) space. However, this weaker hypothesis does not imply the uniqueness of such a geodesic, nor does it imply (ii) of Proposition \ref{propositiongeodesicconvergence}, as shown by the following example:
%%\end{remark}
%%
%%\begin{example}[A proper and uniquely geodesic hyperbolic CAT(0) space which is not regularly geodesic]
%%\label{examplestrip}
%%Let
%%\[
%%X = \{\xx\in\R^2: x_2\in[0,1]\}
%%\]
%%be interpreted as a subspace of $\R^2$ with the usual metric. Then $X$ is hyperbolic, proper, and uniquely geodesic, but is not regularly geodesic.
%%\end{example}
%%\begin{proof}
%%It is hyperbolic since it is roughly isometric to $\R$. It is uniquely geodesic since it is a convex subset of $\R^2$. It is proper because it is a closed subset of $\R^2$. It is not regularly geodesic because if we write $\del X = \{\xi_+,\xi_-\}$, then the two points $\xi_+$ and $\xi_-$ have infinitely many distinct geodesics connecting them: for each $t\in[0,1]$, $\R\times\{t\}$ is a geodesic connecting $\xi_+$ and $\xi_-$.
%%\end{proof}
%%
%%
%%% What follows is an attempted proof which uses the notion of standard parameterization from the start; there is an error though.
%%
%%%The proof of Proposition \ref{propositiongeodesicconvergence} will rely on the notion of the \emph{standard parameterization} of a geodesic. For any pair of distinct points $x,y\in\bord X$, the standard parameterization of %the geodesic $\geo xy$ is the unique isometry $\pi:[t,s]\to \geo xy$ satisfying
%%%\[
%%%r = \lb \pi(r) | y\rb_\zero - \lb \pi(r) | x\rb_\zero \text{ for all $r\in[t,s]$,}
%%%\]
%%%i.e. $\pi$ is the inverse of the isometric embedding $z\mapsto \lb z | y\rb_\zero - \lb z | x\rb_\zero : \geo xy\to [-\infty,+\infty]$. Note that if $x,y\in X$, then $\pi$ is uniquely determined by the requirements that $t = -\lb %\zero|y\rb_x$ and $s = \lb \zero|x\rb_y$ and $\pi(t) = x$, $\pi(s) = y$; nevertheless, if $x,y\in\bord X$, then $\pi$ still makes sense according to the above definition.
%%%
%%%If $\pi:[t,s]\to X$ is an isometric embedding, the \emph{natural extension} of $\pi$ is the map $\w\pi:[-\infty,+\infty]\to X$ defined by the equation
%%%\[
%%%\w\pi(r) = \pi(t\vee r\wedge s).\Footnote{Although the operators $\vee$ and $\wedge$ do not normally associate, they do here, because $t\leq s$.}
%%%\]
%%%\begin{lemma}
%%%\label{lemmaxyz}
%%%Fix $\epsilon > 0$. There exists $\delta = \delta_X(\epsilon) > 0$ such that if $\Delta = \Delta(x,y_1,y_2)$ is a geodesic triangle in $X$ satisfying
%%%\begin{equation}
%%%\label{y1y2delta}
%%%\wbar\Dist(y_1,y_2) \leq \delta,
%%%\end{equation}
%%%and if $\w\pi_1,\w\pi_2:[-\infty,+\infty]\to X$ are the natural extensions of the standard parameterizations of the geodesics $\geo x{y_1}$ and $\geo x{y_2}$, respectively, then
%%%\[
%%%\wbar\Dist(\w\pi_1(r),\w\pi_2(r)) \leq \epsilon \all r\in [-\infty,+\infty].
%%%\]
%%%\end{lemma}
%%%
%%%\begin{proof}
%%%We prove the assertion first for $X = \H^2$ and then in general:
%%%\begin{itemize}
%%%\item[If $X = \H^2$:] Let $\epsilon > 0$, and by contradiction, suppose that for each $\delta_n = \frac 1n > 0$ there exists a triple $(x^{(n)},y_1^{(n)},y_2^{(n)})$ satisfying the hypotheses but not the conclusion of the %lemma. Since $\bord\H^2$ is compact, there exists a convergent subsequence $(x^{(n_k)},y_1^{(n_k)},y_2^{(n_k)}) \to (x,y_1,y_2)\in\left(\bord\H^2\right)^3$.
%%%
%%%Since $\delta_n \to 0$, taking the limit of \eqref{y1y2delta} shows that $\wbar\Dist(y_1,y_2) = 0$, and thus $y_1 = y_2$. For each $k\in\N$ and $i = 1,2$, let $\pi_{k,i}$ be the standard parameterization of the geodesic $\geo {x^{(n_k)}}{y_i^{(n_k)}}$, and let $\w\pi_{k,i}$ be its natural extension. The continuity of the Gromov product implies that $\w\pi_{k,i}\to \w\pi_i$, where $\w\pi_i$ is the natural extension of the standard parameterization of the geodesic $\geo x{y_i}$. Moreover, this convergence is uniform since the family $(\w\pi_{k,i})_{k = 1}^\infty$ is equicontinuous. [More details needed?\internal] But since $y_1 = y_2$, we have $\w\pi_1 = \w\pi_2$. Thus the families $(\w\pi_{k,1})_1^\infty$ and $(\w\pi_{k,2})_1^\infty$ converge uniformly to the same function. For $k$ sufficiently large, this shows that the triple $(x^{(n_k)},y_1^{(n_k)},y_2^{(n_k)})$ satisfies the conclusion of Lemma \ref{lemmaxyz}, a contradiction.
%%%
%%%\item[In general:] Let $\epsilon > 0$, and fix $\w\epsilon > 0$ to be determined, depending on $\epsilon$. Let $\w\delta = \delta_{\H^2}(\w\epsilon)$, and fix $\delta > 0$ to be determined, depending on $\w\delta$. %Now suppose that $\Delta = \Delta(x,y_1,y_2)$ is a geodesic triangle in $X$ satisfying \eqref{y1y2delta}, and let $\pi_1,\pi_2$ be the standard parameterizations of $\geo x{y_1}$ and $\geo x{y_2}$, respectively. Fix $r%\in[-\infty,+\infty]$, and let $z_i = \w\pi_i(r)$. To complete the proof, we must show that $\wbar\Dist(z_1,z_2) \leq \epsilon$.
%%%
%%%By contradiction suppose not, i.e. suppose that $\wbar\Dist(z_1,z_2) > \epsilon$. Then $\Dist(x,z_i) > \epsilon/2$ for some $i = 1,2$; without loss of generality suppose $\Dist(x,z_1) > \epsilon/2$. By Rips' inequality this %implies $\dist(\zero,\w\pi_1(0)) \asymp_{\plus,\epsilon} 0$. Let $w_i = \w\pi_i(0)$. (See Figure \internalmissingref.\internalmakefigure) \comdavid{This figure should have the points $x,y_i,z_i,w_i,\zero$.}
%%%
%%%Now let $\wbar{\Delta} = \Delta(\wbar x,\wbar{y_1},\wbar{y_2})$ be a comparison triangle for $\Delta(x,y_1,y_2)$, and let $\wbar{z_1},\wbar{z_2},\wbar{w_1},\wbar{w_2}$ be the corresponding comparison points. %Without loss of generality, suppose that $\wbar{w_1} = \zero_\H$. It follows from this that $\wbar{z_i} = \wbar\pi_i(r)$ and $\wbar{w_i} = \wbar\pi_i(0)$, where $\wbar\pi_i$ is the natural extension of the standard %parameterization of $\geo{\wbar x}{\wbar{y_i}}$. \alert{ERROR: This is not true.} Moreover, $\dox{y_2} \leq \dox{w_1} + \dist(w_1,y_2) \asymp_{\plus,\epsilon} \dist(\zero_\H,\wbar{y_2})$, and so $\lb y_1|%y_2\rb_\zero \lesssim_{\plus,\epsilon} \lb \wbar{y_1}|\wbar{y_2}\rb_{\zero_\H}$, and thus
%%%\[
%%%\wbar\Dist(\wbar{y_1},\wbar{y_2}) \lesssim_{\times,\epsilon} \wbar\Dist(y_1,y_2) \leq \delta.
%%%\]
%%%Setting $\delta$ equal to $\w\delta$ divided by the implied constant, we have
%%%\[
%%%\wbar\Dist(\wbar{y_1},\wbar{y_2}) \leq \w\delta = \delta_\H(\w\epsilon).
%%%\]
%%%Thus $\wbar\Dist(\wbar{z_1},\wbar{z_2})\leq\w\epsilon$ and $\wbar\Dist(\wbar{w_1},\wbar{w_2})\leq\w\epsilon$.
%%%\begin{itemize}
%%%\item If $\dist(\wbar{z_1},\wbar{z_2})\leq\w\epsilon$, then the CAT(-1) inequality finishes the proof (as long as $\w\epsilon\leq \epsilon$). Thus, suppose that
%%%\begin{equation}
%%%\label{z1z2}
%%%\Dist(\wbar{z_1},\wbar{z_2}) \leq \w\epsilon.
%%%\end{equation}
%%%\item If $\Dist(\wbar{w_1},\wbar{w_2})\leq\w\epsilon$, then $0 = \lb \wbar{w_1}|\wbar{w_2}\rb_{\zero_\H} \geq -\log(\w\epsilon)$, a contradiction for $\w\epsilon$ sufficiently small. Thus, suppose that
%%%\begin{equation}
%%%\label{w1w2}
%%%\dist(\wbar{w_1},\wbar{w_2}) \leq\w\epsilon.
%%%\end{equation}
%%%\end{itemize}
%%%By \eqref{z1z2}, we have $\dist(\zero_\H,\wbar{z_i}) \geq -\log(\w\epsilon)$. Applying \eqref{w1w2} gives
%%%\[
%%%\lb z_i|y_i\rb_{w_i} = \dist(w_i,z_i) = \dist(\wbar{w_i},\wbar{z_i}) \gtrsim_\plus -\log(\w\epsilon).
%%%\]
%%%Applying \eqref{w1w2}, the CAT(-1) inequality, and the asymptotic $\dox{w_1} \asymp_{\plus,\epsilon} 0$, we have
%%%\[
%%%\lb z_i|y_i\rb_\zero \gtrsim_{\plus,\epsilon} -\log(\w\epsilon),
%%%\]
%%%and thus $\wbar\Dist(z_i,y_i)\lesssim_{\times,\epsilon} \w\epsilon$. Using the triangle inequality together with the assumption $\wbar\Dist(y_1,y_2)\leq\delta$, we have
%%%\[
%%%\wbar\Dist(z_1,z_2) \lesssim_{\times,\epsilon} \max(\delta,\w\epsilon).
%%%\]
%%%Setting $\w\epsilon$ equal to $\epsilon$ divided by the implied constant, and decreasing $\delta$ if necessary, completes the proof.
%%%\end{itemize}
%%%\end{proof}
%%%
%%%
%%%\begin{lemma}
%%%\label{lemmageodesicconvergence}
%%%Suppose that $(x_n)_1^\infty$ and $(y_n)_1^\infty$ are sequences in $\bord X$ which converge to points $x_n\to x\in\bord X$ and $y_n\to y\in\bord X$, with $x\neq y$. For each $n$, let $\geo{x_n}{y_n}$ be a geodesic %connecting $x_n$ and $y_n$. Then there exists a geodesic $\geo xy$ connecting $x$ and $y$ such that $\geo{x_n}{y_n} \to \geo xy$ in the Hausdorff metric. If $x = y$, then $\geo{x_n}{y_n}\to \{x\}$ in the Hausdorff %metric.
%%%\end{lemma}
%%%\begin{proof}
%%%For each $n$ let $\pi_n:[t_n,s_n]\to X$ be the standard parameterization of $\geo{x_n}{y_n}$, and for each $m,n\in\N$ let $\pi_{m,n}:[t_{m,n},s_{m,n}]\to X$ be the standard parameterization of $\geo{x_n}{y_m}$. Let $%\w\pi_n$ and $\w\pi_{m,n}$ denote their respective natural extensions. By Lemma \ref{lemmaxyz}, we have
%%%\[
%%%\wbar\Dist(\w\pi_n(r) , \w\pi_{m,n}(r)) \tendsto{m,n} 0 \text{ and } \wbar\Dist(\w\pi_{m,n}(r) , \w\pi_m(r)) \tendsto{m,n} 0
%%%\]
%%%uniformly for $r\in[-\infty,+\infty]$. The triangle inequality gives
%%%\[
%%%\wbar\Dist(\w\pi_n(r) , \w\pi_m(r)) \tendsto{m,n} 0,
%%%\]
%%%i.e. the sequence of functions $(\w\pi_n)_1^\infty$ is uniformly Cauchy. Since $(\bord X,\wbar\Dist)$ is complete, they converge uniformly to a function $\w\pi:[-\infty,+\infty]\to X$.
%%%
%%%Clearly, $\geo{x_n}{y_n} = \w\pi_n([-\infty,+\infty]) \to \w\pi([-\infty,+\infty])$ in the Hausdorff metric. We claim that $\w\pi([-\infty,+\infty])$ is a geodesic connecting $x$ and $y$. Indeed, note that
%%%\[
%%%t_n \tendsto n t := -\lb \zero|y\rb_x \text{ and } s_n \tendsto n s := \lb \zero |x\rb_y.
%%%\]
%%%For all $t < r_1 < r_2 < s$, we have $t_n < r_1 < r_2 < s_n$ for all sufficiently large $n$, which implies that
%%%\[
%%%\dist(\w\pi(r_1) , \w\pi(r_2)) = \lim_{n\to\infty} \dist(\w\pi_n(r_1) , \w\pi_n(r_2)) = \lim_{n\to\infty} (r_2 - r_1) = r_2 - r_1,
%%%\]
%%%i.e. $\w\pi\given(t,s)$ is an isometric embedding. Since $\w\pi$ is continuous (being the uniform limit of continuous functions), $\pi := \w\pi\given[t,s]$ is also an isometric embedding. A similar argument shows that $\w%\pi(r) = \pi(t)$ for all $r\leq t$, and $\w\pi(r) = \pi(s)$ for all $r\geq s$; thus $\w\pi([-\infty,+\infty]) = \pi([t,s])$ is a geodesic. To complete the proof, we must show that $\pi(t) = x$ and $\pi(s) = y$. Indeed,
%%%\[
%%%\pi(t) = \w\pi(-\infty) = \lim_{n\to\infty} \w\pi_n(-\infty) = \lim_{n\to\infty} x_n = x,
%%%\]
%%%and a similar argument shows that $\pi(s) = y$. Thus the geodesic $\pi([t,s])$ connects $x$ and $y$.
%%%\end{proof}
%%%
%%%
%%%Using Lemma \ref{lemmageodesicconvergence}, we prove Proposition \ref{propositiongeodesicconvergence}.
%%%
%%%\begin{proof}[Proof of Proposition \ref{propositiongeodesicconvergence}]~
%%%\begin{itemize}
%%%\item[(i)] For any $x,y\in\bord X$, we may find sequences $X\ni x_n\to x$ and $X\ni y_n\to y$. Applying Lemma \ref{lemmageodesicconvergence} proves the existence of a geodesic connecting $x$ and $y$. To show %uniqueness, suppose that $\geo xy_1$ and $\geo xy_2$ are two geodesics connecting $x$ and $y$. Fix sequences $\geo xy_1\ni x_n^{(1)}\to x$, $\geo xy_2\ni x_n^{(2)}\to x$, $\geo xy_1\ni y_n^{(1)}\to y$,and $\geo %xy_2\ni y_n^{(2)}\to y$. By considering the intertwined sequence, Lemma \ref{lemmageodesicconvergence} shows that the sequences $(\geo{x_n^{(1)}}{y_n^{(1)}})_1^\infty$ and $(\geo{x_n^{(2)}}{y_n^{(2)}})_1^\infty$ %converge in the Hausdorff metric to a common geodesic $\geo xy$. But clearly the former tend to $\geo xy_1$, and the latter tend to $\geo xy_2$; we must have $\geo xy_1 = \geo xy_2$.
%%%
%%%\item[(ii)] is now immediate from Lemma \ref{lemmageodesicconvergence}.
%%%
%%%\end{itemize}
%%%\end{proof}
%%
%%The proof of Proposition \ref{propositiongeodesicconvergence} will proceed through several lemmas, the first of which is as follows:
%%
%%\begin{lemma}
%%\label{lemmaxyz}
%%Fix $\epsilon > 0$. There exists $\delta = \delta_X(\epsilon) > 0$ such that if $\Delta = \Delta(x,y_1,y_2)$ is a geodesic triangle in $X$ satisfying
%%\begin{equation}
%%\label{y1y2delta}
%%\wbar\Dist(y_1,y_2) \leq \delta,
%%\end{equation}
%%then for all $t\in[0,\min_{i = 1}^2 \dist(x,y_i)]$, if $z_i = \geo x{y_i}_t$, then
%%\begin{equation}
%%\label{z1z2epsilon}
%%\wbar\Dist(z_1,z_2) \leq \epsilon.
%%\end{equation}
%%\end{lemma}
%%\begin{proof}
%%We prove the assertion first for $X = \H^2$ and then in general:
%%\begin{itemize}
%%\item[If $X = \H^2$:] Let $\epsilon > 0$, and by contradiction, suppose that for each $\delta = \frac 1n > 0$ there exists a 5-tuple $(x^{(n)},y_1^{(n)},y_2^{(n)},z_1^{(n)},z_2^{(n)})$ satisfying the hypotheses but not the conclusion of the theorem. Since $\bord\H^2$ is compact, there exists a convergent subsequence
%%\[
%%(x^{(n_k)},y_1^{(n_k)},y_2^{(n_k)},z_1^{(n_k)},z_2^{(n_k)}) \to (x,y_1,y_2,z_1,z_2)\in\left(\bord\H^2\right)^5.
%%\]
%%Taking the limit of \eqref{y1y2delta} as $k\to\infty$ shows that $\wbar\Dist(y_1,y_2) = 0$, so $y_1 = y_2$. Conversely, taking the limit of \eqref{z1z2epsilon} shows that $\wbar\Dist(z_1,z_2) \geq \epsilon > 0$, so $z_1\neq z_2$. Write $y = y_1 = y_2$.
%%
%%We will take for granted that Proposition \ref{propositiongeodesicconvergence} holds when $X = \H^2$. (This can be proven using the explicit form of geodesics in this space.) It follows that $z_i\in \geo xy$ if $x\neq y$, and $z_i = x$ if $x = y$. The second case is clearly a contradiction, so we assume that $x\neq y$.
%%
%%Writing $z_i^{(n_k)} = \geo{x^{(n_k)}}{y_i^{(n_k)}}_{t_k}$, we observe that
%%\begin{align*}
%%t_k - \dox{x^{(n_k)}}
%%&= \lb z_i^{(n_k)} | y_i^{(n_k)} \rb_\zero - \lb z_i^{(n_k)} | x^{(n_k)} \rb_\zero - \lb x^{(n_k)} | y_i^{(n_k)} \rb_\zero\\
%%&\tendsto k \lb z_i | y \rb_\zero - \lb z_i | x \rb_\zero - \lb x |y\rb_\zero.
%%\end{align*}
%%Since the left hand side is independent of $i$, so is the right hand side. But the function
%%\[
%%z\mapsto \lb z | y \rb_\zero - \lb z | x \rb_\zero - \lb x | y \rb_\zero
%%\]
%%is an isometric embedding from $\geo xy$ to $[-\infty,+\infty]$; it is therefore injective. Thus $z_1 = z_2$, a contradiction.
%%
%%\item[In general:] Let $\epsilon > 0$, and fix $\w\epsilon > 0$ to be determined, depending on $\epsilon$. Let $\w\delta = \delta_{\H^2}(\w\epsilon)$, and fix $\delta > 0$ to be determined, depending on $\w\delta$. Now suppose that $\Delta = \Delta(x,y_1,y_2)$ is a geodesic triangle in $X$ satisfying \eqref{y1y2delta}, fix $t\geq 0$, and let $z_i = \geo x{y_i}_t$. To complete the proof, we must show that $\wbar\Dist(z_1,z_2) \leq \epsilon$.
%%
%%By contradiction suppose not, i.e. suppose that $\wbar\Dist(z_1,z_2) > \epsilon$. Then $\Dist(x,z_i) > \epsilon/2$ for some $i = 1,2$; without loss of generality suppose $\Dist(x,z_1) > \epsilon/2$. By Proposition \ref{propositionrips} this implies $\dist(\zero,\geo x{z_1}) \asymp_{\plus,\epsilon} 0$; fix $w_1\in\geo x{z_1}$ with $\dox{w_1}\asymp_{\plus,\epsilon} 0$. Let $s = \dist(x,w_1) \leq t$, and let $w_2 = \geo x{z_2}_s$. (See Figure \ref{figurexyz}.)
%%
%%\begin{figure}
%%\begin{center}
%%\begin{tikzpicture}[line cap=round,line join=round,>=triangle 45,x=1.0cm,y=1.0cm]
%%\clip(-1.539,-1.576) rectangle (3.98,1.8);
%%\draw [shift={(4.7207100633754004,-4.251144523671373)}] plot[domain=1.948435252123557:2.6141429865240005,variable=\t]({1.0*5.902960522581992*cos(\t r)+-0.0*5.902960522581992*sin(\t r)},{0.0*5.902960522581992*cos(\t r)+1.0*5.902960522581992*sin(\t r)});
%%\draw [shift={(2.6440776402056905,-5.2367125169303135)}] plot[domain=1.5361702930450398:2.2233807737101374,variable=\t]({1.0*4.980022039672619*cos(\t r)+-0.0*4.980022039672619*sin(\t r)},{0.0*4.980022039672619*cos(\t r)+1.0*4.980022039672619*sin(\t r)});
%%\draw (0.0,-0.0)-- (0.48334729355344397,-0.1414365427938682);
%%\draw [shift={(2.802578610044664,0.5103691534918943)}] plot[domain=1.9130096151383342:4.730441794850227,variable=\t]({1.0*0.7701702590181908*cos(\t r)+-0.0*0.7701702590181908*sin(\t r)},{0.0*0.7701702590181908*cos(\t r)+1.0*0.7701702590181908*sin(\t r)});
%%\draw [shift={(7.3783710087828664,2.3193354663807444)}] plot[domain=3.4843932225971725:3.560165863983011,variable=\t]({1.0*7.320980188110974*cos(\t r)+-0.0*7.320980188110974*sin(\t r)},{0.0*7.320980188110974*cos(\t r)+1.0*7.320980188110974*sin(\t r)});
%%\begin{scriptsize}
%%\draw [fill=black] (0.0,-0.0) circle (1pt);
%%\draw[color=black] (-0.20573038068032204,-0.16484898217738642) node {$\zero$};
%%\draw [fill=black] (-0.38,-1.28) circle (1pt);
%%\draw[color=black] (-0.6011417551376018,-1.450757318050246) node {$x$};
%%\draw [fill=black] (2.544130372359072,1.2358805154383208) circle (1pt);
%%\draw[color=black] (2.8569090399799677,1.294830750435049) node {$y_1$};
%%\draw [fill=black] (2.8164815956355223,-0.25967560810823964) circle (1pt);
%%\draw[color=black] (3.13494327476329,-0.3038660995690469) node {$y_2$};
%%\draw [fill=black] (0.48334729355344397,-0.1414365427938682) circle (1pt);
%%\draw[color=black] (0.458863764973814,0.0892241456218923) node {$w_1$};
%%\draw [fill=black] (0.6894125111628517,-0.6563300275601325) circle (1pt);
%%\draw[color=black] (0.8759151171487972,-0.8641974071794341) node {$w_2$};
%%\draw [fill=black] (1.4043001305961753,0.6321297967948407) circle (1pt);
%%\draw[color=black] (1.3277207486716958,0.8430251189121523) node {$z_1$};
%%\draw [fill=black] (1.7083588447287483,-0.34538858046728116) circle (1pt);
%%\draw[color=black] (1.796903519868552,-0.6861631723961132) node {$z_2$};
%%\end{scriptsize}
%%\end{tikzpicture}
%%\caption{The triangle $\Delta(x,y_1,y_2)$.}
%%\label{figurexyz}
%%\end{center}
%%\end{figure}
%%
%%Now let $\wbar{\Delta} = \Delta(\wbar x,\wbar{y_1},\wbar{y_2})$ be a comparison triangle for $\Delta(x,y_1,y_2)$, and let $\wbar{z_1},\wbar{z_2},\wbar{w_1},\wbar{w_2}$ be the corresponding comparison points. Note that $\wbar{z_i} = \geo{\wbar x}{\wbar{y_i}}_t$ and $\wbar{w_i} = \geo{\wbar x}{\wbar{y_i}}_s$. Without loss of generality, suppose that $\wbar{w_1} = \zero_\H$. Then $\dox{y_2} \leq \dox{w_1} + \dist(w_1,y_2) \asymp_{\plus,\epsilon} \dist(\zero_\H,\wbar{y_2})$, and so $\lb y_1|y_2\rb_\zero \lesssim_{\plus,\epsilon} \lb \wbar{y_1}|\wbar{y_2}\rb_{\zero_\H}$, and thus
%%\[
%%\wbar\Dist(\wbar{y_1},\wbar{y_2}) \lesssim_{\times,\epsilon} \wbar\Dist(y_1,y_2) \leq \delta.
%%\]
%%Setting $\delta$ equal to $\w\delta$ divided by the implied constant, we have
%%\[
%%\wbar\Dist(\wbar{y_1},\wbar{y_2}) \leq \w\delta = \delta_\H(\w\epsilon).
%%\]
%%Thus $\wbar\Dist(\wbar{z_1},\wbar{z_2})\leq\w\epsilon$ and $\wbar\Dist(\wbar{w_1},\wbar{w_2})\leq\w\epsilon$.
%%\begin{itemize}
%%\item If $\dist(\wbar{z_1},\wbar{z_2})\leq\w\epsilon$, then the CAT(-1) inequality finishes the proof (as long as $\w\epsilon\leq \epsilon$). Thus, suppose that
%%\begin{equation}
%%\label{z1z2}
%%\Dist(\wbar{z_1},\wbar{z_2}) \leq \w\epsilon.
%%\end{equation}
%%\item If $\Dist(\wbar{w_1},\wbar{w_2})\leq\w\epsilon$, then $0 = \lb \wbar{w_1}|\wbar{w_2}\rb_{\zero_\H} \geq -\log(\w\epsilon)$, a contradiction for $\w\epsilon$ sufficiently small. Thus, suppose that
%%\begin{equation}
%%\label{w1w2}
%%\dist(\wbar{w_1},\wbar{w_2}) \leq\w\epsilon.
%%\end{equation}
%%\end{itemize}
%%By \eqref{z1z2}, we have $\dist(\zero_\H,\wbar{z_i}) \geq -\log(\w\epsilon)$. Applying \eqref{w1w2} gives
%%\[
%%\lb z_i|y_i\rb_{w_i} = \dist(w_i,z_i) = \dist(\wbar{w_i},\wbar{z_i}) \gtrsim_\plus -\log(\w\epsilon).
%%\]
%%Applying \eqref{w1w2}, the CAT(-1) inequality, and the coarse asymptotic $\dox{w_1} \asymp_{\plus,\epsilon} 0$, we have
%%\[
%%\lb z_i|y_i\rb_\zero \gtrsim_{\plus,\epsilon} -\log(\w\epsilon),
%%\]
%%and thus $\wbar\Dist(z_i,y_i)\lesssim_{\times,\epsilon} \w\epsilon$. Using the triangle inequality together with the assumption $\wbar\Dist(y_1,y_2)\leq\delta$, we have
%%\[
%%\wbar\Dist(z_1,z_2) \lesssim_{\times,\epsilon} \max(\delta,\w\epsilon).
%%\]
%%Setting $\w\epsilon$ equal to $\epsilon$ divided by the implied constant, and decreasing $\delta$ if necessary, completes the proof.
%%\end{itemize}
%%
%%%\ignore{
%%%
%%%We divide the proof into three cases:
%%%
%%%\begin{itemize}
%%%\item[Case 1:] $\wbar\Dist(y_1,y_2) = \dist(y_1,y_2) \leq \delta$. In this case, by Lemma \ref{lemmaxzyeuclidean} and the CAT(-1) inequality, we have $\dist(z_1,z_2)\leq \dist(y_1,y_2)\leq\delta$. Thus letting $\delta\leq\epsilon$ completes the proof in this case.
%%%
%%%\item[Case 2:] $\wbar\Dist(y_1,y_2) = \Dist(y_1,y_2)\leq\delta$. Since $X$ is hyperbolic, by Proposition \internalmissingref\ there exists $y_3\in\geo{y_1}{y_2}$ such that
%%%\[
%%%\dist(y_3,\geo x{y_1}) \asymp_\plus \dist(y_3,\geo x{y_2}) \asymp_\plus 0.
%%%\]
%%%For $i = 1,2$, let $w_i\in\geo x{y_i}$ be such that $\dist(x,w_i) =\dist(x,y_3)$. Note that $\dist(w_1,w_2) \leq \dist(w_1,y_3) + \dist(w_2,y_3) \leq 2\dist(y_3,\geo x{y_1}) + 2\dist(y_3,\geo x{y_2}) \asymp_\plus 0$.
%%%
%%%\begin{itemize}
%%%\item[Case 2A:] $\dist(x,z_1)\leq \dist(x,y_3)$. In this case, consider the comparison triangle $\Delta(\xx,\ww_1,\ww_2)$, and let $\zz_1,\zz_2$ be the corresponding points on the comparison triangles. Without loss of generality, suppose that $\xx = \0$. Then
%%%\[
%%%\|\zz_2 - \zz_1\| = \frac{\|\zz_1\|}{\|\ww_1\|}\|\ww_2 - \ww_1\| = \frac{\dist(x,z_1)}{\dist(x,y_3)} \dist(w_1,w_2).
%%%\]
%%%On the other hand, by (d) of Proposition \ref{propositionbasicidentities} we have
%%%\[
%%%\dist(x,y_3)\geq \lb y_1|y_2\rb_x - \lb y_1|y_2\rb_{y_3} = \lb y_1|y_2\rb_x
%%%\]
%%%and thus
%%%\[
%%%\dist(z_1,z_2) \lesssim_\times \frac{\dist(x,z_1)}{\lb y_1|y_2\rb_x}\cdot
%%%\]
%%%[Continue\internal]
%%%\end{itemize}
%%%\end{itemize}
%%%}% end ignore
%%\end{proof}
%%
%%%\ignore{
%%%\begin{corollary}
%%%If $\epsilon,\delta,\Delta(x,y_1,y_2)$ are as in Lemma \ref{lemmaxyz}, then the Hausdorff distance between $\geo x{y_1}$ and $\geo x{y_2}$ is at most $\epsilon$.
%%%\end{corollary}
%%%\begin{proof}
%%%For every $z_1 = \geo x{y_1}_t\in\geo x{y_1}$, we can find $z_2 = \geo x{y_2}_{t\wedge\dist(x,y_2)}\in\geo x{y_2}$ such that $\wbar\Dist(z_1,z_2) \leq\epsilon$, and vice-versa.
%%%\end{proof}
%%%}% end ignore
%%
%%\begin{notation}
%%\label{notationwpi}
%%If $\pi:[t,s]\to X$ is an isometric embedding, then $\w\pi:[-\infty,+\infty]\to X$ is defined by the equation
%%\[
%%\w\pi(r) = \pi(t\vee r\wedge s).
%%\]
%%\end{notation}
%%
%%\begin{corollary}
%%\label{corollaryxyz}
%%If $\epsilon,\delta,\Delta(x,y_1,y_2)$ are as in Lemma \ref{lemmaxyz}, and if $\pi_1:[t,s_1]\to \geo x{y_1}$ and $\pi_2:[t,s_2]\to\geo x{y_2}$ are isometric embeddings, then
%%\[
%%\wbar\Dist(\w\pi_1(r),\w\pi_2(r)) \lesssim_\times \epsilon \all r\in [-\infty,+\infty].
%%\]
%%\end{corollary}
%%\begin{proof}
%%If $r\leq t$, then $\w\pi_1(r) = x = \w\pi_2(r)$. If $t\leq r \leq \min_{i = 1}^2 s_i$, then $\w\pi_i(r) = \geo x{y_i}_{r - t}$, allowing us to apply Lemma \ref{lemmaxyz} directly. Finally, suppose $r\geq r_0 := \min_{i = 1}^2 s_i$. Without loss of generality suppose that $s_1 \leq s_2$, so that $r_0 = s_1$. Applying the previous case to $r_0$, we have
%%\[
%%\wbar\Dist(y_1,w_2) \leq \epsilon,
%%\]
%%where $w_2 = \pi_2(s_1)$. Now $\w\pi_1(r) = y_1$, and $\w\pi_2(r)\in\geo{w_2}{y_2}$, so Lemma \ref{lemmageodesicdiameter} completes the proof.
%%\end{proof}
%%
%%\begin{lemma}
%%\label{lemmageodesicconvergence}
%%Suppose that $(x_n)_1^\infty$ and $(y_n)_1^\infty$ are sequences in $X$ which converge to points $x_n\to x\in\bord X$ and $y_n\to y\in\bord X$, with $x\neq y$. Then there exists a geodesic $\geo xy$ connecting $x$ and $y$ such that $\geo{x_n}{y_n} \to \geo xy$ in the Hausdorff metric. If $x = y$, then $\geo{x_n}{y_n}\to \{x\}$ in the Hausdorff metric.
%%\end{lemma}
%%\begin{proof}
%%We observe first that if $x = y$, then the conclusion follows immediately from Lemma \ref{lemmageodesicdiameter}. Thus we assume in what follows that $x\neq y$.
%%
%%For any pair $p,q\in X$, we define the \emph{standard parameterization} of the geodesic $\geo pq$ to be the unique isometry $\pi:[-\lb \zero|q\rb_p , \lb \zero|p\rb_q] \to \geo pq$ sending $-\lb \zero|q\rb_p$ to $p$ and $\lb \zero|p\rb_q$ to $q$. For each $n$ let $\pi_n:[t_n,s_n]\to \geo{x_n}{y_n}$ be the standard parameterization, and for each $m,n\in\N$ let $\pi_{m,n}:[t_{m,n},s_{m,n}]\to \geo{x_n}{y_m}$ be the standard parameterization. Let $\w\pi_n:[-\infty,+\infty]\to \geo{x_n}{y_n}$ and $\w\pi_{m,n}:[-\infty,+\infty]\to \geo{x_n}{y_m}$ be as in Notation \ref{notationwpi}. Note that
%%\[
%%t_n - t_{m,n} = \lb \zero|y_m\rb_{x_n} - \lb \zero|y_n\rb_{x_n} = \lb x_n|y_n\rb_\zero - \lb x_n|y_m\rb_\zero \tendsto{m,n} \lb x|y\rb_\zero - \lb x|y\rb_\zero = 0.
%%\]
%%(We have $\lb x|y\rb_\zero < \infty$ since $x\neq y$.) Thus
%%\[
%%\wbar\Dist(\w\pi_n(r), \w\pi_n(r - t_n + t_{m,n})) \leq \dist(\w\pi_n(r), \w\pi_n(r - t_n + t_{m,n})) \leq |t_n - t_{m,n}| \tendsto n 0.
%%\]
%%Here and below, the limit converges uniformly for $r\in[-\infty,+\infty]$. On the other hand, Corollary \ref{corollaryxyz} implies that
%%\[
%%\wbar\Dist(\w\pi_n(r - t_n + t_{m,n}) , \w\pi_{m,n}(r)) \tendsto{m,n} 0,
%%\]
%%so the triangle inequality gives
%%\[
%%\wbar\Dist(\w\pi_n(r) , \w\pi_{m,n}(r)) \tendsto{m,n} 0.
%%\]
%%A similar argument shows that
%%\[
%%\wbar\Dist(\w\pi_{m,n}(r) , \w\pi_m(r)) \tendsto{m,n} 0,
%%\]
%%so the triangle inequality gives
%%\[
%%\wbar\Dist(\w\pi_n(r) , \w\pi_m(r)) \tendsto{m,n} 0,
%%\]
%%i.e. the sequence of functions $(\w\pi_n)_1^\infty$ is uniformly Cauchy. Since $(\bord X,\wbar\Dist)$ is complete, they converge uniformly to a function $\w\pi:[-\infty,+\infty]\to X$.
%%
%%Clearly, $\geo{x_n}{y_n} = \w\pi_n([-\infty,+\infty]) \to \w\pi([-\infty,+\infty])$ in the Hausdorff metric. We claim that $\w\pi([-\infty,+\infty])$ is a geodesic connecting $x$ and $y$. Indeed,
%%\[
%%t_n \tendsto n t := \lb x|y\rb_\zero - \dox x \text{ and } s_n \tendsto n s := \dox y - \lb x|y\rb_\zero.
%%\]
%%For all $t < r_1 < r_2 < s$, we have $t_n < r_1 < r_2 < s_n$ for all sufficiently large $n$, which implies that
%%\[
%%\dist(\w\pi(r_1) , \w\pi(r_2)) = \lim_{n\to\infty} \dist(\w\pi_n(r_1) , \w\pi_n(r_2)) = \lim_{n\to\infty} (r_2 - r_1) = r_2 - r_1,
%%\]
%%i.e. $\w\pi\given(t,s)$ is an isometric embedding. Since $\w\pi$ is continuous (being the uniform limit of continuous functions), $\pi := \w\pi\given[t,s]$ is also an isometric embedding. A similar argument shows that $\w\pi(r) = \pi(t)$ for all $r\leq t$, and $\w\pi(r) = \pi(s)$ for all $r\geq s$; thus $\w\pi([-\infty,+\infty]) = \pi([t,s])$ is a geodesic. To complete the proof, we must show that $\pi(t) = x$ and $\pi(s) = y$. Indeed,
%%\[
%%\pi(t) = \w\pi(-\infty) = \lim_{n\to\infty} \w\pi_n(-\infty) = \lim_{n\to\infty} x_n = x,
%%\]
%%and a similar argument shows that $\pi(s) = y$. Thus the geodesic $\pi([t,s])$ connects $x$ and $y$.
%%\end{proof}
%%
%%%\ignore{
%%%In fact, we demonstrate the following stronger assertion:
%%%
%%%\begin{claim}
%%%\label{claimpirkpir}
%%%For any sequence $[t_k,s_k]\ni r_k \to r\in[t,s]$, we have $\pi_k(r_k) \to \pi(r)$.
%%%\end{claim}
%%%\begin{subproof}
%%%Let $z_k = \pi_k(r_k)$ and let $z = \pi(r)$. If $r\neq\pm\infty$, then
%%%\[
%%%\dist(z_k,z) \leq \dist(z_k,\pi_k(r)) + \dist(\pi_k(r),z) = |r_k - r| + \dist(\pi_k(r),\pi(r)) \tendsto k 0.
%%%\]
%%%On the other hand, suppose $r = \pm\infty$; without loss of generality suppose $r = +\infty$. \comdavid{Question: Is it possible to use the stronger assertion of Lemma \ref{lemmaxyz} to prove this case? Consider $\w\pi(\pm\infty)$.} Fix $N\in\N$. For all $k$ sufficiently large, we have $r_k\geq N$ and therefore
%%%\[
%%%\lb \pi_k(N)|z_k\rb_{\pi_k(0)} = N.
%%%\]
%%%Moreover, since $\pi_k(N)\to\pi(N)$ and $\pi_k(0)\to\pi(0)$, we have
%%%\[
%%%\lb \pi(N)|z_k\rb_{\pi(0)} \geq N - 1
%%%\]
%%%for all $k$ sufficiently large. On the other hand
%%%\[
%%%\lb \pi(N)|z\rb_{\pi(0)} = N,
%%%\]
%%%so by Gromov's inequality we have
%%%\[
%%%\lb z_k|z\rb_{\pi(0)} \gtrsim_\plus N.
%%%\]
%%%Since $N$ was arbitrary, we have $z_k\to z$.
%%%\end{subproof}
%%%
%%%Since $\pi_k(t_k) = x_k \to x$ and $\pi_k(s_k) = y_k \to y$, we have $\pi(t) = x$ and $\pi(s) = y$, i.e. $\pi$ connects $x$ and $y$.
%%%}% end ignore
%%
%%
%%
%%%\ignore{
%%%
%%%[END]
%%%
%%%
%%%
%%%
%%%For each $n$ let
%%%\[
%%%t_n = -\dox{x_n} , \;\; s_n = t_n + \dist(x_n,y_n) = \busemann_{x_n}(y_n,\zero),
%%%\]
%%%and let $\pi_n:[t_n,s_n]\to \geo{x_n}{y_n}$ be an isometry satisfying $\pi_n(t_n) = x_n$ and $\pi_n(s_n) = y_n$. We observe that
%%%\[
%%%t_n \tendsto n t := -\dox x , \;\; s_n \tendsto n s := \begin{cases}
%%%\busemann_x(y,\zero) & y\in X\\
%%%+\infty & y\in\del X
%%%\end{cases}.
%%%\]
%%%(If $y\in\del X$, then
%%%\[
%%%\busemann_{x_n}(y_n,\zero) = \dox{y_n} - 2\lb x_n|y_n\rb_\zero \tendsto n (+\infty) - 2\lb x|y\rb_\zero = +\infty;
%%%\]
%%%$\lb x|y\rb_\zero$ is finite since $x\neq y$.) In particular, note that $t < s$. For each $n$, define $\w\pi_n:[-\infty,+\infty]\to \geo{x_n}{y_n}$ via the formula
%%%\[
%%%\w\pi_n(t) = \pi_n(t_n \vee t \wedge s_n).
%%%\]
%%%\begin{claim}
%%%The sequence $(\w\pi_n)_1^\infty$ converges uniformly to a function $\w\pi:[-\infty,+\infty]\to X$.
%%%\end{claim}
%%%\begin{subproof}
%%%Since $(\bord X,\wbar\Dist)$ is complete, it suffices to show that the sequence $(\w\pi_n)_1^\infty$ is uniformly Cauchy, i.e. that for every $\epsilon > 0$ there exists $N\in\N$ such that for all $m,n\geq N$ and $t\in[-\infty,+\infty]$,
%%%\[
%%%\wbar\Dist(\w\pi_n(t),\w\pi_m(t)) \leq \epsilon.
%%%\]
%%%Indeed, fix such an $\epsilon$, and let $\delta > 0$ be as in Lemma \ref{lemmaxyz}.
%%%
%%%\end{subproof}
%%%
%%%\end{proof}
%%%
%%%
%%%\begin{proof}
%%%Corollary \internalmissingref\ shows that the sequence $(\geo{x_n}{y_n})_1^\infty$ is Cauchy in the Hausdorff metric. Since the Hausdorff metric of a complete metric space is complete [ref?], we have $\geo{x_n}{y_n} \to K$ for some closed set $K\subset \bord X$.
%%%\begin{itemize}
%%%\item[Case 1:] $K\cap X\neq\emptyset$, say $w\in K\cap X$. There exists a sequence $\geo{x_n}{y_n}\ni w_n\to w$. For each $n$ let $t_n = -\dist(w_n,x_n)$ and $s_n = \dist(w_n,y_n)$, so that there exists a parameterization $\pi_n:[t_n,s_n]\to \geo{x_n}{y_n}$ sending $\zero$ to $w_n$.
%%%
%%%Fix $m,n\in\N$, and let's compare $\pi_m$ to $\pi_n$. Let
%%%\[
%%%\w t_m = s_n - \dist(x_m,y_n) , \;\; \w s_m = \w t_m + \dist(x_m,y_n),
%%%\]
%%%so that there exist parameterizations $\pi_{n,m}:[\w t_m,s_n]\to \geo{x_m}{y_n}$ and $\w\pi_m : [\w t_m,\w s_m]\to \geo{x_m}{y_m}$. Applying Lemma \ref{lemmaxyz} twice gives
%%%\[
%%%\wbar\Dist(\pi_n(t_n\vee t\wedge s_n), \w\pi_m(\w t_m\vee t\wedge \w s_m)) \leq 2\epsilon \all t\in\R,
%%%\]
%%%assuming that $m,n$ are large enough so that $\wbar\Dist(x_n,x_m),\wbar\Dist(y_n,y_m) \leq \delta$, where $\delta$ is as in Lemma \ref{lemmaxyz}.
%%%
%%%In particular, letting $t = 0$ gives
%%%\[
%%%\wbar\Dist(w_n,\w\pi_m(0)) \leq 2\epsilon,
%%%\]
%%%so $\w\pi_m(0)\to w$. [This is a very unusual calculation...]
%%%
%%%
%%%\item[Case 2:]
%%%\end{itemize}
%%%\end{proof}
%%%
%%%
%%%
%%%
%%%
%%%
%%%[Previous attempts:]
%%%
%%%\begin{corollary}
%%%\label{corollaryxyz}
%%%Fix $\epsilon > 0$. There exists $\delta > 0$ such that if $\pi_1:[t,s_1]\to X$ is a geodesic connecting $x,y_1\in X$, and $y_2\in X$ is a point satisfying
%%%\[
%%%\wbar\Dist(y_1,y_2) \leq \delta,
%%%\]
%%%then there exists a geodesic $\pi_2:[t,s_2]\to X$ connecting $x$ and $y_2$ satisfying
%%%\[
%%%\wbar\Dist(\pi_1(r),\pi_2(r)) \leq \epsilon \all r\in [t,\min(s_1,s_2)].
%%%\]
%%%\end{corollary}
%%%\begin{proof}
%%%Apply Lemma \ref{lemmaxyz} with $z_i = \pi_i(r)$.
%%%\end{proof}
%%%
%%%[Attempting to ``start over'' again]
%%%
%%%[Note: The main problem with the ``Hausdorff metric'' presentation seems to be that it makes it less obvious that the limit of geodesics is a geodesic.]
%%%\begin{corollary}
%%%Fix $\epsilon > 0$. There exists $\delta > 0$ such that if $x,y_1,y_2\in X$ satisfy
%%%\[
%%%\wbar\Dist(y_1,y_2)\leq\delta,
%%%\]
%%%then the Hausdorff distance between $\geo x{y_1}$ and $\geo x{y_2}$ is at most $\epsilon$.
%%%\end{corollary}
%%%\begin{proof}
%%%Let $\epsilon > 0$, and fix $\delta > 0$ to be determined.
%%%
%%%Fix $z_1\in\geo x{y_1}$, and write $z = \geo x{y_1}_t$ for some $0\leq t\leq \dist(x,y_1)$. If $t\leq \dist(x,y_2)$, then we let $z_2 = \geo x{y_2}_t\in\geo x{y_2}$, and then by Lemma \ref{lemmaxyz}, $\wbar\Dist(z_1,z_2)\leq\epsilon$. Otherwise, let $s = \dist(x,y_2)$, and for $i = 1,2$ let $w_i = \geo x{y_i}_s$. Then $w_2 = y_2$, $\wbar\Dist(w_1,w_2) \leq \epsilon$, and $z_1\in\geo{y_1}{w_1}$. By Lemma \ref{lemmageodesicdiameter},
%%%\[
%%%\wbar\Dist(z_1,y_1) \lesssim_\times \wbar\Dist(y_1,w_1) \leq \wbar\Dist(y_1,y_2) + \wbar\Dist(w_1,w_2) \leq \delta + \epsilon \lesssim_\times \epsilon,
%%%\]
%%%and thus $\wbar\Dist(z_1,y_2) \lesssim_\times\epsilon$. Choosing $\epsilon$ appropriately (and fixing horrible notation), we have $\wbar\Dist(z_1,y_2) \leq \epsilon$.
%%%
%%%Since $z_1$ was arbitrary, we have $\geo x{y_1} \subset \geo x{y_2}^{(\epsilon)}$. A symmetric argument shows that $\geo x{y_2}\subset \geo x{y_1}^{(\epsilon)}$.
%%%\end{proof}
%%%
%%%\begin{proof}[Proof of Proposition \ref{propositiongeodesicconvergence}]
%%%
%%%\end{proof}
%%%
%%%
%%%
%%%
%%%
%%%
%%%[Attempting to ``start over'']
%%%
%%%For any pair $x,y\in X$, we define the \emph{standard parameterization} of the geodesic $\geo xy$ to be the unique isometry $\pi:[-\lb \zero|y\rb_x , \lb \zero|x\rb_y] \to \geo xy$ sending $-\lb \zero|y\rb_x$ to $x$ and $\lb \zero|x\rb_y$ to $y$.
%%%
%%%\begin{corollary}
%%%\label{corollaryxyz}
%%%Fix $\epsilon > 0$. There exists $\delta > 0$ such that if $\geo x{y_1}$ and $\geo x{y_2}$ are geodesics satisfying
%%%\[
%%%\wbar\Dist(y_1,y_2) \leq \delta,
%%%\]
%%%and if $\pi_1:[t_1,s_1]\to\geo x{y_1}$ and $\pi_2:[t_2,s_2]\to\geo x{y_2}$ are their standard parameterizations, then
%%%\[
%%%\wbar\Dist(\pi_1(r),\pi_2(r)) \leq \epsilon \text{ for all applicable $r$.}
%%%\]
%%%\end{corollary}
%%%\begin{proof}
%%%It suffices to show that $|\lb \zero|y_1\rb_x - \lb \zero|y_2\rb_x| \leq \epsilon$. Let $d_i = \lb \zero|y_i\rb_x$; then
%%%\[
%%%|e^{-d_1} - e^{-d_2}| \leq e^{-\lb y_1|y_2\rb_x}.
%%%\]
%%%On the other hand, by (k) of Proposition \ref{propositionbasicidentities}
%%%\[
%%%\lb y_1|y_2\rb_x = \lb y_1|y_2\rb_\zero + \dox x - \lb x|y_1\rb_\zero - \lb x|y_2\rb_\zero \asymp_\plus \lb y_1|y_2\rb_\zero + \dox x
%%%\]
%%%and so
%%%\[
%%%|e^{-d_1} - e^{-d_2}| \lesssim_\times e^{-\dox x}\Dist(y_1,y_2).
%%%\]
%%%Since $d_1,d_2\leq \dox x$, this gives
%%%\[
%%%|d_1 - d_2| \lesssim_\times \Dist(y_1,y_2).
%%%\]
%%%On the other hand, it is clear that $|d_1 - d_2| \leq \dist(y_1,y_2)$.
%%%\end{proof}
%%%
%%%
%%%[END]
%%%
%%%\begin{lemma}
%%%\label{lemmageodesicconvergence}
%%%Fix points $x,y\in \bord X$ and sequences $(x_n)_1^\infty$ and $(y_n)_1^\infty$ in $X$ which converge to $x$ and $y$, respectively. There exists an increasing sequence $(n_k)_1^\infty$ and a geodesic $\geo xy$ connecting $x$ and $y$ such that $\geo{x_{n_k}}{y_{n_k}} \tendsto k \geo xy$ in the Hausdorff metric on $(\bord X,\wbar\Dist)$.% sense of Proposition \ref{propositiongeodesicconvergence}.
%%%\end{lemma}
%%%\begin{proof}
%%%For each $k\in\N$ let $\epsilon_k = 2^{-k}$, and let $\delta_k > 0$ be given by Lemma \ref{lemmaxyz}.
%%%\begin{observation}
%%%There exists an increasing sequences $(n_k)_1^\infty$ satisfying:
%%%\begin{align} \label{ynym}
%%%\wbar\Dist(y_{n_k},y_{n_{k + 1}}) &\leq \min(\delta_k,2^{-k}) &\text{for all $k\in\N$} \\ \label{xnxm}
%%%\wbar\Dist(x_{n_k},x_{n_{k + 1}}) &\leq \min(\delta_k,2^{-k}) &\text{for all $k\in\N$} \\ \label{xkyk1yk}
%%%\dist(x_{n_k},y_{n_{k + 1}}) &\geq \dist(x_{n_k},y_{n_k}) &\text{if $y\in\del X$}\\
%%%\dist(y_{n_{k + 1}},x_{n_{k + 1}}) &\geq \dist(y_{n_{k + 1}},x_{n_k}) &\text{if $x\in\del X$}
%%%\end{align}
%%%\end{observation}
%%%\begin{subproof}
%%%Since $x_n\to x$ and $y_n\to y$, \eqref{ynym} and \eqref{xnxm} can be guaranteed by choosing $n_k$ large enough. [CONTINUE]
%%%
%%%\end{subproof}
%%%
%%%For shorthand let us write $x_k = x_{n_k}$ and $y_k = y_{n_k}$. Let $\pi_1:[t_1,s_1]\to X$ be a geodesic connecting $x_1$ and $y_1$. Define the geodesics $\pi_k:[t_k,s_k]\to X$ and $\pi_k':[t_k,s_{k + 1}]\to X$ inductively using Corollary \ref{corollaryxyz}, so that
%%%\begin{itemize}
%%%\item[(i)] $\pi_k$ connects $x_k$ and $y_k$; $\pi_k'$ connects $x_k$ and $y_{k + 1}$
%%%\item[(ii)]
%%%\begin{equation}
%%%\label{piinduction}
%%%\begin{split}
%%%\wbar\Dist(\pi_k(r),\pi_k'(r)) &\leq \epsilon_k = 2^{-k} \all r\in [t_k,\min(s_k,s_{k + 1})]\\
%%%\wbar\Dist(\pi_k'(r),\pi_{k + 1}(r)) &\leq \epsilon_k = 2^{-k} \all r\in [\max(t_k,t_{k + 1}),s_{k + 1}].
%%%\end{split}
%%%\end{equation}
%%%\end{itemize}
%%%
%%%\begin{claim}
%%%The limits $t = \lim_{k\to\infty} t_k$ and $s = \lim_{k\to\infty} s_k$ both exist. The values $t = -\infty$ and $s = +\infty$ are possible, but $t\neq +\infty$ and $s\neq -\infty$.
%%%\end{claim}
%%%\begin{subproof}
%%%If $y\in X$, then $|s_{k + 1} - s_k| = |\busemann_{x_k}(y_{k + 1},y_k)| \leq \dist(y_{k + 1},y_k) \leq 2^{-k}$ for all $k$, so the sequence $(s_k)_1^\infty$ is Cauchy. On the other hand, if $y\in\del X$, then \eqref{xkyk1yk} implies $s_{k + 1}\geq s_k$ for all $k$. Either way the sequence $(s_k)_1^\infty$ converges in $\OC{-\infty}{+\infty}$; a symmetric argument shows that the sequence $(t_k)_1^\infty$ converges in $\CO{-\infty}{+\infty}$.
%%%\end{subproof}
%%%
%%%Now for all $r\in (t,s)$, we have $r\in[t_k,s_k]$ for all sufficiently large $k$, so by \eqref{piinduction} we have $\wbar\Dist(\pi_k(r),\pi_{k + 1}(r)) \leq 2^{-(k - 1)}$ for all sufficiently large $k$. Thus the sequence $(\pi_k(r))_1^\infty$ is Cauchy, and converges since $X$ is complete.
%%%
%%%Let $\pi:(t,s)\to X$ be the pointwise limit of the sequence $(\pi_k)_1^\infty$. By continuity, $\pi$ is an isometric embedding. Let $\pi:[t,s]\to X$ be the unique continuous extension.
%%%
%%%\begin{claim}
%%%\label{claimpirkpir}
%%%For any sequence $[t_k,s_k]\ni r_k \to r\in[t,s]$, we have $\pi_k(r_k) \to \pi(r)$.
%%%\end{claim}
%%%\begin{subproof}
%%%Let $z_k = \pi_k(r_k)$ and let $z = \pi(r)$. If $r\neq\pm\infty$, then
%%%\[
%%%\dist(z_k,z) \leq \dist(z_k,\pi_k(r)) + \dist(\pi_k(r),z) = |r_k - r| + \dist(\pi_k(r),\pi(r)) \tendsto k 0.
%%%\]
%%%On the other hand, suppose $r = \pm\infty$; without loss of generality suppose $r = +\infty$. Fix $N\in\N$. For all $k$ sufficiently large, we have $r_k\geq N$ and therefore
%%%\[
%%%\lb \pi_k(N)|z_k\rb_{\pi_k(0)} = N.
%%%\]
%%%Moreover, since $\pi_k(N)\to\pi(N)$ and $\pi_k(0)\to\pi(0)$, we have
%%%\[
%%%\lb \pi(N)|z_k\rb_{\pi(0)} \geq N - 1
%%%\]
%%%for all $k$ sufficiently large. On the other hand
%%%\[
%%%\lb \pi(N)|z\rb_{\pi(0)} = N,
%%%\]
%%%so by Gromov's inequality we have
%%%\[
%%%\lb z_k|z\rb_{\pi(0)} \gtrsim_\plus N.
%%%\]
%%%Since $N$ was arbitrary, we have $z_k\to z$.
%%%\end{subproof}
%%%
%%%Since $\pi_k(t_k) = x_k \to x$ and $\pi_k(s_k) = y_k \to y$, we have $\pi(t) = x$ and $\pi(s) = y$, i.e. $\pi$ connects $x$ and $y$.
%%%
%%%To show that $\geo{x_k}{y_k}\to \geo xy$, [A NEW ARGUMENT IS NEEDED] for each $k\in\N$ fix $z_k\in\geo{x_k}{y_k}$. Write $z_k = \pi_k(r_k)$ for some $r_k\in[t_k,s_k]$. Let $(k_\ell)_1^\infty$ be an increasing sequence such that $r_{k_\ell}\to r\in[-\infty,+\infty]$. Then by Claim \ref{claimpirkpir}, we have $z_{k_\ell}\to z := \pi(r)$.
%%%\end{proof}
%%%
%%%}% end ignore
%%
%%Using Lemma \ref{lemmageodesicconvergence}, we prove Proposition \ref{propositiongeodesicconvergence}.
%%
%%\begin{proof}[Proof of Proposition \ref{propositiongeodesicconvergence}]~
%%\begin{itemize}
%%\item[(i)] Given distinct points $x,y\in\bord X$, we may find sequences $X\ni x_n\to x$ and $X\ni y_n\to y$. Applying Lemma \ref{lemmageodesicconvergence} proves the existence of a geodesic connecting $x$ and $y$. To show uniqueness, suppose that $\geo xy_1$ and $\geo xy_2$ are two geodesics connecting $x$ and $y$. Fix sequences $\geo xy_1\ni x_n^{(1)}\to x$, $\geo xy_2\ni x_n^{(2)}\to x$, $\geo xy_1\ni y_n^{(1)}\to y$,and $\geo xy_2\ni y_n^{(2)}\to y$. By considering the intertwined sequences $x_1^{(1)},x_1^{(2)},x_2^{(1)},x_2^{(2)},\ldots$ and $y_1^{(1)},y_1^{(2)},y_2^{(1)},y_2^{(2)},\ldots$, Lemma \ref{lemmageodesicconvergence} shows that the sequences $(\geo{x_n^{(1)}}{y_n^{(1)}})_1^\infty$ and $(\geo{x_n^{(2)}}{y_n^{(2)}})_1^\infty$ converge in the Hausdorff metric to a common geodesic $\geo xy$. But clearly the former tend to $\geo xy_1$, and the latter tend to $\geo xy_2$; we must have $\geo xy_1 = \geo xy_2$.
%%
%%
%%\item[(ii)] Suppose that $\bord X\ni x_n\to x$ and $\bord X\ni y_n\to y$. For each $n$, choose $\what x_n,\what y_n\in \geo{x_n}{y_n}\cap X$ such that $\wbar\Dist(\what x_n,x_n) , \wbar\Dist(\what y_n,y_n) \leq 1/n$. Then $\what x_n\to x$ and $\what y_n\to y$, so by Lemma \ref{lemmageodesicconvergence} we have $\geo{\what x_n}{\what y_n}\to \geo xy$ in the Hausdorff metric, or $\geo{\what x_n}{\what y_n} \to \{x\}$ if $x = y$. To complete the proof it suffices to show that the Hausdorff distance between $\geo{x_n}{y_n}$ and $\geo{\what x_n}{\what y_n}$ tends to zero as $n$ tends to infinity. Indeed, $\geo{\what x_n}{\what y_n} \subset \geo{x_n}{y_n}$, and for each $z\in \geo{x_n}{y_n}$, either $z\in \geo{x_n}{\what x_n}$, $z\in \geo{\what x_n}{\what y_n}$, or $z\in \geo{\what y_n}{y_n}$. In the first case, Lemma \ref{lemmageodesicdiameter} shows that $\wbar\Dist(z,\geo{\what x_n}{\what y_n}) \leq \wbar\Dist(z,\what x_n) \lesssim_\times \wbar\Dist(x_n,\what x_n) \leq 1/n \to 0$; the third case is treated similarly.
%%
%%\end{itemize}
%%\end{proof}
%%
%%
%%Having completed the proof of Proposition \ref{propositiongeodesicconvergence}, in the remainder of this subsection we prove that a version of the CAT(-1) equality holds for ideal triangles.
%%
%%\begin{definition}
%%\label{definitiontriangles}
%%A \emph{geodesic triangle} $\Delta = \Delta(x,y,z)$ consists of three distinct points $x,y,z\in\bord X$ together with the geodesics $\geo xy$, $\geo yz$, and $\geo zx$.
%%
%%A geodesic triangle $\wbar{\Delta} = \Delta(\wbar x,\wbar y,\wbar z)$ is called a \emph{comparison triangle} for $\Delta$ if
%%\[
%%\lb \wbar x|\wbar y\rb_{\wbar z} = \lb x|y\rb_z, \text{ etc.}
%%\]
%%For any point $p\in\geo xy$, its \emph{comparison point} is defined to be the unique point $\wbar p\in\geo {\wbar x}{\wbar y}$ such that
%%\[
%%\lb x|z\rb_p - \lb y|z\rb_p = \lb \wbar x|\wbar z\rb_{\wbar p} - \lb \wbar y|\wbar z\rb_{\wbar p}.
%%\]
%%We say that the geodesic triangle $\Delta$ \emph{satisfies the CAT(-1) inequality} if for all $p,q\in\Delta$ and for any comparison points $\wbar p,\wbar q\in\wbar\Delta$, we have $\dist(p,q) \leq \dist(\wbar p,\wbar q)$.
%%\end{definition}
%%
%%It should be checked that these definitions are consistent with those given in Subsection \ref{subsectionCAT}.
%%
%%\begin{proposition}
%%\label{propositionidealCAT}
%%Any geodesic triangle (including ideal triangles) satisfies the CAT(-1) inequality.
%%\end{proposition}
%%\begin{proof}
%%Let $\Delta = \Delta(x,y,z)$ be a geodesic triangle, and fix $p,q\in \Delta$. Choose sequences $x_n\to x$, $y_n\to y$, and $z_n\to z$. By Proposition \ref{propositiongeodesicconvergence}, we have $\Delta_n = \Delta(x_n,y_n,z_n) \to \Delta$ in the Hausdorff metric, so we may choose $p_n,q_n\in\Delta_n$ so that $p_n\to p$, $q_n\to q$. For each $n$, let $\wbar\Delta_n = \Delta(\wbar x_n,\wbar y_n,\wbar z_n)$ be a comparison triangle for $\Delta_n$. Without loss of generality, we may assume that
%%\begin{equation}
%%\label{zerocenter}
%%\zero\in\geo{\wbar x_n}{\wbar y_n} \text{ and } \lb \wbar x_n|\wbar z_n \rb_\zero = \lb \wbar y_n | \wbar z_n \rb_\zero \asymp_\plus 0.
%%\end{equation}
%%By extracting a convergent subsequence, we may without loss of generality assume that $\wbar x_n\to \wbar x$, $\wbar y_n\to \wbar y$, and $\wbar z_n\to \wbar z$ for some points $\wbar x,\wbar y,\wbar z\in\bord \H^2$. By \eqref{zerocenter}, the points $\wbar x,\wbar y,\wbar z$ are distinct. Thus $\wbar\Delta = \Delta(\wbar x,\wbar y,\wbar z)$ is a geodesic triangle, and is in fact a comparison triangle for $\Delta$. If $\wbar p,\wbar q$ are comparison points for $p,q$, then $\wbar p_n\to \wbar p$ and $\wbar q_n\to q$. It follows that
%%\[
%%\dist(p,q) = \lim_{n\to\infty} \dist(p_n,q_n) \leq \lim_{n\to\infty} \dist(\wbar p_n,\wbar q_n) = \dist(\wbar p,\wbar q).
%%\]
%%\end{proof}
%%
%%
%%%
%%%\begin{definition}
%%%Two embeddings $\pi_1:[t_1,s_1]\to X$ and $\pi_2:[t_2,s_2]\to X$ are \emph{left-aligned} if either
%%%\begin{itemize}
%%%\item $t_1 = t_2 \neq -\infty$ and $\pi_1(t_1) = \pi_2(t_2)$, or
%%%\item $t_1 = t_2 = -\infty$, $\xi := \pi_1(t_1) = \pi_2(t_2)$, and
%%%\[
%%%\busemann_\xi(\pi_1(t),\pi_2(t)) = 0 \all t\leq \min(s_1,s_2).
%%%\]
%%%\end{itemize}
%%%The notion of \emph{right-aligned} embeddings is defined similarly.
%%%\end{definition}
%%%
%%%The key technical lemma in proving Propositions \ref{propositiongeodesicconvergence} and \ref{propositionidealCAT} is the following:
%%%
%%%
%%%
%%%\begin{lemma}
%%%Suppose that $t_n\tendsto n +\infty$ and that for each $n$, $\pi_n:[0,t_n]\to X$ is an isometric embedding. If $\pi_n(0)\tendsto n x\in X$ and $\pi_n(t_n)\tendsto n\eta\in\del X$, then for each $t\in\Rplus$ the sequence $(\pi_n(t))_1^\infty$ is Cauchy (and therefore converges).
%%%\end{lemma}
%%%\begin{proof}
%%%Fix $n,m\in\N$ and let $\wbar{\Delta} = \Delta(\wbar x,\wbar{y_n},\wbar{y_m})$ be a comparison triangle for the geodesic triangle $\Delta = \Delta(x,y_n,y_m)$. Let $z_n = \pi_n(t)$; then
%%%\[
%%%\dist(z_n,z_m) \leq \dist(\wbar{z_n},\wbar{z_m}).
%%%\]
%%%Now we note that
%%%\[
%%%\dist(x,z_n) = \dist(x,z_m) = t
%%%\]
%%%and that
%%%\[
%%%\lb \wbar{y_n}|\wbar{y_m}\rb_{\wbar x} \tendsto{n,m} +\infty \Rightarrow \angle_{\wbar x}(\wbar{y_n},\wbar{y_m}) \tendsto{n,m} 0.
%%%\]
%%%These two equations together with the law of cosines (Proposition \ref{lawofcosines}) imply that $\dist(\wbar{z_n},\wbar{z_m})\tendsto{n,m} 0$. This shows that $(z_n)_1^\infty$ is a Cauchy sequence. \comdavid{Use a compactness argument instead}
%%%\end{proof}
%%%
%%%
%%%\begin{proof}
%%%
%%%Let $x\in X$ and let $(y_n)_1^\infty$ be a sequence in $X$ which converges to a point $\eta\in\del X$. Then we have a sequence of isometric embeddings $\pi_n:[0,t_n]\to X$, where $t_n = \dist(x,y_n)$, satisfying $\pi_n(0) = x$ and $\pi_n(t_n) = y_n$. The idea now is to take a limit of the $\pi_n$s. Fix $t\geq 0$, and we claim that the sequence $(\pi_n(t))_1^\infty$ is Cauchy. Indeed, fix $n,m\in\N$ and let $\wbar{\Delta} = \Delta(\wbar x,\wbar{y_n},\wbar{y_m})$ be a comparison triangle for the geodesic triangle $\Delta = \Delta(x,y_n,y_m)$. Let $z_n = \pi_n(t)$; then
%%%\[
%%%\dist(z_n,z_m) \leq \dist(\wbar{z_n},\wbar{z_m}).
%%%\]
%%%Now we note that
%%%\[
%%%\dist(x,z_n) = \dist(x,z_m) = t
%%%\]
%%%and that
%%%\[
%%%\lb \wbar{y_n}|\wbar{y_m}\rb_{\wbar x} \tendsto{n,m} +\infty \Rightarrow \angle_{\wbar x}(\wbar{y_n},\wbar{y_m}) \tendsto{n,m} 0.
%%%\]
%%%These two equations together with the law of cosines (Proposition \ref{lawofcosines}) imply that $\dist(\wbar{z_n},\wbar{z_m})\tendsto{n,m} 0$. This shows that $(z_n)_1^\infty$ is a Cauchy sequence. \comdavid{Use a compactness argument instead}
%%%
%%%Let $\xi\in\del X$ and let $(y_n)_1^\infty$ be a sequence in $X$ which converges to a point $\eta\in\del X$. Then we have a sequence of isometric embeddings $\pi_n:[0,t_n]\to X$, where $t_n = \busemann_\xi(y_n,\zero)$, satisfying $\pi_n(t_n) = y_n$. The idea now is to take a limit of the $\pi_n$s. Fix $t\in\R$, and we claim that the sequence $(\pi_n(t))_1^\infty$ is Cauchy. Indeed, fix $n,m\in\N$ and let $\wbar{\Delta} = \Delta(\wbar{\xi},\wbar{y_n},\wbar{y_m})$ be a comparison triangle for the geodesic triangle $\Delta = \Delta(\xi,y_n,y_m)$. Let $z_n = \pi_n(t)$; then
%%%\[
%%%\dist(z_n,z_m) \leq \dist(\wbar{z_n},\wbar{z_m}).
%%%\]
%%%Now we note that
%%%\[
%%%\busemann_\xi(z_n,\zero) = \dist(z_m,\zero) = t
%%%\]
%%%and that
%%%\[
%%%\lb \wbar{y_n}|\wbar{y_m}\rb_{\wbar{\xi,\zero}} \tendsto{n,m} +\infty \Rightarrow \angle_{\wbar{\xi,\zero}}(\wbar{y_n},\wbar{y_m}) \tendsto{n,m} 0.
%%%\]
%%%These two equations together with the law of Busemann cosines imply that $\dist(\wbar{z_n},\wbar{z_m})\tendsto{n,m} 0$. This shows that $(z_n)_1^\infty$ is a Cauchy sequence.
%%%
%%%Law of Busemann cosines:
%%%\[
%%%\cosh(a) = \cosh(b - c) + \frac{e^{b + c}}{8}\alpha^2,
%%%\]
%%%where
%%%\begin{align*}
%%%a &= \dist(x,y)\\
%%%b &= \busemann_\xi(x,\zero)\\
%%%c &= \busemann_\xi(y,\zero)\\
%%%\alpha &= \angle_{\xi,\zero}(x,y) :=\lim_{z\to\xi} \frac{\angle_z(x,y)}{\dox z}
%%%\end{align*}
%%%In $\H^{d + 1}$ we have
%%%\begin{align*}
%%%a &= \dist_\H(\xx,\yy)\\
%%%b &= -\log(x_{d + 1})\\
%%%c &= -\log(y_{d + 1})\\
%%%\alpha &= \|\pi[\xx] - \pi[\yy]\|.
%%%\end{align*}
%%%Here $\pi[\xx] = \xx - x_{d + 1}\ee_{d + 1}$. Thus
%%%\begin{align*}
%%%\cosh\dist_\H(\xx,[\yy])
%%%&= \cosh(\log(x_{d + 1}/y_{d + 1})) + \frac{1}{8x_{d + 1}y_{d + 1}}\|\pi[\xx] - \pi[\yy]\|^2\\
%%%&= \frac{x}{2y} + \frac{y}{2x} + \frac{1}{8xy}\|\pi[\xx] - \pi[\yy]\|^2\\
%%%&= \frac{1}{8xy}\left[4x^2 + 4y^2 + \|\pi[\xx] - \pi[\yy]\|^2\right]\\
%%%\frac{\cosh + 1}{\cosh - 1} &= \frac{\|\pi[\xx] - \pi[\yy]\|^2 + 4(x + y)^2}{\|\pi[\xx] - \pi[\yy]\|^2 + 4(x - y)^2}\\
%%%&= \frac{(e^d + 1)^2}{(e^d - 1)^2}
%%%\end{align*}
%%%
%%%
%%%\end{proof}
%%
%%
%%
%%
%%
%%%\ignore{
%%%\begin{proposition}
%%%Let $X$ be a proper hyperbolic metric space satisfying (i) of Proposition \ref{propositiongeodesicconvergence}. Then $X$ is regularly geodesic.
%%%\end{proposition}
%%%\begin{proof}
%%%Suppose that $(x_n)_1^\infty$ and $(y_n)_1^\infty$ are sequences in $\bord X$ which converge to points $x_n\tendsto n x\in\bord X$ and $y_n\tendsto n y\in\bord X$. Suppose that $(z_n)_1^\infty$ is a bounded sequence such that $z_n\in \geo{x_n}{y_n}$.
%%%
%%%Since $X$ is proper, there is a sequence $(n_k)_1^\infty$ and a point $z\in X$ so that $z_{n_k}\tendsto k z$. To complete the proof, we need to show that $z\in\geo xy$.
%%%
%%%\end{proof}
%%%}% end ignore
%%
%%
%%
%%%%%%%%%%%%%%%%%%%%%%%%%%%%%%%%%%%%
%%\bigskip
%%\subsection{The geometry of shadows}\label{subsectionshadows}
%%%%%%%%%%%%%%%%%%%%%%%%%%%%%%%%%%%%
%%
%%\subsubsection{Shadows in regularly geodesic hyperbolic metric spaces}
%%Suppose that $X$ is regularly geodesic. For each $z\in X$ we consider the relation $\pi_z\subset X\times\del X$ defined by
%%\[
%%(x,\xi)\in \pi_z \Leftrightarrow x\in \geo z\xi
%%\]
%%(see Definition \ref{definitiongeodesic} for the definition of $\geo z\xi$). We remark that if $X$ is a \ROSS, then the relation $\pi_z$ is a function when restricted to $X\butnot\{z\}$; in particular, for $\xx\in \B = \B_\F^\alpha$ with $\xx\neq\0$ we have
%%\[
%%\pi_\0(\xx) = \frac{\xx}{\|\xx\|}\cdot
%%\]
%%However, in general the relation $\pi_z$ is not necessarily a function; $\R$-trees provide a good counterexample. The reason is that in an $\R$-tree, there may be multiple ways to extend a geodesic segment to a geodesic ray.
%%
%%For any set $S$, we define its \emph{shadow} with respect to the light source $z$ to be the set
%%\[
%%\pi_z(S) := \{\xi\in\del X:\exists x\in S \;\; (x,\xi)\in \pi_z\}.
%%\]
%%
%%
%%\begin{figure}
%%\begin{center}
%%\begin{tikzpicture}[line cap=round,line join=round,>=triangle 45,x=1.0cm,y=1.0cm]
%%\clip(-3.85,-3.03) rectangle (3.99,3.1);
%%\draw(0,0) circle (3cm);
%%\draw[line width=0.81pt](1.7,1.4) circle (0.44cm);
%%\draw (0,0)-- (1.89,2.33);
%%\draw (0,0)-- (2.65,1.4);
%%\draw [shift={(0,0)},line width=1.2pt] plot[domain=0.49:0.89,variable=\t]({1*3*cos(\t r)+0*3*sin(\t r)},{0*3*cos(\t r)+1*3*sin(\t r)});
%%\draw [shift={(1.83,2.51)}] plot[domain=4.19:4.78,variable=\t]({1*0.96*cos(\t r)+0*0.96*sin(\t r)},{0*0.96*cos(\t r)+1*0.96*sin(\t r)});
%%\begin{scriptsize}
%%\fill [color=black] (0,0) circle (1pt);
%%\draw[color=black] (-0.1,-0.16) node {$z$};
%%\draw[color=black] (0.72,1.66) node {$B(x,\sigma)$};
%%\fill [color=black] (1.89,1.55) circle (1pt);
%%\draw[color=black] (2,1.44) node {$x$};
%%\draw[color=black] (1.59,1.46) node {$\sigma$};
%%\draw[color=black] (3.2,1.98) node {$\pi_z\big(B(x,\sigma)\big)$};
%%\end{scriptsize}
%%\end{tikzpicture}
%%\caption{The set $\pi_z(B(x,\sigma))$. Although this set is not equal to $\Shad_z(x,\sigma)$, they are approximately the same in regularly geodesic spaces by Corollary \ref{corollaryshadowsrelation}. In our drawings, we will draw the set $\pi_z(B(x,\sigma))$ to indicate the set $\Shad_z(x,\sigma)$ (since the latter is hard to draw).}
%%\label{figureshadow}
%%\end{center}
%%\end{figure}
%%
%%
%%\subsubsection{Shadows in hyperbolic metric spaces}
%%In regularly geodesic hyperbolic metric spaces, it is particularly useful to consider $\pi_z(B(x,\sigma))$ where $x\in X$ and $\sigma > 0$. We would like to have an analogue for this set in the Gromov hyperbolic setting.
%%
%%\begin{definition}
%%\label{definitionshadow}
%%For each $\sigma > 0$ and $x,z\in X$, let
%%\[
%%\Shad_z(x,\sigma) = \{\eta\in \del X:\lb z|\eta\rb_x\leq\sigma\}.
%%\]
%%We say that $\Shad_z(x,\sigma)$ is the \emph{shadow cast by $x$ from the light source $z$, with parameter $\sigma$}. For shorthand we will write $\Shad(x,\sigma) = \Shad_\zero(x,\sigma)$.
%%\end{definition}
%%The relation between $\pi_z(B(x,\sigma))$ and $\Shad_z(x,\sigma)$ in the case where $X$ is a regularly geodesic hyperbolic metric space will be made explicit in Corollary \ref{corollaryshadowsrelation} below.
%%
%%Let us establish up front some geometric properties of shadows.
%%
%%\begin{observation}
%%\label{observationshadowsclosed}
%%For each $x,z\in X$ and $\sigma > 0$ the set $\Shad_z(x,\sigma)$ is closed.
%%\end{observation}
%%\begin{proof}
%%This follows directly from Lemma \ref{lemmalowersemicontinuous}.
%%\end{proof}
%%
%%\begin{observation}\label{observationmetricshadow}
%%If $\eta\in \Shad_z(x,\sigma)$, then
%%\[
%%\lb x|\eta\rb_z \asymp_{\plus,\sigma}%\Footnote{Here and from now on $\asymp_{\plus,\sigma}$ means that the implied constant may depend on $\sigma$.}
%% \dist(z,x).
%%%\lb x|\eta\rb_\zero
%%%\asymp_\plus \lb x|\eta\rb_\zero + \lb\zero|\eta\rb_x
%%%\lesssim_\plus \lb x|\eta\rb_\zero + \sigma.
%%\]
%%\end{observation}
%%\begin{proof}
%%Follows directly from (b) of Proposition \ref{propositionbasicidentities} together with the definition of $\Shad_z(x,\sigma)$.
%%\end{proof}
%%
%%
%%
%%
%%
%%\begin{lemma}[Intersecting Shadows Lemma]
%%\label{lemmatau}
%%For each $\sigma > 0$, there exists $\tau = \tau_\sigma > 0$ such that for all $x,y,z\in X$ satisfying $\dist(z,y)\geq \dist(z,x)$ and $\Shad_z(x,\sigma)\cap \Shad_z(y,\sigma)\neq\emptyset$, we have
%%\begin{equation}
%%\label{shadcontainment}
%%\Shad_z(y,\sigma) \subset\Shad_z(x,\tau)
%%\end{equation}
%%and
%%\begin{equation}
%%\label{distbusemann}
%%\dist(x,y) \asymp_{\plus,\sigma} \dist(z,y) - \dist(z,x).
%%\end{equation}
%%\end{lemma}
%%
%%\begin{figure}
%%\begin{center}
%%\definecolor{ttqqqq}{rgb}{0.2,0,0}
%%\definecolor{ffttww}{rgb}{1,0.1,0.1}
%%\begin{tikzpicture}[line cap=round,line join=round,>=triangle 45,x=1.0cm,y=1.0cm]
%%\clip(-4.88,-4.1) rectangle (5.4,4.1);
%%\draw(0,0) circle (4.03cm);
%%\draw(1.84,2.21) circle (0.76cm);
%%\draw(2.92,1.93) circle (0.25cm);
%%\draw [shift={(0,0)},line width=2pt,color=ffttww] plot[domain=0.61:1.14,variable=\t]({1*4.03*cos(\t r)+0*4.03*sin(\t r)},{0*4.03*cos(\t r)+1*4.03*sin(\t r)});
%%\draw [shift={(0,0)},line width=1.2pt,color=yellow] plot[domain=0.66:0.61,variable=\t]({1*4.03*cos(\t r)+0*4.03*sin(\t r)},{0*4.03*cos(\t r)+1*4.03*sin(\t r)});
%%\draw [shift={(0,0)},dash pattern=on 1pt off 2pt,line width=1pt,color=blue] plot[domain=0.51:0.66,variable=\t]({1*4.03*cos(\t r)+0*4.03*sin(\t r)},{0*4.03*cos(\t r)+1*4.03*sin(\t r)});
%%\draw (1.15,2.53)-- (1.67,3.67);
%%\draw (2.28,1.59)-- (3.31,2.31);
%%\draw (2.76,2.13)-- (3.19,2.47);
%%\draw (3.04,1.71)-- (3.51,1.98);
%%\begin{scriptsize}
%%\fill [color=black] (0,0) circle (1pt);
%%\draw[color=black] (-0.08,-0.34) node {$z$};
%%\fill [color=black] (2.07,2.49) circle (1pt);
%%\draw[color=black] (0.5,2.13) node {$B(x,\sigma)$};
%%\fill [color=black] (3.01,1.99) circle (1pt);
%%\draw[color=black] (3.1,1.4) node {$B(y,\sigma)$};
%%\draw[color=ffttww] (3.22,3.36) node {$\Shad_z(x,\sigma)$};
%%\draw[color=blue] (4.38,2.1) node {$\Shad_z(y,\sigma)$};
%%\end{scriptsize}
%%\end{tikzpicture}
%%\caption{In this figure, $\dist(z,y)\geq \dist(z,x)$ and $\Shad_z(x,\sigma)\cap \Shad_z(y,\sigma)\neq\emptyset$. The Intersecting Shadows Lemma (Lemma \ref{lemmatau}) provides a $\tau_\sigma > 0$ such that the shadow cast from $z$ about $B(x,\tau_\sigma)$ will capture $\Shad_z(y,\sigma)$.}
%%\label{figureintersectingshadows}
%%\end{center}
%%\end{figure}
%%
%%\begin{proof}
%%Fix $\eta\in\Shad_z(x,\sigma)\cap \Shad_z(y,\sigma)$, so that by Observation \ref{observationmetricshadow}
%%\[
%%\lb x|\eta\rb_z \asymp_{\plus,\sigma} \dist(z,x)\ \text{ and } \
%%\lb y|\eta\rb_z \asymp_{\plus,\sigma} \dist(z,y) \geq \dist(z,x).
%%\]
%%Gromov's inequality along with (c) of Proposition \ref{propositionbasicidentities} then gives
%%\begin{equation}\label{distbusemann2}
%%\lb x|y\rb_z \asymp_{\plus,\sigma} \dist(z,x).
%%\end{equation}
%%%\ignore{
%%%We then successively deduce the following asymptotics:
%%%\begin{align*}
%%%\frac12(\dist(z,y) + \dist(z,x) - \dist(x,y)) &\asymp_{\plus,\sigma} \dist(z,x)\\
%%%\frac12(\dist(z,y) - \dist(z,x) - \dist(x,y)) &\asymp_{\plus,\sigma} 0\\
%%%\dist(z,y) - \dist(z,x) - \dist(x,y) &\asymp_{\plus,\sigma} 0\\
%%%B_z(y,x) &\asymp_{\plus,\sigma} \dist(x,y).
%%%\end{align*}
%%%}% end ignore
%%Rearranging yields \eqref{distbusemann}. In order to show \eqref{shadcontainment}, fix $\xi\in\Shad_z(y,\sigma)$, so that $\lb y|\xi\rb_z \asymp_{\plus,\sigma} \dist(z,y)\geq \dist(z,x)$. Gromov's inequality and \eqref{distbusemann2} then give
%%\[
%%\lb x|\xi\rb_z \asymp_{\plus,\sigma} \dist(z,x),
%%\]
%%i.e. $\xi\in\Shad_z(x,\tau)$ for some $\tau > 0$ sufficiently large (depending on $\sigma$).
%%\end{proof}
%%
%%\begin{corollary}
%%\label{corollaryshadowsrelation}
%%Suppose that $X$ is regularly geodesic. For every $\sigma > 0$, there exists $\tau = \tau_\sigma > 0$ such that for any $x,z\in X$ we have
%%\begin{equation}
%%\label{shadowsrelation}
%%\pi_z(B(x,\sigma)) \subset \Shad_z(x,\sigma) \subset \pi_z(B(x,\tau)).
%%\end{equation}
%%\end{corollary}
%%\begin{proof}
%%Suppose $\xi\in\pi_z(B(x,\sigma))$. Then there exists a point $y\in B(x,\sigma)\cap \geo z\xi$. By (d) of Proposition \ref{propositionbasicidentities}
%%\begin{align*}
%%\lb z|\xi\rb_x \leq \lb z|\xi\rb_y + \dist(x,y)
%%\leq \lb z|\xi\rb_y + \sigma = \sigma,
%%\end{align*}
%%i.e. $\xi\in\Shad_z(x,\sigma)$. This demonstrates the first inclusion of \eqref{shadowsrelation}. On the other hand, suppose that $\xi\in\Shad_z(x,\sigma)$. Let $y\in\geo z\xi$ be the unique point so that $\dist(z,y) = \dist(z,x)$. Clearly $\xi\in\Shad_z(y,\sigma)$, so $\Shad_z(x,\sigma)\cap\Shad_z(y,\sigma)\neq\emptyset$; by the Intersecting Shadows Lemma \ref{lemmatau} we have
%%\[
%%\dist(x,y) \asymp_{\plus,\sigma} \busemann_\xi(y,x) = 0
%%\]
%%i.e. $\dist(x,y)\leq\tau$ for some $\tau = \tau_\sigma > 0$ depending only on $\sigma$. Then $y\in B(x,\tau)\cap\geo z\xi$, which implies that $\xi = \pi_z(y)\in \pi_z(B(x,\tau))$. This finishes the proof.
%%\end{proof}
%%
%%\begin{lemma}[Bounded Distortion Lemma]
%%\label{lemmaboundeddistortion}
%%Fix $\sigma > 0$. Then for every $g\in\Isom(X)$ and for every $y\in\Shad_{g^{-1}(\zero)}(\zero,\sigma)$ we have
%%\begin{equation}
%%\label{boundeddistortion1}
%%\overline g'(y) \asymp_{\times,\sigma} b^{-\dogo g}.
%%\end{equation}
%%Moreover, for every $y_1,y_2\in \Shad_{g^{-1}(\zero)}(\zero,\sigma)$, we have
%%\begin{equation}
%%\label{boundeddistortion2}
%%\frac{\Dist (g(y_1),g(y_2) )}{\Dist(y_1,y_2)} \asymp_{\times,\sigma} b^{-\dogo g}.
%%\end{equation}
%%\end{lemma}
%%\begin{proof}
%%We have $\overline g'(y) \asymp_\times b^{\busemann_y(\zero,g^{-1}(\zero))}
%%\asymp_\times b^{2\lb g^{-1}(\zero)|y\rb_\zero - \dogo g}
%%\asymp_{\times,\sigma} b^{-\dogo g}$, giving \eqref{boundeddistortion1}.
%%%\begin{align*}
%%%\frac{\Dist (g(y_1),g(y_2) )}{\Dist(y_1,y_2)} &\asymp_\times \frac{b^{- \lb g(y_1) | g(y_2) \rb_\zero}}{b^{- \lb y_1 | y_2 \rb_\zero}} \\
%%%&= b^{-[\lb y_1 | y_2 \rb_{g^{-1}(\zero)} - \lb y_1 | y_2 \rb_\zero ]} \\
%%%&\asymp_\times b^{\frac12 [\busemann_{y_1} (\zero, g^{-1}(\zero)) + \busemann_{y_2} (\zero, g^{-1}(\zero)) ]} \by{the change of veiwpoint formula} \\
%%%&\asymp_{\times,\sigma} b^{- \dist (\zero, g(\zero))} \since{$\busemann_a(b,c) \asymp_\plus \dist(b,c) - 2 \lb a | b \rb_c$} . 
%%%\end{align*}
%%%Note that the last asymptotic is where we used the hypothesis that $y_1,y_2\in\Shad_{g^{-1}(\zero)}(\zero,\sigma)$. Also note that both asymptotics in $\lb a | b \rb_c - \lb a | b \rb_d \asymp_\plus \frac12 \Bigl[ \busemann_a (d,c) + \busemann_b (d,c) \Bigr]$ and $\busemann_a(b,c) \asymp_\plus \dist(b,c) - 2 \lb a | b \rb_c$ hold with equality in CAT(-1) spaces.
%%Now \eqref{boundeddistortion2} follows from \eqref{boundeddistortion1} and the geometric mean value theorem (Proposition \ref{propositionGMVT}).
%%\end{proof}
%%
%%\begin{lemma}[Big Shadows Lemma]
%%\label{lemmabigshadow}
%%For every $\epsilon > 0$, for every $\sigma > 0$ sufficiently large (depending on $\epsilon$), and for every $z\in X$, we have
%%\begin{equation}
%%\label{bigshadow}
%%\Diam(\del X\butnot\Shad_z(\zero,\sigma)) \leq \epsilon.
%%\end{equation}
%%\end{lemma}
%%
%%\begin{figure}
%%\begin{center}
%%\begin{tikzpicture}[line cap=round,line join=round,>=triangle 45,x=1.0cm,y=1.0cm]
%%\clip(-5.14,-4.17) rectangle (5.05,4.2);
%%\draw(0,0) circle (4.12cm);
%%\draw(0,0) circle (3.07cm);
%%\draw [shift={(-3.49,-2.35)}] plot[domain=-0.28:1.95,variable=\t]({1*1.14*cos(\t r)+0*1.14*sin(\t r)},{0*1.14*cos(\t r)+1*1.14*sin(\t r)});
%%\draw [shift={(-1.98,-3.76)}] plot[domain=-0.24:1.94,variable=\t]({1*1.18*cos(\t r)+0*1.18*sin(\t r)},{0*1.18*cos(\t r)+1*1.18*sin(\t r)});
%%\draw [shift={(0,0)},line width=1.2pt] plot[domain=3.46:4.51,variable=\t]({1*4.12*cos(\t r)+0*4.12*sin(\t r)},{0*4.12*cos(\t r)+1*4.12*sin(\t r)});
%%\draw (0,0)-- (2.56,1.7);
%%\begin{scriptsize}
%%\fill [color=black] (0,0) circle (1.0pt);
%%\draw[color=black] (-0.23,0.12) node {$\zero$};
%%\fill [color=black] (-2.5,-2.66) circle (1.0pt);
%%\draw[color=black] (-2.19,-2.84) node {$z$};
%%\draw[color=black] (-3.93,-3.21) node {$\del X\butnot\Shad_z(\zero,\sigma)$};
%%\draw[color=black] (1.97,1.08) node {$\sigma$};
%%\end{scriptsize}
%%\end{tikzpicture}
%%\caption{The Big Shadows Lemma \ref{lemmabigshadow} tells us that for any $\epsilon > 0$, we may choose $\sigma > 0$ sufficiently large so that $\Diam(\del X\butnot\Shad_z(\zero,\sigma)) \leq \epsilon$ for every $z\in X$.}
%%\label{figurebigshadows}
%%\end{center}
%%\end{figure}
%%
%%\begin{proof}
%%If $\xi, \eta\in\del X \setminus \Shad_z(\zero, \sigma)$, then $\lb z | \xi \rb_\zero > \sigma$ and $\lb z | \eta \rb_\zero > \sigma$. Thus by Gromov's inequality we have
%%\[
%%\lb \xi | \eta \rb_\zero \gtrsim_\plus \sigma.
%%\]
%%Exponentiating gives $\Dist(\xi,\eta) \lesssim_\times b^{-\sigma}$. Thus
%%\[
%%\Diam (\del X \setminus\Shad_z (\zero,\sigma)) \lesssim_\times b^{-\sigma} \tendsto\sigma 0,
%%\]
%%and the convergence is uniform in $z$.
%%\end{proof}
%%
%%\begin{figure}
%%\begin{center}
%%\begin{tikzpicture}[line cap=round,line join=round,>=triangle 45,x=1.0cm,y=1.0cm]
%%\clip(-3.54,-3.1) rectangle (3.943,3.38);
%%\draw(0.0,0.0) circle (3.0cm);
%%\draw(1.5,1.5) circle (0.75cm);
%%\draw (0.0,-0.0)-- (1.2343134832984426,2.7343134832984424);
%%\draw (0.0,-0.0)-- (2.7343134832984424,1.2343134832984426);
%%\draw (0.0,-0.0)-- (1.7415599609449397,1.7700977248889884);
%%\draw [shift={(1.317194030436857,2.707806971266285)}] plot[domain=4.204270499810774:5.13736643041908,variable=\t]({1.0*1.0292643361729137*cos(\t r)+-0.0*1.0292643361729137*sin(\t r)},{0.0*1.0292643361729137*cos(\t r)+1.0*1.0292643361729137*sin(\t r)});
%%\draw [shift={(0.0,0.0)},line width=1.2pt] plot[domain=0.4240310394907405:1.146765287304156,variable=\t]({1.0*3.0*cos(\t r)+-0.0*3.0*sin(\t r)},{0.0*3.0*cos(\t r)+1.0*3.0*sin(\t r)});
%%\begin{scriptsize}
%%\draw [fill=black] (0.0,-0.0) circle (.75pt);
%%\draw[color=black] (-0.19797292010819517,-0.09284349186632893) node {$z$};
%%%\draw[color=black] (1.378631811355094,0.9161835362701746) node {$e$};
%%\draw [fill=black] (1.7415599609449397,1.7700977248889884) circle (.75pt);
%%\draw[color=black] (1.8727829942209163,1.50) node {$g(\zero)$};
%%
%%\draw[color=black] (1.5727829942209163,0.71) node {$\Big\uparrow$};
%%\draw[color=black] (1.9727829942209163,0.31) node {$B(g(\zero),\sigma)$};
%%
%%\draw[color=black] (1.2421411016356007,1.814123570799388) node {$\sigma$};
%%\draw[color=black] (3.0115304609745355,2.3508938419017653) node {$\Shad_z(g(\zero),\sigma)$};
%%\end{scriptsize}
%%\end{tikzpicture}
%%\caption{The Diameter of Shadows Lemma \ref{lemmadiameterasymptotic} says that the diameter of $\Shad(g(\zero),\sigma)$ is coarsely asymptotic to $b^{-\dogo g}$.}
%%\end{center}
%%\end{figure}
%%
%%\begin{lemma}[Diameter of Shadows Lemma]
%%\label{lemmadiameterasymptotic}
%%For all $\sigma > 0$ sufficiently large, we have for all $g\in\Isom(X)$ and for all $z\in X$
%%\begin{equation}
%%\label{diameterasymptotic}
%%\Diam_z(\Shad_z(g(\zero),\sigma)) \lesssim_{\times,\sigma} b^{-\dist(z,g(\zero))},
%%\end{equation}
%%with $\asymp$ if $\#(\del X)\geq 3$. Moreover, for every $C > 0$ there exists $\sigma > 0$ such that
%%\begin{equation}
%%\label{ShadcontainsB}
%%B_z(x,C e^{-\dist(z,x)}) \subset \Shad_z(x,\sigma) \all x,z\in X.
%%\end{equation}
%%\end{lemma}
%%\begin{proof}
%%Let $x = g(\zero)$. For any $\xi,\eta\in \Shad_z(x,\sigma)$, we have
%%\begin{align*}
%%\Dist_z(\xi,\eta) \asymp_\times b^{-\lb\xi|\eta\rb_z}
%%\lesssim_{\times\phantom{,\sigma}} b^{-\min\left(\lb x|\xi\rb_z,\lb x|\eta\rb_z\right)}
%%\lesssim_{\times,\sigma} b^{-\dist(z,x)}
%%\end{align*}
%%which demonstrates \eqref{diameterasymptotic}.
%%
%%Now let us prove the converse of \eqref{diameterasymptotic}, assuming $\#(\del X)\geq 3$. Fix $\xi_1,\xi_2,\xi_3\in\del X$, let $\epsilon = \min_{i\neq j}\Dist(\xi_i,\xi_j)/2$, and fix $\sigma > 0$ large enough so that \eqref{bigshadow} holds for every $z\in X$. By \eqref{bigshadow} we have
%%\[
%%\Diam(\del X\butnot \Shad_{g^{-1}(z)}(\zero,\sigma)) \leq \epsilon,
%%\]
%%and thus
%%\[
%%\#\big\{i = 1,2,3: \xi_i\in \Shad_{g^{-1}(z)}(\zero,\sigma)\big\}\geq 2.
%%\]
%%Without loss of generality suppose that $\xi_1,\xi_2\in \Shad_{g^{-1}(z)}(\zero,\sigma)$. By applying $g$, we have $g(\xi_1),g(\xi_2)\in \Shad_z(x,\sigma)$. Then
%%\begin{align*}
%%\Diam_z(\Shad_z(x,\sigma))
%%&\geq_\pt \Dist_z(g(\xi_1),g(\xi_2))\\
%%&\asymp_\times b^{-\lb g(\xi_1)|g(\xi_2)\rb_z}
%%= b^{-\lb \xi_1|\xi_2\rb_{g^{-1}(z)}}\\
%%&\gtrsim_\times b^{-\lb\xi_1|\xi_2\rb_\zero}b^{-\dox{g^{-1}(z)}}
%%\asymp_{\times,\xi_1,\xi_2}b^{-\dist(z,x)}.
%%\end{align*}
%%Finally, given $y\in B_z(x,C b^{-\dist(z,x)})$, we have
%%\[
%%\lb x|y\rb_z \gtrsim_\plus -\log_b(C b^{-\dox x}) \asymp_\plus \dist(z,x)
%%\]
%%and thus $\lb z|y\rb_x \asymp_\plus 0$, demonstrating \eqref{ShadcontainsB}.
%%\end{proof}
%%
%%
%%
%%%%%%%%%%%%%%%%%%%%%%%%%%%%%%%%%%%%
%%\bigskip
%%\subsection{Generalized polar coordinates} \label{subsectionpolar}
%%%%%%%%%%%%%%%%%%%%%%%%%%%%%%%%%%%%
%%
%%Suppose that $X = \E = \E^\alpha$ is the half-space model of a real \ROSS. Fix a point $\xx\in\E$, and consider the numbers $\|\xx\|$ and $\ang(\xx) := \cos^{-1}(x_1 / \|\xx\|)$, i.e. the radial and unsigned angular coordinates of $\xx$. (The angular coordinate is computed with respect to the ray $\{(t,\0) : t\in\Rplus\}$; cf. Figure \ref{figurepolarcoordinates}.) These ``polar coordinates'' of $\xx$ do not completely determine $\xx$, but they are enough to compute certain important quantities depending on $\xx$, e.g. $\dist_\E(\zero,\xx)$, $\busemann_\infty(\zero,\xx)$, and $\busemann_\0(\zero,\xx)$. (We omit the details.) In this subsection we consider a generalization, in a loose sense, of these coordinates to an arbitrary hyperbolic metric space.
%%
%%
%%\begin{figure}
%%\begin{center}
%%\begin{tikzpicture}[line cap=round,line join=round,>=triangle 45,x=1.0cm,y=1.0cm]
%%\clip(-5.2,-0.33) rectangle (5.38,5.43);
%%\draw [shift={(0,0)},color=black,fill=black,fill opacity=0.1] (0,0) -- (42.62:0.79) arc (42.62:90:0.79) -- cycle;
%%\draw (-4,0)-- (4,0);
%%\draw (0,0)-- (0,5);
%%\draw (0,0)-- (2.76,2.54);
%%\draw [dash pattern=on 2pt off 2pt] (2.76,2.54)-- (2.78,0);
%%\begin{scriptsize}
%%\fill [color=black] (0,0) circle (1pt);
%%\draw[color=black] (-0.1,-0.2) node {$\0$};
%%\fill [color=black] (2.76,2.54) circle (1pt);
%%\draw[color=black] (3,2.6) node {$\xx$};
%%\draw[color=black] (1.44,1.82) node {$\|\xx\|$};
%%\draw[color=black] (3.06,1.09) node {$x_1$};
%%\draw[color=black] (0.43,0.9) node {$\ang(\xx)$};
%%\end{scriptsize}
%%\end{tikzpicture}
%%\caption{The quantities $\|\xx\|$ and $\ang(\xx)$ can be interpreted as ``polar coordinates'' of $\xx$.}
%%\label{figurepolarcoordinates}
%%\end{center}
%%\end{figure}
%%
%%Let us note that the isometries of $\E$ which preserve the polar coordinate functions defined above are exactly those of the form $\what T$ where $T\in\O(\EE)$. Equivalently, these are the members of $\Isom(\E)$ which preserve $\0$, $\zero = (1,\0)$, and $\infty$. This suggests that our ``coordinate system'' is fixed by choosing a point in $\E$ and two distinct points in $\del\E$.
%%
%%We now return to the general case of \6\ref{standingassumptions2}. Fix two distinct points $\xi_1,\xi_2\in\del X$.
%%\begin{definition}
%%\label{definitionpolarcoordinates}
%%The \emph{generalized polar coordinate functions} are the functions $r = r_{\xi_1,\xi_2,\zero}$ and $\theta = \theta_{\xi_1,\xi_2,\zero}:X\to\R$ defined by
%%\begin{align*}
%%r(x) &= \frac12\left[\busemann_{\xi_1}(x,\zero) - \busemann_{\xi_2}(x,\zero)\right]\\
%%\theta(x) &= \frac12\left[\busemann_{\xi_1}(x,\zero) + \busemann_{\xi_2}(x,\zero)\right] \asymp_\plus \lb \xi_1|\xi_2\rb_x - \lb \xi_1|\xi_2\rb_\zero.
%%\end{align*}
%%\end{definition}
%%
%%The connection between generalized polar coordinates and classical polar coordinates is given in Proposition \ref{propositionrtheta} below. For now, we list some geometrical facts about generalized polar coordinates. Our first lemma says that the hyperbolic distance from a point to the origin is essentially the sum of the ``radial'' distance and the ``angular'' distance.
%%
%%\begin{lemma}
%%\label{lemmafabs}
%%For all $x\in X$ we have
%%\[
%%\dox x \asymp_{\plus,\zero,\xi_1,\xi_2} \max_{i = 1}^2\busemann_{\xi_i}(x,\zero) = |r(x)| + \theta(x).
%%\]
%%\end{lemma}
%%\begin{proof}
%%The equality is trivial, so we concentrate on the asymptotic. The $\gtrsim$ direction follows directly from (f) of Proposition \ref{propositionbasicidentities}. On the other hand, by Gromov's inequality
%%\begin{align*}
%%\dox x - \max_{i = 1}^2\busemann_{\xi_i}(x,\zero)
%%\asymp_\plus 2\min_{i = 1}^2 \lb x|\xi_i\rb_\zero
%%\lesssim_\plus 2\lb \xi_1|\xi_2\rb_\zero \asymp_{\plus,\zero,\xi_1,\xi_2} 0.
%%\end{align*}
%%\end{proof}
%%
%%Our next lemma describes the effect of isometries on generalized polar coordinates.
%%
%%\begin{lemma}
%%\label{lemmatranslationdistance}
%%Fix $g\in \Isom(X)$ such that $\xi_1,\xi_2\in\Fix(g)$. For all $x\in X$ we have
%%\begin{align} \label{rtranslation}
%%r(g(x)) &\asymp_\plus r(x) + \log_b g'(\xi_1) = r(x) - \log_b g'(\xi_2)\\ \label{thetatranslation}
%%\theta(g(x)) &\asymp_\plus \theta(x),
%%\end{align}
%%with equality if $X$ is strongly hyperbolic. The implied constants are independent of $g$, $\xi_1$, and $\xi_2$.
%%\end{lemma}
%%\begin{proof}
%%\begin{align*}
%%2[r(g(x)) - r(x)]
%%&=_\pt \big[\busemann_{\xi_1}(g(x),\zero) - \busemann_{\xi_2}(g(x),\zero)\big] - \big[\busemann_{\xi_1}(x,\zero) + \busemann_{\xi_2}(x,\zero)\big] \noreason\\
%%&=_\pt \big[\busemann_{\xi_1}(x,g^{-1}(\zero)) - \busemann_{\xi_2}(x,g^{-1}(\zero))\big] - \big[\busemann_{\xi_1}(x,\zero) - \busemann_{\xi_2}(x,\zero)\big] \noreason\\
%%&\asymp_\plus \busemann_{\xi_1}(\zero,g^{-1}(\zero)) - \busemann_{\xi_2}(\zero,g^{-1}(\zero))\\
%%&\asymp_\plus \log_b g'(\xi_1) - \log_b g'(\xi_2). & \by{Proposition \ref{propositionbusemannpreserved}}
%%\end{align*}
%%Now \eqref{rtranslation} follows from Corollary \ref{corollaryproductone}.
%%
%%On the other hand, by (g) of Proposition \ref{propositionbasicidentities},
%%\begin{align*}
%%\theta(g(x)) - \theta(x) &\asymp_\plus \big[\lb \xi_1|\xi_2\rb_{g(x)} - \lb \xi_1|\xi_2\rb_\zero\big] - \big[\lb \xi_1|\xi_2\rb_x - \lb \xi_1|\xi_2\rb_\zero\big]\\
%%&=_\pt \lb g^{-1}(\xi_1)|g^{-1}(\xi_2)\rb_x - \lb \xi_1|\xi_2\rb_x = 0,
%%\end{align*}
%%proving \eqref{thetatranslation}.
%%\end{proof}
%%
%%%Combining with Lemma \ref{lemmafabs} yields the following:
%%%
%%%\begin{corollary}
%%%\label{corollarytranslationdistance}
%%%With $g,\xi_1,\xi_2$ as above,
%%%\[
%%%\dogo{g^n} \asymp_{\plus,\zero,\xi_1,\xi_2} |n\log_b g'(\xi_1)| = |n\log_b g'(\xi_2)|.
%%%\]
%%%\end{corollary}
%%
%%
%%We end this section by describing the relation between generalized polar coordinates and classical polar coordinates.
%%
%%\begin{proposition}
%%\label{propositionrtheta}
%%If $X = \E$, $\zero = (1,\0)$, $\xi_1 = \0$, and $\xi_2 = \infty$, then
%%\begin{align*}
%%r(\xx) &= \log \|\xx\| \\
%%\theta(\xx) &= -\log(x_1/\|\xx\|) = -\log \cos(\ang(\xx)).
%%\end{align*}
%%\end{proposition}
%%Thus the notations $r$ and $\theta$ are slightly inaccurate as they really represent the logarithm of the radius and the negative logarithm of the cosine of the angle, respectively.
%%\begin{proof}[Proof of Proposition \ref{propositionrtheta}]
%%We consider first the case $\|\xx\| = 1$. Let $g(\yy) = \yy/\|\yy\|^2$, and note that $g\in\Isom(\E)$, $g(\zero) = \zero$, and $g(\xi_i) = \xi_{3 - i}$. On the other hand, since $\|\xx\| = 1$ we have $g(\xx) = \xx$, and so
%%\[
%%\busemann_{\xi_1}(\xx,\zero) = \busemann_{g(\xi_1)}(g(\xx),g(\zero)) = \busemann_{\xi_2}(\xx,\zero).
%%\]
%%It follows that $r(\xx) = 0$ and $\theta(\xx) = \busemann_{\xi_2}(\xx,\zero) = \busemann_\infty(\xx,\zero)$. By Proposition \ref{propositionbusemannE}, we have $\busemann_\infty(\xx,\zero) = -\log(x_1/\zero_1) = -\log(x_1/\|\xx\|) = -\log\cos(\ang(\xx))$.
%%
%%The general case follows upon applying Lemma \ref{lemmatranslationdistance} to maps of the form $g_\lambda(\xx) = \lambda\xx$, $\lambda > 0$.
%%\end{proof}
%%
%%
%%
%%%%%%%%%%%%%%%%%%%%%%%%%%%%%%%%%%%%
%%\draftnewpage
%%\section{Discreteness} \label{sectiondiscreteness}
%%%%%%%%%%%%%%%%%%%%%%%%%%%%%%%%%%
%%
%%Let $X$ be a metric space. In this section we discuss several different notions of what it means for a group or semigroup $G\prec\Isom(X)$ to be \emph{discrete}. We show that these notions are equivalent in the Standard Case. Finally, we give examples to show that these notions are no longer equivalent when $X = \H^\infty$.
%%
%%In this section, the standing assumptions that $X$ is a (not necessarly hyperbolic) metric space and that $\zero\in X$ replace the paper's overarching standing assumption that $(X,\zero,b)$ is a Gromov triple (cf. \6\ref{standingassumptions2}). Of course, if $(X,\zero,b)$ is a Gromov triple then $X$ is a metric space and $\zero\in X$, so all theorems in this section can be used in other sections without comment.
%%
%%
%%%%%%%%%%%%%%%%%%%%%%%%%%%%%%%%%%%%
%%\subsection{Topologies on $\Isom(X)$}
%%\label{subsectiontopologies}
%%%%%%%%%%%%%%%%%%%%%%%%%%%%%%%%%%%%
%%
%%In this subsection we discuss different topologies that may be put on the isometry group of the metric space $X$.
%%
%%In the Standard Case, the most natural topology is the \emph{compact-open topology (COT)}, i.e. the topology whose subbasic open sets are of the form
%%\[
%%\GG(K,U) = \{f\in \Isom(X):f(K)\subset U\}
%%\]
%%where $K\subset X$ is compact and $U\subset X$ is open. When we replace $X$ by a metric space which is not proper, it is tempting to replace the compact-open topology with a ``bounded-open'' topology. However, it is hard to define such a topology in a way that does not result in pathologies. It turns out that the compact-open topology is still the ``right'' topology for many applications in an arbitrary metric space. But we are getting ahead of ourselves.
%%
%%Let's start by considering the case where $X$ is a \CROSS\ $X = \H = \H_\F^\alpha$, and figure out what topology or topologies we can put on $\Isom(\H)$. Recall from Theorem \ref{theoremisometries} that
%%\begin{equation}
%%\label{IsomH}
%%\Isom(\H) \equiv \PO^*(\LL;\QQ) \equiv \O^*(\LL;\QQ)/\sim
%%\end{equation}
%%where $\LL = \HH_\F^{\alpha + 1}$, $\QQ$ is the quadratic form \eqref{Qdef}, and $T_1\sim T_2$ means that $[T_1] = [T_2]$ (in the notation of Subsection \ref{subsectionisometries}). Thus $\Isom(\H)$ is isomorphic to a quotient of a subspace of $L(\LL)$, the set of bounded linear maps from $\LL$ to itself. This indicates that to define a topology or topologies on $\Isom(\H)$, it may be best to start from the functional analysis point of view and look for topologies on $L(\LL)$. In particular, we will be interested in the following widely used topologies on $L(\LL)$:
%%\begin{itemize}
%%\item The \emph{uniform operator topology (UOT)} is the topology on $L(\LL)$ which comes from looking it as a metric space with the metric
%%\[
%%\dist(T_1,T_2) = \|T_1 - T_2\| = \sup\{\|(T_1 - T_2)\xx\|:\xx\in\LL,\|\xx\| = 1\}.
%%\]
%%\item The \emph{strong operator topology (SOT)} is the topology on $L(\LL)$ which comes from looking at it as a subspace of the product space $\LL^\LL$. Note that in this topology,
%%\[
%%T_n\tendsto n T \;\;\;\; \Leftrightarrow \;\;\;\; T_n \xx \tendsto n T\xx \all \xx\in\LL.
%%\]
%%The strong operator topology is weaker than the uniform operator topology.
%%\end{itemize}
%%
%%\begin{remark}
%%There are many other topologies used in functional analysis, for example the weak operator topology, which we do not consider here.
%%\end{remark}
%%
%%Starting with either the uniform operator topology or the strong operator topology, we may restrict to the subspace $\O^*(\LL;\QQ)$ and then quotient by $\sim$ to induce a topology on $\Isom(\H)$ using the identification \eqref{IsomH}. For convenience, we will also call these induced topologies the uniform operator topology and the strong operator topology, respectively.
%%
%%We now return to the general case of a metric space $X$. Define the \emph{Tychonoff topology} to be the topology on $\Isom(X)$ inherited from the product topology on $X^X$.
%%
%%\begin{proposition}
%%\label{propositionCOTSOT}~
%%\begin{itemize}
%%\item[(i)] The Tychonoff topology and the compact-open topology on $\Isom(X)$ are identical.
%%\item[(ii)] If $X$ is a \CROSS, then the strong operator topology is identical to the Tychonoff topology (and thus also to the compact-open topology).
%%\end{itemize}
%%\end{proposition}
%%\begin{proof}~
%%\begin{itemize}
%%\item[(i)] Since subbasic sets in the Tychonoff topology take the form $\GG(\{x\},U)$, it is clear that the compact-open topology is at least as fine as the Tychonoff topology. Conversely, suppose that $\GG(K,U)$ is a subbasic open set in the Tychonoff topology, and fix $f\in \GG(K,U)$. Let $\epsilon = \dist(f(K),X\butnot U) > 0$, and let $(x_i)_1^n$ be a set of points in $K$ such that $K\subset \bigcup_1^n B(x_i,\epsilon/3)$. Then
%%\[
%%f\in\UU := \bigcap_{i = 1}^n \GG(\{x_i\},\thicken_{\epsilon/3}(f(K))).\Footnote{Here and elsewhere $\thicken_\epsilon(S) = \{x\in X : \dist(x,S) \leq \epsilon\}$.}
%%\]
%%The set $\UU$ is open in the Tychonoff topology; we claim that $\UU\subset \GG(K,U)$. Indeed, suppose that $\w f\in\UU$. Then for $x\in K$, fix $i$ with $x\in B(x_i,\epsilon/3)$; since $\w f$ is an isometry, $\dist(\w f(x),f(K)) \leq \dist(\w f(x),\w f(x_i)) + \dist(\w f(x_i),f(K)) \leq 2\epsilon/3 < \epsilon$. It follows that $\w f(x)\in U$; since $x\in K$ was arbitrary, $\w f\in \GG(K,U)$.
%%
%%\item[(ii)] It is clear that the strong operator topology is at least as fine as the Tychonoff topology. Conversely, suppose that a set $\UU\subset\Isom(\H)$ is open in the strong operator topology, and fix $[T]\in\UU$. Let $T\in\O^*(\LL;\QQ)$ be a representative of $[T]$. There exist $(\vv_i)_1^n$ in $\LL$ and $\epsilon > 0$ such that for all $\w T\in \O^*(\LL;\QQ)$ satisfying $\|(\w T - T)\vv_i\| \leq \epsilon\all i$, we have $[\w T]\in \UU$. Let $\ff_0 = \ee_0$, and let $V = \lb \ff_0,\vv_1,\ldots,\vv_n\rb$. Extend $\{\ff_0\}$ to an $\F$-basis $\{\ff_0,\ff_1,\ldots,\ff_k\}$ of $V$ with the property that $B_\QQ(\ff_{j_1},\ff_{j_2}) = 0$ for all $j_1\neq j_2$. Without loss of generality, suppose that $k\geq 1$. For each $i = 1,\ldots,n$ we have $\vv_i = \sum_j \ff_j c_{i,j}$ for some $c_{i,j}\in\F$, so there exists $\epsilon_2 > 0$ such that for all $\w T\in\O^*(\LL;\QQ)$ satisfying $\|(\w T - T)\ff_j\| \leq \epsilon_2 \all j$ and $\|\sigma_{\w T} - \sigma_T\| \leq \epsilon_2$, we have $[\w T]\in \UU$.
%%
%%Let
%%\begin{equation}
%%\label{IFdef}
%%I_\F = \begin{cases}
%%\{1\} & \F = \R\\
%%\{1,i\} & \F = \C\\
%%\{1,i,j,k\}& \F = \Q
%%\end{cases},
%%\end{equation}
%%and let
%%\[
%%F = \{\ee_0\} \cup \{\ee_0 \pm (1/2)\ff_1 \ell : j = 1,\ldots,k , \ell\in I_\F\}.
%%\]
%%Fix $\epsilon_3 > 0$ small to be determined, and for the remainder of this proof write $A\sim B$ if $\|A - B\|$ is bounded by a constant which tends to zero as $\epsilon\to 0$. Let
%%\[
%%\VV = \left\{[\w T]\in\Isom(\H): \forall \xx\in F, \;\;\exists \yy_\xx\in [\w T]([\xx]) \text{ such that } \|\yy_\xx - T\xx\| < \epsilon_3\right\}.
%%\]
%%For each $\xx\in F$, we have $[\xx]\in\H$, so the set $\{[\yy]\in\H : \exists \yy\in[\yy] \text{ such that } \|\yy - T\xx\| < \epsilon_3\}$ is open in the natural topology on $\H$. It follows that $\VV$ is open in the Tychonoff topology. Moreover, $[T]\in\VV$. To complete the proof we show that $\VV\subset\UU$. Indeed, fix $[\w T]\in\VV$, and let $\yy = \yy_{\ee_0}$. There exists a representative $\w T\in\O^*(\LL;\QQ)$ such that $\w T\ee_0 = \lambda \yy$ for some $\lambda > 0$. Since
%%\[
%%-1 = \QQ(\ee_0) \sim \QQ(\yy) = \lambda^{-2}\QQ(\lambda\yy) = -\lambda^{-2},
%%\]
%%we have $\lambda\sim 1$ and thus $\w T\ee_0 \sim T\ee_0$.
%%
%%Now for each $\xx\in F\butnot\{\ee_0\}$, there exists $a_\xx\in\F$ such that $\yy_\xx = \w T(\xx a_\xx)$. Fix $j = 1,\ldots,k$ and $\ell\in I_\F$. Writing $a_\pm = a_{\ee_0 \pm (1/2)\ff_j \ell}$, we have
%%\[
%%\left\|T\left(\ee_0 \pm \frac12\ff_j \ell\right) - \w T\left(\left(\ee_0 \pm \frac12\ff_j \ell\right)a_\pm\right)\right\| < \epsilon_3,
%%\]
%%i.e. $T(\ee_0 \pm (1/2)\ff_j \ell) \sim \w T((\ee_0 \pm (1/2)\ff_j \ell) a_\pm)$. Substituting $\pm = +$ and $\pm = -$ and adding the resulting equations gives
%%\[
%%2T\ee_0 \sim \w T(\ee_0 (a_+ + a_-)) + \frac12 \w T(\ff_j \ell (a_+ - a_-));
%%\]
%%using $T\ee_0 \sim \w T\ee_0$ and rearranging gives
%%\[
%%\w T(\ee_0 (2 - a_+ - a_-)) \sim \frac12 \w T(\ff_j \ell (a_+ - a_-)).
%%\]
%%Now by Lemma \ref{lemmaoperatornorm}, we have $\|\w T\|\sim 1$, and thus $\ee_0(2 - a_+ - a_-) \sim (1/2)\ff_j \ell(a_+ - a_-)$. Since $\|\ee_0 a + \ff_j \ell b\| \asymp_\times \max(|a|,|b|)$ for all $a,b\in\F$, it follows that
%%\[
%%2 - a_+ - a_- \sim \ell(a_+ - a_-) \sim 0,
%%\]
%%from which we deduce $a_+\sim a_- \sim 1$. Thus
%%\[
%%T\left(\ee_0 \pm \frac12\ff_j \ell\right) \sim \w T\left(\ee_0 \pm \frac12\ff_j \ell\right).
%%\]
%%Substituting $\pm = +$ and $\pm = -$, subtracting the resulting equations, and using the fact that $T\ee_0\sim \w T\ee_0$ gives
%%\[
%%T(\ff_j \ell) \sim \w T(\ff_j \ell).
%%\]
%%In particular, letting $\ell = 1$ we have $T\ff_j \sim \w T\ff_j$. Thus
%%\[
%%(T\ff_j) (\sigma_T\ell) \sim (\w T\ff_j) (\sigma_{\w T}\ell) \sim (T\ff_j) (\sigma_{\w T}\ell).
%%\]
%%Since this holds for all $\ell\in I_\F$, we have $\sigma_T\sim \sigma_{\w T}$. By the definition of $\sim$, this means that we can choose $\epsilon_3$ small enough so that $\|T\ff_j \ell - \w T\ff_j \ell\| \leq \epsilon_2 \all j$ and $\|\sigma_{\w T} - \sigma_T\| \leq \epsilon_2$. Then $[\w T]\in\UU$, completing the proof.
%%\end{itemize}
%%\end{proof}
%%
%%\begin{proposition}
%%The compact-open topology makes $\Isom(X)$ into a topological group, i.e. the maps
%%\[
%%(g,h) \mapsto gh, \hspace{.5 in}
%%g \mapsto g^{-1}
%%\]
%%are continuous.
%%\end{proposition}
%%\begin{proof}
%%Fix $g_0,h_0\in\Isom(X)$, and let $\GG(\{x\},U)$ be a neighborhood of $g_0 h_0$. For some $\epsilon > 0$, we have $B(g_0 h_0(x),\epsilon)\subset U$. We claim that
%%\[
%%\GG\Big(\{h_0(x)\},B(g_0 h_0(x),\epsilon/2)\Big) \GG\Big(\{x\},B(h_0(x),\epsilon/2)\Big) \subset \GG(\{x\},U).
%%\]
%%Indeed, fix $g\in \GG(\{h_0(x)\},B(g_0 h_0(x),\epsilon/2))$ and $h\in \GG(\{x\},B(h_0(x),\epsilon/2))$. Then
%%\[
%%\dist(gh(x),g_0 h_0(x)) \leq \dist(h(x),h_0(x)) + \dist(gh_0(x),g_0 h_0(x)) \leq \epsilon/2 + \epsilon/2 = \epsilon,
%%\]
%%demonstrating that $gh(x)\in U$, and thus that the map $(g,h) \mapsto gh$ is continuous.
%%
%%Now fix $g_0\in\Isom(X)$, and let $\GG(\{x\},U)$ be a neighborhood of $g_0^{-1}$. For some $\epsilon > 0$, we have $B(g_0^{-1}(x),\epsilon)\subset U$. We claim that
%%\[
%%\GG\big(\{g_0^{-1}(x)\},B(x,\epsilon)\big)^{-1} \subset \GG(\{x\},U).
%%\]
%%Indeed, fix $g\in \GG(\{g_0^{-1}(x)\},B(x,\epsilon))$. Then
%%\[
%%\dist(g^{-1}(x),g_0^{-1}(x)) = \dist(x,g g_0^{-1}(x)) \leq \epsilon,
%%\]
%%demonstrating that $g^{-1}(x)\in U$, and thus that the map $g \mapsto g^{-1}$ is continuous.
%%\end{proof}
%%
%%\begin{remark}[{\cite[9.B(9), p.60]{Kechris}}]
%%\label{remarkpolish}
%%If $X$ is a separable complete metric space, then the group $\Isom(X)$ with the compact-open topology is a Polish space.
%%\end{remark}
%%
%%
%%%%%%%%%%%%%%%%%%%%%%%%%%%%%%%%%%%%
%%\bigskip
%%\subsection{Discrete groups of isometries}
%%\label{subsectiondiscrete}
%%%%%%%%%%%%%%%%%%%%%%%%%%%%%%%%%%%%
%%
%%In this subsection we discuss several different notions of what it means for a group $G\leq\Isom(X)$ to be \emph{discrete}, and then we show that they are equivalent in the Standard Case. However, each of our notions will be distinct when $X = \H = \H_\F^\alpha$ for some infinite cardinal $\alpha$.
%%
%%\begin{definition}
%%\label{definitiondiscreteness}
%%Fix $G\leq\Isom(X)$.
%%\begin{itemize}
%%\item $G$ is called \emph{strongly discrete (SD)} if for every bounded set $B \subset X$, we have
%%\[
%%\#\{g\in G: g(B) \cap B \neq \emptyset\} < \infty.
%%\]
%%\item $G$ is called \emph{moderately discrete (MD)} if for every $x\in X$, there exists an open set $U\ni x$ such that
%%\[
%%\#\{g\in G: g(U) \cap U \neq \emptyset\} < \infty.
%%\]
%%\item $G$ is called \emph{weakly discrete (WD)} if for every $x\in X$, there exists an open set $U \ni x$ such that
%%\[
%%g(U) \cap U \neq \emptyset \Rightarrow g(x) = x.
%%\]
%%\end{itemize}
%%\end{definition}
%%\begin{remark}
%%Strongly discrete groups are known in the literature as \emph{metrically proper}, and moderately discrete groups are known as \emph{wandering}.
%%\end{remark}
%%\begin{remark}
%%\label{remarkdiscretenessequivalent}
%%We may equivalently give the definitions as follows:
%%\begin{itemize}
%%\item $G$ is \emph{strongly discrete (SD)} if for every $R > 0$ and $x\in X$,
%%\begin{equation}
%%\label{SD2}
%%\#\{g\in G: \dist(x,g(x)) \leq R\} < \infty.
%%\end{equation}
%%\item $G$ is \emph{moderately discrete (MD)} if for every $x\in X$, there exists $\epsilon > 0$ such that
%%\begin{equation}
%%\label{MD2}
%%\#\{g\in G: \dist(x,g(x)) \leq \epsilon\} < \infty.
%%\end{equation}
%%\item $G$ is \emph{weakly discrete (WD)} if for every $x\in X$, there exists $\epsilon > 0$ such that
%%\begin{equation}
%%\label{WD2}
%%G(x)\cap B(x,\epsilon) = \{x\}.
%%\end{equation}
%%\end{itemize}
%%\end{remark}
%%
%%As our naming suggests, the condition of strong discreteness is stronger than the condition of moderate discreteness, which is in turn stronger than the condition of weak discreteness.
%%\begin{proposition}
%%\label{propositionSDMDWD}
%%Any strongly discrete group is moderately discrete, and any moderately discrete group is weakly discrete.
%%\end{proposition}
%%\begin{proof}
%%It is clear from the second formulation that strongly discrete groups are moderately discrete. Let $G\leq\Isom(X)$ be a moderately discrete group. Fix $x\in X$, and let $\epsilon > 0$ be such that \eqref{MD2} holds. Letting $\epsilon' = \epsilon \wedge \min\{\dist(x,g(x)) : g(x)\neq x, g(x)\in B(x,\epsilon)\}$, we see that \eqref{WD2} holds.
%%\end{proof}
%%
%%The reverse directions, WD \implies MD and MD \implies SD, both fail in infinite dimensions. Examples \ref{examplevalette} and \ref{exampletranslations}-\ref{examplenotstronglydiscreteBIM} are moderately discrete groups which are not strongly discrete, and Examples \ref{exampleMDWD} and \ref{exampleAutT} are weakly discrete groups which are not moderately discrete.
%%
%%If $X$ is a proper metric space, then the classes MD and SD coincide, but are still distinct from WD. Example \ref{exampleAutT} is a weakly discrete group acting on a proper metric space which is not moderately discrete. We show now that MD $\Leftrightarrow$ SD when $X$ is proper:
%%
%%\begin{proposition}
%%\label{propositionproperMDSD}
%%Suppose that $X$ is proper. Then a subgroup of $\Isom(X)$ is moderately discrete if and only if it is strongly discrete.
%%\end{proposition}
%%\begin{proof}
%%Let $G\leq\Isom(X)$ be a moderately discrete subgroup. Fix $x\in X$, and let $\epsilon > 0$ satisfy \eqref{MD2}. Fix $R > 0$ and let $K = \cl{G(\zero)}\cap B(x,R)$; $K$ is compact since $X$ is proper. The collection $\{B(g(x),\epsilon): g\in G\}$ covers $K$, so there is a finite subcover $\{B(g_i(x),\epsilon): i = 1,\ldots,n\}$. Now
%%\[
%%\#\{g\in G: \dist(x,g(x)\leq R)\} \leq \sum_{i = 1}^n \#\{g\in G: g(x)\in B(g_i(x),\epsilon)\} < \infty,
%%\]
%%i.e. \eqref{SD2} holds.
%%\end{proof}
%%
%%\subsubsection{Parametric discreteness}
%%\begin{definition}
%%\label{definitionparametricdiscreteness}
%%Let $\scrT$ be a topology on $\Isom(X)$. A group $G\leq\Isom(X)$ is \emph{$\scrT$-parametrically discrete ($\scrT$-PD)} if it is discrete as a subspace of $\Isom(X)$ in the topology $\scrT$.
%%\end{definition}
%%
%%Most of the time, we will let $\scrT$ be the compact-open topology (COT). The relation between COT-parametric discreteness and our previous notions of discreteness is as follows:
%%\begin{proposition}
%%\label{propositionparametricdiscreteness}~
%%\begin{itemize}
%%\item[(i)] Any moderately discrete group is COT-parametrically discrete.
%%\item[(ii)] Any weakly discrete group acting on a \CROSS\ is COT-parametrically discrete.%Suppose that $\Fix(g)$ has empty interior for all $g\in\Isom(X)\butnot\{\id\}$. Then any weakly discrete group acting on $X$ is COT-parametrically discrete.
%%\item[(iii)] Any COT-parametrically discrete group acting on a proper metric space is strongly discrete.
%%\end{itemize}
%%\end{proposition}
%%\begin{proof}~
%%\begin{itemize}
%%\item[(i)]
%%Let $G\leq\Isom(X)$ be moderately discrete, and let $\epsilon > 0$ satisfy \eqref{MD2}. Then the set $\UU := \GG(\{\zero\},B(\zero,\epsilon))\subset\Isom(X)$ satisfies $\#(\UU\cap G) < \infty$. But $\UU$ is a neighborhood of $\id$ in the compact-open topology. It follows that $G$ is COT-parametrically discrete.
%%
%%\item[(ii)]
%%Suppose that $X = \H = \H_\F^\alpha$. Let $G\leq\Isom(\H)$ be weakly discrete, and by contradiction suppose it is not COT-parametrically discrete. For any finite set $F\subset\H$, let $\epsilon > 0$ be small enough so that \eqref{WD2} holds for all $x\in F$; since $G$ is not COT-PD, there exists $g = g_F\in G\butnot\{\id\}$ such that $\dist(x,g(x))\leq \epsilon$ for all $x\in F$, and it follows that $g(x) = x$ for all $x\in F$. Now suppose that $J$ is a finite set of indices, and let $F = \{[\ee_0]\} \cup \{[\ee_0 \pm (1/2)\ee_i] \ell : i\in J, \ell\in I_\F\}$, where $I_\F$ is as in \eqref{IFdef}. Then if $T_I$ is a representative of $g_F$ satisfying $T_J\ee_0 = \ee_0$, an argument similar to the proof of Proposition \ref{propositionCOTSOT}(ii) shows that $\sigma_{T_J} = I$ and $T_J \ee_i = \ee_i$ for all $i\in J$.
%%
%%Now define an infinite sequence of indices $(i_n)_1^\infty$ as follows: If $i_1,\ldots,i_{n - 1}$ have been defined, let $T_n = T_{\{i_1,\ldots,i_{n - 1}\}}$, and let $i_n$ be such that $\ee_{i_n}\notin \Fix(T_n)$.
%%
%%Choose a nonnegative summable sequence $(t_n)_1^\infty$, and let $\xx = \ee_0 + \sum_{n = 1}^\infty t_n \ee_{i_n}$. Then $T_n\xx\to \xx$; since $G$ is weakly discrete, it follows that $T_n\xx = \xx$ for all $n$ sufficiently large. Fix such an $n$, and observe that
%%\[
%%0 = T_n\xx - \xx = t_n (T_n(\ee_n) - \ee_n) + \sum_{m > n} t_m (T_n(\ee_m) - \ee_m);
%%\]
%%the triangle inequality gives
%%\[
%%t_n \leq \frac{\sum_{m > n}2t_m}{\|T_n\ee_n - \ee_n\|}\cdot
%%\]
%%By choosing the sequence $(t_n)_1^\infty$ to satisfy
%%\[
%%t_{n + 1} < \frac14 \|T_n\ee_n - \ee_n\| t_n \leq \frac12 t_n ,
%%\]
%%we arrive at a contradiction.
%%
%%%\ignore{
%%%
%%%[Separable case:]
%%%
%%%By contradiction, suppose that we have a sequence $(g_n)_1^\infty$ in $G$ such that $g_n\to\id$ in the compact-open topology.
%%%
%%%Let us view $G$ as a subgroup of $\O^*(\LL;\QQ)$. For each $n$, we have $g_n\neq \id$, so $\Fix_\LL(g_n)$ is a proper subspace of $\LL$. Now the Baire category theorem implies that
%%%\[
%%%\bigcup_{n\in\N}\Fix_\LL(g_n) \neq \LL
%%%\]
%%%so let us choose $\xx\in\LL\butnot \bigcup_n\Fix_\LL(g_n)$. Let $U$ be a neighborhood of $\xx$ such that 
%%%\[
%%%g(\xx)\in U \Rightarrow g(\xx) = \xx.
%%%\]
%%%Since $\xx\notin\Fix_\LL(g_n)$ for all $n$, we have $g_n(\xx)\notin U$ for all $n$. But on the other hand we must have $g_n(\xx)\to\xx$ since $g_n\to\id$ in the compact-open topology. This is a contradiction.
%%%}% end ignore
%%
%%\item[(iii)]
%%Let $G$ be a COT-parametrically discrete group acting by isometries on a proper metric space $X$. By contradiction, suppose that $G$ is not strongly discrete. Then there exists an infinite set $A\subset G$ such that the set $A(\zero)$ is bounded. Without loss of generality we may suppose that $A^{-1} = A$. Note that for each $x\in X$, the set $A(x)$ is bounded and therefore precompact. Now since $X$ is a proper metric space, it is $\sigma$-compact and therefore separable. Let $\SS$ be a countable dense subset of $X$. Then
%%\[
%%\KK := \left(\prod_{q\in\SS}\cl{A(q)}\right)^2
%%\]
%%is a compact metrizable space. For each $g\in A$ let
%%\[
%%\phi_g := \left((g(q))_{q\in\SS},(g^{-1}(q))_{q\in\SS}\right)\in \KK.
%%\]
%%Since $A$ is infinite, there exists an infinite sequence $(g_n)_1^\infty$ in $A$ such that
%%\[
%%\phi_{g_n}\to \left((y_q^{(+)})_{q\in\SS},(y_q^{(-)})_{q\in\SS}\right)\in\KK.
%%\]
%%Thus
%%\[
%%g_n^\pm(q) \tendsto n y_q^{(\pm)} \all q\in\SS.
%%\]
%%The density of $\SS$ and the equicontinuity of the sequences $(g_n)_1^\infty$ and $(g_n^{-1})_1^\infty$ imply that for all $x\in X$, there exist $y_x^{(\pm)}$ such that $g_n^\pm(y)\to y_x^{(\pm)}$. Thus, the sequence $(g_n)_1^\infty$ converges in the Tychonoff topology to some $g^{(+)}\in X^X$. Similarly, the sequence $(g_n^{-1})_1^\infty$ converges to some $g^{(-)}\in X^X$. Again, the equicontinuity of the sequences $(g_n)_1^\infty$ and $(g_n^{-1})_1^\infty$ allows us to take limits and deduce that
%%\[
%%g^{(+)} g^{(-)} = \lim_{n\to\infty} g_n g_n^{-1} = \id.
%%\]
%%Similarly, $g^{(-)}g^{(+)} = \id$. Thus $g^{(+)}$ and $g^{(-)}$ are inverses, and in particular $g^{(+)}\in\Isom(X)$. Since $g_n\to g^{(+)}$ in the compact-open topology, the proof is completed by the following lemma from topological group theory:
%%\begin{lemma}
%%\label{lemmadiscretesubgroup}
%%Let $H$ be a topological group, and let $G$ be a subgroup of $H$. Suppose there is a sequence $(g_n)_1^\infty$ of distinct elements in $G$ which converges to an element of $H$. Then $G$ is not discrete in the topology inherited from $H$.
%%\end{lemma}
%%\begin{subproof}
%%Suppose $g_n\to h\in H$. Then
%%\[
%%g_n g_{n + 1}^{-1} \to h h^{-1} = \id,
%%\]
%%while on the other hand $g_n g_{n + 1}^{-1}\neq \id$ (since the sequence $(g_n)_1^\infty$ consists of distinct elements). This demonstrates that $G$ is not discrete in the inherited topology.
%%\end{subproof}
%%\end{itemize}
%%\end{proof}
%%If $X$ is not a \CROSS, then it is possible for a weakly discrete group to not be COT-parametrically discrete; see Example \ref{exampleAutT}. Conversely, it is possible for a COT-parametrically discrete group to not be weakly discrete; see Examples \ref{examplepoincareextension} amd \ref{examplepoincareextensiontwo}.
%%
%%Now suppose that $X$ is a \CROSS. Recall that UOT denotes the uniform operator topology.
%%\begin{observation}
%%\label{observationCOTUOT}
%%If a subgroup $G\leq\Isom(X)$ is COT-parametrically discrete, then it is also UOT-parametrically discrete.
%%\end{observation}
%%This is because the uniform operator topology is finer, i.e. it has more open sets, and so it is easier for every subset of $G$ to be relatively open in that topology, which is exactly what it means to be discrete.
%%
%%Note that there is an ``order switch'' here; the UOT is finer than the COT, but the condition of being COT-parametrically discrete is stronger than the condition of being UOT-parametrically discrete.
%%
%%A significant example of a group which is UOT-parametrically discrete but not COT-parametrically discrete is given in Example \ref{exampleAutTBIM}.
%%
%%We summarize the relations between our different notions of discreteness in Table \ref{figurediscreteness} below.
%%
%%\subsubsection{Equivalence in finite dimensions}
%%
%%\begin{proposition}\label{propositionfinitedimequivalent}
%%Suppose that $X$ is a finite-dimensional Riemannian manifold. Then the notions of strong discreteness, moderate discreteness, weak discreteness, and COT-parametric discreteness agree. If $X$ is a \CROSS, these notions also agree with the notion of UOT-parametric discreteness.
%%\end{proposition}
%%\begin{proof}
%%By Propositions \ref{propositionSDMDWD} and \ref{propositionparametricdiscreteness}, the conditions of strong discreteness, moderate discreteness, and COT-parametric discreteness agree and imply weak discreteness. Conversely, suppose that $G\leq\Isom(X)$ is weakly discrete, and by contradiction suppose that $G$ is not COT-parametrically discrete. Since $X$ is separable, so is $\Isom(X)$, and thus there exists a sequence $\Isom(X)\butnot\{\id\}\ni g_n\to\id$ in the compact-open topology. For each $n$ let $F_n = \{x\in X : g_n(x) = x\}$. Since $G$ is weakly discrete, $X = \bigcup_1^\infty F_n$, so by the Baire category theorem, $F_n$ has nonempty interior for some $n$. But then $g_n = \id$ on an open set; in particular there exists a point $x_0\in X$ such that $g_n(x_0) = x_0$ and $g_n'(x_0)$ is the identity map on the tangent space of $x_0$. By the naturality of the exponential map, this implies that $g_n$ is the identity map, a contradiction.
%%
%%Finally, suppose $X = \H = \H_\F^\alpha$ is a \CROSS, and let $\LL = \LL_\F^{\alpha + 1}$. Since $\LL$ is finite-dimensional, the SOT and UOT topologies on $L(\LL)$ are equivalent; this demonstrates that the notions of COT-parametric discreteness and UOT-parametric discreteness agree.
%%\end{proof}
%%We shall call a group satisfying any of these equivalent definitions simply \emph{discrete}.
%%
%%\subsubsection{Proper discontinuity}
%%\begin{definition}
%%\label{definitionproperdiscontinuity}
%%A group $G\leq\Isom(X)$ \emph{acts properly discontinuously (PrD)} on $X$ if for every $x\in X$, there exists an open set $U \ni x$ with 
%%\[
%%g(U) \cap U \neq \emptyset \Rightarrow g = \id,
%%\]
%%or equivalently, if
%%\[
%%\dist(x,\{g(x):g\neq\id\}) > 0.
%%\]
%%\end{definition}
%%
%%We note that even in finite dimensions, the notion of proper discontinuity is not the same as the notion of discreteness; it is slightly stronger. We also remark that in finite dimensions Selberg's lemma can be used to pass from a discrete group to a finite-index subgroup that acts properly discontinuously; this is not the case in infinite dimensions.
%%
%%Since $\#\{\id\} = 1 < \infty$, we have the following:
%%\begin{observation}
%%\label{observationproperlydiscontinuous}
%%Any group which acts properly discontinuously is moderately discrete.
%%\end{observation}
%%
%%In particular, if $X$ is proper then any group which acts properly discontinuously is strongly discrete. This provides a connection between our results, in which strong discreteness is often a hypothesis, and many results from the literature in which proper discontinuity and properness are both hypotheses.
%%
%%We summarize the relations between our various notions of discreteness, together with proper discontinuity, in the following table:
%%
%%\begin{table}[h!]
%%\begin{tabular}{ | c | c c c c c c c | }
%% \hline
%% Finite				& SD & $\leftrightarrow$ & MD & $\leftrightarrow$ & WD && \\
%% dimensional	& $\uparrow$ & & & & $\updownarrow$ && \\
%% manifold	 			& PrD & & & & COT-PD & $\leftrightarrow$ & UOT-PD\\
%% \hline
%% General		 		& SD & $\rightarrow$ & MD & $\rightarrow$ & WD && \\
%% metric 		& & $\nearrow$ & & $\searrow$ & && \\
%% space				& PrD & & & & COT-PD && \\
%% \hline
%% Infinite				& SD & $\rightarrow$ & MD & $\rightarrow$ & WD && \\
%% dimensional	& & $\nearrow$ & & & $\downarrow$ && \\
%% \ROSS			& PrD & & & & COT-PD & $\rightarrow$ & UOT-PD \\
%% \hline
%% Proper		 		& SD & $\leftrightarrow$ & MD & $\leftrightarrow$ & COT-PD && \\
%% metric		& $\uparrow$ & & & & $\downarrow$ &&\\
%% space			& PrD & & & & WD && \\
%% \hline
%%\end{tabular}
%%\vskip0.12cm
%%\caption{The relations between different notions of discreteness. All implications not listed have counterexamples; see Section \ref{sectionexamples}.}
%%\label{figurediscreteness}
%%\end{table}
%%
%%%\begin{figure}
%%%\centerline{\mbox{\includegraphics[scale = 0.7]{discreteness.png}}}
%%%\caption{The relations between different notions of discreteness. All implications not listed have counterexamples.}
%%%\label{figurediscreteness}
%%%\end{figure}
%%
%%Observation \ref{observationproperlydiscontinuous} has the following partial converse:
%%\begin{remark}
%%If $X$ is a proper CAT(0) space, then a group acts properly discontinously if and only if it is moderately discrete and torsion free.
%%\end{remark}
%%\begin{proof}
%%Suppose that $G\leq\Isom(X)$ acts properly discontinuously. If $g\in G\butnot\{\id\}$ is a torsion element, then by Cartan's lemma \cite[II.2.8(1)]{BridsonHaefliger}, $g$ has a fixed point, which contradicts that $G$ acts properly discontinuously. Thus $G$ is torsion-free.
%%
%%Conversely, suppose that $G\leq\Isom(X)$ is moderately discrete and torsion-free. Given $x\in X$, let $\epsilon > 0$ be as in \eqref{WD2}, and by contradiction suppose that there exists $g\neq\id$ such that $\dist(x,g(x)) < \epsilon$. By \eqref{WD2}, $g(x) = x$. But then by \eqref{MD2}, the set $\{g^n : n\in\Z\}$ is finite, i.e. $g$ is a torsion element. This is a contradiction, so $G$ acts properly discontinuously.
%%\end{proof}
%%
%%%%%%%%%%%%%%%%%%%%%%%%%%%%%%%%%%
%%%Covering space interpretation
%%%\begin{remark}
%%%Properly discontinuous actions are also known as \emph{covering space actions} and in this language we could also have said a group $G$ acts properly discontinuously on a topological space $X$ if and only if the natural projection
%%%\[
%%%\pi_G : X \to X/G
%%%\]
%%%is a covering map.
%%%\end{remark}
%%
%%
%%\subsubsection{Behavior with respect to restrictions}
%%Fix $G\leq\Isom(X)$, and suppose $Y\subset X$ is a subspace of $X$ preserved by $G$, i.e. $g(Y) = Y$ for all $g\in G$. Then $G$ can be viewed as a group acting on the metric space $(Y,\dist\given_Y)$.
%%
%%\begin{observation}~
%%\label{observationrestrictions}
%%\begin{itemize}
%%\item[(i)] $G$ is strongly discrete $\Leftrightarrow$ $G\given Y$ is strongly discrete
%%\item[(ii)] $G$ is moderately discrete \implies $G\given Y$ is moderately discrete
%%\item[(iii)] $G$ is weakly discrete \implies $G\given Y$ is weakly discrete
%%\item[(iv)] $G$ is $\scrT$-parametrically discrete $\Leftarrow$ $G\given Y$ is $\scrT\given Y$-parametrically discrete
%%\item[(v)] $G$ acts properly discontinuously on $X$ \implies $G$ acts properly discontinuously on $Y$.
%%\end{itemize}
%%\end{observation}
%%In particular, strong discreteness is the only concept which is ``independent of the space being acted on''. It is thus the most robust of all our definitions.
%%
%%Note that for parametric discreteness, the order of implication reverses; restricting to a subspace may cause a group to no longer be discrete. Example \ref{examplepoincareextension} is an example of this phenomenon.
%%
%%\subsubsection{Countability of discrete groups}
%%In finite dimensions, all discrete groups are countable. In general, it depends on what type of discreteness you are considering.
%%
%%\begin{proposition}
%%Fix $G\leq\Isom(X)$, and suppose that either
%%\begin{itemize}
%%\item[(1)] $G$ is strongly discrete, or
%%\item[(2)] $X$ is separable and $G$ is COT-parametrically discrete.
%%\end{itemize}
%%Then $G$ is countable.
%%\end{proposition}
%%\begin{proof}
%%If $G$ is strongly discrete, then
%%\[
%%\#(G) \leq \sum_{n\in\N}\#\{g\in G:\dogo g\leq n\} \leq \sum_{n\in\N}\#(\N) = \#(\N).
%%\]
%%On the other hand, if $X$ is a separable metric space, then by Remark \ref{remarkpolish} $\Isom(X)$ is separable metrizable, so it contains no uncountable discrete subspaces.
%%\end{proof}
%%\begin{remark}
%%An example of an uncountable UOT-parametrically discrete subgroup of $\Isom(\H^\infty)$ is given in Example \ref{exampleAutTBIM}, and an example of an uncountable weakly discrete group acting on a separable $\R$-tree is given in Example \ref{exampleAutT}. An example of an uncountable moderately discrete group acting on a (non-separable) $\R$-tree is given in Remark \ref{remarkMDuncountable}.
%%\end{remark}
%%
%%%\ignore{
%%%% A few definitions from wikipedia!
%%%%
%%%%Faithful (or effective) if for any two distinct g, h in G there exists an x in X such that $g.x = h.x$; or equivalently, if for any $g\neq e$ in $G$ there exists an $x\in X$ such that $g.x = x$. Intuitively, different elements of G induce different permutations of X.
%%%%
%%%%Free (or semiregular) if for any x in X, $g.x = h.x$ implies $g = h$. Equivalently: if there exists an x in X such that g.x = x (that is, if g has at least one fixed point), then g is the identity.
%%%%
%%%%Regular (or simply transitive) if it is both transitive and free; this is equivalent to saying that for any two x, y in X there exists precisely one g in G such that $g.x = y$. In this case, X is known as a principal homogeneous space for G or as a G-torsor.
%%%%
%%%Recall that a group is \emph{torsion-free} if every element of finite order is the identity; it \emph{acts freely} if there are no nontrivial fixed points i.e. $\Fix(g) \neq \emptyset \Rightarrow g = \id$. We summarize the connections between our various notions in the following
%%%
%%%\begin{proposition}\label{propositiondiscrete} 
%%%Let $G$ be a group acting by isometries on a metric space $X$.
%%%\begin{itemize}
%%%\item[(i)] If $G$ is strongly discrete and torsion-free, then $G$ acts freely.
%%%\item[(ii)] $G$ acts properly discontinuously if and only if it is weakly discrete and acts freely.
%%%\item[(iii)] If $G$ is strongly discrete and torsion-free, then $G$ acts properly discontinuously.
%%%\item[(iv)] Suppose that $X$ is a CAT(0) space. Then if $G$ acts freely, then it is torsion-free.
%%%\end{itemize}
%%%\end{proposition}
%%%
%%%\begin{remark}
%%%In finite dimensions, a group acts properly discontinously if and only if it is discrete and torsion free. Indeed, this can be seen from 2, 3, and 4 of the above proposition together with Remark \ref{remarkfinitedimequivalent}.
%%%\end{remark}
%%%
%%%\begin{proof}[Proof of Proposition \ref{propositiondiscrete}]~
%%%\begin{itemize}
%%%\item[(i)] [SD + TF \implies F]
%%%By contradiction suppose that $G$ is not free; then $g(x) = x$ for some $g \neq \id$. Then for every $n$ we have $\dist(x,g^n(x)) = 0$. Now strong discreteness implies that $\#\{g^n : n\in\Z\} < \infty$ which in turn produces an $N\neq 0$ for which $g^N = \id$. This contradicts the assumption that $G$ was torsion-free.
%%%\item[(ii)] [PD $\Leftrightarrow$ WD + F]
%%%This follows straight from the definitions.
%%%\item[(iii)] [SD + TF \implies PD]
%%%This follows from parts 1 and 2, together with Proposition \ref{propositionSDWD}.
%%%\item[(iv)] [If $X$ is a CAT(0) space, then F \implies TF]
%%%By contradiction, suppose that $G$ is not torsion free; then $g^n = \id$ for some $g \neq \id$. Let $H = \lb g \rb$; we have $\#(H) < \infty$. Therefore the orbit $H(\zero)$ is bounded and by \cite[Corollary 2.8]{BridsonHaefliger} we have $\Fix(H) = \Fix(g) \neq\emptyset$. This contradicts that $G$ acts freely.
%%%\item[(v)] [PD \implies COT-PD]
%%%By contradiction, suppose that we have a sequence $(g_n)_1^\infty$ in $G$ such that $g_n\to\id$ in the compact-open topology. Then $g_n(\zero)\to\zero$, which contradicts that $G$ is properly discontinuous.
%%%\end{itemize}
%%%\end{proof}
%%%}% end ignore
%%
%%
%%
%%%%%%%%%%%%%%%%%%%%%%%%%%%%%%%%%%
%%\draftnewpage
%%\section{Classification of isometries and semigroups} \label{sectionclassification}
%%%%%%%%%%%%%%%%%%%%%%%%%%%%%%%%%%
%%
%%In this section we classify subsemigroups $G\prec\Isom(X)$ into six categories, depending on the behavior of the orbit of the basepoint $\zero\in X$. We start by classifying individual isometries, although it will turn out that the category into which an isometry is classified is the same as the category of the cyclic group that it generates.
%%
%%We remark that if $X$ is geodesic and $G\leq\Isom(X)$ is a group, then the main results of this section were proven in \cite{Hamann}. Moreover, our terminology is based on \cite[\63.A]{CCMT}, where a similar classification was given based on \cite[\6 3.1]{Gromov3}.% In fact, the hypothesis that $X$ is geodesic can be removed by embedding $X$ into a geodesic hyperbolic metric space; see \cite[Theorem 4.1]{BonkSchramm} and \cite[Corollary A.10]{BHM}.
%%
%%\bigskip
%%\subsection{Classification of isometries}
%%Fix $g\in\Isom(X)$, and let
%%\[
%%\Fix(g) := \{x\in\bord X:g(x) = x\}.
%%\]
%%Consider $\xi\in\Fix(g)\cap\del X$. Recall that $g'(\xi)$ denotes the dynamical derivative of $g$ at $\xi$ (see \6\ref{subsubsectiondynamicalderivative}). 
%%\begin{definition}
%%\label{definitionderivativefixedpoint}
%%$\xi$ is said to be
%%\begin{itemize}
%%\item a \emph{neutral} or \emph{indifferent} fixed point if $g'(\xi) = 1$,
%%\item an \emph{attracting} fixed point if $g'(\xi) < 1$, and
%%\item a \emph{repelling} fixed point if $g'(\xi) > 1$.
%%\end{itemize}
%%\end{definition}
%%
%%\begin{definition}
%%\label{definitionclassification1}
%%An isometry $g\in\Isom(X)$ is called
%%\begin{itemize}
%%\item \emph{elliptic} if the orbit $\{g^n(\zero):n\in\N\}$ is bounded,
%%\item \emph{parabolic} if it is not elliptic and has a unique fixed point in $\del X$, which is neutral, and
%%\item \emph{loxodromic} if it has exactly two fixed points in $\del X$, one of which is attracting and the other of which is repelling.
%%\end{itemize}
%%\end{definition}
%%\begin{remark}
%%We use the terminology ``loxodromic'' rather than the more common ``hyperbolic'' to avoid confusion with the many other meanings of the word ``hyperbolic''. In particular, when we get to classification of groups it would be a bad idea to call any group ``hyperbolic'' if it is not hyperbolic in the sense of Gromov.
%%\end{remark}
%%
%%The categories of elliptic, parabolic, and loxodromic are clearly mutually exclusive.\Footnote{Proposition \ref{propositionbusemannpreserved} can be used to show that loxodromic isometries are not elliptic.} In the converse direction we have the following:
%%\begin{theorem}
%%\label{classificationofisometries}
%%Any isometry is either elliptic, parabolic, or loxodromic.
%%\end{theorem}
%%
%%The proof of Theorem \ref{classificationofisometries} will proceed through several lemmas.
%%
%%\begin{lemma}[A corollary of {\cite[Proposition 5.1]{Karlsson}}]
%%\label{lemmakarlsson}
%%If $g\in\Isom(X)$ is not elliptic, then $\Fix(g)\cap\del X\neq\emptyset$.
%%\end{lemma}
%%We include the proof for completeness.
%%
%%\begin{proof}
%%For each $t\in\N$, let $n_t$ be the smallest integer such that
%%\[
%%\dogo{g^{n_t}} \geq n_t.
%%\]
%%The sequence $(n_t)_1^\infty$ is nondecreasing. Given $s,t\in\N$ with $s < t$, we have
%%\[
%%\dist(g^{n_s}(\zero) , g^{n_t}(\zero)) = \dogo{g^{n_t - n_s}} < n_t,
%%\]
%%and thus
%%\[
%%\lb g^{n_s}(\zero) | g^{n_t}(\zero)\rb_\zero > \frac12 [ n_s + n_t - n_t] = \frac12 n_s \tendsto{s,t} \infty,
%%\]
%%i.e. $(g^{n_t}(\zero))_t$ is a Gromov sequence. Let $\xi = [(g^{n_t}(\zero))_t]$, and note that
%%\[
%%\lb \xi|g(\xi)\rb_\zero = \lim_{t\to\infty} \lb g^{n_t}(\zero) | g^{n_t + 1}(\zero) \rb \geq \lim_{t\to\infty} \left[\dogo{g^{n_t}} - \dist(g^{n_t}(\zero), g^{n_t + 1}(\zero))\right] = \infty.
%%\]
%%Thus $g(\xi) = \xi$, i.e. $\xi\in\Fix(g)\cap\del X$.
%%\end{proof}
%%
%%\begin{remark}[{\cite[Proposition 5.2]{Karlsson}}]
%%If $g\in\Isom(X)$ is elliptic and if $X$ is CAT(0), then $\Fix(g)\cap X\neq\emptyset$ due to Cartan's lemma (Theorem \ref{Cartan's lemma} below). Thus if $X$ is a CAT(0) space, then any isometry of $X$ has a fixed point in $\bord X$.
%%\end{remark}
%%
%%\begin{lemma}
%%\label{lemmarepellingloxodromic}
%%If $g\in\Isom(X)$ has an attracting or repelling periodic point, then $g$ is loxodromic.
%%\end{lemma}
%%\begin{proof}
%%Suppose that $\xi\in\del X$ is a repelling fixed point for $g\in\Isom(X)$, i.e. $g'(\xi) > 1$. Recall from Proposition \ref{propositioneuclideansimilarity} that
%%\[
%%\Dist_\xi(g^n(y_1), g^n(y_2)) \leq C g'(\xi)^{-n} \Dist_\xi(y_1, y_2) \all y_1,y_2\in\EE_\xi\all n\in\Z
%%\]
%%for some constant $C > 0$. Now let $n$ be large enough so that $g'(\xi)^n > C$; then the above inequality shows that the map $g^n$ is a strict contraction of the complete metametric space $(\EE_\xi,\Dist_\xi)$ (cf. Proposition \ref{propositioneuclideanmetametric}). Then by Theorem \ref{banachcontractionprinciple}, $g$ has a unique fixed point $\eta\in(\EE_\xi)_\refl = \del X\butnot\{\xi\}$. By Corollary \ref{corollaryproductone}, $\eta$ is an attracting fixed point. Corollary \ref{corollaryproductone} also implies that $g$ cannot have a third fixed point. Thus $g$ is loxodromic.
%%
%%On the other hand, if $g$ has an attracting fixed point, then by Proposition \ref{propositionderivativeiteration}, $g^{-1}$ has a repelling fixed point. Thus $g^{-1}$ is loxodromic, so applying Proposition \ref{propositionderivativeiteration} again, we see that $g$ is loxodromic.
%%\end{proof}
%%
%%\begin{proof}[Proof of Theorem \ref{classificationofisometries}]
%%By contradiction suppose that $g$ is not elliptic or loxodromic, and we will show that it is parabolic. By Lemma \ref{lemmakarlsson}, we have $\Fix(g)\cap\del X\neq\emptyset$; on the other hand, by Lemma \ref{lemmarepellingloxodromic}, every fixed point of $g$ in $\del X$ is neutral. It remains to show that $\#(\Fix(g)) = 1$. By contradiction, suppose otherwise. Since $g$ is not elliptic, we clearly have $\Fix(g)\cap X = \emptyset$. Thus we may suppose that there are two distinct neutral fixed points $\xi_1,\xi_2\in\del X$.
%%
%%Now for each $n\in\N$, we have
%%\[
%%\busemann_{\xi_i}(\zero,g^n(\zero)) \asymp_\plus n\log_b(g'(\xi_i)) = 0, \hspace{.5 in} i = 1,2
%%\]
%%by Proposition \ref{propositionbusemannpreserved}. Let $r = r_{\xi_1,\xi_2,\zero}$ and $\theta = \theta_{\xi_1,\xi_2,\zero}$ be as in Subsection \ref{subsectionpolar}. Then by Lemma \ref{lemmatranslationdistance} we have $r(g^n(\zero))\asymp_\plus \theta(g^n(\zero)) \asymp_\plus 0$. Thus by Lemma \ref{lemmafabs} we have
%%\[
%%\dogo{g^n} \asymp_\plus |r(g^n(\zero))| + \theta(g^n(\zero)) \asymp_\plus 0,
%%\]
%%i.e. the sequence $\{g^n(\zero):n\in\N\}$ is bounded. Thus $g$ is elliptic, contradicting our hypothesis.
%%\end{proof}
%%
%%
%%\begin{remark}[Cf. {\cite[Chapter 3, Theorem 1.4]{Chiswell}}]
%%\label{remarkparabolicRtree}
%%For $\R$-trees, parabolic isometries are impossible, so Theorem \ref{classificationofisometries} shows that every isometry is elliptic or loxodromic.
%%\end{remark}
%%\begin{proof}
%%By contradiction suppose that $X$ is an $\R$-tree and that $g\in\Isom(X)$ is a parabolic isometry with fixed point $\xi\in\del X$. Let $x = C(\zero,g(\zero),\xi)\in X$; then $x = \geo\zero\xi_t$ for some $t\geq 0$. Now,
%%\[
%%\dist(g(\zero),x) = \dox x + \busemann_\xi(g(\zero),\zero) = t + 0 = t.
%%\]
%%It follows that $g(x) = \geo{g(\zero)}{\xi}_t = x$. Thus $g$ is elliptic, a contradiction.
%%\end{proof}
%%
%%
%%
%%\subsubsection{More on loxodromic isometries}
%%
%%\begin{notation}
%%Suppose $g\in\Isom(X)$ is loxodromic. Then $g_+$ and $g_-$ denote the attracting and repelling fixed points of $g$, respectively.
%%\end{notation}
%%
%%\begin{theorem}
%%\label{theoremloxodromicextra}
%%Let $g\in\Isom(X)$ be loxodromic. Then
%%\begin{equation}
%%\label{loxodromicinversederivative}
%%g'(g_+) = \frac{1}{g'(g_-)}\cdot
%%\end{equation}
%%Furthermore, for every $x\in \bord X\butnot\{g_-\}$ and for every $n\in\N$ we have
%%\begin{equation}
%%\label{loxodromiccontracting}
%%\Dist(g^n(x), g_+) \lesssim_\times \frac{[g'(g_+)]^n}{\Dist(g_-,g_+)\Dist(x, g_-)},
%%\end{equation}
%%with $\leq$ if $X$ is strongly hyperbolic. In particular
%%\[
%%x\neq g_- \;\;\Rightarrow\;\; g^n(x)\tendsto n g_+,
%%\]
%%and the convergence is uniform over any set whose closure does not contain $g_-$. Finally,
%%\begin{equation}
%%\label{dognoloxodromic}
%%\dogo{g^n} \asymp_\plus |n| \log_b g'(g_-) = |n| \log_b \frac{1}{g'(g_+)}\cdot
%%\end{equation}
%%\end{theorem}
%%\begin{proof}
%%\eqref{loxodromicinversederivative} follows directly from Corollary \ref{corollaryproductone}.
%%
%%To demonstrate \eqref{loxodromiccontracting}, note that
%%\begin{align*}
%%&\lb x|g_-\rb_\zero + \lb g^n(x)|g_+\rb_\zero\\
%%&\gtrsim_\plus \busemann_{g_-}(\zero,x) + \busemann_{g_+}(\zero,g^n(x)) \by{(j) of Proposition \ref{propositionbasicidentities}}\\
%%&\asymp_\plus \busemann_{g_-}(\zero,x) + \busemann_{g_+}(\zero,x) - n \log_b g'(g_+) \by{Proposition \ref{propositionbusemannpreserved}}\\
%%&\asymp_\plus \lb g_-|g_+\rb_x - \lb g_-|g_+\rb_\zero - n \log_b g'(g_+) \by{(g) of Proposition \ref{propositionbasicidentities}}\\
%%&\geq_\pt -\lb g_-|g_+\rb_\zero - n\log_b g'(g_+).
%%\end{align*}
%%Exponentiating and rearranging yields \eqref{loxodromiccontracting}.
%%
%%Finally, \eqref{dognoloxodromic} follows directly from Lemmas \ref{lemmafabs} and \ref{lemmatranslationdistance}.
%%\end{proof}
%%
%%%\comdavid{I commented out the subsection called ``Minimal displacement'' - it seemed like there was no point since we were just repeating the finite-dimensional theory, and the concepts do not seem to be useful for anything.}
%%%\subsection{Minimal displacement}
%%%\comtushar{what should we do with this remark?}
%%%
%%%\begin{remark}
%%%The classification of isometries into these three classes may also be given the following elegant geometric interpretation in terms of \emph{minimal displacement}, see \cite{BridsonHaefliger}. The following definitions can be made %for a metric space $X$ and a mapping $f : X \to X$. The \emph{minimal displacement} of $f$ is defined as
%%%\[
%%%\ell(f) := \inf_{x\in X} \dist(x,f(x)).
%%%\]
%%%The \emph{minimal set} of $f$ is defined as
%%%\[
%%%\text{Min}(f) := \{ x\in X : \dist(x,f(x)) = \ell(f) \} . 
%%%\]
%%%Then given an isometry $g$ of hyperbolic space, exactly one of the following holds:
%%%\begin{itemize}
%%%\item $g$ is parabolic $\Leftrightarrow$ $\text{Min}(g) = \emptyset$.
%%%\item $g$ is elliptic $\Leftrightarrow$ $\text{Min}(g) \neq \emptyset$ and $\ell(f)=0$.
%%%\item $g$ is hyperbolic $\Leftrightarrow$ $\text{Min}(g) \neq \emptyset$ and $\ell(f)>0$.
%%%\end{itemize}
%%%\end{remark}
%%%
%%%[Define asymptotic displacement $\w{\ell}$; they agree in the CAT(-1) case but in general $\w{\ell}$ is more useful]
%%
%%
%%\subsubsection{The story for real ROSSONCTs}
%%If $X$ is a real \ROSS, then we may conjugate each $g\in\Isom(X)$ to a ``normal form'' whose geometrical significance is clearer. The normal form will depend on the classification of $g$ as elliptic, parabolic, or hyperbolic.
%%
%%
%%\begin{proposition}\label{class}
%%Let $X$ be a real \ROSS, and fix $g\in\Isom(X)$.
%%\begin{itemize}
%%\item[(i)] If $g$ is elliptic, then $g$ is conjugate to a map of the form $T \given \B$ for some linear isometry $T\in\O^*(\HH)$.
%%\item[(ii)] If $g$ is parabolic, then $g$ is conjugate to a map of the form $\xx\mapsto \what T \xx + \pp:\E\to\E$, where $T\in\O^*(\BB)$ and $\pp\in\BB$. Here $\BB = \del\E\butnot\{\infty\} = \HH^{\alpha - 1}$.
%%\item[(iii)] If $g$ is hyperbolic, then $g$ is conjugate to a map of the form $\xx\mapsto \lambda \what T \xx: \E \to \E$, where $0< \lambda < 1$ and $T\in\O^*(\BB)$.
%%\end{itemize}
%%\end{proposition}
%%\begin{proof}~
%%\begin{itemize}
%%\item[(i)] If $g$ is elliptic, then by Cartan's lemma (Theorem \ref{Cartan's lemma} below), $g$ has a fixed point $x\in X$. Since $\Isom(X)$ acts transitively on $X$ (Observation \ref{observationtransitivity}), we may conjugate to $\B$ in a way such that $g(\0) = \0$. But then by Proposition \ref{propositionIsomB}, $g$ is of the form (i).
%%\item[(ii)] Let $\xi$ be the neutral fixed point of $g$. Since $\Isom(X)$ acts transitively on $\del X$ (Proposition \ref{propositiontransitivity}), we may conjugate to $\E$ in a way such $g(\infty) = \infty$. Then by Proposition \ref{propositionIsomE} and Example \ref{examplegprimeinfty}, $g$ is of the form (ii).
%%\item[(iii)] Since $\Isom(X)$ acts doubly transitively on $\del X$ (Proposition \ref{propositiontransitivity}), we may conjugate to $\E$ in a way such that $g_+ = \0$ and $g_- = \infty$. Then by Proposition \ref{propositionIsomE} and Example \ref{examplegprimeinfty}, $g$ is of the form (iii). (We have $\pp = \0$ since $\0\in\Fix(g)$.)
%%\end{itemize}
%%\end{proof}
%%
%%
%%%We may summarize our discussion in the following table:\comdavid{Whatever this table is doing, it does not appear to be ``summarizing the discussion'' above. Fix},\comdavid{Depending on the position on the page, the table may not be directly following the colon. We need to be careful.}
%%%{\footnotesize
%%%\begin{table}[h]
%%%\begin{center}
%%%{\renewcommand{\arraystretch}{1.2}
%%%\begin{tabular}{clclclc}
%%%Type & $\#(\Fix(g) \cap \del \H)$ & $\#(\Fix(g) \cap \H)$ & End behavior of $(g^n(x))_{-\infty}^\infty$ \\
%%%\hline
%%%Parabolic &$1$&$0$& Unbounded orbits: $(g^n(x))_{-\infty}^\infty$ either \\ 
%%%&&& converges to or clusters at the fixed point \\
%%%\hline
%%%Hyperbolic &$2$&$0$& Unbounded orbits: \\
%%%&&& $g^n(x) \to g_+$ as $n \to +\infty$ \\
%%%&&& $g^n(x) \to g_-$ as $n \to -\infty$ \\
%%%\hline
%%%&$0$&$1$& \\
%%%Elliptic &$2$&$\#(\R)$& Bounded orbits \\
%%%&$\#(\R)$&$\#(\R)$& \\
%%%\hline
%%%\end{tabular}
%%%}
%%%\end{center}
%%% \caption{Classification of an isometry $g\in\Isom(\H)$.}
%%% \label{tab:classification}
%%%\end{table}
%%%}
%%
%%
%%\begin{remark}
%%If $g\in\Isom(X)$ is elliptic or loxodromic, then the orbit $(g^n(\zero))_1^\infty$ exhibits some ``regularity'' - either it remains bounded forever, or it diverges to the boundary. On the other hand, if $g$ is parabolic then the orbit can oscillate, both accumulating at infinity and returning infinitely often to a bounded region. This is in sharp contrast to finite dimensions, where such behavior is impossible. We discuss such examples in detail in \6\ref{subsubsectionedelstein}.
%%\end{remark}
%%%
%%%\begin{proposition}
%%%\label{propositionparabolicrecurrent}\comdavid{I am writing this subsubsection under the assumption that no one else has explicitly stated Proposition \ref{propositionparabolicrecurrent}. Is this correct?}
%%%There exists a parabolic isometry $g\in\Stab(\Isom(\H);\infty)$ such that the orbit $(g^n(\zero))_1^\infty$ is unbounded but returns infinitely often to a bounded region and in fact accumulates at $\zero$.
%%%\end{proposition}
%%%It is not hard to see that this propostion is equivalent to the following one which was proven by M. Edelstein in 1964:
%%%
%%%\begin{proposition}[{\cite[Theorem 2.1]{Edelstein}}]
%%%\label{propositionedelstein}
%%%There exists a fixed-point-free isometry $g\in\Isom(\ell^2(\N;\C))$ and a sequence $(n_k)_1^\infty$ so that
%%%\[
%%%g^{n_k}(\0) \tendsto k \0.
%%%\]
%%%\end{proposition}
%%%
%%%\begin{proof}[Proof that Proposition \ref{propositionedelstein} implies Proposition \ref{propositionparabolicrecurrent}]\comdavid{Is there any reasonable way to simplify this argument?}
%%%Since $\EE$ and $\ell^2(\N;\C)$ are both separable Hilbert spaces, we may without loss of generality suppose that $g\in\Isom(\EE)$. Let $\what g\in\Stab(\Isom(\H);\infty)$ be the Poincar\'e extension of $g$. Now note that $\what g$ has no fixed points in $\H$; if $\xx$ were such a point, then $(0,x_1,x_2,\ldots)\in\EE$ would be a fixed point of $g$, a contradiction. Thus by Theorem \ref{Cartan's lemma}, $\what g$ is not elliptic. (In particular, by definition the orbit $(g^n(\zero))_1^\infty$ is unbounded.) But $\what g'(\infty) = 1$, so $\what g$ is not loxodromic. Thus $\what g$ is parabolic.
%%%
%%%Recalling that $\zero(\H) = \uparrow$\comdavid{add somewhere}, we have
%%%\[
%%%\what g^{n_k}(\zero) = g^{n_k}(\0) + \uparrow \tendsto k \uparrow \;= \zero,
%%%\]
%%%i.e. the orbit $(g^n(\zero))_1^\infty$ accumulates at $\zero$.
%%%\end{proof}
%%%
%%%For completeness we include Edelstein's proof of Proposition \ref{propositionedelstein}.
%%%
%%%\begin{proof}[Proof of Proposition \ref{propositionedelstein}]
%%%
%%%[Continue\internal]
%%%\end{proof}
%%
%%%%%%%%%%%%%%%%%%%%%%%%%%%%%%%%%%
%%\bigskip
%%\subsection{Classification of semigroups}
%%%%%%%%%%%%%%%%%%%%%%%%%%%%%%%%%%
%%
%%\begin{notation}
%%\label{notationglobalfixedpoints}
%%We denote the set of \emph{global fixed points} of a semigroup $G\prec\Isom(X)$ by
%%\[
%%\Fix(G) := \bigcap_{g\in G}\Fix(g).
%%\]
%%\end{notation}
%%
%%\begin{definition}
%%\label{definitionclassification2}
%%$G$ is
%%\begin{itemize}
%%\item \emph{elliptic} if $G(\zero)$ is a bounded set.
%%\item \emph{parabolic} if $G$ is not elliptic and has a global fixed point $\xi\in \Fix(G)$ such that 
%%\[
%%g'(\xi) = 1 \all g\in G,
%%\]
%%i.e. $\xi$ is neutral with respect to every element of $G$.
%%\item \emph{loxodromic} if it contains a loxodromic isometry.
%%\end{itemize}
%%\end{definition}
%%
%%Below we shall prove the following theorem:
%%
%%\begin{theorem}
%%\label{classificationofsemigroups}
%%Every semigroup of isometries of a hyperbolic metric space is either elliptic, parabolic, or loxodromic.
%%\end{theorem}
%%
%%\begin{observation}
%%An isometry $g$ is elliptic, parabolic, or loxodromic according to whether the cyclic group generated by it is elliptic, parabolic, or loxodromic. A similar statement holds if ``group'' is replaced by ``semigroup''. Thus, Theorem \ref{classificationofisometries} is a special case of Theorem \ref{classificationofsemigroups}.
%%\end{observation}
%%
%%Before proving Theorem \ref{classificationofsemigroups}, let us say a bit about each of the different categories in this classification.
%%
%%\subsubsection{Elliptic semigroups}
%%Elliptic semigroups are the least interesting of the semigroups we consider. Indeed, we observe that any strongly discrete elliptic semigroup is finite. We now consider the question of whether every elliptic semigroup has a global fixed point.
%%
%%\begin{theorem}[Cartan's lemma] \label{Cartan's lemma}
%%If $X$ is a CAT(0) space (and in particular if $X$ is a CAT(-1) space), then every elliptic subsemigroup $G\prec\Isom(X)$ has a global fixed point.
%%\end{theorem}
%%We remark that if $G$ is a group, then this result may be found as \cite[Corollary II.2.8(1)]{BridsonHaefliger}.
%%\begin{proof}
%%Since $G(\zero)$ is a bounded set, it has a unique circumcenter \cite[Proposition II.2.7]{BridsonHaefliger}, i.e. the minimum
%%\[
%%\min_{x\in X}\sup_{g\in G}\dist(x,g(\zero))
%%\]
%%is achieved at a single point $x\in X$. We claim that $x$ is a global fixed point of $G$. Indeed, for each $h\in G$ we have
%%\[
%%\sup_{g\in G}\dist(h^{-1}(x),g(\zero)) = \sup_{g\in G}\dist(x,h g(\zero)) \leq \sup_{g\in G}\dist(x,g(\zero));
%%\]
%%since $x$ is the circumcenter we deduce that $h^{-1}(x) = x$, or equivalently that $h(x) = x$.
%%\end{proof}
%%
%%On the other hand, if we do not restrict to CAT(0) spaces, then it is possible to have an elliptic group with no global fixed point. We have the following simple example:
%%\begin{example}
%%Let $X = \B\butnot B_\B(\0,1)$ and let $g(\xx) = -\xx$. Then $X$ is a hyperbolic metric space, $g$ is an isometry of $X$, and $G = \{\id,g\}$ is an elliptic group with no global fixed point.
%%\end{example}
%%
%%\subsubsection{Parabolic semigroups}
%%Parabolic semigroups will be important in Section \ref{sectionGF} when we consider geometrically finite semigroups. In particular, we make the following definition:
%%\begin{definition}
%%\label{definitionparabolic}
%%Let $G\prec\Isom(X)$. A point $\xi\in\del X$ is a \emph{parabolic fixed point} of $G$ if the semigroup
%%\[
%%G_\xi := \Stab(G;\xi) = \{g\in G:g(\xi) = \xi\}
%%\]
%%is a parabolic semigroup.
%%\end{definition}
%%In particular, if $G$ is a parabolic semigroup then the unique global fixed point of $G$ is a parabolic fixed point.
%%\begin{warning}
%%A parabolic group does not necessarily contain a parabolic isometry; see Example \ref{exampleparabolictorsion}.
%%\end{warning}
%%
%%Note that Proposition \ref{propositioneuclideansimilarity} yields the following observation:
%%\begin{observation}
%%\label{observationuniformlyLipschitz}
%%Let $G\prec\Isom(X)$, and let $\xi$ be a parabolic fixed point of $G$. Then the action of $G_\xi$ on $(\EE_\xi,\Dist_\xi)$ is uniformly Lipschitz, i.e.
%%\[
%%\Dist_\xi(g(y_1),g(y_2)) \asymp_\times \Dist_\xi (y_1,y_2) \all y_1,y_2\in\EE_\xi\all g\in G,
%%\]
%%and the implied constant is independent of $g\in G$. Furthermore, if $X$ is strongly hyperbolic, then $G$ acts isometrically on $\EE_\xi$.
%%\end{observation}
%%
%%\begin{observation}
%%\label{observationeuclideanparabolicasymp}
%%Let $G\prec\Isom(X)$, and let $\xi$ be a parabolic fixed point of $G$. Then for all $g\in G_\xi$,
%%\begin{equation}
%%\label{euclideanparabolicasymp}
%%\Dist_\xi(\zero,g(\zero)) \asymp_\times b^{(1/2)\dogo g},
%%\end{equation}
%%with equality if $X$ is strongly hyperbolic.
%%\end{observation}
%%\begin{proof}
%%This is a direct consequence of \eqref{euclideancomparison}, (h) of Proposition \ref{propositionbasicidentities}, and Proposition \ref{propositionbusemannpreserved}.
%%\end{proof}
%%
%%As a corollary we have the following:
%%\begin{observation}
%%\label{observationparabolic}
%%Let $G\prec\Isom(X)$, and let $\xi$ be a parabolic fixed point of $G$. Then for any sequence $(g_n)_1^\infty$ in $G_\xi$,
%%\[
%%\dogo{g_n}\tendsto n \infty \Leftrightarrow g_n(\zero)\tendsto n\xi.
%%\]
%%\end{observation}
%%\begin{proof}
%%Indeed,
%%\[
%%g_n(\zero)\tendsto n\xi \Leftrightarrow \Dist_\xi(\zero,g_n(\zero))\tendsto n \infty \Leftrightarrow \dogo{g_n}\tendsto n \infty.
%%\]
%%\end{proof}
%%
%%\begin{remark}
%%\label{remarkparabolicRtree2}
%%If $X$ is an $\R$-tree, then any parabolic group must be infinitely generated. This follows from a straightforward modification of the proof of Remark \ref{remarkparabolicRtree}.
%%\end{remark}
%%
%%\subsubsection{Loxodromic semigroups}
%%\label{subsubsectionloxodromic}
%%We now come to loxodromic semigroups, which are the most diverse out of these classes. In fact, they are so diverse that we separate them into three subclasses.
%%\begin{definition}[\cite{CCMT}]
%%\label{definitionCCMT}
%%Let $G\prec\Isom(X)$ be a loxodromic semigroup. $G$ is
%%\begin{itemize}
%%\item \emph{lineal} if $\Fix(g) = \Fix(h)$ for all loxodromic $g,h\in G$.
%%\item \emph{of general type} if it has two loxodromic elements $g,h\in G$ with $\Fix(g)\cap\Fix(h) = \emptyset$.
%%\item \emph{focal} if $\#(\Fix(G)) = 1$.
%%\end{itemize}
%%(We remark that focal groups were called \emph{quasiparabolic} by Gromov \cite[\63, Case 4']{Gromov3}.
%%\end{definition}
%%We observe that any cyclic loxodromic group or semigroup is lineal, so this refined classification does not give any additional information for individual isometries.
%%
%%\begin{proposition}
%%Any loxodromic semigroup is either lineal, focal, or of general type.
%%\end{proposition}
%%\begin{proof}
%%Clearly, $\#(\Fix(G))\leq 2$ for any loxodromic semigroup $G$; moreover, $\#(\Fix(G)) = 2$ if and only if $G$ is lineal. So to complete the proof, it suffices to show that $\#(\Fix(G)) = 0$ if and only if $G$ is of general type. The backward direction is obvious. Suppose that $\#(\Fix(G)) = 0$, but that $G$ is not of general type. Combinatorial considerations show that there exist three points $\xi_1,\xi_2,\xi_3\in\del X$ such that $\Fix(g)\subset\{\xi_1,\xi_2,\xi_3\}$ for all $g\in G$. But then the set $\{\xi_1,\xi_2,\xi_3\}$ would have to be preserved by every element of $g$, which contradicts the definition of a loxodromic isometry.
%%\end{proof}
%%
%%Let $G$ be a focal semigroup, and let $\xi$ be the global fixed point of $G$. The dynamics of $G$ will be different depending on whether or not $g'(\xi) > 1$ for any $g\in G$.
%%\begin{definition}
%%\label{definitionfocal}
%%$G$ will be called \emph{outward focal} if $g'(\xi) > 1$ for some $g\in G$, and \emph{inward focal} otherwise.
%%\end{definition}
%%Note that an inward focal semigroup cannot be a group.
%%
%%
%%\begin{proposition}
%%\label{propositionfocalequivalent}
%%For $G\leq\Isom(X)$, the following are equivalent:
%%\begin{itemize}
%%\item[(A)] $G$ is focal.
%%\item[(B)] $G$ has a unique global fixed point $\xi\in\del X$, and $g'(\xi)\neq 1$ for some $g\in G$.
%%\item[(C)] $G$ has a unique global fixed pont $\xi\in\del X$, and there are two loxodromic isometries $g,h\in G$ so that $g_+ = h_+ = \xi$, but $g_- \neq h_-$.
%%\end{itemize}
%%\end{proposition}
%%\begin{proof}
%%The implications (C) \implies (A) \iff (B) are straightforward. Suppose that $G$ is focal, and let $g\in G$ be a loxodromic isometry. Since $G$ is a group, we may without loss of generality suppose that $g'(\xi) < 1$, so that $g_+ = \xi$. Let $j\in G$ be such that $g_-\notin\Fix(j)$. By choosing $n$ sufficiently large, we may guarantee that $(jg^n)'(\xi) < 1$. Then if $h = jg^n$, then $h$ is loxodromic and $h_+ = \xi$. But $g_-\notin \Fix(h)$, so $g_-\neq h_-$.
%%\end{proof}
%%
%%
%%\bigskip
%%\subsection{Proof of Theorem \ref{classificationofsemigroups}}
%%
%%We begin by recalling the following definition from Subsection \ref{subsectionshadows}:
%%\begin{repdefinition}{definitionshadow}
%%For each $\sigma > 0$ and $x,y\in X$, let
%%\[
%%\Shad_y(x,\sigma) = \{\eta\in\del X:\lb y|\eta\rb_x\leq\sigma\}.
%%\]
%%We say that $\Shad_y(x,\sigma)$ is the \emph{shadow cast by $x$ from the light source $y$, with parameter $\sigma$}. For shorthand we will write $\Shad(x,\sigma) = \Shad_\zero(x,\sigma)$.
%%\end{repdefinition}
%%
%%\begin{lemma}
%%\label{lemmaloxodromic}
%%For every $\sigma > 0$, there exists $r > 0$ such that for every $g\in\Isom(X)$ with $\dogo g \geq r$, if there exists a nonempty closed set
%%\[
%%Z \subset \Shad_{g^{-1}(\zero)}(\zero,\sigma)
%%\]
%%satisfying $g(Z) \subset Z$, then $g$ is loxodromic and $g_+\in Z$.
%%\end{lemma}
%%
%%\begin{proof}
%%%Let $Z = \Shad_{g^{-1}(\zero)}(\zero,\sigma)$, and recall that $Z$ is closed (Observation \ref{observationshadowsclosed}). Note that $Z\neq\emptyset$ since $\zero\in Z$. Our assumption states that
%%%\[
%%%g(Z)\subset Z.
%%%\]
%%Recall from the Bounded Distortion Lemma \ref{lemmaboundeddistortion} that
%%\begin{equation}
%%\label{equationloxodromic}
%%\frac{\Dist(g(y_1),g(y_2))}{\Dist(y_1,y_2)} \leq C b^{-\dogo g} \all y_1,y_2\in Z
%%\end{equation}
%%for some $C > 0$ independent of $g$. Now choose $r > 0$ large enough so that $C b^{-r} < 1$. If $g\in\Isom(X)$ satisfies $\dogo g\geq r$, we can conclude that $g:Z\to Z$ is a strict contraction of the complete metametric space $(Z,\Dist)$. Then by Theorem \ref{banachcontractionprinciple}, $g$ has a unique fixed point $\xi\in Z_\refl = Z\cap \del X$.
%%
%%To complete the proof we must show that $g'(\xi) < 1$ to prove that $g$ is not parabolic and that $\xi = g_+$. Indeed, by the Bounded Distortion Lemma, we have $g'(\xi) \lesssim_\times b^{-\dogo g} \leq b^{-r}$, so choosing $r$ sufficiently large completes the proof.
%%\end{proof}
%%
%%\begin{corollary}
%%\label{corollaryloxodromic}
%%For every $\sigma > 0$, there exists $r = r_\sigma > 0$ such that for every $g\in\Isom(X)$ with $\dogo g\geq r$, if $g$ is not loxodromic, then
%%\begin{equation}
%%\label{nonloxodromic}
%%\lb g(\zero)|g^{-1}(\zero)\rb_\zero \geq \sigma.
%%\end{equation}
%%\end{corollary}
%%\begin{proof}
%%Fix $\sigma > 0$, and let $\sigma' = \sigma + \delta$, where $\delta$ is the implied constant in Gromov's inequality. Apply Lemma \ref{lemmaloxodromic} to get $r' > 0$. Let $r = \max(r',2\sigma')$. Now suppose that $g\in\Isom(X)$ satisfies $\dogo g\geq r\geq r'$ but is not loxodromic. Then by Lemma \ref{lemmaloxodromic}, we have
%%\[
%%\Shad(g(\zero),\sigma')\butnot \Shad_{g^{-1}(\zero)}(\zero,\sigma')\neq \emptyset.
%%\]
%%Let $x$ be a member of this set. By definition this means that
%%\[
%%\lb \zero|x\rb_{g(\zero)} \leq \sigma' < \lb g^{-1}(\zero)|x\rb_\zero.
%%\]
%%Since $\dogo g\geq r\geq 2\sigma'$, we have
%%\[
%%\lb g(\zero)|x\rb_\zero = \dogo g - \lb \zero|x\rb_{g(\zero)} \geq 2\sigma' - \sigma' = \sigma'.
%%\]
%%Now by Gromov's inequality we have
%%\[
%%\lb g(\zero)|g^{-1}(\zero)\rb_\zero \geq \min(\lb g(\zero)|x\rb_\zero,\lb g^{-1}(\zero)|x\rb_\zero) - \delta
%%\geq \sigma' - \delta = \sigma.
%%\]
%%\end{proof}
%%
%%\begin{lemma}\label{keylemma}
%%Let $G\prec\Isom(X)$ be a semigroup which is not loxodromic, and let $(g_n)_1^\infty$ be a sequence in $G$ such that $\dogo{g_n} \to \infty$. Then $(g_n(\zero))_1^\infty$ is a Gromov sequence.
%%\end{lemma}
%%\begin{proof}
%%Fix $\sigma > 0$ large, and let $r = r_\sigma$ be as in Corollary \ref{corollaryloxodromic}. Since $G$ is not loxodromic, \eqref{nonloxodromic} holds for every $g\in G$ for which $\dogo g\geq r$.
%%
%%Fix $n,m\in\N$ with $\dogo{g_n},\dogo{g_m}\geq r$; Corollary \ref{corollaryloxodromic} gives
%%\begin{align} \label{gngn1}
%%\lb g_n(\zero)|g_n^{-1}(\zero)\rb_\zero &\geq \sigma\\ \label{gmgm1}
%%\lb g_m(\zero)|g_m^{-1}(\zero)\rb_\zero &\geq \sigma.
%%\end{align}
%%By contradiction, suppose that $\lb g_n(\zero)|g_m(\zero)\rb_\zero \leq \sigma/2$; then Gromov's inequality together with \eqref{gngn1} gives
%%\begin{equation}
%%\label{gn1gm}
%%\lb g_n^{-1}(\zero)|g_m(\zero)\rb_\zero \asymp_\plus 0.
%%\end{equation}
%%It follows that
%%\[
%%\dogo{g_n g_m} = \dist(g_n^{-1}(\zero),g_m(\zero)) \geq 2r - \lb g_n^{-1}(\zero)|g_m(\zero)\rb_\zero \asymp_\plus 2r.
%%\]
%%Choosing $r$ sufficiently large, we have $\dogo{g_n g_m} \geq r$. So by Corollary \ref{corollaryloxodromic},
%%\begin{equation}
%%\label{gngmgm1gn1}
%%\lb g_n g_m(\zero)|g_m^{-1} g_n^{-1}(\zero)\rb_\zero \geq \sigma.
%%\end{equation}
%%Now
%%\begin{align*}
%%\lb g_n(\zero)|g_n g_m(\zero)\rb_\zero &=_\pt \lb \zero|g_m(\zero)\rb_{g_n^{-1}(\zero)}\\
%%&=_\pt \dogo{g_n} - \lb g_n^{-1}(\zero)|g_m(\zero)\rb_\zero\\
%%&\asymp_\plus \dogo{g_n} \by{\eqref{gn1gm}}\\
%%&\geq_\pt r,
%%\end{align*}
%%i.e.
%%\begin{equation}
%%\label{gngngm}
%%\lb g_n(\zero)|g_n g_m(\zero)\rb_\zero \gtrsim_\plus r.
%%\end{equation}
%%
%%A similar argument yields
%%\begin{equation}
%%\label{gm1gm1gn1}
%%\lb g_m^{-1}(\zero)|g_m^{-1} g_n^{-1}(\zero)\rb_\zero \gtrsim_\plus r.
%%\end{equation}
%%Combining \eqref{gmgm1}, \eqref{gm1gm1gn1}, \eqref{gngmgm1gn1}, and \eqref{gngngm}, together with Gromov's inequality, yields
%%\[
%%\lb g_n|g_m\rb_\zero \gtrsim_\plus \min(\sigma,r).
%%\]
%%This completes the proof.
%%\end{proof}
%%
%%\begin{proof}[Proof of Theorem \ref{classificationofsemigroups}]%\comtushar{this was the first attempt. i'm putting the photos for second attempt on the next page}
%%Suppose that $G$ is neither elliptic nor loxodromic, and we will show that it is parabolic. Since $G$ is not elliptic, there is a sequence $(g_n)_1^\infty$ in $G$ such that $\dogo{g_n} \to \infty$. By Lemma \ref{keylemma}, $(g_n(\zero))_1^\infty$ is a Gromov sequence; let $\xi\in\del X$ be the limit point.
%%
%%Note that $\xi$ is uniquely determined by $G$; if $(h_n(\zero))_1^\infty$ were another Gromov sequence, then we could let
%%\[
%%j_n := \begin{cases}
%%g_{n/2} & n\text{ even}\\
%%h_{(n - 1)/2} & n\text{ odd}
%%\end{cases}.
%%\]
%%The sequence $(j_n(\zero))_1^\infty$ would tend to infinity, so by Lemma \ref{keylemma} it would be a Gromov sequence. But that exactly means that the Gromov sequences $(g_n(\zero))_1^\infty$ and $(h_n(\zero))_1^\infty$ are equivalent. Moreover, it is easy to see that $\xi$ does not depend on the choice of the basepoint $\zero\in X$.
%%
%%In particular, the fact that $\xi$ is canonically determined by $G$ implies that $\xi$ is a global fixed point of $G$. To complete the proof, we need to show that $g'(\xi) = 1$ for all $g\in G$. Suppose we have $g\in G$ such that $g'(\xi)\neq 1$. Then $g$ is loxodromic by Lemma \ref{lemmarepellingloxodromic}, a contradiction.
%%\end{proof}
%%%\ignore{
%%%\noindent Thus $G(\xi) = \xi$. 
%%%
%%%\noindent $G$ is an isometry on $\EE_\xi$.
%%%
%%%\noindent Therefore $G$ is parabolic. 
%%%
%%%\noindent [Note this was left unfinished. See proof of Key Lemma - DSC04237.]
%%%
%%%\comtushar{There was also this observation, but I wasn't sure why we needed it. Also whether we should introduce cross ratios or not. 
%%%\[
%%%b^{[\xi,\eta_0,\eta_1,\eta_2]} = \frac{\Dist_x(\eta_0,\eta_2)}{\Dist_x(\eta_0,\eta_1)} \frac{\Dist_x(\xi,\eta_1)}{\Dist_x(\xi,\eta_0))} = \frac{\Dist_\xi(\eta_0,\eta_2)}{\Dist_\xi(\eta_0,\eta_1)} \cdot
%%%\]
%%%Maybe we needed it to show that $G$ is an isometry on $\EE_\xi$.}
%%%}% end ignore
%%
%%
%%%\comdavid{I am commenting out the ``classification diagram'' since it seems like too much of it depends on $X = \H$.}
%%
%%%\subsection{The classification diagram DSC03592}
%%%
%%%\begin{figure}
%%%\centerline{\mbox{\includegraphics[scale=.12]{photos/DSC03592}}}
%%%\caption{DSC03592}
%%%\label{photoDSC03592}
%%%\end{figure}
%%%
%%%\comtushar{DSC03592 and also 03582-91.}
%%%Given a group $G\leq\Isom(X)$, we may summarize the classification via the following tree of cases:
%%%\begin{enumerate}
%%%\item~ [Case 1: $\sup_{g\in G} \dogo g < \infty .$]
%%%\newline~~ In this case:
%%%\begin{itemize}
%%%\item There exists a global fixed point (GFP) in the interior $X$.
%%%\item All elements of $G$ are elliptic.
%%%\item $G$ embeds into $\Isom(\del \B)$.
%%%\item $\Lambda(G) = \emptyset$ .
%%%\end{itemize}
%%%\item~ [Case 2: $\sup_{g\in G} \dogo g = \infty .$]
%%%\begin{enumerate}
%%%\item~ [Case 2.a: There are no hyperbolic elements in $G$.]
%%%\newline~ \phantom{x} In this case: 
%%%\begin{itemize}
%%%\item The GFP is on the boundary $\del X$.
%%%\item All elements of $G$ are elliptic or parabolic.
%%%\item $G$ embeds into $\O^*(\HH)$.
%%%\item The GFP is unique.
%%%\item $\Lambda(G) = \{ \text{GFP} \}$.
%%%\end{itemize}
%%%\item~ [Case 2.b: There exists a hyperbolic element in $G$.]
%%%\begin{enumerate}
%%%\item~ [Case 2.b.ii: The nonelementary case.]
%%%\newline~ \phantom{x} In this case: 
%%%\begin{itemize}
%%%\item There exists a Schottky subgroup $H \subset G$.
%%%\item $\#(\LambdaG) = \#(\R)$.
%%%\end{itemize}
%%%\item~ [Case 2.b.i: The elementary case.]
%%%\begin{enumerate}
%%%\item~ [Case 2.b.i.A: Fixed point, but not fixed pair.]
%%%\newline~ \phantom{x} In this case: 
%%%\begin{itemize}
%%%\item $G$ is not discrete.
%%%\item $G$ embeds into $\Sim(\HH)$ 
%%%\item The GFP is unique.
%%%\end{itemize}
%%%\item~ [Case 2.b.i.B: Fixed pair $F$.]
%%%\newline~ \phantom{x} In this case we have 
%%%\begin{itemize}
%%%\item There are no other fixed points. In fact, $\Lambda(G) = F$.
%%%\item All elements of $G$ are elliptic or hyperbolic.
%%%\item $G$ embeds into $\Isom(\R) \oplus \Iso (\del \B)$.
%%%\item If $G$ abelian, then the pair $F$ is fixed setwise. 
%%%\end{itemize}
%%%\end{enumerate}
%%%\end{enumerate} 
%%%\end{enumerate} 
%%%\end{enumerate}
%%
%%\bigskip
%%\subsection{Discreteness and focal groups}
%%
%%\begin{proposition}
%%\label{propositionnonfocal}
%%Fix $G\leq\Isom(X)$, and suppose that either
%%\begin{itemize}
%%\item[(1)] $G$ is strongly discrete,
%%\item[(2)] $X$ is CAT(-1) and $G$ is moderately discrete, or
%%\item[(3)] $X$ admits unique geodesic extensions (e.g. $X$ is a \CROSS) and $G$ is weakly discrete.
%%\end{itemize}
%%Then $G$ is not focal.
%%\end{proposition}
%%%\begin{remark}
%%%The proposition fails for semigroups; see e.g. Example \internalmissingref.
%%%\end{remark}
%%\begin{proof}[Strongly discrete case]
%%Suppose that $G$ is a focal group. Let $\xi\in\del X$ be its global fixed point, and let $g,h\in G$ be as in (C) of Proposition \ref{propositionfocalequivalent}. Since $h^{-n}(\zero)\to h_- \neq \xi$, we have
%%\[
%%\lb h^{-n}(\zero)|\xi\rb_\zero \asymp_{\plus,h} 0
%%\]
%%and thus
%%\[
%%\lb h^n(\zero)|\xi\rb_\zero \asymp_\plus \dogo{h^n} - \lb\zero|\xi\rb_{h^n(\zero)} \asymp_{\plus,h} \dogo{h^n}.
%%\]
%%Applying $g$ we have
%%\[
%%\lb gh^n(\zero)|\xi\rb_\zero
%%= \lb h^n(\zero)|\xi\rb_{g^{-1}(\zero)}
%%\asymp_{\plus,g} \lb h^n(\zero)|\xi\rb_\zero
%%\asymp_{\plus,h} \dogo{h^n}
%%\]
%%and applying Gromov's inequality we have
%%\[
%%\lb h^n(\zero) | gh^n(\zero)\rb_\zero \asymp_{\plus,g,h} \dogo{h^n} \asymp_{\plus,g} \dogo{g h^n}.
%%\]
%%Now
%%\begin{align*}
%%\dogo{h^{-n} gh^n}
%%&= \dist(h^n(\zero),gh^n(\zero))\\
%%&= \dogo{h^n} + \dogo{g h^n} - 2\lb h^n(\zero),gh^n(\zero)\rb_\zero
%%\asymp_{\plus,g,h} 0.
%%\end{align*}
%%Since $G$ is strongly discrete, this implies that the collection $\{h^{-n} gh^n:n\in\N\}$ is finite, and so for some $n_1 < n_2$ we have
%%\[
%%h^{-n_1} gh^{n_1} = h^{-n_2} gh^{n_2}
%%\]
%%or
%%\[
%%h^{n_2 - n_1} g = g h^{n_2 - n_1},
%%\]
%%i.e. $h^{n_2 - n_1}$ commutes with $g$. But then $h^{n_2 - n_1}(g_-) = g_-$, contradicting that $g_-\neq h_-$.
%%\end{proof}
%%
%%\begin{proof}[Moderately discrete case]
%%Suppose that $G$ is a focal group. Let $\xi\in\del X$ be its global fixed point, and let $g,h\in G$ be as in (C) of Proposition \ref{propositionfocalequivalent}. Let
%%\[
%%k = [g,h] = g^{-1} h^{-1} gh\in G.
%%\]
%%We observe first that
%%\begin{equation}
%%\label{kprimeequalsone}
%%k'(\xi) = \frac{1}{g'(\xi)}\frac{1}{h'(\xi)}g'(\xi) h'(\xi) = 1.
%%\end{equation}
%%Note that strong hyperbolicity is necessary to deduce equality in \eqref{kprimeequalsone} rather than merely a coarse asymptotic.
%%
%%Next, we claim that $k(g_-)\neq g_-$. Indeed, $g_-\notin\Fix(h)$, so $h(g_-)\neq g_-$. This in turn implies that $h(g_-)\notin\Fix(g)$, so $gh(g_-)\neq h(g_-)$. Now applying $g^{-1} h^{-1}$ to both sides shows that $k(g_-)\neq g_-$.
%%
%%\begin{claim}
%%\label{claimfocalMD}
%%$g^{-n} k g^n(\zero)\to \zero$.
%%\end{claim}
%%\begin{subproof}
%%Indeed,
%%\[
%%\dogo{g^{-n} k g^n} = \dist(g^n(\zero),k g^n(\zero)).
%%\]
%%Let
%%\begin{align*}
%%x &= \xi\\
%%y &= \zero\\
%%z &= k(\zero)\\
%%p_n &= g^n(\zero)\\
%%q_n &= k g^n(\zero).
%%\end{align*}
%%(See Figure \ref{figurexyzpq}.) Then $p_n,q_n\in\Delta := \Delta(x,y,z)$. By Proposition \ref{propositionidealCAT} $\dist(p_n,q_n) \leq \dist(\wbar p_n,\wbar q_n)$, where $\wbar p_n,\wbar q_n$ are comparison points for $p_n,q_n$ on the comparison triangle $\wbar \Delta = \Delta(\wbar x,\wbar y,\wbar z)$. Now notice that
%%\[
%%\busemann_{\wbar x}(\wbar p_n,\wbar q_n) = \busemann_\xi(g^n(\zero),k g^n(\zero)) = 0
%%\]
%%by Proposition \ref{propositionbusemannpreserved} and \eqref{kprimeequalsone}. On the other hand, $\wbar p_n,\wbar q_n\to \wbar x$. An easy calculation based on \eqref{distE} and Proposition \ref{propositionbusemannE} (letting $\wbar x = \infty$) shows that $\dist(\wbar p_n,\wbar q_n)\to 0$, and thus that $\dogo{g^{-n} k g^n}\to 0$ i.e. $g^{-n} k g^n(\zero)\to \zero$.
%%\end{subproof}
%%Since $G$ is moderately discrete, this implies that the collection $\{g^{-n} k g^n:n\in\N\}$ is finite. As before (in the proof of the strongly discrete case), this implies that $g^n$ and $k$ commute for some $n\in\N$. But $(g^n)_- = g_-$, and $k(g_-) \neq g_-$, which contradicts that $g^n$ and $k$ commute.
%%\end{proof}
%%
%%\begin{figure}
%%\begin{center}
%%\begin{tikzpicture}[line cap=round,line join=round,>=triangle 45,x=1.0cm,y=1.0cm]
%%\clip(-5.13,-0.16) rectangle (6.57,8.26);
%%\draw (0,1.08)-- (0,6.6);
%%\draw (2,1.1)-- (2,6.6);
%%\draw [shift={(1,0.14)}] plot[domain=1.32:1.82,variable=\t]({1*3.98*cos(\t r)+0*3.98*sin(\t r)},{0*3.98*cos(\t r)+1*3.98*sin(\t r)});
%%\draw [shift={(1,-0.03)}] plot[domain=0.84:2.31,variable=\t]({1*1.49*cos(\t r)+0*1.49*sin(\t r)},{0*1.49*cos(\t r)+1*1.49*sin(\t r)});
%%\draw (-4,0)-- (6,0);
%%\begin{scriptsize}
%%\fill [color=black] (2,1.1) circle (1pt);
%%\draw[color=black] (2.78,1.2) node {$z=k(\zero)$};
%%\fill [color=black] (1.03,6.98) circle (1pt);
%%\draw[color=black] (1.43,7.15) node {$x=\xi$};
%%\fill [color=black] (0,4) circle (1pt);
%%\draw[color=black] (-0.89,4.14) node {$p_n = g^n(\zero)$};
%%\fill [color=black] (2,4) circle (1pt);
%%\draw[color=black] (3.07,4.06) node {$q_n = kg^n(\zero)$};
%%\fill [color=black] (0,1.08) circle (1pt);
%%\draw[color=black] (-0.3,1.2) node {$\zero$};
%%\end{scriptsize}
%%\end{tikzpicture}
%%\caption{The higher the point $g^n(\zero)$ is, the smaller its displacement under $k$ is.}
%%\label{figurexyzpq}
%%\end{center}
%%\end{figure}
%%
%%\begin{proof}[Weakly discrete case]
%%Suppose that $G$ is a focal group. Let $\xi$, $g$, $h$, and $k$ be as above. Without loss of generality, supposet that $\zero\in \geo{g_-}{\xi}$.
%%
%%\begin{claim}
%%$g^{-n} k g^n(\zero)\neq \zero$ for all $n\in\N$.
%%\end{claim}
%%\begin{subproof}
%%Fix $n\in\N$. As observed above, $k(g_-)\neq g_-$. On the other hand, $k(\xi) = \xi$, and $g^n(\zero)\in \geo{g_-}{\xi}$. Since $X$ admits unique geodesic extensions, it follows that $k(g^n(\zero))\neq g^n(\zero)$, or equivalently that $g^{-n} k g^n(\zero)\neq \zero$.
%%\end{subproof}
%%Together with Claim \ref{claimfocalMD}, this contradicts that $G$ is weakly discrete.
%%\end{proof}
%%
%%
%%%%%%%%%%%%%%%%%%%%%%%%%%%%%%%%
%%\draftnewpage
%%\section{Limit sets} \label{sectionlimitsets}
%%%%%%%%%%%%%%%%%%%%%%%%%%%%%%%%%%
%%
%%Throughout this section, we fix a subsemigroup $G\prec\Isom(X)$. We define the \emph{limit set} of $G$, along with various subsets. We then define several concepts in terms of the limit set including elementariness and compact type, while relating other concepts to the limit set, such as the quasiconvex core and irreducibility of a group action. We also prove that the limit set is minimal in an approprate sense (Proposition \ref{propositionminimal} - Proposition \ref{propositionminimal2}).
%%
%%
%%
%%
%%%%%%%%%%%%%%%%%%%%%%%%%%%%%%%%%%%%
%%\bigskip
%%\subsection{Modes of convergence to the boundary}
%%\label{subsectionmodes}
%%%%%%%%%%%%%%%%%%%%%%%%%%%%%%%%%%%%
%%We recall (Observation \ref{observationboundaryconvergence}) that a sequence $(x_n)_1^\infty$ in $X$ converges to a point $\xi\in\del X$ if and only if
%%\[
%%\lb x_n|\xi\rb_\zero \tendsto n \infty.
%%\]
%%In this subsection we define more restricted modes of convergence. To get an intuition let us consider the case where $X = \E = \E^\alpha$ is the half-space model of a real \ROSS. Consider a sequence $(\xx_n)_1^\infty$ in $\E$ which converges to a point $\xi\in\BB := \del\E\butnot\{\infty\} = \HH^{\alpha - 1}$. We say that $\xx_n \to \xi$ \emph{conically} if there exists $\theta > 0$ such that if we let
%%\[
%%C(\xi,\theta) = \{\xx\in\E: x_1 \geq \sin(\theta)\|\xx - \xi\|\}
%%\]
%%then $\xx_n\in C(\xi,\theta)$ for all $n\in\N$. We call $C(\xi,\theta)$ the \emph{cone centered at $\xi$ with inclination $\theta$}; see Figure \ref{figurecone}.
%%
%%
%%
%%\definecolor{qqwuqq}{rgb}{0,0.39,0}
%%\begin{figure}
%%\begin{center}
%%\begin{tikzpicture}[line cap=round,line join=round,>=triangle 45,x=1.0cm,y=1.0cm]
%%\clip(-5.69,-0.26) rectangle (5.35,3.56);
%%\draw [shift={(0,0)},color=qqwuqq,fill=qqwuqq,fill opacity=0.1] (0,0) -- (147.02:0.44) arc (147.02:180:0.44) -- cycle;
%%\draw (0,0)-- (4.26,2.76);
%%\draw (-4.26,2.76)-- (0,0);
%%\draw (-5,0)-- (5,0);
%%\begin{scriptsize}
%%\fill [color=black] (0,0) circle (1.2pt);
%%\draw[color=black] (0.09,-0.17) node {$\xi$};
%%\draw[color=black] (-0.60,0.18) node {$\theta$};
%%\fill [color=black] (-1.62,2.69) circle (0.81pt);
%%\fill [color=black] (0.35,1.2) circle (0.81pt);
%%\fill [color=black] (-1.67,1.77) circle (0.81pt);
%%\fill [color=black] (-0.26,1.24) circle (0.81pt);
%%\fill [color=black] (0.04,0.64) circle (0.81pt);
%%\draw[color=black] (0.35,2.2) node {$C(\xi,\theta)$};
%%\end{scriptsize}
%%\end{tikzpicture}
%%\caption{A sequence converging conically to $\xi$. For each point $\xx$, the height of $\xx$ is greater than $\sin(\theta)$ times the distance from $\xx$ to $\xi$.}
%%\label{figurecone}
%%\end{center}
%%\end{figure}
%%
%%
%%%\begin{figure}
%%%\centerline{\mbox{\includegraphics[scale=.5]{cone.png}}}
%%%\caption{A sequence converging conically to $\xi$. For each point $\xx$, the height of $\xx$ is greater than $\sin(\theta)$ times the distance from $\xx$ to $\xi$.}
%%%\label{figurecone}
%%%\end{figure}
%%
%%
%%\begin{proposition}
%%\label{propositionradialconvergence}
%%Let $(\xx_n)_1^\infty$ be a sequence in $\E$ converging to a point $\xi\in\BB$. Then the following are equivalent:
%%\begin{itemize}
%%\item[(A)] $(\xx_n)_1^\infty$ converges conically to $\xi$.
%%\item[(B)] The sequence $(\xx_n)_1^\infty$ lies within a bounded distance of the geodesic ray $\geo\zero\xi$.
%%\item[(C)] There exists $\sigma > 0$ such that for all $n\in\N$,
%%\[
%%\lb\zero|\xi\rb_{\xx_n}\leq\sigma,
%%\]
%%or equivalently
%%\begin{equation}
%%\label{sigmaradially}
%%\xi\in\Shad(\xx_n,\sigma).
%%\end{equation}
%%\end{itemize}
%%Moreover, the equivalence of (B) and (C) holds in all geodesic hyperbolic metric spaces.
%%\end{proposition}
%%\begin{proof}
%%The equivalence of (B) and (C) follows directly from (i) of Proposition \ref{propositionrips}. Moreover, conditions (B) and (C) are clearly independent of the basepoint $\zero$. Thus, in proving (A) \iff (B) we may without loss of generality suppose that $\xi = \0$ and $\zero = (1,\0)$. Note that if $\theta > 0$ is fixed, then
%%\[
%%C(\0,\theta) = \{\xx\in\E : \ang(\xx) \leq \pi/2 - \theta\}
%%= \{\xx\in\E : \theta(\xx) \leq -\log\cos(\pi/2 - \theta)\},
%%\]
%%where $\theta = \theta_{\0,\infty,\zero}$ is as in Proposition \ref{propositionrtheta}. Since $-\log\cos(\pi/2 - \theta) \to \infty$ as $\theta \to 0$, we have (A) if and only if the sequence $(\theta(\xx_n))_1^\infty$ is bounded. But
%%\begin{align*}
%%\theta(\xx_n) = \lb \0 | \infty \rb_{\xx_n}
%%&\asymp_\plus \dist(\xx_n,\geo\0\infty) \by{(i) of Proposition \ref{propositionrips}}\\
%%&=_\pt \dist(\xx_n,\geo\zero\0), &\text{(for $n$ sufficiently large)}
%%\end{align*}
%%which completes the proof.
%%\end{proof}
%%
%%Condition (B) of Proposition \ref{propositionradialconvergence} motivates calling this kind of convergence \emph{radial}; we shall use this terminology henceforth. However, condition (C) is best suited to a general hyperbolic metric space.
%%
%%\begin{definition}
%%\label{definitionradialconvergence}
%%Let $(x_n)_1^\infty$ be a sequence in $X$ converging to a point $\xi\in\del X$. We will say that $(x_n)_1^\infty$ converges to $\xi$
%%\begin{itemize}
%%\item \emph{$\sigma$-radially} if \eqref{sigmaradially} holds for all $n\in\N$,
%%\item \emph{radially} if it converges $\sigma$-radially for some $\sigma > 0$,
%%\item \emph{$\sigma$-uniformly radially} if it converges $\sigma$-radially, $x_1 = \zero$, and
%%\[
%%\dist(x_n,x_{n + 1})\leq \sigma \all n\in\N,
%%\]
%%\item \emph{uniformly radially} if it converges $\sigma$-uniformly radially for some $\sigma > 0$.
%%\end{itemize}
%%\end{definition}
%%Note that a sequence can converge $\sigma$-radially and uniformly radially without converging $\sigma$-uniformly radially.
%%
%%We next define horospherical convergence. Again, we motivate the discussion by considering the case of a real \ROSS\ $X = \E = \E^\alpha$. This time, however, we will let $\xi = \infty$, and we will say that a sequence $(\xx_n)_1^\infty$ converges \emph{horospherically} to $\xi$ if
%%\[
%%\height(\xx_n) \tendsto n \infty,
%%\]
%%where the \emph{height} of a point $\xx\in\E$ is its first coordinate $x_1$. This terminology comes from defining a \emph{horoball centered at $\infty$} to be a set of the form $H_{\infty,t} = \{\xx : \height(\xx) > e^t\}$; then $\xx_n\to\infty$ horospherically if and only if for every horoball $H_{\infty,t}$ centered at infinity, we have $\xx_n\in H_{\infty,t}$ for all sufficiently large $n$. (See also Definition \ref{definitionhoroball} below.)
%%
%%Recalling (cf. Proposition \ref{propositionbusemannE}) that
%%\[
%%\height(x) = b^{\busemann_\infty(\zero,x)},
%%\]
%%the above discussion motivates the following definition:
%%\begin{definition}
%%\label{definitionhorosphericalconvergence}
%%A sequence $(x_n)_1^\infty$ in $X$ converges \emph{horospherically} to a point to $\xi\in\del X$ if
%%\[
%%\busemann_\xi(\zero,x_n)\tendsto n +\infty.
%%\]
%%\end{definition}
%%
%%\begin{observation}
%%If $x_n\to\xi$ radially, then $x_n\to\xi$ horospherically.
%%\end{observation}
%%\begin{proof}
%%Indeed,
%%\[
%%\busemann_\xi(\zero,x_n) \asymp_\plus \dox{x_n} - 2\lb\zero|\xi\rb_{x_n} \asymp_\plus \dox{x_n} \tendsto n \infty.
%%\]
%%\end{proof}
%%
%%The converse is false, as illustrated in Figure \ref{figurehorospherical}.
%%
%%
%%\begin{figure}
%%\begin{center}
%%\begin{tikzpicture}[line cap=round,line join=round,>=triangle 45,x=1.0cm,y=1.0cm]
%%\clip(-5.69,-0.39) rectangle (5.35,3.76);
%%\draw (0,0)-- (4.26,2.76);
%%\draw (-4.26,2.76)-- (0,0);
%%\draw (-5,0)-- (5,0);
%%\draw(0.03,1.86) circle (1.86cm);
%%\begin{scriptsize}
%%\fill [color=black] (0,0) circle (1.2pt);
%%\draw[color=black] (0.12,-0.26) node {$\xi$};
%%\fill [color=black] (-1.61,2.26) circle (0.81pt);
%%\fill [color=black] (-0.91,0.42) circle (0.81pt);
%%\fill [color=black] (-1.61,1.39) circle (0.81pt);
%%\fill [color=black] (-1.32,0.75) circle (0.81pt);
%%\fill [color=black] (-0.5,0.19) circle (0.81pt);
%%\end{scriptsize}
%%\end{tikzpicture}
%%\caption{A sequence converging horospherically but not radially to $\xi$.}
%%\label{figurehorospherical}
%%\end{center}
%%\end{figure}
%%
%%%\begin{figure}
%%%\centerline{\mbox{\includegraphics[scale=.5]{horosphere.png}}}
%%%\caption{A sequence converging horospherically but not radially to $\xi$.}
%%%\label{figurehorospherical}
%%%\end{figure}
%%
%%\begin{observation}
%%\label{observationdependenceonbasepoint}
%%The concepts of convergence, radial convergence, uniformly radial convergence, and horospherical convergence are independent of the basepoint $\zero$, whereas the concepts of $\sigma$-radial convergence and $\sigma$-uniformly radial convergence depend on the basepoint. (Regarding $\sigma$-radial convergence, this dependence on basepoint is not too severe; see Proposition \ref{propositionnearinvariance} below.)
%%\end{observation}
%%
%%%%%%%%%%%%%%%%%%%%%%%%%%%%%%%%%%%%
%%\bigskip
%%\subsection{Limit sets}
%%%%%%%%%%%%%%%%%%%%%%%%%%%%%%%%%%%%
%%
%%We define the limit set of $G$, a subset of $\del X$ which encodes geometric information about $G$. We also define a few important subsets of the limit set.
%%
%%\begin{definition}
%%\label{definitionlimitset}
%%Let
%%\begin{align*}
%%\Lambda(G) &:= \{\eta\in\del X: g_n(\zero)\to \eta \text{ for some sequence $(g_n)_1^\infty$ in $G$}\}\\
%%\Lr(G) &:= \{\eta\in\del X: g_n(\zero)\to \eta\text{ radially for some sequence $(g_n)_1^\infty$ in $G$}\}\\
%%\Lrsigma(G) &:= \{\eta\in\del X: g_n(\zero)\to \eta\text{ $\sigma$-radially for some sequence $(g_n)_1^\infty$ in $G$}\}\\
%%\Lur(G) &:= \{\eta\in\del X: g_n(\zero)\to \eta\text{ uniformly radially for some sequence $(g_n)_1^\infty$ in $G$}\}\\
%%\Lursigma(G) &:= \{\eta\in\del X: g_n(\zero)\to \eta\text{ $\sigma$-uniformly radially for some sequence $(g_n)_1^\infty$ in $G$}\}\\
%%\Lh(G) &:= \{\eta\in\del X: g_n(\zero)\to \eta\text{ horospherically for some sequence $(g_n)_1^\infty$ in $G$}\}.
%%\end{align*}
%%These sets are respectively called the \emph{limit set}, \emph{radial limit set}, \emph{$\sigma$-radial limit set}, \emph{uniformly radial limit set}, \emph{$\sigma$-uniformly radial limit set}, and \emph{horospherical limit set} of the semigroup $G$.
%%\end{definition}
%%Note that
%%\begin{align*}
%%\Lr &= \bigcup_{\sigma > 0}\Lrsigma\\
%%\Lur &= \bigcup_{\sigma > 0}\Lursigma\\
%%\Lur &\subset \Lr \subset \Lh \subset \Lambda.
%%\end{align*}
%%
%%\begin{observation}
%%\label{observationlimitsets}
%%The sets $\Lambda$, $\Lr$, $\Lur$, and $\Lh$ are invariant\Footnote{By \emph{invariant} we always mean \emph{forward invariant}.} under the action of $G$, and are independent of the basepoint $\zero$. The set $\Lambda$ is closed.
%%\end{observation}
%%\begin{proof}
%%The first assertion follows from Observation \ref{observationdependenceonbasepoint} and the second follows directly from the definition of $\Lambda$ as the intersection of $\del X$ with the set of accumulation points of the set $G(\zero)$.
%%\end{proof}
%%
%%
%%\begin{proposition}[Near-invariance of the sets $\Lrsigma$]
%%\label{propositionnearinvariance}
%%For every $\sigma > 0$, there exists $\tau > 0$ such that for every $g\in G$, we have
%%\begin{equation}
%%\label{nearinvariance}
%%g(\Lrsigma) \subset \Lrtau.
%%\end{equation}
%%If $X$ is strongly hyperbolic, then \eqref{nearinvariance} holds for all $\tau > \sigma$.
%%\end{proposition}
%%\begin{proof}
%%Fix $\xi\in\Lrsigma$. There exists a sequence $(h_n)_1^\infty$ so that $h_n(\zero)\to \xi$ $\sigma$-radially, i.e.
%%\[
%%\lb \zero|\xi\rb_{h_n(\zero)} \leq \sigma\all n\in\N
%%\]
%%and $h_n(\zero)\to \xi$. Now
%%\[
%%\lb \zero|g^{-1}(\zero)\rb_{h_n(\zero)} \geq \dogo{h_n} - \dogo{g^{-1}} \tendsto n +\infty.
%%\]
%%Thus, for $n$ sufficiently large, Gromov's inequality gives
%%\begin{equation}
%%\label{ghnsigma}
%%\lb g^{-1}(\zero)|\xi\rb_{h_n(\zero)} \lesssim_\plus \sigma
%%\end{equation}
%%i.e.
%%\[
%%\lb \zero|g(\xi)\rb_{gh_n(\zero)} \lesssim_\plus \sigma.
%%\]
%%So $gh_n(\zero)\to g(\xi)$ $\tau$-radially, where $\tau$ is equal to $\sigma$ plus the implied constant of this asymptotic. Thus, $g(\xi)\in\Lrtau$.
%%
%%If $X$ is strongly hyperbolic, then by using \eqref{gromovforCAT} instead of Gromov's inequality, the implied constant of \eqref{ghnsigma} can be made arbitrarily small. Thus $\tau$ may be taken arbitrarily close to $\sigma$.
%%\end{proof}
%%
%%%%%%%%%%%%%%%%%%%%%%%%%%%%%%%%%%%%
%%\bigskip
%%\subsection{Cardinality of the limit set}
%%%%%%%%%%%%%%%%%%%%%%%%%%%%%%%%%%%%
%%
%%In this section we characterize the cardinality of the limit set according to the classification of the semigroup $G$.
%%
%%\begin{proposition}[Cardinality of the limit set by classification]
%%\label{propositioncardinalitylimitset}
%%Fix $G\prec\Isom(X)$.
%%\begin{itemize}
%%\item[(i)] If $G$ is elliptic, then $\LambdaG = \emptyset$.
%%\item[(ii)] If $G$ is parabolic or inward focal with global fixed point $\xi$, then $\LambdaG = \{\xi\}$.
%%\item[(iii)] If $G$ is lineal with fixed pair $\{\xi_1,\xi_2\}$, then $\LambdaG \subset \{\xi_1,\xi_2\}$, with equality if $G$ is a group.
%%\item[(iv)] If $G$ is outward focal or of general type, then $\#(\LambdaG) \geq \#(\R)$. Equality holds if $X$ is separable.
%%\end{itemize}
%%\end{proposition}
%%Case (i) is immediate, while case (iv) requires the theory of Schottky groups and will be proven in Section \ref{sectionschottky} (see Proposition \ref{propositionnonelementaryequivalent}).
%%\begin{proof}[Proof of \text{(ii)}]
%%For $g\in G$, $g'(\xi) \leq 1$, so by Proposition \ref{propositionbusemannpreserved}, we have $\busemann_\xi(g(\zero),\zero) \lesssim_\plus 0$. In particular, by (h) of Proposition \ref{propositionbasicidentities} we have
%%\[
%%\lb x | \xi \rb_\zero \gtrsim_\plus \frac12 \dox x \all x\in G(\zero).
%%\]
%%This implies that $x_n\to \xi$ for any sequence $(x_n)_1^\infty$ in $G(\zero)$ satisfying $\dox{x_n}\to \infty$. It follows that $\LambdaG = \{\xi\}$.
%%\end{proof}
%%\begin{proof}[Proof of \text{(iii)}]
%%By Lemma \ref{lemmatranslationdistance} we have
%%\[
%%\theta(g(\zero)) \asymp_\plus \theta(\zero) = \zero \all g\in G,
%%\]
%%where $\theta = \theta_{\xi_1,\xi_2,\zero} = \theta_{\xi_2,\xi_1,\zero}$ is as in Subsection \ref{subsectionpolar}. Thus
%%\[
%%\lb \xi_1 | \xi_2 \rb_x \asymp_\plus 0 \all x\in G(\zero).
%%\]
%%Fix a sequence $G(\zero)\ni x_n \to \xi\in\LambdaG$. By Gromov's inequality, there exists $i = 1,2$ such that
%%\[
%%\lb \zero | \xi_i \rb_{x_n} \asymp_\plus 0 \text{ for infinitely many $n$.}
%%\]
%%It follows that $x_n \to \xi_i$ radially along some subsequence, and in particular $\xi = \xi_i$. Thus $\LambdaG \subset \{\xi_1,\xi_2\}$.
%%\end{proof}
%%
%%\begin{definition}
%%\label{definitionelementary}
%%Fix $G\prec\Isom(X)$. $G$ is called \emph{elementary} if $\#(\Lambda) < \infty$ and \emph{nonelementary} if $\#(\Lambda) = \infty$.
%%\end{definition}
%%
%%Thus, according to Proposition \ref{propositioncardinalitylimitset}, elliptic, parabolic, lineal, and inward focal semigroups are elementary while outward focal semigroups and semigroups of general type are nonelementary.
%%
%%\begin{remark}
%%In the Standard Case, some authors (e.g. \cite[\65.5]{Ratcliffe}) define a subgroup of $\Isom(X)$ to be elementary if there is a global fixed point or a global fixed geodesic line. According to this definition, focal groups are considered elementary. By contrast, we follow \cite{CCMT} and others in considering them to be nonelementary.
%%
%%Another common definition in the Standard Case is that a group is elementary if it is virtually abelian. This agrees with our definition, but beyond the Standard Case this equivalence no longer holds (cf. Observation \ref{observationmargulislemma} and Remark \ref{remarkmargulislemma}).
%%\end{remark}
%%
%%%%%%%%%%%%%%%%%%%%%%%%%%%%%%%%%%%%
%%\bigskip
%%\subsection{Minimality of the limit set}
%%%%%%%%%%%%%%%%%%%%%%%%%%%%%%%%%%%%
%%
%%Observation \ref{observationlimitsets} identified the limit set $\Lambda$ as a closed $G$-invariant subset of the Gromov boundary $\del X$. In this section, we give a characterization of $\Lambda$ depending on the classification of $G$.
%%
%%\begin{proposition}[Cf. {\cite[Th\'eor\`eme 5.1]{Coornaert}}]
%%\label{propositionminimal}
%%Fix $G\prec\Isom(X)$. Then any closed $G$-invariant subset of $\del X$ containing at least two points contains $\Lambda$.
%%\end{proposition}
%%\begin{proof}
%%We begin with the following lemma, which will also be useful later:
%%\begin{lemma}
%%\label{lemmaminimal}
%%Let $(x_n)_1^\infty$, $(y_n^{(1)})_1^\infty$, $(y_n^{(2)})_1^\infty$ be sequences in $\bord X$ satisfying
%%\[
%%\lb y_n^{(1)} | y_n^{(2)} \rb_{x_n} \asymp_\plus 0
%%\]
%%and
%%\[
%%x_n \to \xi\in\del X.
%%\]
%%Then $\xi\in\cl{\{y_n^{(i)} : n\in\N , i = 1,2\}}$.
%%\end{lemma}
%%\begin{subproof}
%%For $n\in\N$ fixed, by Gromov's inequality there exists $i_n = 1,2$ such that
%%\[
%%\lb \zero | y_n^{(i_n)} \rb_{x_n} \asymp_\plus 0.
%%\]
%%It follows that
%%\[
%%\lb x_n | y_n^{(i_n)}\rb_\zero \asymp_\plus \dox{x_n} \tendsto n \infty.
%%\]
%%On the other hand
%%\[
%%\lb x_n | \xi\rb_\zero \tendsto n \infty,
%%\]
%%so by Gromov's inequality
%%\[
%%\lb y_n^{(i_n)} | \xi \rb_\zero \tendsto n \infty,
%%\]
%%i.e. $y_n^{(i_n)} \to \xi$.
%%\end{subproof}
%%Now let $F$ be a closed $G$-invariant subset of $\del X$ containing two points $\xi_1\neq\xi_2$, and let $\eta\in\Lambda$. Then there exists a sequence $(g_n)_1^\infty$ so that $g_n(\zero)\to\eta$. Applying Lemma \ref{lemmaminimal} with $x_n = g_n(\zero)$ and $y_n^{(i)} = g_n(\xi_i)\in F$ completes the proof.
%%
%%%Let $F$ be a closed $G$-invariant subset of $\del X$ containing two points $\xi_1\neq\xi_2$, and let $\eta\in\Lambda$. Then there exists a sequence $(g_n)_1^\infty$ so that $g_n(\zero)\to\eta$. We have
%%%\[
%%%0 \asymp_{\plus,\xi_1,\xi_2} \lb \xi_1|\xi_2\rb_\zero
%%%\gtrsim_\plus \min\left(\lb g_n^{-1}(\zero)|\xi_1\rb_\zero,\lb g_n^{-1}(\zero)|\xi_2\rb_\zero\right).
%%%\]
%%%Let $i_n$ be $1$ or $2$, depending on which term in the minimum is smaller. We then have 
%%%\[
%%%\lb g_n^{-1}(\zero)|\xi_{i_n}\rb_\zero \asymp_{\plus,\xi_1,\xi_2} 0,
%%%\]
%%%or equivalently
%%%\[
%%%\lb g_n(\zero)|g_n(\xi_{i_n})\rb_\zero \asymp_{\plus,\xi_1,\xi_2} \dogo{g_n} \tendsto n \infty.
%%%\]
%%%On the other hand, since $g_n(\zero)\to\eta$ we have
%%%\[
%%%\lb g_n(\zero)|\eta\rb_\zero\tendsto n \infty.
%%%\]
%%%Applying Gromov's inequality yields
%%%\[
%%%\lb g_n(\xi_{i_n})|\eta\rb_\zero \tendsto n \infty
%%%\]
%%%and thus $g_n(\xi_{i_n})\to\eta$. But $g_n(\xi_{i_n})\in F$ for all $n$ by the $G$-invariance of $F$; since $F$ is closed we have $\eta\in F$.
%%\end{proof}
%%The proof of Proposition \ref{propositionminimal} may be compared to the proof of \cite[Theorem 3.1]{FSU4}, where a quantitative convergence result is proven assuming that $\eta$ is in the radial limit set (and assuming that $G$ is a group).
%%
%%\begin{corollary}
%%\label{corollaryminimal}
%%Let $G\prec\Isom(X)$ be nonelementary.
%%\begin{itemize}
%%\item[(i)] If $G$ is outward focal with global fixed point $\xi$, then $\LambdaG$ is the smallest closed $G$-invariant subset of $\del X$ which contains a point other than $\xi$.
%%\item[(ii)] (Cf. {\cite[Theorem 5.3.7]{Beardon_book}}) If $G$ is of general type, then $\LambdaG$ is the smallest nonempty closed $G$-invariant subset of $\del X$.
%%\end{itemize}
%%\end{corollary}
%%\begin{proof}
%%Any $G$-invariant set containing a point which is not fixed by $G$ contains two points.
%%\end{proof}
%%
%%\begin{corollary}
%%Let $G\prec\Isom(X)$ be nonelementary. Then
%%\[
%%\Lambda = \cl{\Lr} = \cl{\Lur}.
%%\]
%%\end{corollary}
%%\begin{proof}
%%The implications $\supset$ are clear. On the other hand, for each loxodromic $g\in G$ we have $g_+\in\Lur$. Thus $\Lur\neq\emptyset$, and $\Lur\nsubset\{\xi\}$ if $G$ is outward focal with global fixed point $\xi$. By Proposition \ref{propositioncardinalitylimitset}, $G$ is either outward focal or of general type. Applying Corollary \ref{corollaryminimal}, we have $\cl{\Lur}\supset\LambdaG$.
%%\end{proof}
%%
%%\begin{remark}
%%If $G$ is elementary, it is easily verified that $\Lambda = \Lr = \Lur$ unless $G$ is parabolic, in which case $\Lr = \Lur = \emptyset \propersubset \Lambda$.
%%\end{remark}
%%
%%If $G$ is a nonelementary group, then Corollary \ref{corollaryminimal} immediately implies that the set of loxodromic fixed points of $G$ is dense in $\Lambda$. However, if $G$ is not a group then this conclusion does not follow, since the set of attracting loxodromic fixed points is not necessarily $G$-invariant. (The set of attracting fixed points is the right set to consider, since the set of repelling fixed points is not necessarily a subset of $\Lambda$.) Nevertheless, we have the following:
%%
%%\begin{proposition}
%%\label{propositionminimal2}
%%Let $G\prec\Isom(X)$ be nonelementary. Then the set
%%\[
%%\Lambda_+ := \{g_+ : g\in G \text{ is loxodromic}\}.
%%\]
%%is dense in $\Lambda$.
%%\end{proposition}
%%\begin{proof}
%%First note that it suffices to show that $\cl{\Lambda_+}$ contains all elements of $\Lambda$ which are not global fixed points. Indeed, if this is true, then $\cl{\Lambda_+}$ is $G$-invariant, and applying Corollary \ref{corollaryminimal} completes the proof.
%%
%%Fix $\xi\in\Lambda$ which is not a global fixed point of $G$, and choose $h\in G$ such that $h(\xi)\neq \xi$. Fix $\epsilon > 0$ small enough so that $\Dist(B,h(B)) > \epsilon$, where $B = B(\xi,\epsilon)$. Let $\sigma > 0$ be large enough so that the Big Shadows Lemma \ref{lemmabigshadow} holds. Since $\xi\in\Lambda$, there exists $g\in G$ such that
%%\[
%%\Shad(g(\zero),\sigma) \subset B.
%%\]
%%Let $Z = g^{-1}(\Shad(g(\zero),\sigma)) = \Shad_{g^{-1}(\zero)}(\zero,\sigma)$. Then by Lemma \ref{lemmabigshadow}, $\Diam(\del X\butnot Z)\leq\epsilon$. Thus $\del X\butnot Z$ can intersect at most one of the sets $B$, $h(B)$. So $B\subset Z$ or $h(B)\subset Z$. If $B\subset Z$ then
%%\[
%%g(B) \subset B \text{ and } B \subset \Shad_{g^{-1}(\zero)}(\zero,\sigma),
%%\]
%%whereas if $h(B)\subset Z$ then
%%\[
%%gh(B) \subset B \text{ and } B \subset \Shad_{(gh)^{-1}(\zero)}(\zero,\sigma + \dogo h).
%%\]
%%So by Lemma \ref{lemmaloxodromic}, we have $j_+\in B$, where $j = g$ or $j = gh$ is a loxodromic isometry.
%%%Fix $\sigma > 0$ large to be determined, and let $r$ be large enough so that Corollary \ref{corollaryloxodromic} holds. Since $\xi\in\Lambda$, there exists $g\in G$ with $\lb g(\zero)|\xi\rb_\zero \geq \sigma$. By contradiction suppose that $g$ and $hg$ are both not loxodromic. Then
%%%\begin{align*}
%%%\Dist(\xi,h(\xi)) &\leq_{\phantom{\times,h}} \Dist(g(\zero),\xi) + \Dist(g(\zero),g^{-1}(\zero)) + \Dist(g^{-1}(\zero),g^{-1}h^{-1}(\zero))\\
%%%&\hspace{.5 in} + \Dist(g^{-1}h^{-1}(\zero), hg(\zero)) + \Dist(hg(\zero),h(\xi))\\
%%%&\lesssim_{\times,h} b^{-\sigma} \tendsto\sigma 0.
%%%\end{align*}
%%%By choosing $\sigma$ large enough, this is a contradiction. Thus at least one of $g,hg$ is loxodromic. Let $j$ denote this loxodromic element. The proof of Lemma \ref{lemmaloxodromic} shows that [Continue\internal]
%%\end{proof}
%%
%%The following improvement over Proposition \ref{propositionminimal2} has a quite intricate proof:
%%
%%\begin{proposition}[Cf. {\cite[Theorem 5.3.8]{Beardon_book}}, {\cite[p.349]{Koebe}}]
%%\label{propositionminimal3}
%%Let $G\prec\Isom(X)$ be of general type. Then
%%\[
%%\{(g_+,g_-) : g\in G \text{ is loxodromic}\}
%%\]
%%is dense in $\Lambda(G)\times\Lambda(G^{-1})$. Here $G^{-1} = \{g^{-1} : g\in G\}$.
%%\end{proposition}
%%\begin{proof}
%%
%%\begin{claim}
%%\label{claimi04}
%%Let $g$ be a loxodromic isometry and fix $\epsilon > 0$. There exists $\delta = \delta(\epsilon,g)$ such that for all $\xi_1,\xi_2\in\del X$ with $\Dist(\xi_2,\Fix(g)) \geq \epsilon$,
%%\[
%%\#\{i = 0,\ldots,4 : \Dist(g^i(\xi_1),\xi_2) \leq \delta\} \leq 1.
%%\]
%%\end{claim}
%%\begin{subproof}
%%Suppose that $\Dist(g^i(\xi_1),\xi_2) \leq \delta$ for two distinct values of $i$. Then $\Dist(g^{i_1}(\xi_1),g^{i_2}(\xi_1)) \leq 2\delta$. For every $n$, we have
%%\[
%%\Dist(g^{n + i_1}(\xi_1),g^{n + i_2}(\xi_1)) \lesssim_\times b^{\dogo{g^n}}\delta
%%\]
%%and thus by the triangle inequality
%%\[
%%\Dist(g^{i_1}(\xi_1),g^{n(i_2 - i_1) + i_1}(\xi_1)) \lesssim_{\times,n} \delta.
%%\]
%%By Theorem \ref{theoremloxodromicextra}, if $n$ is sufficiently large then $\Dist(g^{n(i_2 - i_1) + i_1}(\xi_1), g_+) \leq \epsilon/2$, which implies that
%%\[
%%\epsilon/2 \leq \Dist(\xi_2,\Fix(g)) - \Dist(g^{n(i_2 - i_1) + i_1}(\xi_1), g_+) \leq \Dist(\xi_2, g^{n(i_2 - i_1) + i_1}(\xi_1)) \lesssim_{\times,n} \delta,
%%\]
%%which is a lower bound on $\delta$ independent of $\xi_1,\xi_2$. Choosing $\delta$ less than this lower bound yields a contradiction.
%%\end{subproof}
%%
%%\begin{claim}
%%\label{claimxi1234}
%%There exist $\epsilon,\rho > 0$ such that for all $\xi_1,\xi_2,\xi_3,\xi_4\in\Lambda$, there exists $j\in G$ such that
%%\begin{equation}
%%\label{xi1234}
%%\Dist(j(\xi_k),\xi_\ell) \geq \epsilon \all k = 1,2 \all \ell = 3,4 \text{ and } \dogo j \leq \rho.
%%\end{equation}
%%\end{claim}
%%\begin{subproof}
%%Fix $g,h\in G$ loxodromic with $\Fix(g)\cap\Fix(h) = \emptyset$, and let
%%\[
%%\rho = \max_{i = 0}^4 \max_{j = 0}^4 \dogo{g^i h^j}.
%%\]
%%Now fix $\xi_1,\xi_2,\xi_3,\xi_4\in\Lambda$. By Claim \ref{claimi04}, for each $k = 1,2$ and $\eta\in\Fix(g)$, we have
%%\[
%%\#\{j = 0,\ldots,4 : \Dist(h^j(\xi_k),\eta) \leq \delta_1 := \delta(\Dist(\Fix(g),\Fix(h)),h)\} \leq 1.
%%\]
%%It follows that there exists $j\in\{0,\ldots,4\}$ such that $\Dist(h^j(\xi_k),\eta) \geq \delta_1$ for all $k = 1,2$ and $\eta\in\Fix(g)$. Applying Claim \ref{claimi04} again, we see that for each $k = 1,2$ and $\ell = 3,4$, we have
%%\[
%%\#\{i = 0,\ldots,4 : \Dist(g^{-i}(\xi_\ell),h^j(\xi_k)) \leq \delta_2 := \delta(\delta_1,g^{-1})\} \leq 1.
%%\]
%%It follows that there exists $i\in\{0,\ldots,4\}$ such that $\Dist(g^{-i}(\xi_\ell),h^j(\xi_k)) \geq \delta_2$ for all $k = 1,2$ and $\ell = 3,4$. But then
%%\[
%%\Dist(g^i h^j(\xi_k) , \xi_\ell) \gtrsim_\times \delta_2,
%%\]
%%completing the proof.
%%\end{subproof}
%%
%%Now fix $\xi_+\in\Lambda$, $\xi_-\in\Lambda(G^{-1})$ distinct, and fix $\delta > 0$ arbitrarily small. By the definition of $\Lambda$, there exist $g,h\in G$ such that
%%\[
%%\Dist(g(\zero),\xi_+) , \Dist(h^{-1}(\zero),\xi_-) \leq \delta.
%%\]
%%Let $\sigma > 0$ be large enough so that the Big Shadows Lemma \ref{lemmabigshadow} holds for $\epsilon = \epsilon/2$, where $\epsilon$ is as in Claim \ref{claimxi1234}. Then
%%\[
%%\Diam(\del X\butnot\Shad_{g^{-1}(\zero)}(\zero,\sigma)) , \Diam(\del X\butnot \Shad_{h(\zero)}(\zero,\sigma)) \leq \epsilon/2.
%%\]
%%On the other hand,
%%\[
%%\Diam(\Shad(g(\zero),\sigma)) , \Diam(\Shad(h^{-1}(\zero),\sigma)) \lesssim_\times \delta.
%%\]
%%Since $\Shad(g(\zero),\sigma)$ is far from $h^{-1}(\zero)$ and $\Shad(h^{-1}(\zero),\sigma)$ is far from $g(\zero)$, the Bounded Distortion Lemma \ref{lemmaboundeddistortion} gives
%%\[
%%\Diam(h(\Shad(g(\zero),\sigma))) , \Diam(g^{-1}(\Shad(h^{-1}(\zero),\sigma))) \lesssim_\times \delta.
%%\]
%%Choose $\xi_1\in h(\Shad(g(\zero),\sigma))$, $\xi_2\in g^{-1}(\Shad(h^{-1}(\zero),\sigma))$, $\xi_3\in \del X\butnot\Shad_{g^{-1}(\zero)}(\zero,\sigma)$ and $\xi_4\in \del X\butnot\Shad_{h(\zero)}(\zero,\sigma)$. By Claim \ref{claimxi1234}, there exists $j\in G$ such that \eqref{xi1234} holds. Then
%%\[
%%\Diam(jh(\Shad(g(\zero),\sigma))) , \Diam(j^{-1} g^{-1}(\Shad(h^{-1}(\zero),\sigma))) \lesssim_\times \delta,
%%\]
%%and by choosing $\delta$ sufficiently small, we can make these diameters less than $\epsilon/2$. It follows that
%%\[
%%jh(\Shad(g(\zero),\sigma)) \subset \Shad_{g^{-1}(\zero)}(\zero,\sigma) \text{ and } j^{-1} g^{-1}(\Shad(h^{-1}(\zero),\sigma)) \subset \Shad_{h(\zero)}(\zero,\sigma)
%%\]
%%or equivalently
%%\[
%%gjh(\Shad(g(\zero),\sigma)) \subset \Shad(g(\zero),\sigma) \text{ and } (gjh)^{-1}(\Shad(h^{-1}(\zero),\sigma)) \subset \Shad(h^{-1}(\zero),\sigma).
%%\]
%%By Lemma \ref{lemmaloxodromic}, it follows that $gjh$ is a loxodromic isometry satisfying
%%\[
%%(gjh)_+\in\Shad(g(\zero),\sigma) ,\; (gjh)_-\in\Shad(h^{-1}(\zero),\sigma).
%%\]
%%In particular $\Dist((gjh)_+,\xi_+) , \Dist((gjh)_-,\xi_-) \lesssim_\times \delta$. Since $\delta$ was arbitrary, this completes the proof.
%%\end{proof}
%%
%%%%%%%%%%%%%%%%%%%%%%%%%%%%%%%%%%%%
%%\bigskip
%%\subsection{Convex hulls}
%%\label{subsectionconvexhulls}
%%%%%%%%%%%%%%%%%%%%%%%%%%%%%%%%%%%%
%%
%%In this subsection, we assume that $X$ is regularly geodesic (see Subsection \ref{subsectiongeodesicsCAT}). Recall that for $x,y\in\bord X$, $\geo xy$ denotes the geodesic segment, line, or ray joining $x$ and $y$.
%%
%%\begin{definition}
%%\label{definitionconvexhull}
%%Given $S\subset\bord X$, let
%%\begin{align*}
%%\Hull_1(S) &:= \bigcup_{x,y\in S} \geo xy,\\
%%\Hull_n(S) &:= \underbrace{\Hull_1\cdots \Hull_1}_{\text{$n$ times}}(S)\\
%%\Hull_\infty(S) &:= \cl{\bigcup_{n\in\N}\Hull_n(S)}.
%%\end{align*}
%%The set $\Hull_n(S)$ will be called the \emph{$n$th convex hull} of $S$. Moreover, $\Hull_\infty(S)$ will be called the \emph{convex hull} of $S$, and $\Hull_1(S)$ will be called the \emph{quasiconvex hull} of $S$.
%%\end{definition}
%%
%%The terminology ``convex hull'' comes from the following fact:
%%
%%\begin{proposition}
%%\label{propositionconvexhull}
%%$\Hull_\infty(S)$ is the smallest closed set $F\subset \bord X$ such that $S\subset F$ and
%%\begin{equation}
%%\label{convex}
%%\geo xy\subset F \all x,y\in F.
%%\end{equation}
%%\end{proposition}
%%A set $F$ satisfying \eqref{convex} will be called \emph{convex}.
%%\begin{proof}
%%It is clear that $S\subset\Hull_\infty(S)\subset\bord X$. To show that $\Hull_\infty(S)$ is convex, fix $x,y\in\Hull_\infty(S)$. Then there exist sequences $A\ni x_n\to x$ and $A\ni y_n\to y$, where $A = \bigcup_{n\in\N}\Hull_n(S)$. For each $n$, $\geo{x_n}{y_n}\subset A\subset \Hull_\infty(S)$. But since $X$ is regularly geodesic, $\geo{x_n}{y_n}\to \geo xy$ in the Hausdorff metric on $\bord X$. Since $\Hull_\infty(S)$ is closed, it follows that $\geo xy\subset\Hull_\infty(S)$.
%%
%%Conversely, if $S\subset F\subset \bord X$ is a closed convex set, then an induction argument shows that $F\supset \Hull_n(S)$ for all $n$. Since $F$ is closed, we have $F\supset \Hull_\infty(S)$.
%%\end{proof}
%%
%%Another connection between the operations $\Hull_1$ and $\Hull_\infty$ is given by the following proposition:
%%
%%\begin{proposition}
%%\label{propositionhull}
%%Suppose that $X$ is a \ROSS. Then there exists $\tau > 0$ such that for every set $S\subset\bord X$ we have
%%\[
%%X\cap \Hull_1(S) \subset X\cap \Hull_\infty(S) \subset \thicken_\tau(X\cap\Hull_1(S)).
%%\]
%%(Recall that $\thicken_\tau(S)$ denotes the $\tau$-thickening of a set with respect to the hyperbolic metric $\dist$.)
%%\end{proposition}
%%\begin{proof}
%%The proof will proceed using the ball model $X = \B = \B_\F^\alpha$. We will need the following lemma:
%%\begin{lemma}
%%\label{lemmaF}
%%There exists a closed convex set $F\propersubset\bord\B$ whose interior intersects $\del\B$.
%%\end{lemma}
%%\begin{subproof}
%%If $\alpha < \infty$, this is a consequence of \cite[Theorem 3.3]{Anderson}.
%%
%%We will use the finite-dimensional case to prove the infinite-dimensional case. Suppose that $\alpha$ is infinite. Let $Y = \B_\F^3 \subset X$. Then by the $\alpha < \infty$ case of Lemma \ref{lemmaF}, there exists a closed convex set $F_2 \propersubset\bord Y$ whose interior intersects $\del Y$, say $\xi\in \interior(F_2)\cap\del Y$. Choose $\epsilon > 0$ such that $B_Y(\xi,\epsilon) \subset F_2$. Then
%%\[
%%F_1 := \Hull_\infty(B_Y(\xi,\epsilon)) \subset F_2 \propersubset \bord Y
%%\]
%%by Proposition \ref{propositionconvexhull}. On the other hand, $F_1$ is invariant under the action of the group
%%\[
%%G_1 := \{g\in\Isom(Y) : g(\0) = \0, g(\xi) = \xi\}.
%%\]
%%Let
%%\[
%%G = \{g\in\Isom(X) : g(\0) = \0, g(\xi) = \xi\},
%%\]
%%and note that $G(\bord Y) = \bord X$. Let $F = G(F_1)$, and note that $F\cap \bord Y = F_1$. We claim that $F$ is convex. Indeed, suppose that $x,y\in F$; then there exists $g\in G$ such that $g(x),g(y)\in\bord Y$. (Note that in this step, we need all three dimensions of $Y$.) Then $g(x),g(y)\in F\cap \bord Y = F_1$, so by the convexity of $F_1$ we have $g(\geo xy) = \geo{g(x)}{g(y)} \subset F_1 \subset F$. Since $F$ is $G$-invariant, we have $\geo xy\subset F$.
%%
%%In addition to being convex, $F$ is also closed and contains the set $G(B_Y(\xi,\epsilon)) = B_X(\xi,\epsilon)$. Thus, $\xi\in\interior(F)$. Finally, since $F\cap \bord Y = F_1 \propersubset \bord Y$, it follows that $F\propersubset\bord X$.
%%\end{subproof}
%%Let $F$ be as in Lemma \ref{lemmaF}. Since $F\propersubset\bord\B$ is a closed set, it follows that $\B\butnot F\neq\emptyset$. By the transitivity of $\Isom(\B)$ (Observation \ref{observationtransitivity}), we may without loss of generality assume that $\0\in\B\butnot F$. By the transitivity of $\Stab(\Isom(\B);\0)$ on $\del\B$, we may without loss of generality assume that $\ee_1\in\interior(F)$. Fix $\epsilon > 0$ such that $B(\ee_1,\epsilon) \subset F$.
%%
%%We now proceed with the proof of Proposition \ref{propositionhull}. It is clear from the definitions that $\B\cap\Hull_1(S) \subset \B\cap\Hull_\infty(S)$. To prove the second inclusion, fix $\zz\in\B \butnot \thicken_\tau(\Hull_1(S))$ and we will show that $\zz\notin\Hull_\infty(S)$. By the transitivity of $\Isom(\B)$, we may without loss of generality assume that $\zz = \0$. Now for every $\xx, \yy\in S$, we have $\zz = \0 \notin \thicken_\tau(\geo\xx\yy)$. By (i) of Proposition \ref{propositionrips}, we have
%%\[
%%\lb \xx | \yy \rb_\0 \gtrsim_\plus \tau
%%\]
%%and thus by \eqref{euclideanball1},
%%\[
%%\|\yy - \xx\| \lesssim_\times e^{-\tau}.
%%\]
%%By choosing $\tau$ sufficiently large, this implies that
%%\[
%%\|\yy - \xx\| \leq \epsilon/2 \all \xx,\yy\in S.
%%\]
%%Moreover, since $\dist(\0,\xx) = 2\lb \xx | \xx \rb_\0 \gtrsim_\plus \tau$, by choosing $\tau$ sufficiently large we may guarantee that
%%\[
%%\|\xx\| \geq 1 - \epsilon/2 \all \xx\in S.
%%\]
%%Since the claim is trivial if $S = \emptyset$, assume that $S\neq \emptyset$ and choose $\xx\in S$. Without loss of generality, assume that $\xx = \lambda \ee_1$ for some $\lambda \geq 2/3$. Then $S \subset B_\EE(\xx,\epsilon/2) \subset B(\ee_1,\epsilon) \subset F$. But then $F$ is a closed convex set containing $S$, so by Proposition \ref{propositionconvexhull} $\Hull_\infty(S) \subset F$. Since $\zz = \0 \notin F$, it follows that $\zz\notin\Hull_\infty(S)$.
%%\end{proof}
%%\begin{corollary}
%%\label{corollaryanderson}
%%Suppose that $X$ is a \ROSS. Then for every closed set $S\subset\bord X$, we have
%%\[
%%\Hull_\infty(S) \cap \del X = S\cap \del X.
%%\]
%%\end{corollary}
%%\begin{proof}
%%The inclusion $\supset$ is immediate. Suppose that $\xi\in\Hull_\infty(S)\cap\del X$, and find a sequence $X\cap\Hull_\infty(S)\ni x_n\to \xi$. By Proposition \ref{propositionhull}, for each $n$ there exist $y_n^{(1)},y_n^{(2)}\in S$ such that $x_n\in\thicken_\tau(\geo{y_n^{(1)}}{y_n^{(2)}})$; by Proposition \ref{propositionrips} we have $\lb y_n^{(1)} | y_n^{(2)}\rb_{x_n} \asymp_\plus 0$. Applying Lemma \ref{lemmaminimal} gives $\xi\in S$.
%%\end{proof}
%%
%%\begin{remark}
%%Corollary \ref{corollaryanderson} was proven for the case where $X$ is a pinched (finite-dimensional) Hadamard manifold and $S\subset\del X$ by M. T. Anderson \cite[Theorem 3.3]{Anderson}. It was conjectured to hold whenever $X$ is ``strictly convex'' by Gromov \cite[p.11]{Gromov1}, who observed that it holds in the Standard Case. However, this conjecture was proven to be false independently by A. Ancona \cite[Corollary C]{Ancona} and A. Borb\'ely \cite[Theorem 1]{Borbely}, who each constructed a three-dimensional CAT(-1) manifold $X$ and a point $\xi\in\del X$ such that for every neighborhood $U$ of $\xi$, $\Hull_\infty(U) = \bord X$.% As late as 2006, it was stated by N. Monod \cite[p.789]{Monod} that Corollary \ref{corollaryanderson} was open for infinite-dimensional complex hyperbolic spaces. In particular, in the notation of \cite{Monod}, Corollary \ref{corollaryanderson} shows that $\scrT_c \given \del X = \scrT \given\del X$ whenever $X$ is a \ROSS\ (in particular, including the case where $X$ is an infinite-dimensional complex hyperbolic space).
%%
%%
%%%If $X$ is not assumed to be a \ROSS, then the validity of Proposition \ref{propositionhull} is far from clear. Its proof depended on the existence of ``large'' convex sets (i.e. those whose interior contains a point of $\del X$) which are nonetheless not equal to all of $\bord X$. The existence of such sets holds for many CAT(-1) spaces which it is easy to think of, and in particular for complex and quaternionic hyperbolic space, $\R$-trees, and for two-dimensional Riemannian manifolds with curvature $\leq -1$, but it is by no means clear that it holds for all regularly geodesic hyperbolic metric spaces, or even for all CAT(-1) spaces.
%%
%%Thus, although the $\infty$-convex hull has more geometric and intuitive appeal based on Proposition \ref{propositionconvexhull}, without more hypotheses there is no way to restrain its geometry. The $1$-convex hull is thus more useful for our applications. Proposition \ref{propositionhull} indicates that in the case of a \ROSS, we are not losing too much by the change.
%%\end{remark}
%%
%%
%%\begin{definition}
%%\label{definitionconvexcore}
%%The \emph{convex core} of a semigroup $G\prec\Isom(X)$ is the set
%%\[
%%\CC_\Lambda := X\cap \Hull_\infty(\LambdaG),
%%\]
%%and the \emph{quasiconvex core} is the set
%%\[
%%\CC_\zero := X\cap \cl{\Hull_1(G(\zero))}.
%%\]
%%\end{definition}
%%
%%\begin{observation}
%%The convex core and quasiconvex core are both closed $G$-invariant sets. The quasiconvex core depends on the distinguished point $\zero$. However:
%%\end{observation}
%%
%%\begin{proposition}
%%\label{propositionC0comparison}
%%Fix $x,y\in X$. Then
%%\[
%%\CC_x \subset \thicken_R(\CC_y)
%%\]
%%for some $R > 0$.
%%\end{proposition}
%%\begin{proof}
%%Fix $z\in\CC_x$. Then $z\in \geo{g(x)}{h(x)}$ for some $g,h\in G$. It follows that
%%\[
%%\lb g(y)|h(y)\rb_z \asymp_\plus \lb g(x)|h(x)\rb_z = 0.
%%\]
%%So by Proposition \ref{propositionrips}, $\dist(z,\geo{g(y)}{h(y)}) \asymp_\plus 0$. But $\geo{g(y)}{h(y)} \subset \CC_y$, so $\dist(z,\CC_y) \asymp_\plus 0$. Letting $R$ be the implied constant completes the proof.
%%\end{proof}
%%
%%\begin{remark}
%%In many cases, we can get information about the action of $G$ on $X$ by looking just at its restriction to $\CC_\Lambda$ or to $\CC_\zero$. We therefore also remark that if $X$ is a CAT(-1) space, then $\CC_\Lambda$ is also a CAT(-1) space.
%%\end{remark}
%%
%%In the sequel the following notation will be useful:
%%\begin{notation}
%%For a set $S\subset\bord X$ let
%%\begin{equation}
%%\label{primenotation}
%%S' = \cl S \cap \del X.
%%\end{equation}
%%\end{notation}
%%
%%\begin{observation}
%%\label{observationboundaryofconvexcore}
%%$(\CC_\zero)' = \LambdaG$.
%%\end{observation}
%%\begin{proof}
%%Since $\LambdaG = (G(\zero))'$ and $G(\zero) \subset\CC_\zero$, we have $(\CC_\zero)' \supset \LambdaG$. Suppose that $\xi\in(\CC_\zero)'$, and let $\CC_\zero\ni x_n\to \xi$. By definition, for each $n$ there exist $y_n^{(1)},y_n^{(2)}\in G(\zero)$ such that $x_n\in \geo{y_n^{(1)}}{y_n^{(2)}}$. Lemma \ref{lemmaminimal} completes the proof.
%%\end{proof}
%%
%%%%%%%%%%%%%%%%%%%%%%%%%%%%%%%%%%%%
%%\bigskip
%%\subsection{Semigroups which act irreducibly on \ROSSs}
%%\label{subsectionirreducibly}
%%%%%%%%%%%%%%%%%%%%%%%%%%%%%%%%%%%%
%%
%%\begin{definition}
%%\label{definitionreducibly}
%%Suppose that $X$ is a \ROSS, and fix $G\prec\Isom(X)$. We shall say that $G$ \emph{acts reducibly} on $X$ if there exists a nontrivial totally geodesic $G$-invariant subset $S\propersubset \bord X$. Otherwise, we shall say that $G$ \emph{acts irreducibly} on $X$.
%%\end{definition}
%%
%%\begin{remark}
%%A parabolic or focal subsemigroup of $\Isom(X)$ may act either reducibly or irreducibly on $X$.
%%\end{remark}
%%
%%%\comdavid{May need some rewording with our new def of $\CC_\zero$}
%%%
%%%\comdavid{Also, we should include the parabolic case to be compatible with Proposition \ref{propositionparabolicdiscrete}.}
%%%
%%%\comdavid{Actually, I disagree, I like the idea of all actions being reducible to an irreducible action.}
%%
%%\begin{proposition}
%%\label{propositionactsreducibly}
%%Let $G\prec\Isom(X)$ be nonelementary. Then the following are equivalent:
%%\begin{itemize}
%%\item[(A)] $G$ acts reducibly on $X$.
%%\item[(B)] There exists a nontrivial totally geodesic subset $S\propersubset \bord X$ such that $\LambdaG\subset S$.
%%\item[(C)] There exists a nontrivial totally geodesic subset $S\propersubset \bord X$ such that $\CC_\Lambda\subset S$.
%%\item[(D)] There exists a nontrivial totally geodesic subset $S\propersubset \bord X$ such that $\CC_\zero\subset S$ for some $\zero\in X$.
%%\end{itemize}
%%\end{proposition}
%%\begin{proof}[Proof of \text{(A) \implies (B)}]
%%Let $S\propersubset \bord X$ be a nontrivial totally geodesic $G$-invariant subset. Fix $\zero\in S\cap X$. Then $\LambdaG\subset\cl{G(\zero)}\subset S$.
%%\end{proof}
%%\begin{proof}[Proof of \text{(B) \implies (C)}]
%%If $S$ is any totally geodesic set which contains $\LambdaG$, then $S$ is a closed convex set containing $\LambdaG$, so by Proposition \ref{propositionconvexhull}, $\CC_\Lambda\subset S$.
%%\end{proof}
%%\begin{proof}[Proof of \text{(C) \implies (D)}]
%%Since $G$ is nonelementary, $\CC_\Lambda \neq \emptyset$. Fix $\zero\in\CC_\Lambda$; then $\CC_\zero\subset \CC_\Lambda$.
%%\end{proof}
%%\begin{proof}[Proof of \text{(D) \implies (A)}]
%%Let $S$ be the smallest totally geodesic subset of $X$ which contains $\CC_\zero$, i.e.
%%\[
%%S := \bigcap\{W:W\supseteq \CC_\zero\text{ totally geodesic}\}.
%%\]
%%Then our hypothesis implies that $S\propersubset \bord X$. Since $\zero\in S$, $S$ is nontrivial. It is obvious from the definition that $S$ is $G$-invariant. This completes the proof.
%%\end{proof}
%%
%%\begin{remark}
%%If $G\prec\Isom(X)$ is nonelementary, then Proposition \ref{propositionactsreducibly} gives us a way to find a nontrivial totally geodesic set on which $G$ acts reducibly; namely, the smallest totally geodesic set containing $\LambdaG$, or equivalently $\CC_G$, will have this property (cf. Lemma \ref{lemmatotallygeodesic}). On the other hand, there exists a parabolic group $G\leq\Isom(\H^\infty)$ such that $G$ does not act irreducibly on any nontrivial totally geodesic subset $S\subset\bord \H^\infty$ (Remark \ref{remarkparabolictorsion}).
%%\end{remark}
%%
%%
%%
%%%%%%%%%%%%%%%%%%%%%%%%%%%%%%%%%%%%
%%\bigskip
%%\subsection{Semigroups of compact type}
%%%%%%%%%%%%%%%%%%%%%%%%%%%%%%%%%%%%
%%
%%\begin{definition}
%%\label{definitioncompacttype}
%%We say that a semigroup $G\prec\Isom(X)$ is of \emph{compact type} if its limit set $\Lambda$ is compact.
%%\end{definition}
%%
%%\begin{proposition}
%%\label{propositioncompacttype}
%%For $G\prec\Isom(X)$, the following are equivalent:
%%\begin{itemize}
%%\item[(A)] $G$ is of compact type.
%%\item[(B)] Every sequence $(x_n)_1^\infty$ in $G(\zero)$ with $\dox{x_n}\to \infty$ has a convergent subsequence.
%%\end{itemize}
%%Furthermore, if $X$ is regularly geodesic, then \text{(A)-(B)} are equivalent to:
%%\begin{itemize}
%%\item[(C)] The set $\CC_\zero$ is a proper metric space.
%%\end{itemize}
%%and if $X$ is a \ROSS, then they are equivalent to:
%%\begin{itemize}
%%\item[(D)] The set $\CC_\Lambda$ is a proper metric space.
%%\end{itemize}
%%\end{proposition}
%%\begin{proof}[Proof of \text{(A)} \implies \text{(B)}]
%%Fix a sequence $(g_n)_1^\infty$ in $G$ with $\dogo{g_n}\to \infty$. The existence of such a sequence implies that $G$ is not elliptic. If $G$ is parabolic or inward focal, then the proof of Proposition \ref{propositioncardinalitylimitset}(ii) shows that $g_n(\zero)\to \xi$, where $\Lambda = \{\xi\}$. So we may assume that $G$ is lineal, outward focal, or of general type, in which case Proposition \ref{propositioncardinalitylimitset} gives $\#(\Lambda)\geq 2$.
%%
%%Fix distinct $\xi_1,\xi_2\in\LambdaG$, and let $(n_k)_1^\infty$ be a sequence such that $(g_{n_k}(\xi_i))_1^\infty$ converges for $i = 1,2$, and such that
%%\[
%%\lb g_{n_k}^{-1}(\zero)|\xi_1\rb_\zero \leq \lb g_{n_k}^{-1}(\zero)|\xi_2\rb_\zero
%%\]
%%for all $k$. (If this is not possible, switch $\xi_1$ and $\xi_2$.) We have
%%\[
%%0 \asymp_{\plus,\xi_1,\xi_2} \lb \xi_1|\xi_2\rb_\zero
%%\gtrsim_\plus \min\left(\lb g_{n_k}^{-1}(\zero)|\xi_1\rb_\zero,\lb g_{n_k}^{-1}(\zero)|\xi_2\rb_\zero\right)
%%= \lb g_{n_k}^{-1}(\zero)|\xi_1\rb_\zero
%%\]
%%and thus
%%\[
%%\lb g_{n_k}(\zero)|g_{n_k}(\xi_1)\rb_\zero \asymp_{\plus,\xi_1,\xi_2} \dogo{g_{n_k}} \tendsto n \infty.
%%\]
%%On the other hand, there exists $\eta\in\LambdaG$ such that $g_{n_k}(\xi_1)\tendsto k\eta$, and thus
%%\[
%%\lb g_{n_k}(\xi_1)|\eta\rb_\zero\tendsto n \infty.
%%\]
%%Applying Gromov's inequality yields
%%\[
%%\lb g_{n_k}(\zero)|\eta\rb_\zero \tendsto n \infty
%%\]
%%and thus $g_{n_k}(\zero)\tendsto k\eta$. This completes the proof.
%%\end{proof}
%%\begin{proof}[Proof of \text{(B)} \implies \text{(A)}]
%%Fix a sequence $(\xi_n)_1^\infty$ in $\LambdaG$. For each $n\in\N$, choose $g_n\in G$ with
%%\[
%%\lb g_n(\zero)|\xi_n\rb_\zero \geq n.
%%\]
%%In particular $\dogo{g_n}\geq n\tendsto n \infty$. Thus by our hypothesis, there exists a convergent subsequence $g_{n_k}(\zero)\tendsto k\eta\in\LambdaG$. Now
%%\[
%%\Dist(\xi_{n_k},\eta) \leq \Dist(g_{n_k}(\zero),\xi_{n_k}) + \Dist(g_{n_k}(\zero),\eta)
%%\lesssim_\times b^{-n_k} + \Dist(g_{n_k}(\zero),\eta) \tendsto k 0,
%%\] 
%%i.e. $\xi_{n_k}\tendsto k\eta$.
%%\end{proof}
%%\begin{proof}[Proof of \text{(A)} \implies \text{(C)}]
%%Let $(x_n)_1^\infty$ be a bounded sequence in $\CC_\zero$. For each $n\in\N$, there exist $y_n^{(1)},y_n^{(2)}\in G(\zero)$ such that $x_n\in\thicken_{1/n}(\geo{y_n^{(1)}}{y_n^{(2)}})$. Choose a sequence $(n_k)_1^\infty$ on which $y_{n_k}^{(1)}\tendsto k\alpha$ and $y_{n_k}^{(2)}\tendsto k\beta$. Since $X$ is regularly geodesic we have
%%\[
%%\geo{y_{n_k}^{(1)}}{y_{n_k}^{(2)}}\tendsto k\geo{y^{(1)}}{y^{(2)}}.
%%\]
%%For each $k$, choose $z_k\in \geo{y_{n_k}^{(1)}}{y_{n_k}^{(2)}}$ with $\dist(x_{n_k},z_k)\leq 1/n_k$. Since the sequence $(z_k)_1^\infty$ is bounded, it must have a subsequence which converges to a point in $\geo{y^{(1)}}{y^{(2)}}$; it follows that the corresponding subsequence of $(x_{n_k})_1^\infty$ is also convergent. Thus every bounded sequence in $\CC_\zero$ has a convergent subsequence, so $\CC_\zero$ is proper.
%%\end{proof}
%%\begin{proof}[Proof of \text{(C)} \implies \text{(B)}]
%%Obvious since $G(\zero)\subset\CC_\zero$.
%%\end{proof}
%%\begin{proof}[Proof of \text{(A)} \implies \text{(D)}]
%%Note first of all that we cannot get (A) \implies (D) immediately from Proposition \ref{propositionhull}, since the $\tau$-thickening of a compact set is no longer compact.
%%
%%By \cite[Proposition 1.5]{Bowditch_convexity}, there exists a metric $\rho$ on $\bord X$ compatible with the topology such that the map $F\mapsto \Hull_1(F)$ is a semicontraction with respect to the Hausdorff metric of $(\bord X,\rho)$. (Finite-dimensionality is not used in any crucial way in the proof of \cite[Proposition 1.5]{Bowditch_convexity},\Footnote{One should keep in mind that the Cartan--Hadamard theorem \cite[IX, Theorem 3.8]{Lang_differential_geometry} can be used as a substitute for the Hopf--Rinow theorem in most circumstances.} and in any case for \ROSSs\ it can be proven by looking at finite-dimensional subsets, as we did in the proof of Proposition \ref{propositionhull}.) We remark that if $\F = \R$, then such a metric $\rho$ can be prescribed explicitly: if $X = \B$ is the ball model, then the Euclidean metric on $\bord\B\subset\HH$ has this property, due to the fact that geodesics in the ball model are line segments in $\HH$ (cf. \eqref{geodesic}).\Footnote{Recall that our ``ball model'' $\amsbb B$ is the Klein model rather than the Poincar\'e model.}
%%
%%Now let us demonstrate (D). It suffices to show that $\bord\CC_\Lambda = \Hull_\infty(\Lambda)$ is compact. Since $\Hull_\infty(\Lambda)$ is by definition closed, it suffices to show that $\Hull_\infty(\Lambda)$ is totally bounded with respect to the $\rho$ metric. Indeed, fix $\epsilon > 0$. Since $\Lambda$ is compact, there is a finite set $F_\epsilon\subset\Lambda$ such that
%%\[
%%\Lambda\subset \thicken_{\epsilon/2}(F_\epsilon).
%%\]
%%(In this proof, all neighborhoods are taken with respect to the $\rho$ metric.) Let $X_\epsilon \subset X$ be a finite-dimensional totally geodesic set containing $F_\epsilon$. Then $\Lambda\subset \thicken_{\epsilon/2}(X_\epsilon)$. On the other hand, since $X_\epsilon$ is compact, there exists a finite set $F_\epsilon'\subset X_\epsilon$ such that $X_\epsilon\subset \thicken_{\epsilon/2}(F_\epsilon')$.
%%
%%Now, our hypothesis on $\rho$ implies that
%%\[
%%\Hull_1(\thicken_{\epsilon/2}(X_\epsilon)) \subset \thicken_{\epsilon/2}(\Hull_1(X_\epsilon)) = \thicken_{\epsilon/2}(X_\epsilon),
%%\]
%%and thus that $\thicken_{\epsilon/2}(X_\epsilon)$ is convex. But $\Lambda\subset \thicken_{\epsilon/2}(X_\epsilon)$, so $\Hull_\infty(\Lambda) \subset \thicken_{\epsilon/2}(X_\epsilon)$. Thus
%%\[
%%\Hull_\infty(\Lambda)\subset \thicken_{\epsilon/2}(X_\epsilon) \subset \thicken_\epsilon(F_\epsilon').
%%\]
%%Since $\epsilon$ was arbitrary, this shows that $\Hull_\infty(\Lambda)$ is totally bounded, completing the proof.
%%\end{proof}
%%\begin{proof}[Proof of \text{(D)} \implies \text{(B)}]
%%Since property (B) is clearly basepoint-independent, we may without loss of generality suppose $\zero\in\CC_\Lambda$. Then (D) \implies (C) \implies (B).
%%\end{proof}
%%
%%As an example of an application we prove the following corollary.
%%\begin{corollary}
%%Suppose that $X$ is regularly geodesic. Then any moderately discrete subgroup of $\Isom(X)$ of compact type is strongly discrete.
%%\end{corollary}
%%\begin{proof}
%%If $G$ is a moderately discrete group, then $G\given\CC_\zero$ is moderately discrete by Observation \ref{observationrestrictions}, and therefore strongly discrete by Propositions \ref{propositionproperMDSD} and \ref{propositioncompacttype}. Thus by Observation \ref{observationrestrictions}, $G$ is strongly discrete.
%%\end{proof}
%%
%%
%%A well-known characterization of the complement of the limit set in the Standard Case is that it is the set of points where the action of $G$ is discrete. We extend this characterization to hyperbolic metric spaces for groups of compact type:
%%
%%
%%\begin{proposition}
%%\label{propositionSDonclX}
%%Let $G\leq\Isom(X)$ be a strongly discrete group of compact type. Then the action of $G$ on $\bord X\butnot\LambdaG$ is strongly discrete in the following sense: For any set $S\subset\bord X\butnot\LambdaG$ satisfying
%%\begin{equation}
%%\label{Lambdabounded}
%%\Dist(S,\LambdaG) > 0,
%%\end{equation}
%%we have
%%\[
%%\#\{g\in G : g(S)\cap S \neq \emptyset\} < \infty.
%%\]
%%\end{proposition}
%%%\comdavid{Note that in the case where $G$ is parabolic with fixed point $\xi$, the condition \eqref{Lambdabounded} reduces to $S$ being $\xi$-bounded.}
%%\begin{proof}
%%By contradiction, suppose that there exists a sequence of distinct $(g_n)_1^\infty$ such that $g_n(S)\cap S\neq \emptyset$ for all $n\in\N$. Since $G$ is strongly discrete, we have $\dogo{g_n}\to\infty$, and since $G$ is of compact type there exist an increasing sequence $(n_k)_1^\infty$ and $\xi_+,\xi_-\in\LambdaG$ such that $g_{n_k}(\zero)\to \xi_+$ and $g_{n_k}^{-1}(\zero)\to\xi_-$. In the remainder of the proof we restrict to this subsequence, so that $g_n(\zero)\to \xi_+$ and $g_n^{-1}(\zero)\to\xi_-$.
%%
%%For each $n$, fix $x_n\in g_n^{-1}(g_n(S)\cap S)$, so that $x_n,g_n(x_n)\in S$. Then
%%\[
%%\Dist(x_n,\xi_-), \Dist(g_n(x_n),\xi_+) \geq \Dist(S,\LambdaG) \asymp_\times 1,
%%\]
%%and so
%%\[
%%\lb x_n|\xi_-\rb_\zero , \lb g_n(x_n)|\xi_+ \rb_\zero \asymp_\plus 0.
%%\]
%%On the other hand, $\lb g_n^{-1}(\zero)|\xi_-\rb_\zero, \lb g_n(\zero)|\xi_+\rb_\zero \to \infty$. Applying Gromov's inequality gives
%%\[
%%\lb x_n|g_n^{-1}(\zero)\rb_\zero , \lb g_n(x_n)|g_n(\zero) \rb_\zero \asymp_\plus 0
%%\]
%%for all $n$ sufficiently large. But then
%%\[
%%\dogo{g_n} = \lb g_n(x_n)|\zero\rb_{g_n(\zero)} + \lb g_n(x_n)|g_n(\zero) \rb_\zero \asymp_\plus 0,
%%\]
%%a contradiction.
%%\end{proof}
%%
%%\endinput